\documentclass[UTF8,fontset=fandol,openany,11pt]{ctexbook}
\usepackage[a4paper,inner=25mm,outer=23mm,top=24mm,bottom=24mm,headheight=15pt]{geometry}
\usepackage{amsmath,amssymb,capt-of,needspace}
\usepackage{xcolor,graphicx,longtable,booktabs,array,enumitem,fancyhdr,fvextra,xurl,hyperref,tikz,float}
\usetikzlibrary{arrows.meta,positioning,calc}
\definecolor{amber}{HTML}{BA872A}
\definecolor{bluegray}{HTML}{517E9B}
\definecolor{forest}{HTML}{207746}
\definecolor{leaf}{HTML}{58A92E}
\definecolor{mist}{HTML}{EAF3E8}
\definecolor{ink}{HTML}{23342D}
\definecolor{linegray}{HTML}{CBD5CE}
\hypersetup{colorlinks=true,linkcolor=forest,urlcolor=forest,citecolor=forest,bookmarksnumbered=true,pdftitle={Go 程序剖析与诊断},pdfauthor={GoLab 性能学院}}
\setcounter{tocdepth}{1}
\setcounter{secnumdepth}{2}
\setlength{\parskip}{4pt}
\setlength{\parindent}{2em}
\setlength{\emergencystretch}{3em}
\setlist{nosep,topsep=5pt,leftmargin=2em}
\ctexset{chapter={format=\Huge\bfseries\color{forest},beforeskip=8pt,afterskip=24pt},section={format=\Large\bfseries\color{forest}},subsection={format=\large\bfseries\color{ink}}}
\DefineVerbatimEnvironment{codeblock}{Verbatim}{breaklines=true,breakanywhere=true,fontsize=\footnotesize,frame=single,rulecolor=\color{linegray},framesep=5pt,tabsize=4}
\pagestyle{fancy}
\fancyhf{}
\fancyhead[L]{\small\color{forest}GoLab\quad Go 程序剖析与诊断}
\fancyhead[R]{\small\thepage}
\fancyfoot[C]{\scriptsize 原理 · 工具 · 实验 · 研究}
\renewcommand{\headrulewidth}{0.3pt}
\begin{document}
\begin{titlepage}
\pagecolor{mist}\color{ink}
\vspace*{15mm}
{\large\bfseries\color{forest}GOLAB / GRADUATE COURSE\par}
\vspace{24mm}
{\fontsize{34}{44}\selectfont\bfseries Go 程序剖析\par 与诊断\par}
\vspace{8mm}
{\Large 原理、工具、实验与研究前沿\par}
\vspace{20mm}
{\large pprof\quad execution trace\quad Delve\par}
\vspace{4mm}
{\large runtime / debug\quad perf / eBPF\quad PGO\par}
\vfill
\noindent\color{forest}\rule{\linewidth}{1.5pt}\par
\vspace{5mm}
{\large 中文研究生课程教材\quad ·\quad 机制深研增订版\par}
\vspace{4mm}
源码语义基线：Go 1.25.0\par
文献窗口：2021-10-02 至 2026-10-02\par
原教材检索日期：2026 年 10 月 2 日\par
机制深研查证：2026 年 10 月 4 日；增订：2026 年 10 月 5 日\par
\vspace{10mm}
\end{titlepage}
\nopagecolor\color{black}
\frontmatter
\chapter*{前言与使用说明}
\addcontentsline{toc}{chapter}{前言与使用说明}
本教材面向具有 Go 编程、操作系统和基础统计知识的研究生及工程师。学习目标是形成可复查的诊断与研究能力：从性能症状提出可反驳假设，通过采样、追踪、调试和系统证据定位机制，再以保持正确性的干预与重复实验验证结论。

本增订版保留原教材第 1–16 章，并在其后追加第 17–22 章机制深研。原教材按“基础—工具—应用—实现—研究”组织：第十五章研究成熟开源项目的实现，第十六章研究近五年前沿工作；这些章节的内容体系完整保留，并依据勘误记录修正 PGO 空输入边界及 trace 注解表述。第十七章重建测量契约，第十八至二十章分别深入 eBPF、pprof 与 Perfetto 的实现，第二十一章完成同一请求的跨工具实验，第二十二章形成组合诊断方法。

原教材各章包含学习目标、先修要求、正文、实验、测验解析和分析题答案，基础课程建议 16 周、每周 6 小时；项目评分量规见第十四章。完成基础课程后，可继续按六章深研的顺序学习，结合固定源码、图解、交互实验与推导练习复核工具机制。

涉及 Go 运行时与标准库的语义以 Go 1.25.0 源码为基线，必要处列明引入版本；第三方项目以各章注明的版本为准。该基线用于复查，不代表推荐部署旧补丁版本；实际部署应采用受支持的补丁版本。不同操作系统、架构、内核和运行时版本的功能支持需按具体环境核对。

第十六章的检索窗口为 2021 年 10 月 2 日至 2026 年 10 月 2 日。入选论文、工程发布和标准状态分开标注，不把标准快照当作经过同行评议的研究，也不把其他语言或网络系统的结果当作 Go 的直接实证。原始来源和固定版本见参考资料与附带研究证据文件。

\textbf{阅读与实践：}HTML 学习网站与本书由同一内容源生成。网站优先呈现机制图解和分步动画，并保留可展开的完整推导；支持离线答题、错题复习和本地进度保存。第二版逐项审计重要断言，并依据原始源码及论文的方法、评价和局限重写第十五、十六章。2.1 校订版在此基础上逐章复核正文、图注、实验、题目及答案，修正技术细节与歧义表达；逐项差异和证据见配套校订报告。实验代码位于配套目录；第十四章还内嵌完整基础服务源码。增订版统一了 22 章的目录、64 张同源机制图和 168 条参考来源；HTML 与 PDF 使用相同章节顺序。原教材的 39 张机制图保留，新增 25 张 SVG/TikZ 同源图及四项交互实验。所有图均为教学机制示意，原理实验中的合成输入不代表生产采样。本文件为独立 LaTeX 源码，无需外部图片、章节文件或 BibTeX 项目。

\textbf{验证边界：}教材中的假设数值和模型计算不作为生产实测结果。配套验证报告单独记录本次实际运行环境、通过项目和未运行能力；工具的构建与运行检查不能代替重复统计研究、线上灰度或生产收益评估。
\tableofcontents
\listoffigures
\mainmatter
\part*{第一篇 · 测量基础}
\addcontentsline{toc}{part}{第一篇 · 测量基础}
\chapter{性能剖析的科学方法}
\label{ch:01}
\noindent{\large\color{forest}先形成可证伪的假设，再选择观察工具。}\par\medskip
\noindent\textbf{先修要求：}具备 Go 函数、接口、goroutine 与基本 Linux 使用经验。

\noindent\textbf{完成本章后，应能够：}
\begin{itemize}
\item 区分工作量、吞吐、延迟、资源消耗与成本
\item 建立从症状到机制再到干预验证的证据链
\item 运用 Amdahl 定律与 Little 定律限制优化预期
\item 定义可复现、可反驳的性能结论
\end{itemize}
\section{问题定义：剖析究竟回答什么}
程序剖析研究资源消耗在代码、时间和请求类别之间如何分布。调试研究某个状态为什么出现；追踪研究事件发生的先后关系；指标描述聚合趋势。四者共同服务于因果解释，但任何单一工具都不能直接证明“这行代码导致了线上慢请求”。Go 官方诊断指南按 profiling、tracing、debugging 与 runtime statistics 组织能力，本书沿用这一分类。\cite{go-diagnostics}

先把问题写成实验陈述：“在固定版本、固定输入分布和相同资源约束下，吞吐从 A 下降到 B；假设是分配率上升导致 GC CPU 增长。”陈述至少包含对象、负载、观察区间、比较基线、成功指标和可推翻条件。只有“CPU 很高”并不充分：批处理充分利用 CPU 可能恰恰是好事。

\begingroup\small\setlength{\tabcolsep}{4pt}
\begin{longtable}{>{\raggedright\arraybackslash}p{0.298\linewidth}>{\raggedright\arraybackslash}p{0.298\linewidth}>{\raggedright\arraybackslash}p{0.298\linewidth}}
\toprule
\textbf{量} & \textbf{定义与单位} & \textbf{常见误读} \\
\midrule
\endfirsthead
\toprule
\textbf{量} & \textbf{定义与单位} & \textbf{常见误读} \\
\midrule
\endhead
延迟 & 一次操作从进入到完成的时间，ms & 把平均值等同于 p99 \\
\addlinespace[3pt]
成功吞吐 & 成功完成操作数/观测时间，req/s & 以发出数替代成功完成数 \\
\addlinespace[3pt]
CPU 用量 & 窗口内所有线程累计的 CPU 秒/墙钟秒，core & 以单核百分比解释整机百分比 \\
\addlinespace[3pt]
分配率 & 窗口内累计分配字节的增量/窗口时长，B/s & 等同于当前活跃堆 \\
\addlinespace[3pt]
资源效率 & CPU 秒/成功请求、分配字节/成功请求 & 比较不同业务工作量 \\
\addlinespace[3pt]
错误与超时 & 失败数/尝试数 & 删除失败请求后只看快请求 \\
\addlinespace[3pt]
\bottomrule
\end{longtable}
\endgroup
可复查的性能结论必须说明统计范围、单位和分母。将吞吐提高 10\% 写成“性能提升 10\%”会掩盖错误率、成本或延迟的变化。

\section{时间分解与排队模型}
\begin{figure}[H]
\centering
\begin{tikzpicture}[x=1cm,y=1cm]
\node[anchor=east,font=\small] at (1,0.0) {泳道1};
\draw[linegray] (1.2,-0.35) -- (14.799999999999999,-0.35);
\path[fill=leaf!60] (1.200,-0.250) rectangle (3.920,0.250);
\node[align=center,font=\small,text width=2.62cm] at (2.560,0.750) {计算};
\path[fill=bluegray!60] (3.920,-0.250) rectangle (12.080,0.250);
\node[align=center,font=\small,text width=8.06cm] at (8.000,0.750) {等待下游};
\path[fill=leaf!60] (12.080,-0.250) rectangle (14.800,0.250);
\node[align=center,font=\small,text width=2.62cm] at (13.440,0.750) {编码};
\draw[-{Stealth[length=2mm]},forest] (1.2,-1.15) -- (14.999999999999998,-1.15);
\draw[forest] (1.20,-1.15) -- (1.20,-1.25);
\node[below,font=\scriptsize] at (1.20,-1.25) {0};
\draw[forest] (3.92,-1.15) -- (3.92,-1.25);
\node[below,font=\scriptsize] at (3.92,-1.25) {20};
\draw[forest] (6.64,-1.15) -- (6.64,-1.25);
\node[below,font=\scriptsize] at (6.64,-1.25) {40};
\draw[forest] (9.36,-1.15) -- (9.36,-1.25);
\node[below,font=\scriptsize] at (9.36,-1.25) {60};
\draw[forest] (12.08,-1.15) -- (12.08,-1.25);
\node[below,font=\scriptsize] at (12.08,-1.25) {80};
\draw[forest] (14.80,-1.15) -- (14.80,-1.25);
\node[below,font=\scriptsize] at (14.80,-1.25) {100};
\node[anchor=east,font=\scriptsize] at (14.799999999999999,-1.9) {时间 / ms};
\end{tikzpicture}
\caption{CPU 时间和请求时间不是一把尺}\label{fig:01-time}
\begin{minipage}{0.96\linewidth}\footnotesize 教学示意：用于解释机制，不是运行时的实测数据或完整实现。

\textbf{读图顺序：}1. 开始计算\quad 2. 挂起等待\quad 3. 恢复执行\quad 4. 按同一边界核算。
\textbf{机制依据：}\cite{go-diagnostics}。
\end{minipage}
\end{figure}
墙钟时间包含执行、等待 CPU、等待锁、I/O、定时器以及业务主动休眠。对于串行且互斥的阶段，可用 \texttt{Twall = Tcpu + Trunnable + Tsync + Tio + Tother} 建立账本；并行请求不能把所有 goroutine 的 CPU 和等待直接相加当作该请求延迟，因为这些区间可能重叠。

稳定系统中 Little 定律为 \texttt{L = \ensuremath{\lambda}W}：平均在途数等于到达率乘平均逗留时间。例如，在没有失败或丢弃、到达与完成速率平衡的教学场景中，完成速率为 1000 req/s，平均逗留时间为 0.04 s，则平均在途请求数约为 40。到达率、在途数和逗留时间必须对应同一系统边界与同一类请求；若存在失败请求，不能把总在途数与成功完成速率混用。队列持续增长时，也不能直接套用上述稳态估计。

当利用率逼近容量，排队延迟可能急剧增大。因此 CPU 从 70\% 增至 90\% 并不意味着 p99 会按相同比例增加。实验应同时展示吞吐—延迟曲线，确定目标负载位于容量曲线的哪个区段。把并发数设为固定值的闭环压测会在服务变慢时自动降低发压，可能掩盖过载；开环压测则须记录调度延迟、丢失请求和发生器是否饱和。

\section{优化收益的上界与归一化}
Amdahl 定律：若可优化部分占原总运行时间的比例为 f，该部分加速 s 倍，则整体加速比为 \texttt{S = 1 /\allowbreak{} ((1\ensuremath{-}f)+f/\allowbreak{}s)}。如果编码阶段占 40\%，局部加速 2 倍，模型给出的整体加速比为 1.25；“热点函数的处理速度提高一倍”并不等于业务吞吐翻倍。这个模型假定其他成本不变；缓存、竞争和并发的非线性相互作用会使真实结果偏离。

对比 profile 时先问工作是否相等。若两个版本的观测时长相同，新版本成功吞吐翻倍、累计 CPU 时间增加 20\%，则 CPU/成功请求下降了 40\%：\texttt{1.\allowbreak{}2/\allowbreak{}2=0.\allowbreak{}6}。应报告原始消耗和归一化消耗，并说明是否按请求数、输入字节或记录数归一化。只看占比还存在分母陷阱：某函数占比下降，可能只是其他函数变贵。

性能预算比全局排名更可执行。将 100 ms 延迟预算分给排队、计算、下游与余量，或将 0.5 CPU-ms/请求拆到编码、业务和框架；预算是规划工具，实测仍需核对阶段重叠与未知项。

\section{证据阶梯：从相关性走向干预}
\begin{figure}[H]
\centering
\begin{tikzpicture}[x=1cm,y=1cm]
\node[draw=linegray!75!black,fill=linegray!13,rounded corners=3pt,text width=2.58cm,minimum height=1.12cm,align=center,font=\small] (p0-symptom) at (1.50,0.00) {延迟上升};
\node[draw=linegray!75!black,fill=linegray!13,rounded corners=3pt,text width=2.58cm,minimum height=1.12cm,align=center,font=\small] (p0-cpu) at (4.50,0.97) {CPU 假设};
\node[draw=leaf!75!black,fill=leaf!13,rounded corners=3pt,text width=2.58cm,minimum height=1.12cm,align=center,font=\small] (p0-wait) at (4.50,-0.97) {等待假设};
\node[draw=linegray!75!black,fill=linegray!13,rounded corners=3pt,text width=2.58cm,minimum height=1.12cm,align=center,font=\small] (p0-probe) at (7.50,0.00) {采样＋时序};
\node[draw=linegray!75!black,fill=linegray!13,rounded corners=3pt,text width=2.58cm,minimum height=1.12cm,align=center,font=\small] (p0-act) at (10.50,0.00) {单变量干预};
\node[draw=leaf!75!black,fill=leaf!13,rounded corners=3pt,text width=2.58cm,minimum height=1.12cm,align=center,font=\small] (p0-check) at (13.50,0.00) {复查与反证};
\draw[-{Stealth[length=2mm]},forest!70!black] (p0-symptom.east) -- (p0-cpu.west);
\draw[-{Stealth[length=2mm]},forest!70!black] (p0-symptom.east) -- (p0-wait.west);
\draw[-{Stealth[length=2mm]},forest!70!black] (p0-cpu.east) -- (p0-probe.west);
\draw[-{Stealth[length=2mm]},forest!70!black] (p0-wait.east) -- (p0-probe.west);
\draw[-{Stealth[length=2mm]},forest!70!black] (p0-probe.east) -- (p0-act.west);
\draw[-{Stealth[length=2mm]},forest!70!black] (p0-act.east) -- (p0-check.west);
\end{tikzpicture}
\caption{证据怎样一步步排除错误解释}\label{fig:01-evidence}
\begin{minipage}{0.96\linewidth}\footnotesize 教学示意：用于解释机制，不是运行时的实测数据或完整实现。

\textbf{读图顺序：}1. 先写症状\quad 2. 保留竞争假设\quad 3. 选择能区分的观测\quad 4. 施加可回退干预\quad 5. 检查可反驳预测。
\textbf{机制依据：}\cite{go-diagnostics}。
\end{minipage}
\end{figure}
建议使用六步闭环：症状量化 \ensuremath{\rightarrow}\allowbreak{} 机制假设 \ensuremath{\rightarrow}\allowbreak{} 工具选择 \ensuremath{\rightarrow}\allowbreak{} 定位与解释 \ensuremath{\rightarrow}\allowbreak{} 单变量干预 \ensuremath{\rightarrow}\allowbreak{} 回归与部署观察。每个阶段都保存事实、推断、未知项，避免把函数名当作根因。

例如：CPU profile 中 \texttt{runtime.\allowbreak{}mallocgc} 的排名升高是一项观测；“业务对象分配增加”是候选解释；“改用池即可解决”尚是未经验证的干预。应进一步用累计分配 profile 定位分配栈，用基准测试确认 B/op，检查存活对象和逃逸原因，再衡量池化带来的持有内存、同步与正确性成本。

强证据包括：相同工作量下重复观测到变化；移除候选机制后症状按预测缓解；把改动撤回后症状复现；解释能预测另一个输入规模的结果。负结果同样有价值：热点消失而延迟不变，可能说明真正限制来自下游或排队。它不应被作为“测试失败”删除。

\section{观察者效应与伦理边界}
采样、栈展开、事件记录、符号解析都会占用资源。不要把启用与关闭剖析器时的数据直接拼接成一组可比测量。可设置 A=原版本且不开启采集、B=原版本且开启采集、C=优化版本且不开启采集：A 与 B 用于估计采集扰动，A 与 C 用于估计优化收益。如果还需要比较优化前后的 profile，应另采集优化版本的样本，并保持采集配置一致。高频采集应先在副本上测量开销。

profile、trace 和 core 可能包含函数名、路径、命令行、请求标签或内存中的敏感值。生产采集需要鉴权、最小权限、低基数标签、短留存和审计。开源工具能采到数据，不代表可以将生产样本上传公共页面。调试器暂停进程、core 生成和高密度事件采集都可能影响服务可用性。

本书所有示例服务默认仅监听回环地址。若需要远程操作，应在已有授权通道内做端口转发，并限制访问范围。不要通过把监听地址改为 \texttt{0.\allowbreak{}0.\allowbreak{}0.\allowbreak{}0} 来省略诊断面访问控制。

\section{全书路线与结业标准}
第 02–04 章建立运行时与统计基础；第 05–10 章形成工具使用能力；第 11–14 章完成优化、生产落地与研究型实验；第 15 章从成熟开源实现理解机制，第 16 章追踪近五年研究。每章均包含目标、实验、诊断选择题和有参考答案的分析题。

建议用 16 周修读：每周一章，理论 2 小时、实验 3 小时、阅读 1 小时；第 13–14 周完成小组项目，最后两周进行源码报告和研究复现。若已有 Go 生产经验，可先做第 13 章综合案例，再按证据缺口回读。

结业能力不是记住命令，而是提交可以被另一位工程师复查的证据包：代码与构建标识、负载描述、环境约束、原始 profile/trace、统计方法、关键观察、替代解释、优化补丁、正确性检查和适用范围。

\section{实验任务}
\subsection*{建立一份性能调查卡}
\begin{enumerate}
\item 选择一个函数或服务，记录版本、输入分布、CPU/内存限制及目标 SLO。

\item 分别写出一个 CPU、一个等待、一个内存假设；为每个假设写明能推翻它的观测。

\item 定义相同工作量：例如处理同一批 100000 条记录，而不是简单运行相同时长。

\item 填写计划采集的指标、profile 与 trace；明确采集开销和停止条件。

\end{enumerate}
\section{自测与解析}
\noindent\textbf{1. CPU 消耗相同、p99 延迟显著增大时，首选判断是什么？}
\begin{enumerate}[label=\Alph*.]
\item CPU profile 已足以定位
\item 应检查等待、排队与负载分布
\item 增加 CPU 采样率必能定位
\item 这一定是 GC
\end{enumerate}
\noindent{\color{forest}\textbf{答案 B。}}CPU profile 主要反映线程实际占用 CPU 时的执行。排队、等待和长尾延迟还需要结合 trace、指标及负载分布解释。

\noindent\textbf{2. 占总时间 30\% 的阶段加速 3 倍，整体理想加速约为多少？}
\begin{enumerate}[label=\Alph*.]
\item 3 倍
\item 1.25 倍
\item 1.8 倍
\item 无法计算任何上界
\end{enumerate}
\noindent{\color{forest}\textbf{答案 B。}}在可优化部分占比和工作量固定、其他成本不变的模型中，整体加速比为 \(\frac{1}{0.7+0.3/3}=1.25\)。

\noindent\textbf{3. 哪个结论最可复查？}
\begin{enumerate}[label=\Alph*.]
\item 新版本更快
\item 火焰图变窄
\item 固定输入下 CPU-ms/请求下降，重复实验并保留失败率
\item 排名第一的函数被删掉
\end{enumerate}
\noindent{\color{forest}\textbf{答案 C。}}必须给出指标、分母、控制条件和重复观测。

\section{分析题与参考答案}
\noindent\textbf{题 1：}在到达与完成速率平衡的服务中，同一类请求的完成速率为 2000 req/s，平均在途数为 100。估计平均逗留时间，并指出一个使该估计不适用的条件。

\noindent\textbf{参考答案：}在同一稳态系统边界内，W=L/\ensuremath{\lambda}=100/2000=0.05 s，即 50 ms。若队列在观测窗口内持续增长，或在途数包含另一类请求而完成速率不包含，则不能直接使用这一估计。

\noindent\textbf{题 2：}在相同工作量和可比采集配置下，某函数的 CPU 占比从 20\% 降到 10\%，但总 CPU 时间增至原来的 2 倍。这个函数的绝对 CPU 成本下降了吗？

\noindent\textbf{参考答案：}没有下降。若旧版本总 CPU 时间为 100 秒，该函数占 20 秒；新版本总 CPU 时间为 200 秒，该函数仍占 20 秒。若工作量也发生变化，则还需按完成工作量归一化，不能仅凭占比判断效率。

\medskip\noindent\textbf{本章主要阅读：}\cite{go-diagnostics}。
\chapter{Go 运行时与系统模型}
\label{ch:02}
\noindent{\large\color{forest}把火焰图和时间线中的运行时符号，与具体的执行机制对应起来。}\par\medskip
\noindent\textbf{先修要求：}第 01 章；Go 并发和基本操作系统知识。

\noindent\textbf{完成本章后，应能够：}
\begin{itemize}
\item 解释 G、M、P 与 OS 线程之间的关系
\item 区分 runnable、running、waiting 对 profile 的影响
\item 理解栈、堆、逃逸与 GC 的成本来源
\item 把语言内存模型与性能诊断分开使用
\end{itemize}
\section{G、M、P：语言并发不是操作系统线程数}
\begin{figure}[H]
\centering
\begin{tikzpicture}[x=1cm,y=1cm]
\node[draw=bluegray!75!black,fill=bluegray!13,rounded corners=3pt,text width=7.08cm,minimum height=1.12cm,align=center,font=\small] (p0-g1) at (3.75,1.95) {G1\\{\scriptsize 等待}};
\node[draw=leaf!75!black,fill=leaf!13,rounded corners=3pt,text width=7.08cm,minimum height=1.12cm,align=center,font=\small] (p0-g2) at (3.75,0.00) {G2\\{\scriptsize 运行}};
\node[draw=leaf!75!black,fill=leaf!13,rounded corners=3pt,text width=7.08cm,minimum height=1.12cm,align=center,font=\small] (p0-g3) at (3.75,-1.95) {G3\\{\scriptsize 运行}};
\node[draw=amber!75!black,fill=amber!13,rounded corners=3pt,text width=7.08cm,minimum height=1.12cm,align=center,font=\small] (p0-m1) at (11.25,-1.95) {M1 /\allowbreak{} P1\\{\scriptsize 持有}};
\node[draw=amber!75!black,fill=amber!13,rounded corners=3pt,text width=7.08cm,minimum height=1.12cm,align=center,font=\small] (p0-m2) at (11.25,0.00) {M2 /\allowbreak{} P2\\{\scriptsize 持有}};
\node[draw=bluegray!75!black,fill=bluegray!13,rounded corners=3pt,text width=7.08cm,minimum height=1.12cm,align=center,font=\small] (p0-net) at (11.25,1.95) {网络等待\\{\scriptsize 等待}};
\draw[-{Stealth[length=2mm]},forest!70!black] (p0-g2.east) -- (p0-m2.west);
\draw[-{Stealth[length=2mm]},forest!70!black] (p0-g3.east) -- (p0-m1.west);
\draw[-{Stealth[length=2mm]},forest!70!black] (p0-g1.east) -- (p0-net.west);
\end{tikzpicture}
\caption{G 可以很多，执行 Go 代码仍需要 P}\label{fig:02-gmp}
\begin{minipage}{0.96\linewidth}\footnotesize 教学时刻：G1 等待网络时让出 P，G3 在 M1/P1 上运行；仅绘制本步真实关系。唤醒后 G1 先变为 runnable，并不立即获得 P。M/P 配对仅为简化示意。

\textbf{读图顺序：}1. 有工作可运行\quad 2. 两个 G 开始运行\quad 3. G1 等待 I/O\quad 4. 就绪不等于已运行。
\textbf{机制依据：}\cite{go-scheduler}。
\end{minipage}
\end{figure}
G 表示 goroutine，M 表示执行 Go 或系统代码的操作系统线程，P 持有执行 Go 代码所需的调度资源。\texttt{GOMAXPROCS} 决定可供并行执行 Go 代码的 P 数量，并不限制程序总线程数、goroutine 数量或所有 cgo 活动。调度器在本地运行队列、全局队列、网络轮询器和工作窃取之间协调工作，具体策略属于版本相关实现细节。\cite{go-scheduler}

在 Go 调度器中，running G 已被安排到 M 上执行；runnable G 已具备运行条件，正在等待调度；waiting G 则可能在 channel、锁、网络或计时器上等待。Go 的 running 状态不保证对应的 OS 线程此刻正在 CPU 上运行，线程仍可能被操作系统抢占。因此，CPU profile 主要反映实际占用 CPU 的执行，而执行追踪可以帮助分析 Go 层面的可运行延迟和同步唤醒。CPU 空闲而 goroutine 很多可能意味着大家都在等 I/O，不能仅据数量判断调度器瓶颈。

本书源码分析固定 Go 1.25.0；实验核心程序兼容 Go 1.22，但 FlightRecorder、testing.B.Loop 等能力有单独版本门槛。Go 1.25 引入的容器感知默认 GOMAXPROCS 会考虑 Linux cgroup CPU 带宽和亲和性等约束，但当主模块声明的 Go 语言版本为 1.24 或更早时，\texttt{containerm\allowbreak{}axprocs} 与 \texttt{updatemaxp\allowbreak{}rocs} 的兼容默认值仍为 0。本书 lab 的 go.mod 声明 go 1.22，因此只换成 Go 1.25 工具链并不会自动启用这两项新默认；对照实验须明确记录模块语言版本与 GODEBUG。不要把宿主机核数直接当成容器可以持续使用的 CPU 容量。

Go 1.25 的默认计算通常取逻辑核数、CPU 亲和性数量和配额容量中的较小者；非整数配额向上取整，并且除非逻辑核或亲和性不足 2，否则不会因小配额自动降到少于 2 个 P。通过环境变量显式设置正数 \texttt{GOMAXPROCS}，或调用 \texttt{runtime.\allowbreak{}GOMAXPROCS(n)}（n > 0），会关闭自动更新；\texttt{runtime.\allowbreak{}GOMAXPROCS(0)} 只是查询。\texttt{runtime.\allowbreak{}SetDefault\allowbreak{}GOMAXPROCS()} 可恢复默认策略，但仍尊重 GODEBUG 开关。\cite{tools-go125}

\section{阻塞、系统调用与网络轮询}
当 goroutine 阻塞于可调度的同步操作时，运行时可以让 M 执行其他 G。进入可能长时间阻塞的系统调用时，P 也可以被重新交给其他线程。网络轮询器让大量网络连接在就绪之前不必各占一个阻塞的 OS 线程。这解释了“百万 goroutine”与“百万 OS 线程”成本的差别，也解释了为何 CPU profile 通常不记录网络等待本身。

并非所有等待都由 Go 运行时完整观察。cgo、内核调度、缺页异常、磁盘阻塞、容器 CPU 节流可能要求 perf、strace、eBPF 或 cgroup 指标。每跨一层就需要重新明确时间边界：系统调用耗时不等同于设备服务时间，off-CPU 时间也不全是锁等待。

第 08 章学习把唤醒事件与实际获得 CPU 连接起来；第 10 章补齐内核侧证据。跨层关联最好保留 PID、构建标识、时间区间与相同实验标签，不能仅根据相似函数名拼接结论。

\section{栈、堆、逃逸与对象生命期}
\begin{figure}[H]
\centering
\begin{tikzpicture}[x=1cm,y=1cm]
\node[align=center,text width=7cm,font=\small\bfseries] at (3.55,2.05) {根保持可达};
\node[draw=leaf!75!black,fill=leaf!13,rounded corners=3pt,text width=1.95cm,minimum height=1.12cm,align=center,font=\scriptsize] (p0-root) at (1.18,-0.70) {全局根\\{\scriptsize 可达}};
\node[draw=leaf!75!black,fill=leaf!13,rounded corners=3pt,text width=1.95cm,minimum height=1.12cm,align=center,font=\scriptsize] (p0-slice) at (3.55,0.12) {切片视图（100 B）\\{\scriptsize 可达}};
\node[draw=leaf!75!black,fill=leaf!13,rounded corners=3pt,text width=1.95cm,minimum height=1.12cm,align=center,font=\scriptsize] (p0-big) at (5.92,-0.70) {10 MiB 数组\\{\scriptsize 可达}};
\node[dashed,draw=linegray!75!black,fill=linegray!13,rounded corners=3pt,text width=1.95cm,minimum height=1.12cm,align=center,font=\scriptsize] (p0-copy) at (3.55,-1.52) {独立 100 B\\{\scriptsize 不纳入}};
\draw[-{Stealth[length=2mm]},forest!70!black] (p0-root.east) -- (p0-slice.west);
\draw[-{Stealth[length=2mm]},forest!70!black] (p0-slice.east) -- (p0-big.west);
\node[align=center,text width=7cm,font=\small\bfseries] at (11.25,2.05) {复制并替换引用};
\node[draw=leaf!75!black,fill=leaf!13,rounded corners=3pt,text width=1.95cm,minimum height=1.12cm,align=center,font=\scriptsize] (p77-root) at (8.88,-0.70) {全局根\\{\scriptsize 可达}};
\node[dashed,draw=amber!75!black,fill=amber!13,rounded corners=3pt,text width=1.95cm,minimum height=1.12cm,align=center,font=\scriptsize] (p77-slice) at (11.25,-1.52) {切片视图（100 B）\\{\scriptsize 待GC}};
\node[dashed,draw=amber!75!black,fill=amber!13,rounded corners=3pt,text width=1.95cm,minimum height=1.12cm,align=center,font=\scriptsize] (p77-big) at (13.62,-0.70) {10 MiB 数组\\{\scriptsize 待GC}};
\node[draw=leaf!75!black,fill=leaf!13,rounded corners=3pt,text width=1.95cm,minimum height=1.12cm,align=center,font=\scriptsize] (p77-copy) at (11.25,0.12) {独立 100 B\\{\scriptsize 可达}};
\draw[-{Stealth[length=2mm]},forest!70!black] (p77-root.east) -- (p77-copy.west);
\end{tikzpicture}
\caption{小切片为何能留住大数组}\label{fig:02-reachability}
\begin{minipage}{0.96\linewidth}\footnotesize 100 B 表示切片可访问的数据长度，不是切片头的大小；切片节点仅示意引用关系，不断言它是独立的堆分配。10 MiB 为底层数组大小。“失去引用”“GC 回收”“向 OS 归还页”是不同阶段，图中的节点数不等于真实 heap 对象计数。

\textbf{读图顺序：}1. 先分配大数组\quad 2. 根保持可达\quad 3. 复制并替换引用\quad 4. GC 与 RSS 仍不同步。
\textbf{连线语义：}全局根\ensuremath{\rightarrow}\allowbreak{}切片视图（100 B）：引用；切片视图（100 B）\ensuremath{\rightarrow}\allowbreak{}10 MiB 数组：底层数组；全局根\ensuremath{\rightarrow}\allowbreak{}独立 100 B：复制后。
\textbf{机制依据：}\cite{go-memory}。
\end{minipage}
\end{figure}
goroutine 栈可以增长和收缩。局部变量是否分配在栈上由编译器逃逸分析与实现限制决定，语法上出现 \texttt{new} 不保证上堆，返回指针也不能替代对编译器结果的检查。接口装箱、闭包捕获、对象尺寸、对象是否需要在调用返回后继续存活等因素，都可能影响分配。

\begin{codeblock}
go build -gcflags='-m=2' ./...
go test -run='^$' -bench=BenchmarkEncode -benchmem ./...
\end{codeblock}

逃逸诊断用于解释某一构建的决定；\texttt{B/\allowbreak{}op} 和 \texttt{allocs/\allowbreak{}op} 用于实测动态代价。诊断输出中出现 escape 不意味着它一定是热点；某处启动时分配一次，可能无须优化。反之，极小对象每请求分配数百次，累积的分配和 GC 成本可能显著。

对象的逻辑生命周期、GC 可达生命周期和操作系统驻留生命周期不相同。一个小切片保留大数组、全局 map 未删除键、goroutine 闭包捕获对象，都可能使可达性超出业务所需。heap profile 记录对象的分配位置；要判断“谁还持有它”，通常还需分析代码中的引用关系或进行受控实验，不能把分配栈直接当作持有者的引用链。

\section{GC：吞吐、延迟与内存预算}
Go 使用并发、追踪式垃圾回收。GC 仍存在短暂停顿、后台标记、写屏障、扫描与分配辅助成本。大部分 GC 时间并不表现为全局 STW；“停顿很短”不能推出“GC 开销很低”。应同时观察分配率、存活堆、扫描量、GC CPU 与请求延迟。\cite{go-memory}

简化模型为 \texttt{目标堆大小 \ensuremath{\approx} 存活堆大小 + (存活堆大小 + GC 根扫描量) \ensuremath{\times} GOGC/\allowbreak{}100}。其中各项均以字节计，GC 根扫描量包括需要扫描的 goroutine 栈和全局变量等。该式用于理解现代 Go 的堆目标，忽略了最小堆目标和内存限制等约束，也不是 RSS 的上界。提高 GOGC 通常以更多堆空间换取较少 GC 周期；降低 GOGC 可能减少堆增长但增加 GC CPU。设置 \texttt{GOMEMLIMIT} 约束运行时管理内存的软预算，需要为进程其他内存、cgo、映射与容器余量留空间。压力过高时运行时可为避免 GC 抖动而超出软限制。

因此调参不能替代根因修复。若每请求制造大量临时字符串，减少分配可能同时改善 CPU 和内存；若有效工作集本就超过预算，只调低 GOGC 往往加重回收压力。详细采集与验证见第 06 章。

\section{缓存、NUMA 与共享状态}
在多核系统中，算法复杂度相同不意味着运行时间相同。连续数组和指针链在缓存局部性上差异很大；由不同核频繁写入、但位于同一缓存行的相邻变量可能引起伪共享，使缓存行反复失效；锁的争用成本可能来自排队、缓存行在多个核之间迁移和内核调度的共同作用。

NUMA 系统中，访问本地与远端内存也可能具有不同成本，需要结合线程调度、内存放置与机器拓扑分析。Go 的 CPU profile 通常不能直接证明某个热点由“缓存未命中”或远端内存访问造成。这需要硬件计数器、固定频率/负载对照、数据结构变体和平台支持。硬件事件具有型号、权限与采样偏差限制，不宜把一台机器上的 cycles/instruction 阈值推广为通用标准。

诊断共享状态时，先验证程序正确性，再考虑分片、批处理、只读快照或缩短临界区。增加 goroutine 数可能提高并发，也可能扩大共享资源竞争；并发不是必然加速。

\section{正确性优先：内存模型与动态检测}
Go 内存模型以单个 goroutine 内的 sequenced-before 顺序和跨 goroutine 的 synchronized-before 关系为基础，取它们的传递闭包，定义 happens-before 关系。互斥锁、channel 和原子操作等提供具体的同步保证。\cite{go-model} 把 map 改成无锁共享访问以消除 mutex profile 热点，可能只是引入数据竞争。性能优化必须保持 API 契约、错误处理、取消和资源释放行为。

\begin{codeblock}
go test -race ./...
go vet ./...
\end{codeblock}

race detector 只检查真实执行过的路径，未报告竞争不等于证明程序无竞争。它会明显改变时间、内存与调度，因此开启 race 的运行用于正确性检查，常规性能结论应使用独立的正式构建。\cite{go-race}

死锁、活锁、逻辑竞态和 goroutine 泄漏也不都属于 data race。诊断应分别使用栈转储、执行追踪、取消测试和生命周期审计。每个优化结果都应回答：快的原因是否只是少做了必要的工作？

\section{实验任务}
\subsection*{构建运行时状态地图}
\begin{enumerate}
\item 在 lab 目录启动示例服务，分别访问 /cpu、/wait、/lock 和 /retain。

\item 分别采集 CPU profile、goroutine profile 和 trace，指出运行、可运行和等待状态的差异。

\item 用 GOMAXPROCS=1 与 GOMAXPROCS=2 运行相同 CPU 工作量；记录吞吐和调度等待，不预设一定线性加速。

\item 开启 -race 单独运行正确性检查，并把该结果与性能测量分开记录。

\end{enumerate}
\section{自测与解析}
\noindent\textbf{1. GOMAXPROCS=2 直接限制什么？}
\begin{enumerate}[label=\Alph*.]
\item 总 goroutine 数
\item 进程 OS 线程总数
\item 可供并行执行 Go 代码的 P 数量
\item 全部网络连接数
\end{enumerate}
\noindent{\color{forest}\textbf{答案 C。}}OS 线程数和 goroutine 数都可能超过 2；cgo 调用及线程在 Go 代码之外的活动也不能只按 P 数量推断。

\noindent\textbf{2. heap profile 中归因权重最高的分配栈，一定能指出谁仍在持有这些对象吗？}
\begin{enumerate}[label=\Alph*.]
\item 是
\item 否，分配位置与引用持有者不同
\item 只要开启 gc=1 就是
\item 只对切片成立
\end{enumerate}
\noindent{\color{forest}\textbf{答案 B。}}profile 记录的是分配栈，而不是完整的对象引用图；对象的分配者与持有者可能不同。

\noindent\textbf{3. race 检查未报错意味着什么？}
\begin{enumerate}[label=\Alph*.]
\item 证明没有并发问题
\item 执行过的路径中未检测到该类数据竞争
\item 没有死锁
\item 可直接用该次运行做性能结论
\end{enumerate}
\noindent{\color{forest}\textbf{答案 B。}}动态检测只覆盖实际执行的路径，而且 race detector 会改变执行时间、内存占用和调度。

\section{分析题与参考答案}
\noindent\textbf{题 1：}服务的 goroutine 数量很高，但 CPU 用量较低。请给出三个可以分别检验的假设。

\noindent\textbf{参考答案：}候选假设包括：goroutine 在等待下游 I/O；在等待共享锁或 channel；任务未及时取消，导致挂起的 goroutine 积累。可分别结合 trace 中的网络等待、block/mutex profile，以及 goroutine 栈和生命周期实验加以区分；这些现象可能同时存在，不能直接据此建议调高 GOMAXPROCS。

\noindent\textbf{题 2：}一个 100 字节切片来自 10 MiB 数组并长期保留，会发生什么？

\noindent\textbf{参考答案：}若切片仍被后续代码使用，它可能使整个底层数组保持可达。因此，业务只需要 100 字节，不等于只保留了 100 字节内存。将所需区间复制到独立的小数组可以解除对大数组的引用，但仍须结合 heap 趋势、复制成本和正确性测试评估是否值得。

\medskip\noindent\textbf{本章主要阅读：}\cite{go-scheduler}\cite{go-memory}\cite{go-model}\cite{go-race}\cite{tools-go125}。
\chapter{基准测试与实验设计}
\label{ch:03}
\noindent{\large\color{forest}从单次测量的快慢判断，走向可复现的统计比较。}\par\medskip
\noindent\textbf{先修要求：}第 01–02 章；基本概率统计。

\noindent\textbf{完成本章后，应能够：}
\begin{itemize}
\item 明确基准测试的计时边界，避免准备工作混入目标测量或结果被优化消除
\item 选择正确的重复单元和负载生成方式
\item 解释采样误差、置信区间与实际显著性
\item 建立版本、环境、工作量一致的比较协议
\end{itemize}
\section{可运行基准与测量边界}
Go benchmark 使用 \texttt{go test -\allowbreak{}bench} 驱动。下面给出兼容 Go 1.22 的基准函数模式；将计算结果保存在包级变量中，以降低计算被编译器消除的风险。该代码片段需放在导入 testing 且定义了 encodeSlow 的测试文件中。初始化是否在计时区内应由研究问题决定，不能默认一律排除。\cite{go-testing}

\begin{codeblock}
var result string
func BenchmarkEncodeSlow(b *testing.B) {
    input := []byte("profiling-with-go")
    b.ReportAllocs()
    b.SetBytes(int64(len(input)))
    b.ResetTimer()
    for i := 0; i < b.N; i++ {
        result = encodeSlow(input)
    }
}
\end{codeblock}

若研究完整请求，序列化、内存分配与错误检查应被包含；若研究一个纯算法，可在计时前构造输入。\texttt{ResetTimer} 清零累计时间、分配统计和用户报告的指标，但不改变计时器的启停状态；\texttt{StopTimer}/\texttt{StartTimer} 应克制使用，频繁切换本身会造成扰动。Go 1.24 引入 \texttt{B.\allowbreak{}Loop}，能简化常见计时模式并通过保活循环体内函数调用的参数与结果，降低整个调用被优化消除的风险；这不意味着被调用函数内部的优化被全部关闭。不要把该 API 写成所有 Go 版本通用，也不要混用 \texttt{b.\allowbreak{}N} 和 \texttt{b.\allowbreak{}Loop()} 控制同一个基准。

基准值 \texttt{ns/\allowbreak{}op} 的 op 是作者定义的一次迭代。若一次迭代处理 1000 条记录，必须说明它不是每条记录时间。\texttt{RunParallel} 的 ns/op 是整个并发 benchmark 的墙钟时间除以全部迭代数，即吞吐的倒数尺度；它不是各 goroutine 逐次延迟的算术平均，更不是单请求 p99。假设两个 worker 各执行一次 10 ms 操作且完全重叠，报告约 5 ms/op，而每次操作延迟仍约 10 ms。

\section{命令、构建与原始数据}
本书附带的 lab 是完整 Go 模块。先用功能测试确认慢、快编码器输出相同，再开展性能测量。

\begin{codeblock}
go test ./...
go test -run='^$' -bench=BenchmarkEncode -benchmem -count=10 > bench.txt
go test -run='^$' -bench=BenchmarkEncodeSlow -benchtime=3s   -cpuprofile=cpu.pb.gz -memprofile=mem.pb.gz
go tool pprof -top cpu.pb.gz
\end{codeblock}

profile 文件名后缀只表示约定，不改变实际编码。\texttt{go test} 在生成 profile 时会保留测试二进制；应保留匹配的二进制，尤其涉及符号化、反汇编和不同构建。记录 \texttt{go version}、\texttt{go env GOOS GOARCH}、构建参数、commit、CPU 型号、容器配额、GOMAXPROCS、GOGC 和输入校验值。

同时采集 CPU 与内存可能有额外开销。用于定位的运行和最终报告数值的运行应分开，除非研究问题就是“启用某采集配置的系统性能”。尽量保存原始输出，而不是只复制平均值到表格。

\texttt{-\allowbreak{}cpuprofile} 的采集生命周期也不等于 benchmark 的计时器：testing 在运行测试或基准之前启动 CPU 采集，并在结束阶段停止；\texttt{b.\allowbreak{}StopTimer()}、\texttt{ResetTimer()} 并不会暂停 CPU profiler。因此未计入 ns/op 的准备、校准或收尾工作仍可能出现在 CPU profile 中，不能仅按火焰图百分比拆分 ns/op。

\section{重复、随机化与环境噪声}
\begin{figure}[H]
\centering
\begin{tikzpicture}[x=1cm,y=1cm]
\node[anchor=east,font=\small] at (1,0.0) {泳道1};
\draw[linegray] (1.2,-0.35) -- (14.799999999999999,-0.35);
\path[fill=forest!60] (1.200,-0.250) rectangle (4.600,0.250);
\node[align=center,font=\small,text width=3.30cm] at (2.900,0.750) {A 基线};
\path[fill=forest!60] (4.600,-0.250) rectangle (8.000,0.250);
\node[align=center,font=\small,text width=3.30cm] at (6.300,0.750) {B 改动};
\path[fill=forest!60] (8.000,-0.250) rectangle (11.400,0.250);
\node[align=center,font=\small,text width=3.30cm] at (9.700,0.750) {B 改动};
\path[fill=forest!60] (11.400,-0.250) rectangle (14.800,0.250);
\node[align=center,font=\small,text width=3.30cm] at (13.100,0.750) {A 基线};
\draw[-{Stealth[length=2mm]},forest] (1.2,-1.15) -- (14.999999999999998,-1.15);
\draw[forest] (1.20,-1.15) -- (1.20,-1.25);
\node[below,font=\scriptsize] at (1.20,-1.25) {0};
\draw[forest] (3.92,-1.15) -- (3.92,-1.25);
\node[below,font=\scriptsize] at (3.92,-1.25) {0.8};
\draw[forest] (6.64,-1.15) -- (6.64,-1.25);
\node[below,font=\scriptsize] at (6.64,-1.25) {1.6};
\draw[forest] (9.36,-1.15) -- (9.36,-1.25);
\node[below,font=\scriptsize] at (9.36,-1.25) {2.4};
\draw[forest] (12.08,-1.15) -- (12.08,-1.25);
\node[below,font=\scriptsize] at (12.08,-1.25) {3.2};
\draw[forest] (14.80,-1.15) -- (14.80,-1.25);
\node[below,font=\scriptsize] at (14.80,-1.25) {4};
\node[anchor=east,font=\scriptsize] at (14.799999999999999,-1.9) {时间 / 实验顺序};
\end{tikzpicture}
\caption{用 ABBA 顺序识别环境漂移}\label{fig:03-experiment}
\begin{minipage}{0.96\linewidth}\footnotesize 教学示意：用于解释机制，不是运行时的实测数据或完整实现。

\textbf{读图顺序：}1. 第一次基线\quad 2. 切换到改动\quad 3. 重复改动\quad 4. 回到基线。
\textbf{机制依据：}\cite{go-benchstat}。
\end{minipage}
\end{figure}
一次进程内的循环迭代并非独立实验：它们共享缓存、堆状态和频率。统计推断需要明确重复单元，是独立进程、机器、输入批次还是时间窗。\texttt{-\allowbreak{}count=10} 会在同一个测试进程中重复基准，产生多份结果；这些结果仍共享进程、主机及环境趋势，不能宣称已代表所有生产部署。

以 ABBA 或随机次序运行旧/新版本，可减轻热机、频率、后台负载和温度漂移。保持 CPU/NUMA 配置、容器额度、输入分布和依赖版本一致。若服务依赖下游，记录下游耗时和错误；若只评估本地代码，可使用明确声明的固定响应下游，但结论限于该模型。

\texttt{benchstat} 是 Go 生态常用的基准比较工具，具体输出与统计方法随工具版本变化。安装时固定所用版本并记录帮助信息，不将 p 值当作改动收益的充分证据。\cite{go-benchstat}

\begin{codeblock}
benchstat before.txt after.txt
\end{codeblock}

应同时给出效应大小、波动区间、样本数和业务阈值。即使 0.5\% 的差异通过了统计检验，也可能不值得为此引入复杂缓存；即使均值改善了 10\%，若 p99 或错误率恶化，也可能不符合业务目标。

\section{采样误差：估计不是计数真相}
\begin{figure}[H]
\centering
\begin{tikzpicture}[x=1cm,y=1cm]
\path[fill=leaf!25] (0.00,0.00) rectangle (13.97,0.78);
\node[align=center,text width=13.92cm,font=\small] at (7.00,0.40) {总样本};
\path[fill=bluegray!38] (0.00,1.05) rectangle (1.38,1.83);
\node[font=\scriptsize] at (0.70,1.45) {2};
\path[fill=bluegray!38] (1.40,1.05) rectangle (5.57,1.83);
\node[align=center,text width=4.12cm,font=\small] at (3.50,1.45) {散列};
\path[fill=leaf!38] (5.60,1.05) rectangle (13.97,1.83);
\node[align=center,text width=8.32cm,font=\small] at (9.80,1.45) {其他};
\node[anchor=north west,align=left,text width=14cm,font=\scriptsize] at (0,-.45) {1=总样本: flat 0 /\allowbreak{} cum 1000；2=编码: flat 100 /\allowbreak{} cum 100；3=散列: flat 300 /\allowbreak{} cum 300；4=其他: flat 600 /\allowbreak{} cum 600};
\end{tikzpicture}
\caption{样本变多后，热点估计才逐渐稳定}\label{fig:03-sampling}
\begin{minipage}{0.96\linewidth}\footnotesize 构造的采样序列，仅用于讲解估计量；不是实际收敛曲线，也不保证样本独立。

\textbf{读图顺序：}1. 样本少\quad 2. 继续观察\quad 3. 更长窗口\quad 4. 区分误差来源。
\textbf{机制依据：}\cite{go-pprof}。
\end{minipage}
\end{figure}
先固定“命中”的定义，例如仅统计叶端属于该函数的样本。若将各样本是否命中近似为独立、同分布的 Bernoulli 事件，N 个样本中命中 k 次，则 \texttt{p\_\allowbreak{}hat=k/\allowbreak{}N}，标准误近似 \texttt{SE=sqrt(p\_\allowbreak{}hat(1\ensuremath{-}p\_\allowbreak{}hat)/\allowbreak{}N)}。若观测到 p\_hat=0.1、N=1000，则估计标准误约为 0.0095，粗略的双侧 95\% 正态近似区间为 10\%\ensuremath{\pm}1.86 个百分点。这里的数字是教学模型，不是 profiler 承诺的置信区间。

真实栈样本存在时间相关性，周期性工作会与采样周期混叠；函数归属也受内联、栈展开和采样丢失影响。极小比例、很少样本或存在相关性时，正态近似不可靠。零命中不证明零开销；在独立、同分布的 Bernoulli 试验模型下，零命中对应的单侧 95\% 置信上界为 \texttt{1\ensuremath{-}0.\allowbreak{}05\textasciicircum{}(1/\allowbreak{}N)}；N 较大时可近似为 \texttt{3/\allowbreak{}N}。这约束的是该模型中的命中概率，并不能排除采集器的系统性漏采。

增加时长通常比随意增加频率更稳妥，但会混合更多工作阶段。应先稳定负载，再重复多个可比较时间窗，报告哪些热点在各窗口中持续出现。若目标是偶发慢请求，更长时间的聚合 CPU profile 可能反而稀释事件，此时应考虑 trace 或飞行记录器。

\section{延迟分布与协调遗漏}
延迟分布常有长尾。报告 p50、p95、p99 以及成功数和错误数；不要把多个实例的 p99 再简单平均称为全局 p99。全局分位数应从合适的原始样本或可合并分布结构估计，并说明桶边界和误差。

闭环负载发生器通常等收到响应后才发出下一个请求。如果研究目标是外部持续到达的请求流，服务阻塞时，闭环发生器减少发压，原本应到达的请求及其等待时间便不会进入已观测延迟，这会造成协调遗漏。如果真实业务本来就由固定数量、等待响应的用户组成，闭环模型仍可能合适；关键是明确要模拟的到达过程。开环发生器按照外部时间表发压，应分别记录预定到达时刻、实际发送时刻和完成时刻，区分发生器调度滞后与服务耗时，并验证发生器自身的容量。

为了比较实现机制，建议固定工作量；为了评估 SLO，建议画多负载点的响应曲线，并固定到达过程。二者回答不同问题。不要在新版本跑更多请求后只比较 profile 总量，也不要在闭环过载时宣称“p99 仍很稳定，因此容量充足”。

\section{实验报告的最小可复现协议}
建议报告按“问题—假设—方法—结果—解释—限制”写作。方法章节至少覆盖自变量、响应变量、控制变量、随机化、重复次数、预热、异常剔除规则和采集开销。异常值若来自真实 GC、节流或尾部排队，正是待解释现象，不能为了曲线好看而删除。

因果主张要控制替代解释：性能改变是否由输入更小、缓存更热、失败更早或正确性降级导致？运行一组撤回实验、输入规模扫描和关闭 profiler 的对照实验，可以增强结论可信度。最后以真实工作负载或至少另一个独立负载验证外部有效性。

课程评分不奖励预设“成功优化”。明确解释无收益、给出原因并诚实保留数据，同样是合格研究结果。第 14 章提供完整评分量规。

\section{实验任务}
\subsection*{完成一次严谨的编码器比较}
\begin{enumerate}
\item 进入 lab，阅读 encodeSlow 与 encodeFast，先运行 TestEncoders。

\item 将两种基准各测量 10 次，交错或随机安排比较顺序，保存原始输出、Go 版本、CPU 信息及容器配额；若只使用 -count=10，应注明这些重复来自同一测试进程。

\item 检查 ns/op、B/op 和 allocs/op；在相同负载下比较关闭与开启 CPU 采集的结果，估计观察者效应。

\item 将输入长度改为 16、256、4096 字节等多个规模，检验结论是否稳定。

\item 对相同工作量采集 profile，解释 CPU 差异的来源；将尚未验证的生产环境外推写入适用范围与限制。

\end{enumerate}
\section{自测与解析}
\noindent\textbf{1. 为何不能把 benchmark 的每次内层迭代都当成独立实验？}
\begin{enumerate}[label=\Alph*.]
\item Go 不支持统计分析
\item 共享缓存、堆与环境会使样本相关
\item 迭代一定没有执行
\item 独立样本必须来自不同语言
\end{enumerate}
\noindent{\color{forest}\textbf{答案 B。}}独立性是统计假设，不由循环次数自动保证。

\noindent\textbf{2. 提高采样率的正确态度是？}
\begin{enumerate}[label=\Alph*.]
\item 越高一定越真实
\item 先测开销、平台支持与采样相关性
\item 设为每纳秒一次
\item 用它代替所有 trace
\end{enumerate}
\noindent{\color{forest}\textbf{答案 B。}}频率提升会增加扰动，且不消除系统性偏差。

\noindent\textbf{3. 能否将多个实例的 p99 取算术平均，得到全局 p99？}
\begin{enumerate}[label=\Alph*.]
\item 能
\item 不能，需可合并分布或原始样本
\item 只有实例数偶数时能
\item 只需查看平均 CPU 用量
\end{enumerate}
\noindent{\color{forest}\textbf{答案 B。}}分位数不是可直接线性平均的统计量。

\section{分析题与参考答案}
\noindent\textbf{题 1：}在独立、同分布样本及正态近似下，预期占比为 10\% 的热点若要达到约 \ensuremath{\pm}1 个百分点的双侧 95\% 误差范围，需要多少个样本？

\noindent\textbf{参考答案：}解不等式 1.96\ensuremath{\times}sqrt(0.1\ensuremath{\times}0.9/N)\ensuremath{\leq}0.01，得到 N\ensuremath{\geq}3457.44，因此至少需要 3458 个样本。这只是一项模型内的样本量规划，不包含时间相关性、栈展开错误或负载偏差；实际采集仍需用多个可比较时间窗检验稳定性。

\noindent\textbf{题 2：}优化后报告的平均延迟下降 8\%，但超时比例从 0.1\% 升到 2\%，应如何解释？

\noindent\textbf{参考答案：}首先检查平均延迟是否只统计成功请求。报告应同时给出尝试总数、成功数、超时数、超时阈值及失败请求已观测到的耗时；超时请求的真实完成时长通常不可知，应说明超时截止后数据的处理方式。若最慢的请求因超时被排除，成功请求的均值就可能看似改善。还需比较容量曲线与业务 SLO，不能仅凭均值判定优化成功。

\medskip\noindent\textbf{本章主要阅读：}\cite{go-testing}\cite{go-benchstat}。
\chapter{pprof 的数据模型与读图方法}
\label{ch:04}
\noindent{\large\color{forest}理解样本值、调用路径、符号化和查询过滤，才能避免读错图。}\par\medskip
\noindent\textbf{先修要求：}第 01–03 章；调用栈和采样基础。

\noindent\textbf{完成本章后，应能够：}
\begin{itemize}
\item 解释 profile.proto 中 sample、location、function 与 mapping
\item 正确选择 sample\_index 并读懂 flat/cum
\item 区分火焰图、调用图和时间线
\item 在一致工作量下解释差分 profile 与标签过滤
\end{itemize}
\section{一个 profile 的结构}
\begin{figure}[H]
\centering
\begin{tikzpicture}[x=1cm,y=1cm]
\node[draw=linegray!75!black,fill=linegray!13,rounded corners=3pt,text width=3.33cm,minimum height=1.12cm,align=center,font=\small] (p0-sample) at (1.88,0.00) {Sample\\{\scriptsize location\_id=[7,3]\\value=[6,60000000]}};
\node[draw=linegray!75!black,fill=linegray!13,rounded corners=3pt,text width=3.33cm,minimum height=1.12cm,align=center,font=\small] (p0-loc) at (5.62,0.00) {Location 7\\{\scriptsize address=0x1040\\line[0].function\_id=12}};
\node[draw=linegray!75!black,fill=linegray!13,rounded corners=3pt,text width=3.33cm,minimum height=1.12cm,align=center,font=\small] (p0-func) at (9.38,0.97) {Function 12\\{\scriptsize name=encode}};
\node[draw=linegray!75!black,fill=linegray!13,rounded corners=3pt,text width=3.33cm,minimum height=1.12cm,align=center,font=\small] (p0-mapping) at (9.38,-0.97) {Mapping\\{\scriptsize 文件 /\allowbreak{} build ID}};
\node[draw=leaf!75!black,fill=leaf!13,rounded corners=3pt,text width=3.33cm,minimum height=1.12cm,align=center,font=\small] (p0-report) at (13.12,0.00) {报告\\{\scriptsize 选择 cpu 时间}};
\draw[-{Stealth[length=2mm]},forest!70!black] (p0-sample.east) -- (p0-loc.west);
\draw[-{Stealth[length=2mm]},forest!70!black] (p0-loc.east) -- (p0-func.west);
\draw[-{Stealth[length=2mm]},forest!70!black] (p0-loc.east) -- (p0-mapping.west);
\draw[-{Stealth[length=2mm]},forest!70!black] (p0-func.east) -- (p0-report.west);
\draw[-{Stealth[length=2mm]},forest!70!black] (p0-sample.north) -- (1.88,2.12) -- (13.12,2.12) -- (p0-report.north);
\end{tikzpicture}
\caption{pprof 用 ID 连接样本、位置和函数}\label{fig:04-format}
\begin{minipage}{0.96\linewidth}\footnotesize 教学格式：sample\_type 按顺序为 samples/count 和 cpu/nanoseconds；value 是 int64 向量，所以 60000000 表示 60 ms，不能把带单位字符串写进原始字段。

\textbf{读图顺序：}1. 先选值类型\quad 2. 解析位置\quad 3. 解析函数和映射\quad 4. 按权重聚合。
\textbf{连线语义：}Sample\ensuremath{\rightarrow}\allowbreak{}Location 7：局部ID；Location 7\ensuremath{\rightarrow}\allowbreak{}Function 12：行记录；Location 7\ensuremath{\rightarrow}\allowbreak{}Mapping：地址映射；Function 12\ensuremath{\rightarrow}\allowbreak{}报告：符号；Sample\ensuremath{\rightarrow}\allowbreak{}报告：权重。
\textbf{机制依据：}\cite{go-proto}。
\end{minipage}
\end{figure}
pprof 既指 Go 的采集接口，也指分析工具和通用 profile 数据格式。Go 程序可以通过 runtime/pprof 或 HTTP 诊断端点导出 profile，再交给 go tool pprof 解析。不同语言采集器也可产生相同格式，但相同格式不保证相同统计语义。\cite{go-proto}\cite{go-pprof}

profile.proto 中，\texttt{sample\_\allowbreak{}type} 定义各项样本值的类型与单位；每个 \texttt{sample.\allowbreak{}value} 都必须具有相同的长度，并与 \texttt{sample\_\allowbreak{}type} 逐项对应。\texttt{sample.\allowbreak{}location\_\allowbreak{}id} 引用栈中的 location，叶端位于索引 0。location 的 \texttt{line} 记录引用 function，并保存对应源码行号；function 保存函数名、文件名与函数起始行；\texttt{mapping\_\allowbreak{}id} 则关联二进制映射。\texttt{string\_\allowbreak{}table} 通过共享字符串表减少重复存储，标签为样本附加字符串或数值维度。

\begingroup\small\setlength{\tabcolsep}{4pt}
\begin{longtable}{>{\raggedright\arraybackslash}p{0.298\linewidth}>{\raggedright\arraybackslash}p{0.298\linewidth}>{\raggedright\arraybackslash}p{0.298\linewidth}}
\toprule
\textbf{结构} & \textbf{作用} & \textbf{需要检查的条件} \\
\midrule
\endfirsthead
\toprule
\textbf{结构} & \textbf{作用} & \textbf{需要检查的条件} \\
\midrule
\endhead
sample\_type & 例如 alloc\_objects/count、alloc\_space/bytes & 选择的统计量是否与问题一致 \\
\addlinespace[3pt]
sample.value & 对应各 sample\_type 的数值向量 & 不是任意原始事件计数 \\
\addlinespace[3pt]
location & 地址、行号和可能的内联帧 & 是否正确展开与符号化 \\
\addlinespace[3pt]
function & 函数名、文件名、起始行 & 是否匹配该构建 \\
\addlinespace[3pt]
mapping & 地址区间、文件与 build ID & 二进制与调试符号是否正确 \\
\addlinespace[3pt]
label & 样本属性或业务分类 & 是否低基数、是否包含敏感数据 \\
\addlinespace[3pt]
\bottomrule
\end{longtable}
\endgroup
Go 的 CPU profile 通常同时包含样本计数和估计的 CPU 时间；heap profile 则包含多个内存计量维度。优先检查 \texttt{go tool pprof -\allowbreak{}raw file.\allowbreak{}pb.\allowbreak{}gz}，而不是依靠图形默认值猜测。

\section{flat 与 cum：相同样本的两种归因}
\begin{figure}[H]
\centering
\begin{tikzpicture}[x=1cm,y=1cm]
\path[fill=bluegray!25] (0.00,0.00) rectangle (13.97,0.78);
\node[align=center,text width=13.92cm,font=\small] at (7.00,0.40) {main};
\path[fill=leaf!38] (0.00,1.05) rectangle (12.57,1.83);
\node[align=center,text width=12.52cm,font=\small] at (6.30,1.45) {serve};
\path[fill=leaf!51] (0.00,2.10) rectangle (8.38,2.88);
\node[align=center,text width=8.32cm,font=\small] at (4.20,2.50) {encode};
\path[fill=bluegray!51] (8.40,2.10) rectangle (12.58,2.88);
\node[align=center,text width=4.12cm,font=\small] at (10.50,2.50) {hash};
\path[fill=bluegray!38] (12.60,1.05) rectangle (13.97,1.83);
\node[font=\scriptsize] at (13.30,1.45) {5};
\node[anchor=north west,align=left,text width=14cm,font=\scriptsize] at (0,-.45) {1=main: flat 0 /\allowbreak{} cum 100；2=serve: flat 0 /\allowbreak{} cum 90；3=encode: flat 60 /\allowbreak{} cum 60；4=hash: flat 30 /\allowbreak{} cum 30；5=cleanup: flat 10 /\allowbreak{} cum 10};
\end{tikzpicture}
\caption{一次采样同时贡献给多个 cum}\label{fig:04-flat-cum}
\begin{minipage}{0.96\linewidth}\footnotesize 100 个等权教学样本；无递归、内联重归因或栈截断。真实工具还需处理这些情况。

\textbf{读图顺序：}1. 60 次落在 encode\quad 2. 再加入 hash 的 30 次\quad 3. 加入 cleanup 的 10 次\quad 4. 选择归因问题。
\textbf{机制依据：}\cite{go-proto}\cite{go-pprof}。
\end{minipage}
\end{figure}
flat 将所选样本权重归给报告中保留的叶端位置，具体归因还受内联展开、聚合和栈节点过滤影响；cum 则统计包含该报告节点的样本权重，涵盖它自身及其下游调用路径。对于 CPU profile，cum 不等于函数独占的 CPU 时间，也不能与其子函数的 cum 直接相加。一个业务包装函数的 flat 很小、cum 很大，通常意味着成本主要发生在其下游调用中。

教学例：假设有 100 个等权样本，60 个栈为 main\ensuremath{\rightarrow}\allowbreak{}serve\ensuremath{\rightarrow}\allowbreak{}encode，30 个为 main\ensuremath{\rightarrow}\allowbreak{}serve\ensuremath{\rightarrow}\allowbreak{}hash，10 个为 main\ensuremath{\rightarrow}\allowbreak{}cleanup。serve 的 flat 为 0、cum 为 90；encode 的 flat 与 cum 均为 60。若将 main 的 100、serve 的 90 和 encode 的 60 相加，就会重复统计共享样本。

\begin{codeblock}
go tool pprof -top cpu.pb.gz
go tool pprof -top -cum cpu.pb.gz
go tool pprof -list=encodeSlow cpu.pb.gz
go tool pprof -raw cpu.pb.gz
\end{codeblock}

调用图的边表示调用关系与样本归因，默认会裁剪权重较低的节点；递归、内联、多父调用和聚合会让图不同于静态调用图。节点未出现可能是过滤或阈值所致，并不能证明函数没有执行。

递归还要避免另一种重复：在同一报告聚合节点中，同一条样本重复经过该节点时，cum 只计该样本一次。例如一条权重 10 的 A\ensuremath{\rightarrow}\allowbreak{}A\ensuremath{\rightarrow}\allowbreak{}B 样本，在按函数聚合时 A 的 cum 为 10，而非 20；火焰图展开出的不同递归深度也不能再相加解释为额外 CPU。\cite{tools-pprof-cli}

\section{火焰图不是时间线}
在常规的非负权重 profile 中，火焰图的框宽反映对应调用路径的所选样本权重，纵向表示调用栈层次。左右位置一般来自排序和布局，不表示开始时间；颜色往往只是区分或编码类别，必须看具体工具图例。“左边先发生，右边后发生”是常见误读。

差分模式需要另外解读。本书核对的 pprof 网页实现中，使用 -diff\_base 时，框宽反映子树中增加量与减少量绝对值的合计，净变化另以阴影表示。在文档的示例中，一个子分支减少 150、另一个增加 200，框宽对应 350，净增加为 50，不能把框宽直接读作净增量。\cite{tools-pprof-cli}

pprof 网页火焰图聚合的是整个 profile 中的栈和权重，适合分析成本分布；go tool trace 提供时间轴，适合分析事件先后和等待。持续消耗 CPU 的函数适合用 CPU 采样定位；发生在 5 ms 内的偶发锁竞争或排队，则通常需要事件追踪补充证据。工具之间的差异来自数据模型，不是界面优劣。

\begin{codeblock}
go tool pprof -http=127.0.0.1:8081 cpu.pb.gz
\end{codeblock}

图形视图中部分调用图输出依赖 Graphviz，终端的 top/list 输出则不需要。建议保留命令行摘要，即使浏览器图形环境暂时不可用，核心证据仍然可复查。

\section{选择统计量与过滤范围}
堆 profile 的四个常见 sample index 为 \texttt{alloc\_\allowbreak{}objects}、\texttt{alloc\_\allowbreak{}space}、\texttt{inuse\_\allowbreak{}objects}、\texttt{inuse\_\allowbreak{}space}。累计分配适合分析分配压力，存活量适合分析保留内存；对象数和字节数则区分大量小对象与少量大对象。Go 的 heap profile 默认通常选择 inuse\_space，allocs profile 则默认选择 alloc\_space。研究报告应显式记录所选统计量，避免因默认值不同而比较了不同指标。

\begin{codeblock}
go tool pprof -sample_index=alloc_space -top mem.pb.gz
go tool pprof -sample_index=inuse_space -top mem.pb.gz
go tool pprof -focus='encode|hash' -top cpu.pb.gz
go tool pprof -hide='runtime\.' -top cpu.pb.gz
\end{codeblock}

focus 与 ignore 以整条样本为单位筛选：只要栈中包含匹配位置，样本就可能被选入或排除。hide 与 show 则主要过滤栈中的报告节点；例如，上面的 hide 命令隐藏 runtime 节点，不会仅因栈中含有 runtime 就删除整条样本。如果一条样本的所有栈节点都被隐藏，工具会移除这条空栈样本。若误用 \texttt{-\allowbreak{}ignore='runtime\textbackslash{}.\allowbreak{}'}，大量包含运行时栈帧的业务样本也可能被排除。报告应记录过滤条件、active filters 和百分比分母。隐藏运行时节点只是一种观察方式，并不能证明 GC 或调度与业务成本无关。

标签过滤也不能让不支持标签的 profile 突然具有请求维度。Go 的 pprof labels 主要供 CPU 和 goroutine profile 使用，标签传播及继承应按具体版本检查；不能默认 heap、block 和 mutex profile 都支持按标签细分。\cite{go-pprof}

默认 \texttt{relative\_\allowbreak{}percentages=false} 时，pprof 在 focus/ignore 等过滤前确定总量，过滤后的表不一定重新归一到 100\%。显式启用 \texttt{-\allowbreak{}relative\_\allowbreak{}percentages} 才先过滤再取报告分母；hide/show 主要改变栈节点显示，不等同删除整条样本。比较图表时应保留这个开关以及 active filters。

\section{差分、基线与代表性}
在相同工作量、相同采样配置、可比构建下，差分能提示成本迁移。\texttt{-\allowbreak{}base} 常用于同一进程的累计快照相减；在正常单调累计的情形，百分比相对总量之差。\texttt{-\allowbreak{}diff\_\allowbreak{}base} 适合前后对比，百分比相对基线总量并保留有意义的负变化。一般有正负样本的情况分母处理更复杂，应同时报告绝对权重并核对工具帮助。不要把不同时间长度的两个累计 profile 直接相减后宣布优化成功。

\begin{codeblock}
go tool pprof -top -base=before.pb.gz after.pb.gz
go tool pprof -top -diff_base=before.pb.gz after.pb.gz
\end{codeblock}

比较周期内若请求数不同，可同时报告总量和每请求归一化结果。可使用工具归一化能力作形状比较，但按总 profile 权重归一化会掩盖总成本变化，不能替代业务分母。更稳妥的实验是完成相同输入批次后分别采集，并把独立基准结果与 profile 的解释配对。

对长期运行进程的累计分配、block 和 mutex profile，先后快照之差可近似表示窗口内的增量；进程重启、采样率变化或语义不同会破坏可比性。heap 中的存活量估计与 GC 周期及对象回收有关，可能自然减少；它的差分表示保留量的净变化，不能解释为窗口内的累计分配。

\section{符号化、优化编译与证据保存}
分析样本需要从程序计数器恢复函数和源码。ASLR、内联、裁剪符号、二进制不匹配、源码版本漂移都会影响结果。发布时应保留构建 ID、完整二进制或对应符号产物及源码 commit；对 cgo 和原生库还需保留相应库映射与调试信息。

为提高 Delve 对变量和控制流的可见性，使用 \texttt{-\allowbreak{}gcflags='all=-\allowbreak{}N -\allowbreak{}l'} 会禁用优化和内联，适用于调试构建；不能把它的 profile 当作正常生产优化构建的性能特征。先剖析真实构建，再为具体控制流问题使用调试构建，明确两者的边界。

证据包至少包含 profile 文件、采集命令、采集时间、duration、sample index、构建标识、输入或流量统计、环境和解读命令。分享截图时仍保留原始文件；截图无法恢复过滤、采样权重和符号化细节。

\section{实验任务}
\subsection*{同一文件的四种内存视角}
\begin{enumerate}
\item 在 lab 中运行第 03 章生成 mem.pb.gz 的基准命令，保留 profile 文件及匹配的测试二进制。

\item 使用 -raw 检查 sample\_type，然后分别选择四种内存 sample\_index，记录各自的类型与单位。

\item 比较 flat、cum 和 list 输出，写出每种视图能回答的问题。

\item 打开网页火焰图，选择一个父节点，解释其宽度为何不能与子节点相加。

\item 选取一个未在 top 中出现的函数，检查过滤条件和显示阈值，并查看 list 输出；区分“未采到样本”与“采到了但被显示规则裁剪”，不要把零命中直接判为零成本。

\end{enumerate}
\section{自测与解析}
\noindent\textbf{1. 火焰图从左到右通常表示什么？}
\begin{enumerate}[label=\Alph*.]
\item 执行先后
\item 按栈聚合后的布局，通常不含时序
\item CPU 核心编号
\item 网络包到达次序
\end{enumerate}
\noindent{\color{forest}\textbf{答案 B。}}横向宽度代表选定的样本权重，左右布局通常不表示事件先后；要分析时序，应查看 trace。

\noindent\textbf{2. 为何不能将父函数的 cum 与子函数的 cum 直接相加？}
\begin{enumerate}[label=\Alph*.]
\item 单位不同
\item 共享样本会重复计数
\item cum 永远为 0
\item 只因递归存在
\end{enumerate}
\noindent{\color{forest}\textbf{答案 B。}}父函数的 cum 已包含经过下游函数的样本权重，相加会重复计数。

\noindent\textbf{3. 分析临时小对象分配压力应优先比较什么？}
\begin{enumerate}[label=\Alph*.]
\item 只查看 inuse\_space
\item 同时比较 alloc\_objects 与 alloc\_space
\item 只查看 RSS 截图
\item 源码行数
\end{enumerate}
\noindent{\color{forest}\textbf{答案 B。}}临时对象可能很快被回收，不再体现在存活堆中；累计分配对象数和字节数有助于揭示分配压力。

\section{分析题与参考答案}
\noindent\textbf{题 1：}在 100 个等权样本中，调用栈 A\ensuremath{\rightarrow}\allowbreak{}B 出现 40 次，A\ensuremath{\rightarrow}\allowbreak{}C 出现 60 次。求按函数聚合后 A、B 的 flat 与 cum。

\noindent\textbf{参考答案：}A 没有作为叶端出现，因此 A 的 flat=0、cum=100；B 的 flat=40、cum=40。这里假设没有内联、递归、过滤或其他改变归因的因素。

\noindent\textbf{题 2：}旧版本的 CPU profile 采集了 30 秒，新版本采集了 60 秒。能否直接用两者的 CPU 总权重判断效率变化？

\noindent\textbf{参考答案：}不能直接作出效率结论。应固定完成工作量，或按同一口径的成功请求数归一化，同时报告采集时长、吞吐及失败率。仅按 profile 总权重归一化，通常只能比较成本分布，可能掩盖总消耗的变化。

\medskip\noindent\textbf{本章主要阅读：}\cite{go-proto}\cite{go-pprof}\cite{go-diagnostics}\cite{tools-pprof-cli}。
\part*{第二篇 · 工具与机制}
\addcontentsline{toc}{part}{第二篇 · 工具与机制}
\chapter{CPU 剖析}
\label{ch:05}
\noindent{\large\color{forest}把采样热点转化为可检验的计算成本解释}\par\medskip
\noindent\textbf{先修要求：}第 02–04 章；能够构建 Go 程序、运行基准测试并理解调用栈。

\noindent\textbf{完成本章后，应能够：}
\begin{itemize}
\item 区分墙钟时间、累计 CPU 时间和 CPU profile 样本权重
\item 从 flat、cum、调用图及源码视图逐层定位计算热点
\item 正确传播低基数 labels，并识别未覆盖的异步工作
\item 用独立复测判断 CPU 优化是否转化为业务收益
\end{itemize}
\section{测量对象：CPU 时间不是请求耗时}
\begin{figure}[H]
\centering
\begin{tikzpicture}[x=1cm,y=1cm]
\node[anchor=east,font=\small] at (1,0.0) {泳道1};
\draw[linegray] (1.2,-0.35) -- (14.799999999999999,-0.35);
\path[fill=leaf!60] (1.200,-0.250) rectangle (5.280,0.250);
\node[align=center,font=\small,text width=3.98cm] at (3.240,0.750) {哈希};
\path[fill=bluegray!60] (5.280,-0.250) rectangle (12.080,0.250);
\node[align=center,font=\small,text width=6.70cm] at (8.680,0.750) {sleep};
\path[fill=leaf!60] (12.080,-0.250) rectangle (14.800,0.250);
\node[align=center,font=\small,text width=2.62cm] at (13.440,0.750) {哈希};
\draw[-{Stealth[length=2mm]},forest] (1.2,-1.15) -- (14.999999999999998,-1.15);
\draw[forest] (1.20,-1.15) -- (1.20,-1.25);
\node[below,font=\scriptsize] at (1.20,-1.25) {0};
\draw[forest] (3.92,-1.15) -- (3.92,-1.25);
\node[below,font=\scriptsize] at (3.92,-1.25) {20};
\draw[forest] (6.64,-1.15) -- (6.64,-1.25);
\node[below,font=\scriptsize] at (6.64,-1.25) {40};
\draw[forest] (9.36,-1.15) -- (9.36,-1.25);
\node[below,font=\scriptsize] at (9.36,-1.25) {60};
\draw[forest] (12.08,-1.15) -- (12.08,-1.25);
\node[below,font=\scriptsize] at (12.08,-1.25) {80};
\draw[forest] (14.80,-1.15) -- (14.80,-1.25);
\node[below,font=\scriptsize] at (14.80,-1.25) {100};
\node[anchor=east,font=\scriptsize] at (14.799999999999999,-1.9) {时间 / ms};
\end{tikzpicture}
\caption{当前 G 的等待时间不计作 CPU 执行}\label{fig:05-cpu-wait}
\begin{minipage}{0.96\linewidth}\footnotesize 教学示意：用于解释机制，不是运行时的实测数据或完整实现。

\textbf{读图顺序：}1. 执行中可能被采到\quad 2. 等待区间\quad 3. 恢复执行\quad 4. 两个不同问题。
\textbf{机制依据：}\cite{tools-pprof}。
\end{minipage}
\end{figure}
CPU profile 回答的是“采样期间，处理器执行哪些调用栈”，而非“每个函数从进入到返回用了多久”。设墙钟窗口为 T，进程内各线程消耗的 CPU 时间之和为 C，则平均使用的 CPU 核数近似为 C/T。四个核充分繁忙时，10 秒窗口可以产生约 40 CPU 秒；窗口内等待网络的请求即使持续 8 秒，也可能只贡献很少 CPU 样本。报告中的百分比通常以所选样本类型的总权重为分母，不能直接读作单核利用率或请求延迟占比。\cite{tools-pprof}

Go 1.25 的 \texttt{runtime/\allowbreak{}pprof.\allowbreak{}StartCPUPr\allowbreak{}ofile} 固定设置名义 100 Hz 采样率，样本通常带有约 10 ms 的 CPU 时间权重。它并不每 10 ms 遍历全部 goroutine，也不准确记录每一次函数调用。操作系统、调度、cgo、信号处理、采样缓存丢失及栈展开都会影响实际覆盖。一次 2 ms 的函数调用可能完全没有被采到；大量重复的小调用则可以凭总体占比稳定出现。应观察足够的工作量，而非把采样周期当成最小可优化函数耗时。\cite{tools-pprof}

把 CPU、on-CPU、off-CPU 与 wall time 分开记账。CPU 时间包含执行所消耗的时间；内核性能工具的 on-CPU 事件能覆盖特定配置下的用户态和内核态执行；off-CPU 是线程不在 CPU 上的时间集合，可能包含睡眠、I/O 等待或等待调度。Go 中处于 runnable 状态的 goroutine 还可能在运行队列中等待调度；与此同时，持有 P 的线程可能正在执行其他 goroutine。只有把运行时状态和操作系统状态关联起来，才可能解释尾延迟。第 07、08、10 章分别补足同步等待、时序与系统边界。

名义 100 Hz 是 \texttt{StartCPUPr\allowbreak{}ofile} 内部固定值，不是 HTTP 参数可调的默认选项。不要先调用 \texttt{runtime.\allowbreak{}SetCPUProf\allowbreak{}ileRate} 再期待 pprof 自动采用该频率：两者管理同一运行时采集器，活动期间改频率会被拒绝，混用还可能使生命周期不一致。教学动画中的采样间隔变化应理解为统计模型，不是标准 pprof 的调频接口。

\section{采集：选择代表性窗口和正确二进制}
短程序可使用 \texttt{runtime/\allowbreak{}pprof}，基准测试可使用 \texttt{go test -\allowbreak{}cpuprofile}，服务可使用第 10 章的诊断 HTTP 端点。保留实际参与采集的二进制、源码提交、构建参数、Go 版本、负载参数和起止时间。源码行号与反汇编需要匹配的构建产物；相同函数名不足以证明二进制匹配。不要为了采 CPU profile 默认关闭内联与优化，否则测到的是调试构建的性能。

\begin{codeblock}
go test -run='^$' -bench='BenchmarkEncode$' -benchtime=10s   -cpuprofile=cpu.pprof -o encode.test .
go tool pprof -top ./encode.test cpu.pprof
go tool pprof -http=127.0.0.1:8081 ./encode.test cpu.pprof
\end{codeblock}

上述命令假设当前包确实含有 \texttt{BenchmarkE\allowbreak{}ncode}。通过 \texttt{-\allowbreak{}run='\textasciicircum{}\$'} 避开普通测试，通过明确的 benchmark 正则固定采样工作负载。服务端采集则必须在施加稳定负载的同时进行：

\begin{codeblock}
curl --fail --max-time 40 -o cpu.pprof   'http://127.0.0.1:6060/debug/pprof/profile?seconds=30'
go tool pprof -top ./service cpu.pprof
\end{codeblock}

默认 CPU 窗口为 30 秒，但窗口长度应取决于事件频度及稳定性。长窗口可以降低随机误差，也可能把冷启动、稳态和故障恢复混合成无法解释的平均值。先划分阶段，再选择窗口。进程只允许一个活动的 CPU profile；HTTP 采集与程序内部的 \texttt{StartCPUPr\allowbreak{}ofile} 会竞争同一全局资源。必须检查开始采集的错误，并在退出前调用 \texttt{StopCPUPro\allowbreak{}file}；后者返回时才保证 profile 写出完成。\cite{tools-http-pprof}\cite{tools-pprof}

\section{从 flat 与 cum 走向调用路径}
\texttt{flat} 是样本直接归属于某函数的权重，\texttt{cum} 是该函数及其下游调用所覆盖的权重。若入口 \texttt{Handle} 的 flat 为 0.1 秒、cum 为 8 秒，而 \texttt{sha256.\allowbreak{}block} 的 flat 为 6 秒，则入口是主要工作路径，哈希实现是主要执行地点。不能将所有函数的 cum 相加；同一份样本沿调用栈进入多个祖先节点。递归、内联和多调用方会让“每层加一次”的直觉更加危险，具体累计以工具的调用图聚合规则为准；官方说明同一样本重复出现同一位置时仅计入一次。\cite{tools-pprof-cli}

\begin{codeblock}
go tool pprof -top -cum ./service cpu.pprof
go tool pprof -list='Encode|Marshal' ./service cpu.pprof
go tool pprof -disasm='main.hotLoop' ./service cpu.pprof
\end{codeblock}

按 flat 排序适合寻找自身执行成本较高的函数，按 cum 排序适合寻找值得重新设计的调用路径。接着用 callers/callees 与源码视图回答：热点被谁调用、输入规模是否异常、是否重复计算、是否出现不必要的格式转换。如果编码库位居榜首，仅替换一条循环语句未必有效；上层重复序列化可能才是原因。

火焰图的宽度表示聚合后的样本权重，高度表示调用深度，横向相邻的两个块通常不代表时间先后。火焰图不能证明“先发生 GC，再发生请求超时”。源码行权重也可能由于编译器移动指令、内联和指令地址映射而集中到某一行；反汇编帮助检查实现，但不是逐语句精确计时器。

初学者可分三步缩小范围：先找到贡献最大的业务路径，再在该路径内寻找成本最高的阶段，最后用基准测试或输入规模实验检验机制。每一步都应保留原始总量与当前过滤范围；如果视图按筛选后的集合重新归一化，一个占原始总量很小的函数也可能显示为 100\%。pprof 的默认百分比分母及 \texttt{relative\_\allowbreak{}percentages} 开关见第 04 章。

还有一类热点来自算法输入规模。若排序耗时随输入项数增长，应该用几组不同规模的数据测量增长趋势，而不仅仅替换一个比较函数。将问题表述为“每个请求排序几次、每次多少元素”，再用调用计数或业务指标确认；调用次数本身通常不能从普通 CPU 样本精确恢复。找到高累计入口后，检查重复解析、无效重试和不必要的跨格式转换，往往比优化叶子函数常数因子更有效。

在数学上，假设可优化部分占总成本比例为 p，局部加速倍数为 s，且其余工作与负载保持不变，则整体加速比模型为 \texttt{1 /\allowbreak{} ((1-\allowbreak{}p) + p/\allowbreak{}s)}。如果热点占四成，即使将这部分耗时降至零，总耗时仍为原来的六成，整体加速比上限约为 1.67。这个模型帮助筛选投资方向，但请求存在并行、排队或外部等待时，需要先判断该 CPU 成本是否位于关键路径。

\section{使用 labels 归因请求类别及其限制}
当不同业务共用编码、压缩或加密代码时，函数栈无法直接说明成本属于哪类请求。\texttt{pprof.\allowbreak{}Do} 将传入 context 的标签与 \texttt{Labels} 参数合并，并安装到当前 goroutine；回调结束后恢复传入 context 所携带的标签。回调内新建的 goroutine 继承创建者当时的 labels。\texttt{WithLabels} 只返回新的 context，并不会独自改变当前 goroutine 的采样标签。Go 1.25 明确只有 CPU 与 goroutine profiles 使用这些 labels，不能据此承诺 heap、block 或 mutex 已实现按租户归因。\cite{tools-labels}\cite{tools-label-runtime}

以下为集成片段，\texttt{encode} 与 \texttt{payload} 来自业务代码：

\begin{codeblock}
pprof.Do(ctx, pprof.Labels("route", "search", "phase", "encode"),
    func(ctx context.Context) {
        encode(payload)
    })
\end{codeblock}

\begin{codeblock}
go tool pprof -tags ./service cpu.pprof
go tool pprof -top -tagfocus='route=search' ./service cpu.pprof
\end{codeblock}

常驻 worker pool 不会因为接收了带标签的 context 就自动更换 goroutine labels。应在 worker 处理每个任务时调用 \texttt{pprof.\allowbreak{}Do}，并让任务携带相应的 context。新 goroutine 的继承也不等于生命周期自动结束：若请求启动永久后台 goroutine，其标签可能长久保留。把请求级工作、后台维护和共享批处理分别建模；共享 GC 或批量任务的成本通常不能简单归还到某个请求。

标签维度应是有限的路由、操作类型或服务等级，避免每请求唯一的 ID、完整 URL、用户输入和令牌。高基数会削弱聚合效果，并增加存储、处理和隐私保护成本。labels 用于按类别拆分聚合数据，不能替代分布式调用链；关联跨进程因果关系仍需 span 与请求上下文。

\section{采样误差、编译器效应与常见误判}
\begin{figure}[H]
\centering
\begin{tikzpicture}[x=1cm,y=1cm]
\path[fill=leaf!25] (0.00,0.00) rectangle (13.97,0.78);
\node[align=center,text width=13.92cm,font=\small] at (7.00,0.40) {全部样本};
\path[fill=bluegray!38] (0.00,1.05) rectangle (8.38,1.83);
\node[align=center,text width=8.32cm,font=\small] at (4.20,1.45) {workload=A};
\path[fill=bluegray!51] (0.00,2.10) rectangle (8.38,2.88);
\node[align=center,text width=8.32cm,font=\small] at (4.20,2.50) {encode};
\path[fill=bluegray!38] (8.40,1.05) rectangle (13.97,1.83);
\node[align=center,text width=5.52cm,font=\small] at (11.20,1.45) {workload=B};
\path[fill=bluegray!51] (8.40,2.10) rectangle (13.97,2.88);
\node[align=center,text width=5.52cm,font=\small] at (11.20,2.50) {hash};
\node[anchor=north west,align=left,text width=14cm,font=\scriptsize] at (0,-.45) {1=全部样本: flat 0 /\allowbreak{} cum 100；2=workload=A: flat 0 /\allowbreak{} cum 60；3=workload=B: flat 0 /\allowbreak{} cum 40；4=encode: flat 60 /\allowbreak{} cum 60；5=hash: flat 40 /\allowbreak{} cum 40};
\end{tikzpicture}
\caption{标签筛选改变了观察的样本集合}\label{fig:05-labels}
\begin{minipage}{0.96\linewidth}\footnotesize 教学示意：用于解释机制，不是运行时的实测数据或完整实现。

\textbf{读图顺序：}1. 混合负载\quad 2. 只选 A\quad 3. 只选 B\quad 4. 检查支持边界。
\textbf{机制依据：}\cite{tools-labels}\cite{tools-label-runtime}。
\end{minipage}
\end{figure}
在理想独立采样模型中，若总样本数为 n、某热点占比为 p，则比例估计的标准误差近似为 \texttt{sqrt(p \ensuremath{\times} (1-\allowbreak{}p) /\allowbreak{} n)}。例如假设 n=100、p=0.1，标准误差约为 3 个百分点；仅凭 10\% 降为 8\% 无法证明优化。真实程序存在阶段性和自相关，独立模型通常过于乐观。用不同采集窗口、不同进程的重复实验评估稳定性，比只看一次火焰图更可靠。

较少样本不一定说明程序快：可能没有施加负载、任务提前退出、采集窗口落在空闲阶段，或耗时主要发生在外部服务；采集器对部分内核或非 Go 执行路径的覆盖不足，也可能影响结果。\texttt{runtime.\allowbreak{}mallocgc} 或 GC 标记占比高通常指向分配压力；应转到第 06 章查找分配源。\texttt{runtime.\allowbreak{}futex} 等系统边界符号则需要结合调用上下文与第 10 章工具分析，不能仅凭名字认定锁竞争。

内联让逻辑函数与机器栈不再一一对应，去虚拟化和逃逸分析会改变工作位置。Go 的符号化可以恢复许多内联信息，但行级热点仍需结合优化报告。运行 \texttt{go build -\allowbreak{}gcflags='all=-\allowbreak{}m=2'} 时输出可能很大，应限定到相关包；输出本身只是编译器决策证据，不是优化收益证据。

比较前后 profile 时，必须同时比较每请求 CPU、吞吐和延迟。假设总 CPU 从 10 秒涨到 12 秒，但处理请求从 1000 增至 2000，则单位请求 CPU 从 10 ms 降至 6 ms。反之，某函数占比减半可能只是其他工作增多。第 11 章将把这种归一化用于 PGO 与算法优化决策。

复测时应将诊断采集与正式性能测量分开运行：一组保留诊断证据，另一组关闭采集测量收益。测试顺序可随机交替，避免温度、频率、缓存预热与邻居负载随时间变化造成系统性偏差。若优化只改变少数长尾输入，还应分别报告各输入类别与整体加权结果；在类别分布改变后得到的平均下降，不能归功于代码修改。

最终结论最好写成可反驳的命题：某路径对同样大小的输入重复执行三次编码，删除两次后单位请求 CPU 下降，同时输出与错误语义保持一致。与“火焰图更漂亮”相比，这样的结论明确了机制、条件和验证指标。

\section{可运行实验：让计算与等待分离}
下面是完整程序，保存为 \texttt{cpu.\allowbreak{}go}。程序先做带标签的哈希计算，再做 600 ms 等待。哈希结果参与最终输出，避免把实验写成可被编译器删除的无用计算。该程序用于展示测量语义，不代表真实服务基准。哈希循环以约 3 秒的墙钟截止时间为结束条件，实际哈希次数随机器、调度和采集开销而变化，因而不是固定工作量基准。输出的 \texttt{wall} 还包含停止采集、等待 profile 写完及关闭文件的时间，不应预先认定恰好等于 3.6 秒。

\begin{codeblock}
package main

import (
	"context"
	"crypto/sha256"
	"fmt"
	"os"
	"runtime/pprof"
	"time"
)

func main() {
	f, err := os.Create("cpu.pprof")
	if err != nil {
		panic(err)
	}
	if err := pprof.StartCPUProfile(f); err != nil {
		f.Close()
		panic(err)
	}
	started := time.Now()
	var sum [32]byte
	pprof.Do(context.Background(), pprof.Labels("phase", "hash"),
		func(context.Context) {
			until := time.Now().Add(3 * time.Second)
			for time.Now().Before(until) {
				for i := 0; i < 10000; i++ {
					sum = sha256.Sum256(sum[:])
				}
			}
		})
	pprof.Do(context.Background(), pprof.Labels("phase", "wait"),
		func(context.Context) { time.Sleep(600 * time.Millisecond) })
	pprof.StopCPUProfile()
	if err := f.Close(); err != nil {
		panic(err)
	}
	fmt.Printf("wall=%v result=%x\n", time.Since(started), sum[:4])
}
\end{codeblock}

\begin{codeblock}
go build -o cpu-lab cpu.go
./cpu-lab
go tool pprof -top ./cpu-lab cpu.pprof
go tool pprof -tags ./cpu-lab cpu.pprof
\end{codeblock}

应检验的预期是哈希阶段贡献主要 CPU 样本，而 600 ms 的 \texttt{Sleep} 不会变成同等 CPU 样本；实际数字由本机测量获得。扩大等待时间会增加墙钟耗时而不相应增加计算成本。再把哈希阶段改成两个并发 worker，并保持各自状态独立。若比较优化收益，应先改为固定哈希总次数，并在两种实现之间保持总工作量相同；若仍让每个 worker 运行约 3 秒，则这是并发容量演示。累计 CPU 时间与墙钟时间的关系取决于可用核数、GOMAXPROCS 及配额，不能预先写成固定两倍。

\section{实验任务}
\subsection*{CPU profile 到优化假设的闭环}
\begin{enumerate}
\item 构建并执行 5.6 的完整程序，保存 go version、二进制和初始 profile；为每次运行使用独立输出目录，避免 cpu.pprof 被覆盖。

\item 将等待改为 2 秒，在相同机器上重新构建采集；比较墙钟与样本总量，解释差别。

\item 在原程序上分别运行 flat、cum、tags 视图，写出热点调用路径及标签覆盖范围。

\item 选择一个实际业务 benchmark，记录代表性输入和每次操作的 CPU 时间与墙钟时间；提出一项减少重复计算的修改并独立复测。

\item 提交原始 profile、命令、至少三个独立进程的复测结果及汇总统计与结论；若结果不稳定，报告不确定性而不宣称优化成立。

\end{enumerate}
\section{自测与解析}
\noindent\textbf{1. 10 秒采集窗口得到 28 CPU 秒，最合理的解释是什么？}
\begin{enumerate}[label=\Alph*.]
\item profile 文件损坏
\item 进程平均使用约 2.8 个 CPU 核
\item 单个请求延迟为 28 秒
\item 所有 goroutine 均等待 18 秒
\end{enumerate}
\noindent{\color{forest}\textbf{答案 B。}}CPU profile 聚合多线程计算时间；CPU 秒可以超过墙钟秒，C/T 估计平均使用的核数。

\noindent\textbf{2. 入口函数 cum 很高而 flat 很低，下一步应该做什么？}
\begin{enumerate}[label=\Alph*.]
\item 删除入口函数
\item 沿调用图查找下游热点及重复工作
\item 只优化入口处的变量赋值
\item 把 cum 与所有子函数 cum 相加
\end{enumerate}
\noindent{\color{forest}\textbf{答案 B。}}cum 说明该调用路径承载大量成本，flat 低说明执行成本主要在下游；cum 包含关系不能重复求和。

\noindent\textbf{3. 任务进入早已存在的 worker pool 后，CPU labels 不见了，最可能缺少什么？}
\begin{enumerate}[label=\Alph*.]
\item 给 context 增加更多任意字段
\item 在 worker 处理任务时调用 pprof.Do
\item 开启 heap 的 debug=2
\item 将所有路由改成唯一 UUID
\end{enumerate}
\noindent{\color{forest}\textbf{答案 B。}}传递 context 不会自动更新既有 goroutine 的 labels，需要在实际执行工作的位置安装并恢复标签。

\noindent\textbf{4. Go 1.25 的 pprof labels 原生用于哪些 profile？}
\begin{enumerate}[label=\Alph*.]
\item 所有 profile
\item 仅 heap 和 allocs
\item CPU 和 goroutine
\item 仅 block 和 mutex
\end{enumerate}
\noindent{\color{forest}\textbf{答案 C。}}固定版本的 runtime/pprof 文档明确限于 CPU 与 goroutine；其他采集器的扩展能力不能当成标准库承诺。

\section{分析题与参考答案}
\noindent\textbf{题 1：}假设优化前后两个观测窗口时长相同。优化前处理 6000 个请求，共有 30 CPU 秒，其中 Encode 为 12 CPU 秒；优化后处理 9000 个请求，共有 36 CPU 秒，Encode 为 9 CPU 秒。计算总 CPU/请求及 Encode CPU/请求，并评价只看总 CPU 的结论。

\noindent\textbf{参考答案：}优化前总成本为 30/6000=5 CPU-ms/请求，Encode 为 2 CPU-ms/请求；优化后总成本为 36/9000=4 CPU-ms/请求，Encode 为 1 CPU-ms/请求。单位请求总 CPU 下降 20\%，编码成本下降 50\%。总 CPU 上升是吞吐也上升后的结果。还需控制请求分布与成功请求的统计口径，排除输入和错误率变化等替代解释，并检查延迟是否满足目标，才能作因果归因。

\noindent\textbf{题 2：}一个 15 秒请求包含约 20 ms 编码、14 秒网络等待和约 980 ms 排队。CPU profile 能否说明 15 秒都耗在哪里？设计最小证据链。

\noindent\textbf{参考答案：}不能。CPU profile 主要能归因编码及其他实际执行，网络等待需要 trace 的 net 阻塞事件或请求级 span，排队需要 runnable 调度时延或应用队列时间。应先明确 980 ms 是调度队列还是业务队列，再在 trace task/region 中标注阶段，并用与采样窗口一致的请求计时验证总账。

\medskip\noindent\textbf{本章主要阅读：}\cite{tools-pprof}\cite{tools-http-pprof}\cite{tools-labels}\cite{tools-label-runtime}\cite{tools-pprof-cli}。
\chapter{内存、分配与垃圾回收}
\label{ch:06}
\noindent{\large\color{forest}把分配流量、存活堆和操作系统驻留内存分开解释}\par\medskip
\noindent\textbf{先修要求：}第 02–05 章；理解可达性、垃圾回收周期和 Go 切片与 map 的基本语义。

\noindent\textbf{完成本章后，应能够：}
\begin{itemize}
\item 正确选择 heap profile 的四种 sample type
\item 解释字节抽样、GC 可见性延迟和差分 profile 的含义
\item 区分逃逸、保留、分配抖动和 RSS 增长
\item 用 GC 指标与分配栈检验内存优化的因果链
\end{itemize}
\section{四种 sample type：存量与流量的坐标系}
\begin{figure}[H]
\centering
\begin{tikzpicture}[x=1cm,y=1cm]
\node[align=center,text width=7cm,font=\small\bfseries] at (3.55,2.25) {分配两批对象};
\node[draw=leaf!75!black,fill=leaf!13,rounded corners=3pt,text width=1.95cm,minimum height=1.12cm,align=center,font=\scriptsize] (p0-root) at (1.18,0.12) {业务根\\{\scriptsize 可达}};
\node[draw=leaf!75!black,fill=leaf!13,rounded corners=3pt,text width=1.95cm,minimum height=1.12cm,align=center,font=\scriptsize] (p0-temp) at (1.18,-1.52) {临时 1 MiB\\{\scriptsize 可达}};
\node[draw=leaf!75!black,fill=leaf!13,rounded corners=3pt,text width=1.95cm,minimum height=1.12cm,align=center,font=\scriptsize] (p0-keep) at (3.55,-0.70) {保留 1 MiB\\{\scriptsize 可达}};
\node[draw=leaf!75!black,fill=leaf!13,rounded corners=3pt,text width=1.95cm,minimum height=1.12cm,align=center,font=\scriptsize] (p0-alloc) at (5.92,0.12) {alloc\_space\\{\scriptsize 值：2}};
\node[draw=linegray!75!black,fill=linegray!13,rounded corners=3pt,text width=1.95cm,minimum height=1.12cm,align=center,font=\scriptsize] (p0-live) at (5.92,-1.52) {inuse\_space\\{\scriptsize 值：2}};
\draw[-{Stealth[length=2mm]},forest!70!black] (p0-root.east) -- (p0-keep.west);
\draw[-{Stealth[length=2mm]},forest!70!black] (p0-temp.north) -- (1.18,1.27) -- (5.92,1.27) -- (p0-alloc.north);
\draw[-{Stealth[length=2mm]},forest!70!black] (p0-keep.east) -- (p0-alloc.west);
\draw[-{Stealth[length=2mm]},forest!70!black] (p0-keep.east) -- (p0-live.west);
\node[align=center,text width=7cm,font=\small\bfseries] at (11.25,2.25) {后续 GC 与 profile 更新};
\node[draw=leaf!75!black,fill=leaf!13,rounded corners=3pt,text width=1.95cm,minimum height=1.12cm,align=center,font=\scriptsize] (p77-root) at (8.88,0.12) {业务根\\{\scriptsize 可达}};
\node[dashed,draw=linegray!75!black,fill=linegray!13,rounded corners=3pt,text width=1.95cm,minimum height=1.12cm,align=center,font=\scriptsize] (p77-temp) at (8.88,-1.52) {临时 1 MiB\\{\scriptsize 已回收}};
\node[draw=leaf!75!black,fill=leaf!13,rounded corners=3pt,text width=1.95cm,minimum height=1.12cm,align=center,font=\scriptsize] (p77-keep) at (11.25,-0.70) {保留 1 MiB\\{\scriptsize 可达}};
\node[draw=linegray!75!black,fill=linegray!13,rounded corners=3pt,text width=1.95cm,minimum height=1.12cm,align=center,font=\scriptsize] (p77-alloc) at (13.62,0.12) {alloc\_space\\{\scriptsize 值：2}};
\node[draw=leaf!75!black,fill=leaf!13,rounded corners=3pt,text width=1.95cm,minimum height=1.12cm,align=center,font=\scriptsize] (p77-live) at (13.62,-1.52) {inuse\_space\\{\scriptsize 值：1}};
\draw[-{Stealth[length=2mm]},forest!70!black] (p77-root.east) -- (p77-keep.west);
\draw[-{Stealth[length=2mm]},linegray,dashed] (p77-temp.north) -- (8.88,1.27) -- (13.62,1.27) -- (p77-alloc.north);
\draw[-{Stealth[length=2mm]},forest!70!black] (p77-keep.east) -- (p77-alloc.west);
\draw[-{Stealth[length=2mm]},forest!70!black] (p77-keep.east) -- (p77-live.west);
\end{tikzpicture}
\caption{累计分配与当前存活是两种账本}\label{fig:06-alloc-live}
\begin{minipage}{0.96\linewidth}\footnotesize 以 MiB 为教学计量，省略采样缩放和运行时自身分配。失去引用不立即改变 profile；GC 及其记账可见性完成后才展示下降，真实数据最多可滞后两个 GC 周期。

\textbf{读图顺序：}1. 分配两批对象\quad 2. 临时对象失去引用\quad 3. 后续 GC 与 profile 更新\quad 4. 再分配并回收一批。
\textbf{连线语义：}业务根\ensuremath{\rightarrow}\allowbreak{}保留 1 MiB：引用；临时 1 MiB\ensuremath{\rightarrow}\allowbreak{}alloc\_space：分配；保留 1 MiB\ensuremath{\rightarrow}\allowbreak{}alloc\_space：分配；保留 1 MiB\ensuremath{\rightarrow}\allowbreak{}inuse\_space：仍可达。
\textbf{机制依据：}\cite{tools-mprof}\cite{tools-pprof}。
\end{minipage}
\end{figure}
Go 的 \texttt{heap} 与 \texttt{allocs} 使用同一组底层内存采样记录，区别主要是默认展示的 sample type：heap 默认看 \texttt{inuse\_\allowbreak{}space}，allocs 默认看 \texttt{alloc\_\allowbreak{}space}。四种视角分别沿“存量或累计分配”与“对象数或字节数”展开。\cite{tools-pprof}

\begingroup\small\setlength{\tabcolsep}{4pt}
\begin{longtable}{>{\raggedright\arraybackslash}p{0.298\linewidth}>{\raggedright\arraybackslash}p{0.298\linewidth}>{\raggedright\arraybackslash}p{0.298\linewidth}}
\toprule
\textbf{sample type} & \textbf{问题} & \textbf{主要用途} \\
\midrule
\endfirsthead
\toprule
\textbf{sample type} & \textbf{问题} & \textbf{主要用途} \\
\midrule
\endhead
inuse\_space & 采样可见状态下仍在使用多少字节 & 存活堆、长期保留、大对象 \\
\addlinespace[3pt]
inuse\_objects & 仍在使用多少堆分配块 & 存活分配块数量、小块保留 \\
\addlinespace[3pt]
alloc\_space & 进程运行以来累计分配多少字节 & 分配流量、GC 压力 \\
\addlinespace[3pt]
alloc\_objects & 累计分配多少堆分配块 & 小块频繁分配、分配次数估计 \\
\addlinespace[3pt]
\bottomrule
\end{longtable}
\endgroup
这里的 objects 使用采样器的堆分配块口径，并不保证与语言层对象一一对应；tiny allocator 和采样外推的边界见下一节。

“inuse”不等于操作系统 RSS，“alloc”不等于内存泄漏。一个反复分配临时缓冲区、用完后不再持有它们的服务，alloc\_space 可以持续增长，而 GC 后的 inuse\_space 保持稳定；一个缓存每次只增加一点数据的服务，分配速率不高却可能长期保留更多对象。先观察趋势，再确定要用 profile 回答的问题。

\begin{codeblock}
go tool pprof -top -sample_index=inuse_space ./service heap.pprof
go tool pprof -top -sample_index=inuse_objects ./service heap.pprof
go tool pprof -top -sample_index=alloc_space ./service heap.pprof
go tool pprof -top -sample_index=alloc_objects ./service heap.pprof
\end{codeblock}

若同一栈在对象数视图很大、字节数视图很小，可能是大量细碎对象；反过来则可能是少量大缓冲区。平均大小可以辅助判断，但采样校正、不同分配尺寸和聚合栈意味着它不是精确类型布局。profile 的栈说明对象在哪里分配，并不直接给出谁正在持有它。

\section{字节采样与 GC 的可见性边界}
Go 1.25 默认 \texttt{runtime.\allowbreak{}MemProfile\allowbreak{}Rate=512*1024}，目标是每分配这么多字节平均采到一次分配，并非每 512 KiB 给整个堆拍照，也不是均匀抽取对象。大对象更容易被抽中；在近似泊松字节抽样模型下，尺寸 s 的对象被选中概率约为 \texttt{1-\allowbreak{}exp(-\allowbreak{}s/\allowbreak{}R)}，工具利用采样率校正权重。这些数字是外推估计量：Go 1.25 按桶内平均分配尺寸计算逆概率权重，最后把对象数和字节数转换为 int64 写入 profile。它并不输出精确的原始事件清单，也不是浮点对象计数；界面上的缩放单位不能改变这一点。\cite{tools-mprof}

\texttt{MemProfile\allowbreak{}Rate=1} 可记录采集器覆盖的每个堆分配块，适用于隔离实验；代价可能很高。它不包含栈上对象，tiny allocator 还会将多个小对象装入同一堆块，因此不等于逐一记录所有语言级对象。空间按分配槽大小记账，可包含内部碎片，而非精确等于请求的切片长度。官方要求若修改该值，应在程序生命周期尽早只改一次，因为分析工具假设整段记录使用恒定采样率。采完后随意改回默认值会令一份 profile 混合不同抽样机制。

heap profile 的对外语义以最近完成的 GC 为基准，避免近期尚未回收的垃圾压倒存活对象。底层 \texttt{runtime.\allowbreak{}MemProfile} 文档还提醒结果可能最多滞后两个 GC 周期：分配记录与释放记录需要协调发布，尤其释放随清扫发生，不能把这一实现细节当成即时对象清单。二者并不矛盾：它们共同强调 heap profile 不是读取瞬间精确堆快照。\cite{tools-pprof}\cite{tools-mprof}

\begin{codeblock}
curl --fail -o heap.pprof   'http://127.0.0.1:6060/debug/pprof/heap?gc=1'
\end{codeblock}

\texttt{gc=1} 先请求一次 GC，再取 profile，适合有意对齐 GC 状态的实验；它会扰动服务，不能无条件放进高频采集任务。\texttt{runtime.\allowbreak{}GC()} 也不承诺把全部闲置堆立刻返还给操作系统。没有发生 GC 时，heap profile 有报告已知分配的特殊规则，分析关闭 GC 的实验尤其要说明这一点。\cite{tools-http-pprof}

参数也不能机械组合：Go 1.25 的 heap handler 遇到 \texttt{seconds=N} 时先进入差分分支，不再执行后面的 \texttt{gc} 处理。因此，\texttt{heap?seconds=5\&gc=1} 中的 \texttt{gc=1} 不会触发该 handler 的强制 GC 分支；这不排除采集期间因其他原因发生 GC。若研究要求 GC 对齐，应分别取明确的 \texttt{heap?gc=1} 快照，再按所选 sample type 比较。

\section{差分分析与泄漏假设}
内存泄漏诊断应先问“在相同稳态工作下，经历 GC 后的存活集是否持续增长”。在固定工作批次前后，采集同一进程的两份 heap profile；确认期间未重启、采样率未改变，并比较相同 sample type、相同构建及相近 GC 状态。仅比较进程运行 1 分钟与 1 小时的 alloc\_space，通常只证明程序一直在工作。

\begin{codeblock}
curl --fail -o before.pprof   'http://127.0.0.1:6060/debug/pprof/heap?gc=1'
# 在另一个终端执行固定批次的代表性请求。
curl --fail -o after.pprof   'http://127.0.0.1:6060/debug/pprof/heap?gc=1'
go tool pprof -top -sample_index=alloc_space   -base=before.pprof ./service after.pprof
go tool pprof -top -sample_index=inuse_space   -diff_base=before.pprof ./service after.pprof
\end{codeblock}

第一种差分估计该批次新增的累计分配字节；若要计算分配速率，还须除以两次快照之间的时长。第二种差分关注存活字节的变化。存活差分可能为负，意味着相应分配栈的保留量减少，不能将它当作坏数据删掉。\texttt{-\allowbreak{}diff\_\allowbreak{}base} 以基线样本总量作为百分比分母；累计 profile 使用 \texttt{-\allowbreak{}base} 时通常以总量之差为分母。解释时保留绝对字节变化，并阅读本地工具帮助。\cite{tools-pprof-cli}

HTTP 的 \texttt{heap?seconds=N} 也是快照差分，并不是另一种“持续采集实时存活堆”的机制。\cite{tools-http-pprof}

若字符串/切片引用使一个很小的子片段保留整个底层大数组，分配栈常指向原始大缓冲区，持有者却是下游缓存。此时需要检查对象所有权、容器增长、生命周期和取消路径。常见模式包括无界 map、未退出的 goroutine 捕获请求数据、未释放的计时资源及错误的缓存淘汰。pprof 提供候选分配点，不能自动展示完整保留引用图。

证据强度依次为：监控发现趋势，profile 指出增长来源，代码找出保留路径，受控实验切断该路径，重复工作后存活集恢复有界。仅看到某个 map 分配较多，还不足以宣布泄漏。

泄漏实验还要控制系统是否已到达稳态。缓存从空状态装入常用键、连接池逐步建立连接、程序首次使用某种类型触发初始化，都会造成合理的单调增长。可以先完成预热，再施加多个相同工作批次，在每批请求结束后，确认相关后台任务已完成或正确取消，再观察 GC 后的存活集。如果增长逐渐趋平，倾向于有界缓存；如果增长与请求总数或某类错误次数保持稳定斜率，则更值得追查生命周期遗漏。

分配源和持有者的分离尤其重要。一个通用解码函数可能为所有请求分配对象，却只在某条异常分支被缓存永久引用。优化解码函数的对象大小只能减缓增长，修复异常分支的所有权才可能消除无界保留。

\section{逃逸分析、对象复用与成本转移}
\begin{figure}[H]
\centering
\begin{tikzpicture}[x=1cm,y=1cm]
\node[draw=linegray!75!black,fill=linegray!13,rounded corners=3pt,text width=3.33cm,minimum height=1.12cm,align=center,font=\small] (p0-objects) at (1.88,0.97) {Go 活对象};
\node[draw=linegray!75!black,fill=linegray!13,rounded corners=3pt,text width=3.33cm,minimum height=1.12cm,align=center,font=\small] (p0-heap) at (5.62,0.00) {堆页/\allowbreak{}栈等};
\node[draw=linegray!75!black,fill=linegray!13,rounded corners=3pt,text width=3.33cm,minimum height=1.12cm,align=center,font=\small] (p0-other) at (1.88,-0.97) {cgo/\allowbreak{}映射};
\node[draw=linegray!75!black,fill=linegray!13,rounded corners=3pt,text width=3.33cm,minimum height=1.12cm,align=center,font=\small] (p0-rss) at (9.38,0.00) {进程 RSS};
\node[draw=leaf!75!black,fill=leaf!13,rounded corners=3pt,text width=3.33cm,minimum height=1.12cm,align=center,font=\small] (p0-cgroup) at (13.12,0.00) {容器记账};
\draw[-{Stealth[length=2mm]},forest!70!black] (p0-objects.east) -- (p0-heap.west);
\draw[-{Stealth[length=2mm]},forest!70!black] (p0-heap.east) -- (p0-rss.west);
\draw[-{Stealth[length=2mm]},forest!70!black] (p0-other.north) -- (1.88,2.12) -- (9.38,2.12) -- (p0-rss.north);
\draw[-{Stealth[length=2mm]},forest!70!black] (p0-rss.east) -- (p0-cgroup.west);
\end{tikzpicture}
\caption{Go 堆、运行时内存、RSS 与容器使用量}\label{fig:06-memory-layers}
\begin{minipage}{0.96\linewidth}\footnotesize 边界示意而非严格集合包含图：共享页、RSS 与 cgroup 的计量规则可能不同，箭头不表示数值相等。

\textbf{读图顺序：}1. 观察对象\quad 2. 对象之外还有页和栈\quad 3. 加入进程其他内存\quad 4. 再检查容器边界。
\textbf{连线语义：}Go 活对象\ensuremath{\rightarrow}\allowbreak{}堆页/栈等：可达对象；堆页/栈等\ensuremath{\rightarrow}\allowbreak{}进程 RSS：驻留页；cgo/映射\ensuremath{\rightarrow}\allowbreak{}进程 RSS：进程其他内存；进程 RSS\ensuremath{\rightarrow}\allowbreak{}容器记账：不同统计边界。
\textbf{机制依据：}\cite{tools-gc-guide}\cite{tools-runtime-debug-api}。
\end{minipage}
\end{figure}
逃逸分析决定对象能否放在栈上，它是编译器证明结果，不是“用了指针就进堆”的语法规则。返回局部对象地址可能逃逸，也可能在内联后被消除；接口装箱、闭包捕获、容器保存及对象大小都可能影响分配。先使用 benchmark 的 \texttt{-\allowbreak{}benchmem} 判断每操作分配，再用编译器诊断定位原因。

\begin{codeblock}
go test -run='^$' -bench='BenchmarkEncode$' -benchmem -count=5 .
go build -gcflags='-m=2' .
\end{codeblock}

未写包匹配模式时，\texttt{-\allowbreak{}gcflags='-\allowbreak{}m=2'} 只作用于命令行直接列出的包；此处的 \texttt{.\allowbreak{}} 表示当前包。确有必要时再用 \texttt{all=-\allowbreak{}m=2} 扩展到依赖包。诊断中的“escapes to heap”解释潜在存放位置，profile 则反映真实运行路径和频率，两者缺一不可。热路径中减少一次分配，可能比冷路径完全零分配更有价值。

预分配容量可以减少增长和拷贝，但过度预分配会增加峰值与保留内存。复用 \texttt{[]byte} 可能减少 alloc\_space，却因为一个过大的容量长期留在池中而提高 inuse\_space。\texttt{sync.\allowbreak{}Pool} 是临时对象复用机制，元素可随时被运行时移除，不是具有容量和持久性保证的缓存；对象再次交给请求之前，必须确保残留状态已正确清理或覆盖，避免把上一请求的数据带入下一请求。

零拷贝也有权衡。复制一个小片段可能增加短期 alloc\_space，却释放巨大的底层数组，降低长期 inuse\_space。反过来，减少复制可能降低 CPU，却增加生命周期耦合。应把吞吐、延迟、分配字节、存活字节和代码复杂性一起列入决策。任何使用 unsafe 的“零分配”技巧还需要独立证明内存与并发安全，不能由性能收益替代正确性论证。

\section{GC 指标、内存限制与 RSS 的解释}
GC 成本受到分配速度、存活对象规模、指针密度、根集合及可用 CPU 的共同影响。其他条件相近时，较高的分配速率会更快耗尽本周期的可分配空间，从而提高 GC 频率；累计 alloc\_space 较大本身不能说明当前 GC 更频繁。可扫描的存活对象、指针及 goroutine 栈根增多，通常会增加标记工作；无指针的大缓冲区并不对应等比例的指针扫描量。CPU profile 中的 GC worker 与 mark assist、运行时 GC 指标、延迟分布需要一起观察，不能只看一次暂停时间。

\texttt{GOGC} 控制以存活堆和根集合为基础的增长目标，属于吞吐与内存之间的交换。官方模型中，目标堆近似为“上次存活堆 + (上次存活堆 + 根集合) \ensuremath{\times} GOGC/100”，实际目标还受最小堆、pacer 与内存限制等因素影响。增大 GOGC 往往用更多内存换取较低 GC 频率；减小则反向交换。\texttt{GOMEMLIMIT} 是 Go 运行时的软内存限制，不是容器 RSS 的绝对上限，也不能阻止 cgo 库或 mmap 使用额外内存。\cite{tools-gc-guide}

运行时可管理内存近似由 \texttt{runtime/\allowbreak{}metrics} 的 \texttt{/\allowbreak{}memory/\allowbreak{}classes/\allowbreak{}total:\allowbreak{}bytes} 减去 \texttt{/\allowbreak{}memory/\allowbreak{}classes/\allowbreak{}heap/\allowbreak{}released:\allowbreak{}bytes} 表征；RSS 还受非 Go 内存、共享映射、页驻留和内核统计口径影响。堆已回收的页可能仍由 Go 保留以供再次分配，因此 HeapAlloc 降低而 RSS 暂时不变并不自动表示泄漏。需要同时检查 heap in-use、heap free/released、goroutine 栈、cgo/mmap 和操作系统映射。\cite{tools-metrics}

诊断顺序可以是：先观察是否触及容器限制，再区分运行时管理内存与外部内存，然后判断是分配抖动还是存活增长，最后再调整 GOGC/内存限制。先调参数而不理解存活集，可能把正常分配变成持续 GC 抖动。

估算峰值时不能只取上次 GC 后存活堆。服务还需要容纳正在构建的新对象、尚未完成的并发请求、缓存扩容时新旧存储并存、运行时元数据和外部库分配。假设每个并发请求额外需要 10 MiB 工作集，额外并发 20 个请求就可能再增加约 200 MiB；这类模型虽然粗糙，却能指导负载阶梯实验和并发上限设计。

GC 优化报告应同时给出资源交换。减少对象数可能降低扫描和分配管理成本，压缩数据结构可能降低存活字节但增加解码 CPU；降低内存限制可能缓解容器风险，却放大回收频率。任何单一指标改善都需要放回容量约束和业务目标中评价。

\section{可运行实验：临时分配与长期保留}
以下是完整的隔离实验，保存为 \texttt{memory.\allowbreak{}go}。为突出采样语义，它在 main 开始设置 \texttt{MemProfile\allowbreak{}Rate=1}，因此只能作为教学实验，不能直接照搬到生产服务。\texttt{churn} 产生临时对象，\texttt{retained} 保存少量长期对象；具体字节总量仍以本机结果为准。

\begin{codeblock}
package main

import (
	"os"
	"runtime"
	"runtime/pprof"
)

var retained [][]byte
var transient []byte

//go:noinline
func churn() {
	for i := 0; i < 4000; i++ {
		transient = make([]byte, 16*1024)
		transient[0] = byte(i)
	}
	transient = nil
}

//go:noinline
func keep() {
	for i := 0; i < 64; i++ {
		retained = append(retained, make([]byte, 64*1024))
	}
}

func main() {
	runtime.MemProfileRate = 1
	churn()
	keep()
	runtime.GC()
	runtime.GC()
	f, err := os.Create("heap.pprof")
	if err != nil {
		panic(err)
	}
	if err := pprof.WriteHeapProfile(f); err != nil {
		f.Close()
		panic(err)
	}
	if err := f.Close(); err != nil {
		panic(err)
	}
	runtime.KeepAlive(retained)
}
\end{codeblock}

\begin{codeblock}
go build -o memory-lab memory.go
./memory-lab
go tool pprof -top -alloc_space ./memory-lab heap.pprof
go tool pprof -top -inuse_space ./memory-lab heap.pprof
\end{codeblock}

两次 GC 用于让隔离实验中的记录更充分发布，不是生产端点必须执行两次 GC 的规则。预期临时分配在 alloc\_space 中明显，而保留缓冲区在 inuse\_space 中明显。把 \texttt{retained=nil} 放在 GC 前再运行，应检验其存活量是否下降；把采样率恢复为默认值则应在新进程中完成，再比较估计噪声和采集开销。

\section{实验任务}
\subsection*{分配、保留与 GC 的三维对照}
\begin{enumerate}
\item 分别运行原始 memory.go、清空 retained 的版本、默认采样率的版本，每次使用独立进程，并在不同输出目录保存二进制和 heap.pprof。

\item 为每个版本生成四种 sample type 的 top 表，解释同一分配栈在不同视角的差异。

\item 在一个可重复的业务 benchmark 中记录 B/op 与 allocs/op，并选择一项复用或复制策略修改。

\item 比较修改前后分配量、稳态存活量、吞吐和延迟，避免仅以 allocs/op 为唯一成功标准。

\item 提交 GC 对齐方式与采样率，并明确哪些结果来自模型预期，哪些来自实际运行。

\end{enumerate}
\section{自测与解析}
\noindent\textbf{1. 服务 alloc\_space 随时间增长，但 GC 后 inuse\_space 稳定，最合理的第一判断是什么？}
\begin{enumerate}[label=\Alph*.]
\item 已经证明内存泄漏
\item 累计分配量增加，尚未证明长期保留增长
\item 操作系统必定 OOM
\item heap profile 不工作
\end{enumerate}
\noindent{\color{forest}\textbf{答案 B。}}alloc\_space 是累计分配量，正常工作的程序也会持续增长；泄漏假设需要稳态存活集与保留路径证据。

\noindent\textbf{2. 小切片导致大数组长期保留，heap profile 最直接显示什么？}
\begin{enumerate}[label=\Alph*.]
\item 所有持有该切片的对象关系
\item 原始数组的分配栈及估计存活量
\item 准确的内核 RSS
\item 每个对象的释放时间
\end{enumerate}
\noindent{\color{forest}\textbf{答案 B。}}标准 heap profile 按分配栈归因，不提供完整引用保留图，持有者需要通过代码和生命周期分析定位。

\noindent\textbf{3. 为什么不应在一个进程里反复切换 MemProfileRate？}
\begin{enumerate}[label=\Alph*.]
\item Go 语法禁止赋值
\item 工具假设内存 profile 使用恒定采样率
\item 它会自动停用 GC
\item 只有 CPU profile 支持该变量
\end{enumerate}
\noindent{\color{forest}\textbf{答案 B。}}混合采样率会破坏校正假设；官方建议在进程早期至多调整一次。

\noindent\textbf{4. GOMEMLIMIT 的正确理解是什么？}
\begin{enumerate}[label=\Alph*.]
\item 容器 RSS 的严格上限
\item 所有 cgo 与 mmap 分配总额的硬限制
\item Go 运行时管理内存的软限制
\item 使所有分配改用栈
\end{enumerate}
\noindent{\color{forest}\textbf{答案 C。}}运行时会通过 GC 与内存回收努力接近软限制，但不覆盖所有外部内存，也不是绝对 RSS 保证。

\section{分析题与参考答案}
\noindent\textbf{题 1：}假设两个 GC 对齐快照的 alloc\_space 分别为 1.2 GiB 和 2.0 GiB，inuse\_space 分别为 80 MiB 和 84 MiB；间隔处理 200000 个请求。估算每请求分配，并说明 4 MiB 增长是否证明泄漏。

\noindent\textbf{参考答案：}区间分配约 0.8 GiB，即约 8.59\ensuremath{\times}10\textasciicircum{}8 字节；除以 200000，得约 4295 字节/请求，也就是约 4.2 KiB/请求。存活增长 4 MiB 可能是缓存预热、并发差异或采样波动；需在稳态相同负载下多次 GC 对齐采样，检查持续斜率并验证保留路径。

\noindent\textbf{题 2：}某优化将每请求分配从 64 KiB 降到 8 KiB，但稳态堆从 200 MiB 升到 900 MiB，容器限制为 1 GiB。如何解释并重新设计？

\noindent\textbf{参考答案：}可能用无界对象池或保留大容量缓冲区换取了较低分配流量。先比较 inuse\_space 分配栈、池中容量分布和并发数，确认持有路径。可限制缓存容量、不把过大缓冲区放回对象池、分级复用并清除跨请求引用。成功标准应同时包括分配率、稳态和峰值内存、GC CPU、尾延迟及容器安全余量。

\medskip\noindent\textbf{本章主要阅读：}\cite{tools-pprof}\cite{tools-mprof}\cite{tools-http-pprof}\cite{tools-gc-guide}\cite{tools-metrics}\cite{tools-pprof-cli}。
\chapter{并发、阻塞与锁竞争}
\label{ch:07}
\noindent{\large\color{forest}从等待者、释放者和调度状态重建同步瓶颈}\par\medskip
\noindent\textbf{先修要求：}第 02–06 章；熟悉 channel、Mutex、WaitGroup、context 取消和 Go 内存模型基础。

\noindent\textbf{完成本章后，应能够：}
\begin{itemize}
\item 区分 block 与 mutex 的归因位置、采样机制和时间单位
\item 解释累计等待为何可以超过墙钟时间
\item 识别 Go 1.22 前后 mutex profile 的语义变化
\item 联合 goroutine 快照、trace 和受控实验验证并发瓶颈
\end{itemize}
\section{并发瓶颈的四种候选机制}
“goroutine 很多”“CPU 很低”“锁函数很热”都不是完整诊断。需要先区分四类问题：计算饱和、同步资源竞争、外部资源等待，以及工作已就绪却迟迟不能运行。它们可能具有相同的吞吐下降或尾延迟上升，但需要不同的证据。CPU profile 主要观察执行；block profile 观察特定 Go 同步操作阻塞；mutex profile 观察锁竞争的释放路径；trace 把状态变化放回时间轴。

\texttt{goroutine} 诊断提供当前调用栈的快照，可发现大量相似的 channel receive、锁等待或网络读；其中 \texttt{debug=2} 文本还会显示 goroutine 状态。快照中的数量不等于频率，也不等于各状态累计时间。某一时刻看到 1000 个 worker 在 channel receive，可能只是空闲池正常等任务；重复快照显示相同请求 goroutine 长期等待、数量持续增长，才更支持生命周期或资源泄漏假设。

\begin{codeblock}
curl --fail -o goroutines.txt   'http://127.0.0.1:6060/debug/pprof/goroutine?debug=2'
curl --fail -o goroutine.pprof   'http://127.0.0.1:6060/debug/pprof/goroutine'
go tool pprof -top ./service goroutine.pprof
\end{codeblock}

二进制 profile 适合聚合，文本快照适合阅读状态和调用链。采集时刻、请求负载和取消策略必须一起记录。若故障属于数据竞态，应在独立正确性实验中使用 race detector；race 检查会改变执行开销与调度，不宜用其吞吐代表正常构建。

\section{Block profile：等待发生在哪里}
\begin{figure}[H]
\centering
\begin{tikzpicture}[x=1cm,y=1cm]
\node[align=center,text width=7cm,font=\small\bfseries] at (3.55,3.70) {三个等待者};
\node[draw=amber!75!black,fill=amber!13,rounded corners=3pt,text width=3.13cm,minimum height=1.12cm,align=center,font=\scriptsize] (p0-mu) at (5.32,-0.70) {共享锁\\{\scriptsize 持有}};
\node[draw=amber!75!black,fill=amber!13,rounded corners=3pt,text width=3.13cm,minimum height=1.12cm,align=center,font=\scriptsize] (p0-g1) at (1.77,1.77) {G1\\{\scriptsize 持有}};
\node[draw=bluegray!75!black,fill=bluegray!13,rounded corners=3pt,text width=3.13cm,minimum height=1.12cm,align=center,font=\scriptsize] (p0-g2) at (1.77,0.12) {G2\\{\scriptsize 等待}\\{\scriptsize 值：10}};
\node[draw=bluegray!75!black,fill=bluegray!13,rounded corners=3pt,text width=3.13cm,minimum height=1.12cm,align=center,font=\scriptsize] (p0-g3) at (1.77,-1.53) {G3\\{\scriptsize 等待}\\{\scriptsize 值：10}};
\node[draw=bluegray!75!black,fill=bluegray!13,rounded corners=3pt,text width=3.13cm,minimum height=1.12cm,align=center,font=\scriptsize] (p0-g4) at (1.77,-3.17) {G4\\{\scriptsize 等待}\\{\scriptsize 值：10}};
\draw[-{Stealth[length=2mm]},forest!70!black] (p0-g1.east) -- (p0-mu.west);
\draw[-{Stealth[length=2mm]},forest!70!black] (p0-g2.east) -- (p0-mu.west);
\draw[-{Stealth[length=2mm]},forest!70!black] (p0-g3.east) -- (p0-mu.west);
\draw[-{Stealth[length=2mm]},forest!70!black] (p0-g4.east) -- (p0-mu.west);
\node[align=center,text width=7cm,font=\small\bfseries] at (11.25,3.70) {解锁并交接};
\node[draw=amber!75!black,fill=amber!13,rounded corners=3pt,text width=3.13cm,minimum height=1.12cm,align=center,font=\scriptsize] (p77-mu) at (13.02,-0.70) {共享锁\\{\scriptsize 持有}};
\node[draw=forest!75!black,fill=forest!13,rounded corners=3pt,text width=3.13cm,minimum height=1.12cm,align=center,font=\scriptsize] (p77-g1) at (9.47,1.77) {G1\\{\scriptsize 完成}};
\node[draw=amber!75!black,fill=amber!13,rounded corners=3pt,text width=3.13cm,minimum height=1.12cm,align=center,font=\scriptsize] (p77-g2) at (9.47,0.12) {G2\\{\scriptsize 持有}};
\node[draw=bluegray!75!black,fill=bluegray!13,rounded corners=3pt,text width=3.13cm,minimum height=1.12cm,align=center,font=\scriptsize] (p77-g3) at (9.47,-1.53) {G3\\{\scriptsize 等待}};
\node[draw=bluegray!75!black,fill=bluegray!13,rounded corners=3pt,text width=3.13cm,minimum height=1.12cm,align=center,font=\scriptsize] (p77-g4) at (9.47,-3.17) {G4\\{\scriptsize 等待}};
\draw[-{Stealth[length=2mm]},linegray,dashed] (p77-g1.east) -- (p77-mu.west);
\draw[-{Stealth[length=2mm]},forest!70!black] (p77-g2.east) -- (p77-mu.west);
\draw[-{Stealth[length=2mm]},forest!70!black] (p77-g3.east) -- (p77-mu.west);
\draw[-{Stealth[length=2mm]},forest!70!black] (p77-g4.east) -- (p77-mu.west);
\end{tikzpicture}
\caption{锁内 10 ms 怎样累积出多份等待}\label{fig:07-lock}
\begin{minipage}{0.96\linewidth}\footnotesize 教学排队模型，不保证锁交接顺序；真实 Go mutex 含自旋、饥饿模式和调度路径。

\textbf{读图顺序：}1. G1 进入临界区\quad 2. 三个等待者\quad 3. 解锁并交接\quad 4. 增加并发未扩大容量。
\textbf{机制依据：}\cite{tools-mprof}\cite{tools-go122}。
\end{minipage}
\end{figure}
Block profile 记录 goroutine 在特定同步原语上阻塞的时间，包括 Mutex/RWMutex、WaitGroup、Cond 以及 channel send/receive/select。栈归因到发生阻塞的位置，例如 \texttt{sync.\allowbreak{}Mutex.\allowbreak{}Lock}。它不等于所有 off-CPU 时间：网络 I/O、任意 syscall、睡眠和 runnable 排队不能统统期待从 block profile 中找到。\cite{tools-pprof}

\begin{codeblock}
// 集成片段：启动时从受控配置选择，默认不开启。
runtime.SetBlockProfileRate(1_000_000)
\end{codeblock}

这里的参数单位是纳秒，目标是平均每累计这么多阻塞时间采一个事件；它不是“每一百万次阻塞采一次”，也不是“每毫秒截取所有 goroutine”。Go 1.25 对持续时间短于 \texttt{rate} 阈值的事件按与时长成正比的概率抽样，对达到阈值的事件直接采集，再做逆概率权重校正；因此不同长度的事件并非相同概率。\texttt{rate=1} 记录所有符合条件的阻塞事件，\texttt{rate<=0} 关闭。全量记录适合小型隔离实验，不能假定开销恒定且可忽略。\cite{tools-mprof}

\begin{codeblock}
curl --fail --max-time 35 -o block.pprof   'http://127.0.0.1:6060/debug/pprof/block?seconds=30'
go tool pprof -top -sample_index=delay ./service block.pprof
go tool pprof -top -sample_index=contentions ./service block.pprof
\end{codeblock}

标准记录在阻塞结束等记录点上更新。一个一直没有解除的等待可能尚未进入完整的阻塞统计；此时快照和 trace 比等待 profile 更重要。\texttt{delay} 是累计估计等待量，\texttt{contentions} 是按采样机制校正的事件量，二者的比值可作粗略平均等待，但不能恢复分布、p99 或最慢请求。

HTTP \texttt{block?seconds=N} 是两个累计快照之差，不会把每个等待裁剪为与窗口的交集。某个等待在窗口开始前已持续 50 秒、窗口内又等 1 秒才完成，其本次新增记录可能包含约 51 秒等待。反之，到窗口结束仍未完成的等待可能尚未发布；因此短窗口的 delay/T 只能在边界效应可忽略时近似平均等待人数。

\section{Mutex profile：谁结束了造成竞争的临界区}
\begin{figure}[H]
\centering
\begin{tikzpicture}[x=1cm,y=1cm]
\node[anchor=east,font=\small] at (1,0.0) {泳道1};
\draw[linegray] (1.2,-0.35) -- (14.799999999999999,-0.35);
\node[anchor=east,font=\small] at (1,-1.7) {泳道2};
\draw[linegray] (1.2,-2.05) -- (14.799999999999999,-2.05);
\path[fill=amber!60] (1.200,-0.250) rectangle (8.618,0.250);
\node[align=center,font=\small,text width=7.32cm] at (4.909,0.750) {G1 持锁};
\path[fill=leaf!60] (8.618,-0.250) rectangle (9.360,0.250);
\draw[linegray] (8.989,0.250) -- (8.989,0.680);
\node[circle,fill=mist,inner sep=2pt,font=\scriptsize] at (8.989,0.850) {2};
\path[fill=amber!60] (2.436,-1.950) rectangle (8.618,-1.450);
\node[align=center,font=\small,text width=6.08cm] at (5.527,-0.950) {G2 阻塞};
\path[fill=forest!60] (8.618,-1.950) rectangle (11.091,-1.450);
\node[align=center,font=\small,text width=2.37cm] at (9.855,-0.950) {G2 可运行};
\path[fill=leaf!60] (11.091,-1.950) rectangle (14.800,-1.450);
\node[align=center,font=\small,text width=3.61cm] at (12.945,-0.950) {G2 实际运行};
\draw[-{Stealth[length=2mm]},forest] (1.2,-2.8499999999999996) -- (14.999999999999998,-2.8499999999999996);
\draw[forest] (1.20,-2.8499999999999996) -- (1.20,-2.9499999999999997);
\node[below,font=\scriptsize] at (1.20,-2.9499999999999997) {0};
\draw[forest] (3.92,-2.8499999999999996) -- (3.92,-2.9499999999999997);
\node[below,font=\scriptsize] at (3.92,-2.9499999999999997) {11};
\draw[forest] (6.64,-2.8499999999999996) -- (6.64,-2.9499999999999997);
\node[below,font=\scriptsize] at (6.64,-2.9499999999999997) {22};
\draw[forest] (9.36,-2.8499999999999996) -- (9.36,-2.9499999999999997);
\node[below,font=\scriptsize] at (9.36,-2.9499999999999997) {33};
\draw[forest] (12.08,-2.8499999999999996) -- (12.08,-2.9499999999999997);
\node[below,font=\scriptsize] at (12.08,-2.9499999999999997) {44};
\draw[forest] (14.80,-2.8499999999999996) -- (14.80,-2.9499999999999997);
\node[below,font=\scriptsize] at (14.80,-2.9499999999999997) {55};
\node[anchor=east,font=\scriptsize] at (14.799999999999999,-3.5999999999999996) {时间 / ms};
\node[align=left,anchor=north west,text width=14cm,font=\scriptsize] at (1,-3.8499999999999996) {2=G1 解锁};
\end{tikzpicture}
\caption{同一次争用：阻塞栈与解锁栈}\label{fig:07-attribution}
\begin{minipage}{0.96\linewidth}\footnotesize 教学示意：用于解释机制，不是运行时的实测数据或完整实现。

\textbf{读图顺序：}1. 等待发生\quad 2. 持有者结束临界区\quad 3. 已就绪但尚未执行\quad 4. 开始执行。
\textbf{机制依据：}\cite{tools-pprof}\cite{tools-trace-doc}。
\end{minipage}
\end{figure}
Mutex profile 的栈通常落在释放锁的调用路径，例如 \texttt{sync.\allowbreak{}Mutex.\allowbreak{}Unlock}。它将竞争成本归到结束临界区的一侧，帮助定位持锁太久或过于集中的共享资源。其时间值是其他 goroutine 等待该锁的近似累计时间，不是 Unlock 指令的 CPU 耗时，也不是持锁者单次占用锁的准确时长。五个等待者同时等 1 秒，相关释放栈可以记录约 5 秒竞争。\cite{tools-pprof}

\begin{codeblock}
// 集成片段：平均抽取 1/10 的竞争事件。
previous := runtime.SetMutexProfileFraction(10)
// 完成诊断窗口后恢复 previous。
_ = previous
\end{codeblock}

\texttt{rate=1} 采集所有符合条件的竞争事件，\texttt{rate=0} 关闭，\texttt{rate<0} 只读取当前设置。无竞争的锁获取并不是 mutex contention 事件。采样按事件抽取，并按比例扩大计数和等待权重，这与 block 的时长抽样不同。运行时内部锁也可出现；无法完整归因的事件可能显示为 \texttt{\_\allowbreak{}LostConten\allowbreak{}dedRuntime\allowbreak{}Lock}，不能仅凭该符号把责任归给某个业务 Mutex。\cite{tools-mprof}

Go 1.22 是重要分界：官方发行说明指出 mutex profile 开始按等待 goroutine 的数量缩放竞争量，并加入运行时内部锁竞争。例子中 100 个 goroutine 各等 10 ms，新语义记录约 1 秒，而此前可能只表现为 10 ms。因此跨 Go 1.21 与 1.22 的原始 profile 总 delay 不能直接当作性能退化证据。不同版本对内部锁栈的完整程度还会继续演进，应记录运行时版本并读其发布说明。\cite{tools-go122}

同一个 Mutex 在不同路径中解锁时，竞争量可能分散到多个栈；多个锁实例具有相同栈时也会聚合到一起。Mutex profile 并不天然提供每个锁地址的精细账本。需结合业务分片键、临界区实现和 trace 确认资源身份。

\section{将等待量转化为队列与吞吐解释}
设观察窗口为 T，多个 goroutine 累计同步等待时间为 W，则 \texttt{W/\allowbreak{}T} 可近似解释为该窗口平均同时等待的 goroutine 数，前提是采样覆盖和窗口边界足够合理。它不是 CPU 利用率，也不是每请求延迟。假设 10 秒内产生 80 goroutine 秒等待，可解释为平均约 8 个等待者；即使有 16 核空闲，它也可能发生在一个不可并行的共享锁上。

把临界区看成串行服务站，可用排队模型构造假设：若每请求平均持锁 s 秒，该锁串行服务能力上限大致为 \texttt{1/\allowbreak{}s}，到达率逼近上限时排队往往迅速增加。这个模型忽略批处理、公平性、锁交接、缓存与分布，不能直接预测 p99，但能解释为什么再加 goroutine 不一定提高吞吐。

常见有效改变包括缩短临界区、把 I/O 或昂贵格式化移出锁内、按独立数据分片、降低重复获取频率、使用不可变快照，以及通过所有权把共享状态移交给单一执行者。把 Mutex 改成 channel 可能只是换一种排队位置；把有界队列改成无界队列可能提高短期接收速度，却把问题转化为内存与尾延迟。

每个方案都应说明不变量。例如将读写锁改成快照，要证明读者观察的版本一致性；把批量更新移出锁内，要证明顺序与原子性仍满足业务要求。优化必须首先保持正确性。第 13 章将把这类证据链扩展到完整服务。

平均等待并发的解释来自时间积分：每个等待者在时间轴上贡献一条等待区间，将这些区间与观测窗口交集的长度相加，相当于对窗口内的同时等待人数进行积分。除以窗口长度后得到平均人数。若一个等待跨越窗口边界，或统计在事件结束时才发布，短窗口会出现偏差；采样校正也会引入方差。因此应同时报告窗口和绝对等待量，避免把瞬时人数与长期平均混为一谈。

锁的公平性与吞吐之间也存在取舍。高负载下可能发生频繁交接，缓存行在核之间迁移，或者一个请求多次抢不到锁。即便平均 delay 改善，少数任务仍可能饥饿；需要查看等待分布与任务完成性，而不只是总量。

\section{联合证据与诊断陷阱}
Block 与 mutex 往往描述同一等待的不同视角，不能简单相加。前者可能显示 \texttt{cache.\allowbreak{}Get} 的 Lock 等待，后者显示 \texttt{cache.\allowbreak{}Refresh} 的 Unlock 路径；两者共同支持“刷新临界区阻塞读取”的假设。此时应检查刷新是否持锁执行序列化或远程调用，并在 trace 中观察读者何时被唤醒。

仅在栈快照中看到大量 goroutine 处于等待，并不能认定程序死锁：服务空闲时，主 goroutine 或网络轮询器可能正在正常等待。若运行时实际输出 \texttt{fatal error:\allowbreak{} all goroutines are asleep -\allowbreak{} deadlock!}，则已触发全局死锁检测并终止进程，不能再用“服务正常空闲”解释这条错误。反过来，未触发该检测也不保证没有死锁；涉及 cgo、外部系统或局部子系统的逻辑死锁可能不在其覆盖范围内。若只是某类请求永久等待，需检查 channel 关闭、WaitGroup 计数、取消传播和生产者是否仍可能运行。

另一种误判是将 runnable 队列等待归因到 Mutex。goroutine 已具备继续执行的条件，却因 P 不足、CPU 配额节流或其他 goroutine 长时间运行而尚未得到调度时，主要证据是 trace 的 ready-to-running 间隔、\texttt{/\allowbreak{}sched/\allowbreak{}latencies:\allowbreak{}seconds} 与操作系统调度信息。它们对应的层次不同，互相印证比互相替代更可靠。\cite{tools-metrics}

采集开销会改变高争用系统的调度。先用低成本指标定位异常窗口，再短时开启高精度 profile；对比开启和关闭时吞吐、p99 与 CPU。若采集本身让争用恶化，必须下调采样率或缩小实验范围。完成记录后恢复设置，并保存采样参数以便解释权重。

对于工作池，还应分开测量提交、排队、执行和交付结果四个阶段。队列满时生产者被阻塞，有时是正常背压；无限排队则让生产者看似迅速完成，但端到端延迟和内存继续增长。是否应阻塞、拒绝、限流或取消，应由服务契约决定，不能仅以消除 block profile 热点为目标。

检查取消路径时，为每个可能长期阻塞的操作说明谁负责终止它。发送结果时调用方已经超时退出，发送者若无人接收且没有取消分支，就可能永久挂起；这种 goroutine 还可能持有请求缓冲区，形成第 06 章的内存保留问题。同步剖析与内存剖析之间因此存在具体的因果联系。

用压力实验发现瓶颈后，还要用较低负载验证修复没有引入空转。把阻塞等待改成反复轮询可以让 block 数值下降，却消耗更多 CPU。正确优化应减少等待原因或控制需求，而不是让等待从一种图表中消失。

排查报告应保留锁保护的数据不变量，以及恢复采样配置的实际时间。

\section{可运行实验：把同一竞争看成等待与释放}
下面是完整程序，保存为 \texttt{locks.\allowbreak{}go}。故意在锁内睡眠，仅用于制造容易观察的同步竞争。四个 worker 共享 Mutex；WaitGroup 自身也可能在 block profile 中出现，这反映主 goroutine 等待工作完成的真实同步事件；启动用的 \texttt{start} channel 等待也可能被记录。分析时应区分这些协调开销与共享 Mutex 的竞争。

\begin{codeblock}
package main

import (
	"os"
	"runtime"
	"runtime/pprof"
	"sync"
	"time"
)

func writeProfile(name string) {
	f, err := os.Create(name + ".pprof")
	if err != nil {
		panic(err)
	}
	if err := pprof.Lookup(name).WriteTo(f, 0); err != nil {
		f.Close()
		panic(err)
	}
	if err := f.Close(); err != nil {
		panic(err)
	}
}

func main() {
	runtime.SetBlockProfileRate(1)
	runtime.SetMutexProfileFraction(1)
	var mu sync.Mutex
	var wg sync.WaitGroup
	start := make(chan struct{})
	for n := 0; n < 4; n++ {
		wg.Add(1)
		go func() {
			defer wg.Done()
			<-start
			for i := 0; i < 100; i++ {
				mu.Lock()
				time.Sleep(500 * time.Microsecond)
				mu.Unlock()
			}
		}()
	}
	close(start)
	wg.Wait()
	runtime.SetBlockProfileRate(0)
	runtime.SetMutexProfileFraction(0)
	writeProfile("block")
	writeProfile("mutex")
}
\end{codeblock}

\begin{codeblock}
go build -o locks-lab locks.go
./locks-lab
go tool pprof -top ./locks-lab block.pprof
go tool pprof -top ./locks-lab mutex.pprof
\end{codeblock}

先核对调用栈归因，再把 Sleep 移到 Unlock 后面。比较总完成时间和两个 profile 的变化，并说明语义是否仍可代表原工作：本实验没有受保护的真实状态，所以移动是安全的；真实程序不能机械地把业务操作移出锁。实际 sleep 时长与调度粒度相关，不预设输出恰好是 200 ms。

\section{实验任务}
\subsection*{同步竞争的四种证据}
\begin{enumerate}
\item 运行 7.6 实验，保存 Go 版本和原始 block/mutex profile，并使用外部计时工具记录墙钟完成时间；程序本身未打印耗时。各组产物放在独立目录，避免覆盖。

\item 分别将 worker 数改为 1、4、16；原程序每个 worker 执行 100 次，因此总操作数分别为 100、400、1600。按总完成操作数除以墙钟时间计算吞吐，再比较每操作等待量；若研究固定总工作量，应把同一总迭代数分配给不同数量的 worker。

\item 将 Sleep 移出临界区，用相同参数重新采集，解释 Mutex 与 WaitGroup 在 block 中的不同角色。

\item 为原实验增加一个短 trace，比较等待、唤醒与重新运行的间隔；不要把 block 与 mutex delay 相加。

\item 写出一个保留正确性的真实锁优化方案，列出不变量、反例与复测指标。

\end{enumerate}
\section{自测与解析}
\noindent\textbf{1. mutex profile 在 Unlock 栈显示 6 秒 delay，说明什么？}
\begin{enumerate}[label=\Alph*.]
\item Unlock 指令消耗 6 CPU 秒
\item 该释放路径被归因了约 6 goroutine 秒竞争等待
\item 单次持锁必定持续 6 秒
\item CPU profile 一定丢失了 6 秒
\end{enumerate}
\noindent{\color{forest}\textbf{答案 B。}}mutex delay 是竞争等待的近似累计归因，可聚合多个等待者与多个事件，既不是指令时间也不是单次持锁时间。

\noindent\textbf{2. SetBlockProfileRate(1000000) 的参数主要控制什么？}
\begin{enumerate}[label=\Alph*.]
\item 每百万次阻塞采一次
\item 每秒最多一个样本
\item 按阻塞时间抽样，目标每累计 1 ms 阻塞时间平均一个事件
\item 每 1 ms 全部 goroutine 的快照
\end{enumerate}
\noindent{\color{forest}\textbf{答案 C。}}单位为纳秒，抽样与事件持续时间相关，并非等概率按事件序号取样。

\noindent\textbf{3. 为什么 Go 1.21 与 1.22 的 mutex 总 delay 不宜直接比较？}
\begin{enumerate}[label=\Alph*.]
\item 文件扩展名不同
\item 1.22 改为考虑等待 goroutine 数，并纳入内部锁竞争
\item 1.22 完全取消采样
\item 1.21 只运行单线程
\end{enumerate}
\noindent{\color{forest}\textbf{答案 B。}}计量与覆盖语义变化会改变总量，即使应用逻辑完全相同也可能出现明显差别。

\noindent\textbf{4. 长期卡住但尚未解除的 channel 等待，首选补充什么证据？}
\begin{enumerate}[label=\Alph*.]
\item 只增加 alloc\_space 采样
\item goroutine 栈快照与执行 trace
\item 把 block 的现有计数视为所有等待
\item 仅查看函数汇编
\end{enumerate}
\noindent{\color{forest}\textbf{答案 B。}}尚未完成的等待可能还未发布完整阻塞事件；快照提供当前位置，trace 提供状态变化与唤醒关系。

\section{分析题与参考答案}
\noindent\textbf{题 1：}假设同一进程在 20 秒窗口前后的快照差分显示，某锁栈的 delay 增加了 100 goroutine 秒；窗口内完成 10000 次请求。估算平均等待并发与每请求归因等待，并说明限制。

\noindent\textbf{参考答案：}若采样覆盖充分、跨窗口记账偏差可忽略，平均等待并发约为 100/20=5 个 goroutine；按请求归一化的归因等待量为 100/10000=10 goroutine-ms/请求。后者不是请求 p99，也不一定是每个请求真正等待 10 ms；可能少数请求等待很多、一个请求创建多个等待者，或栈同时服务后台任务。尤其要检查是否包含窗口前已经开始的等待，以及窗口结束时尚未记账的等待；若这些影响显著，不能把 delay/T 直接解释为窗口内平均等待人数。

\noindent\textbf{题 2：}CPU 仅使用 1 核、机器有 16 核、mutex delay 很高。增加 worker 后吞吐不变而延迟变差，给出模型与验证方案。

\noindent\textbf{参考答案：}先核对 GOMAXPROCS、亲和性与 CPU 配额，避免把机器的 16 个核都当作进程可用容量。一个共享锁可能将关键工作串行化，其理想服务率上限近似为 1/每请求平均持锁时间；增加 worker 只扩大队列。验证释放栈中的临界区，观察 trace 中被阻塞的等待者和唤醒节奏，并在保持不变量的前提下缩短或分片临界区。若争用下降且相同到达率的尾延迟改善，才支持该解释。

\medskip\noindent\textbf{本章主要阅读：}\cite{tools-pprof}\cite{tools-mprof}\cite{tools-go122}\cite{tools-metrics}。
\chapter{执行追踪与时序分析}
\label{ch:08}
\noindent{\large\color{forest}还原 goroutine 状态转换、关键路径和罕见事件之前的现场}\par\medskip
\noindent\textbf{先修要求：}第 02、05、07 章；理解 goroutine、线程、P、GC 与同步机制。

\noindent\textbf{完成本章后，应能够：}
\begin{itemize}
\item 用状态转换区分同步等待、网络等待和调度延迟
\item 采集与分析 Go 1.22+ trace，并理解代际分段设计
\item 用 task 关联跨 goroutine 工作，并在各 goroutine 内正确使用 region 与 log
\item 在 Go 1.25 使用 FlightRecorder，并解释时间与容量的非硬保证
\end{itemize}
\section{从聚合调用栈进入状态时间轴}
Profile 将相似栈汇总，trace 保留事件的相对顺序及状态转换。两段请求可以具有相近 CPU profile，却因一个先排队再计算、另一个立刻计算而产生不同延迟。运行时 trace 记录 goroutine 创建、运行、阻塞、解除阻塞、系统调用、GC 和用户标注等事件，让分析者定位“资源什么时候可用，工作什么时候真正继续”。它不是每条机器指令的记录，也不包含所有业务语义。\cite{tools-trace-doc}

分析一条 goroutine 生命周期时，先区分 running、runnable、waiting 等状态。等待 channel 时它还不能执行；被发送者唤醒之后变成 runnable，直到某个 P 调度它才重新 running。这段 ready-to-running 时间属于调度延迟。若将二者合并成“锁慢”，会把 CPU 供给不足错误归因到同步实现。

运行时中的“正在运行”也不完全等于操作系统一直让该线程占用 CPU。线程可能受到内核调度或容器配额影响；Go trace 与 OS 调度信息的结合见第 10 章。跨进程服务延迟需要分布式 tracing 提供 RPC 边界，不能指望一份进程内 trace 自动还原全部网络链路。

阅读顺序应先从总体 CPU/GC/线程视图识别异常时段，再选择相关 task 或 goroutine，最后查看阻塞栈与唤醒关系。直接放大整份 trace 寻找“红色区域”容易失去分母，也容易把正常等待误认成故障。

\section{采集和 Go 1.22 重写的意义}
\begin{figure}[H]
\centering
\begin{tikzpicture}[x=1cm,y=1cm]
\node[anchor=east,font=\small] at (1,0.0) {泳道1};
\draw[linegray] (1.2,-0.35) -- (14.799999999999999,-0.35);
\node[anchor=east,font=\small] at (1,-1.7) {泳道2};
\draw[linegray] (1.2,-2.05) -- (14.799999999999999,-2.05);
\path[fill=leaf!60] (1.200,-0.250) rectangle (2.632,0.250);
\draw[linegray] (1.916,0.250) -- (1.916,0.680);
\node[circle,fill=mist,inner sep=2pt,font=\scriptsize] at (1.916,0.850) {1};
\path[fill=bluegray!60] (2.632,-0.250) rectangle (12.653,0.250);
\node[align=center,font=\small,text width=9.92cm] at (7.642,0.750) {父 G 等结果};
\path[fill=bluegray!60] (2.632,-1.950) rectangle (11.221,-1.450);
\node[align=center,font=\small,text width=8.49cm] at (6.926,-0.950) {子 G 后端区间};
\path[fill=leaf!60] (11.221,-1.950) rectangle (12.653,-1.450);
\draw[linegray] (11.937,-1.450) -- (11.937,-1.020);
\node[circle,fill=mist,inner sep=2pt,font=\scriptsize] at (11.937,-0.850) {4};
\path[fill=leaf!60] (12.653,-0.250) rectangle (14.800,0.250);
\draw[linegray] (13.726,0.250) -- (13.726,0.680);
\node[circle,fill=mist,inner sep=2pt,font=\scriptsize] at (13.726,0.850) {5};
\draw[-{Stealth[length=2mm]},forest] (1.2,-2.8499999999999996) -- (14.999999999999998,-2.8499999999999996);
\draw[forest] (1.20,-2.8499999999999996) -- (1.20,-2.9499999999999997);
\node[below,font=\scriptsize] at (1.20,-2.9499999999999997) {0};
\draw[forest] (3.92,-2.8499999999999996) -- (3.92,-2.9499999999999997);
\node[below,font=\scriptsize] at (3.92,-2.9499999999999997) {19};
\draw[forest] (6.64,-2.8499999999999996) -- (6.64,-2.9499999999999997);
\node[below,font=\scriptsize] at (6.64,-2.9499999999999997) {38};
\draw[forest] (9.36,-2.8499999999999996) -- (9.36,-2.9499999999999997);
\node[below,font=\scriptsize] at (9.36,-2.9499999999999997) {57};
\draw[forest] (12.08,-2.8499999999999996) -- (12.08,-2.9499999999999997);
\node[below,font=\scriptsize] at (12.08,-2.9499999999999997) {76};
\draw[forest] (14.80,-2.8499999999999996) -- (14.80,-2.9499999999999997);
\node[below,font=\scriptsize] at (14.80,-2.9499999999999997) {95};
\node[anchor=east,font=\scriptsize] at (14.799999999999999,-3.5999999999999996) {时间 / ms};
\node[align=left,anchor=north west,text width=14cm,font=\scriptsize] at (1,-3.8499999999999996) {1=父 G 发起；4=子 G 发送结果；5=父 G 完成};
\end{tikzpicture}
\caption{一条请求的 task 可以跨 goroutine}\label{fig:08-trace}
\begin{minipage}{0.96\linewidth}\footnotesize 教学示意：用于解释机制，不是运行时的实测数据或完整实现。

\textbf{读图顺序：}1. 建立 task\quad 2. 子 goroutine 开始 region\quad 3. 发送和唤醒\quad 4. 请求结束。
\textbf{机制依据：}\cite{tools-trace-annotation}。
\end{minipage}
\end{figure}
最简单的采集方式是使用测试命令或诊断端点。采集器必须在问题发生前开启；故障结束后才启动 trace，通常会缺少关键因果事件。下面的测试命令假定当前包确有 TestPipeline；若使用其他测试，应替换正则并确认目标测试实际执行。两个查看器分别打开不同文件，前一个会持续占用 8082 端口；切换前应退出前一个进程，或为后一个选择其他端口。

\begin{codeblock}
go test -run='TestPipeline$' -trace=trace.out .
go tool trace -http=127.0.0.1:8082 trace.out
curl --fail --max-time 10 -o service.trace   'http://127.0.0.1:6060/debug/pprof/trace?seconds=5'
go tool trace -http=127.0.0.1:8082 service.trace
\end{codeblock}

Go 1.21 的栈回溯优化显著降低了常态采集 CPU 开销；Go 1.22 继续重写执行追踪器，常称为 trace v2。其关键改变包括定期分成可处理的代际片段、改进启动停止延迟、补全系统调用持续时间，并记录 goroutine 所在线程。多数平台采用操作系统时钟，便于与更底层的同主机信息关联；这并不把所有时间戳直接变成可跨主机比较的 UTC。\cite{tools-go122}\cite{tools-trace2024}

分段意味着工具不必把整个历史看成一个不可拆分、无限增长的结构，也为保留最近窗口和流式处理奠定基础。需要区分“格式支持分段处理”与“特定版本查看器已经做到恒定内存”：2024 年官方文章明确当时 \texttt{go tool trace} 仍加载完整 trace。教材不将该历史实现状态推广到所有后续版本，也不承诺任意大小文件均可无代价打开。

尽量使用与生成版本相同或明确兼容的工具链打开 trace，保留采集端 Go 版本。正常 \texttt{trace.\allowbreak{}Start} 同时只能有一个消费者；HTTP trace 与程序内普通 trace 可能发生冲突。读取开始采集的错误，保证 \texttt{trace.\allowbreak{}Stop} 在文件关闭之前完成。\cite{tools-trace-doc}

\section{Task、region、log 的语义与生命周期}
Task 描述一个用户定义的逻辑操作，可跨越多个 goroutine；\texttt{trace.\allowbreak{}NewTask} 返回携带任务信息的新 context 与 Task，\texttt{Task.\allowbreak{}End} 标记逻辑完成。Region 描述同一 goroutine 内的一段活动，开始和结束必须保持正确的 goroutine 与嵌套关系。Log 是带类别的瞬时事件。它们补足“当前业务正在做什么”，而不是替代运行时自动记录的调度事件。\cite{tools-trace-annotation}

以下为集成片段，\texttt{encode} 与 \texttt{send} 表示业务函数：

\begin{codeblock}
ctx, task := trace.NewTask(ctx, "request")
defer task.End()
trace.WithRegion(ctx, "encode", func() { encode() })
trace.Log(ctx, "cache", "miss")
done := make(chan struct{})
go func() {
    defer close(done)
    trace.WithRegion(ctx, "send", func() { send() })
}()
<-done
\end{codeblock}

这里 task 等待子 goroutine 完成才结束；若在创建子 goroutine 后立即返回，task 的生命周期就可能错误地短于实际工作。一个 Region 不能在父 goroutine Start 后交给子 goroutine End；应由子 goroutine 自己创建 region，并通过 context 关联同一 task。

task type、region type 和日志类别应是有限集合；具体请求参数不应放进类型字符串。日志消息也应控制长度并去掉敏感数据。\texttt{trace.\allowbreak{}Logf} 在未启用 trace 时会避免其内部格式化，但 Go 会先求值调用实参，因此昂贵的 \texttt{buildSummary()} 仍可能执行；必要时在外层使用 \texttt{trace.\allowbreak{}IsEnabled()}。注解本身也有开销，特别是每个极小循环迭代创建任务，可能改变分配与时序。

\section{关键路径、时序证据与导出 profile}
请求延迟由依赖图中的关键路径决定，不是所有 goroutine 工作时间之和。一个父请求并行执行两个 100 ms 子任务，如果两者完全重叠，完成延迟接近 100 ms 加协调成本，而累计工作可以接近 200 ms。trace 的价值在于查看重叠、依赖和何时唤醒；只按函数累计耗时排序会漏掉这些结构。

实用分析步骤是：选中异常 task；检查持续时间较长的 region，并结合依赖确认它是否处于关键路径；区分它内部的 running、同步/网络等待和 runnable 排队；追到解除阻塞的 goroutine；检查唤醒之前发生了什么；最后把假设映射回业务依赖。例如消费者早已等待，生产者却在持锁写日志，就应优先验证生产者关键路径，而非优化消费者等待函数。

\begin{codeblock}
go tool trace -pprof=sched service.trace > sched.pprof
go tool trace -pprof=sync service.trace > sync.pprof
go tool trace -pprof=net service.trace > net.pprof
go tool trace -pprof=syscall service.trace > syscall.pprof
go tool pprof -top sched.pprof
\end{codeblock}

这些文件从 trace 事件推导出按调用栈聚合的事件计数与累计延迟，与运行时直接采集的 CPU/block/mutex profile 不是同一数据来源。它们不保留每个事件的完整延迟分布，不能仅凭导出结果恢复事件延迟的 p99。\texttt{sched} 关注调度延迟，\texttt{sync} 关注同步阻塞，\texttt{net} 关注网络阻塞，\texttt{syscall} 关注系统调用阻塞。导出后时间顺序被聚合掉，宜用于定位候选栈，再回到时间轴验证。\cite{tools-trace-cmd}

观察相关性还不是因果证明。GC 恰好与慢请求重叠，不等于 GC 导致全部延迟；检查请求是否确实暂停、是否有 assist、是否同时存在队列或下游慢调用。受控改变一个机制并重现时序变化，才形成更强的因果证据。

任务分析还需区分业务完成与资源清理。请求返回客户端之后，后台写日志、清理缓冲区或异步复制可能继续执行；如果这些工作是必要承诺的一部分，应该使用独立任务或正确延长逻辑生命周期。反过来，把永久后台循环放入某个请求任务会使任务迟迟不结束，失去可比较的延迟分布。

对于扇出请求，可以画一个很小的依赖图：入口、各子任务、汇合和响应。先标注真实依赖，再把时间区间映射进去。若本可并行的子任务直到前一子任务完成后才启动，或汇合后仍存在不必要的串行处理，就有潜在的优化空间。已经完成的子任务稍后才被读取结果，本身未必增加总延迟；如果汇合必须等待最慢依赖，则应优先处理那条路径。

\section{Go 1.25 FlightRecorder：保留异常前的窗口}
Go 1.25 将 FlightRecorder 纳入标准库 \texttt{runtime/\allowbreak{}trace}；2024 年的 \texttt{golang.\allowbreak{}org/\allowbreak{}x/\allowbreak{}exp/\allowbreak{}trace} 实验接口属于历史阶段，不能与标准库 API 混写。FlightRecorder 持续维护最近执行轨迹，在慢请求、超时或错误率异常触发后，将此前窗口写出，解决“发现问题时已经来不及 Start”的难题。\cite{tools-go125}\cite{tools-flight-blog}

\begin{codeblock}
// Go 1.25+ 集成片段：生命周期由单一管理者控制。
fr := trace.NewFlightRecorder(trace.FlightRecorderConfig{
    MinAge: 10 * time.Second,
    MaxBytes: 32 << 20,
})
if err := fr.Start(); err != nil { return err }
defer fr.Stop()
// 触发后由单一写入者创建文件并调用 fr.WriteTo(file)。
\end{codeblock}

必须精确理解边界：\texttt{MinAge} 表示希望窗口覆盖到至少这么久以前的事件，但容量配置优先；\texttt{MaxBytes} 是窗口大小的提示，并不保证运行时内存绝不超出该值，也不保证 \texttt{WriteTo} 输出字节数受同样硬限制。事件速率过高时，不能同时承诺固定 10 秒历史和固定 32 MiB 占用。需要根据本机事件速率、异常窗口与内存预算实测规划。\cite{tools-flight-source}

Go 1.25 允许一个活动 FlightRecorder 与一个使用 \texttt{trace.\allowbreak{}Start} 的普通消费者同时存在；不能同时启用多个 FlightRecorder。\texttt{WriteTo} 只能由一个 goroutine 同时执行，并发调用会返回错误。\texttt{WriteTo} 期望快照覆盖到调用时刻的最新数据，但源码明确说明这不是硬保证。触发器应做去重、冷却、文件限额和错误处理。生命周期操作由明确的管理者串行控制，不要让多个 HTTP 请求随意 Start/Stop。

Flight recording 把持续写盘变成按需导出，并未让持续采集免费。它仍维护事件缓冲区，触发快照还涉及刷新与 I/O。不要把网上某个负载的“低开销”数字作为生产保证；首先分别测量仅 recorder、触发导出以及普通 trace 三种情形的吞吐、p99、CPU 和内存。

\section{完整实验：带 task 的跨 goroutine 请求与飞行记录}
\begin{figure}[H]
\centering
\begin{tikzpicture}[x=1cm,y=1cm]
\node[anchor=east,font=\small] at (1,0.0) {泳道1};
\draw[linegray] (1.2,-0.35) -- (14.799999999999999,-0.35);
\node[anchor=east,font=\small] at (1,-1.7) {泳道2};
\draw[linegray] (1.2,-2.05) -- (14.799999999999999,-2.05);
\path[fill=forest!60] (1.200,-0.250) rectangle (4.222,0.250);
\node[align=center,font=\small,text width=2.92cm] at (2.711,0.750) {已淘汰片段};
\path[fill=forest!60] (4.222,-0.250) rectangle (13.289,0.250);
\node[align=center,font=\small,text width=8.97cm] at (8.756,0.750) {保留窗口};
\path[fill=leaf!60] (12.987,-1.950) rectangle (13.289,-1.450);
\draw[linegray] (13.138,-1.450) -- (13.138,-1.020);
\node[circle,fill=mist,inner sep=2pt,font=\scriptsize] at (13.138,-0.850) {3};
\path[fill=forest!60] (13.289,-1.950) rectangle (14.800,-1.450);
\draw[linegray] (14.044,-1.450) -- (14.044,-1.020);
\node[circle,fill=mist,inner sep=2pt,font=\scriptsize] at (14.044,-0.850) {4};
\draw[-{Stealth[length=2mm]},forest] (1.2,-2.8499999999999996) -- (14.999999999999998,-2.8499999999999996);
\draw[forest] (1.20,-2.8499999999999996) -- (1.20,-2.9499999999999997);
\node[below,font=\scriptsize] at (1.20,-2.9499999999999997) {0};
\draw[forest] (3.92,-2.8499999999999996) -- (3.92,-2.9499999999999997);
\node[below,font=\scriptsize] at (3.92,-2.9499999999999997) {1.8};
\draw[forest] (6.64,-2.8499999999999996) -- (6.64,-2.9499999999999997);
\node[below,font=\scriptsize] at (6.64,-2.9499999999999997) {3.6};
\draw[forest] (9.36,-2.8499999999999996) -- (9.36,-2.9499999999999997);
\node[below,font=\scriptsize] at (9.36,-2.9499999999999997) {5.4};
\draw[forest] (12.08,-2.8499999999999996) -- (12.08,-2.9499999999999997);
\node[below,font=\scriptsize] at (12.08,-2.9499999999999997) {7.2};
\draw[forest] (14.80,-2.8499999999999996) -- (14.80,-2.9499999999999997);
\node[below,font=\scriptsize] at (14.80,-2.9499999999999997) {9};
\node[anchor=east,font=\scriptsize] at (14.799999999999999,-3.5999999999999996) {时间 / s};
\node[align=left,anchor=north west,text width=14cm,font=\scriptsize] at (1,-3.8499999999999996) {3=异常触发；4=导出快照};
\end{tikzpicture}
\caption{触发时保留过去，而不是只从现在开始}\label{fig:08-flight}
\begin{minipage}{0.96\linewidth}\footnotesize 窗口图为教学示意。MaxBytes 优先于 MinAge 但属于目标提示，不保证全部内存或文件大小硬上限。

\textbf{读图顺序：}1. 持续滚动记录\quad 2. 在异常处触发\quad 3. 写出可分析的快照\quad 4. 审查覆盖限制。
\textbf{机制依据：}\cite{tools-flight-source}。
\end{minipage}
\end{figure}
下面是完整的 Go 1.25+ 程序，保存为 \texttt{flight.\allowbreak{}go}。它只使用 FlightRecorder，顺序执行 20 个带任务标记的请求后导出快照；每个请求内部由父、子 goroutine 配合完成。短 sleep 用来制造可见等待，数字是实验输入而非性能测量结论。

\begin{codeblock}
package main

import (
	"context"
	"os"
	"runtime/trace"
	"time"
)

func request() {
	ctx, task := trace.NewTask(context.Background(), "request")
	defer task.End()
	done := make(chan struct{})
	go func() {
		defer close(done)
		trace.WithRegion(ctx, "backend", func() {
			time.Sleep(40 * time.Millisecond)
			trace.Log(ctx, "backend", "completed")
		})
	}()
	trace.WithRegion(ctx, "wait-result", func() { <-done })
}

func main() {
	fr := trace.NewFlightRecorder(trace.FlightRecorderConfig{
		MinAge:   2 * time.Second,
		MaxBytes: 8 << 20,
	})
	if err := fr.Start(); err != nil {
		panic(err)
	}
	defer fr.Stop()
	for i := 0; i < 20; i++ {
		request()
	}
	f, err := os.Create("flight.trace")
	if err != nil {
		panic(err)
	}
	if _, err := fr.WriteTo(f); err != nil {
		f.Close()
		panic(err)
	}
	if err := f.Close(); err != nil {
		panic(err)
	}
}
\end{codeblock}

\begin{codeblock}
go version
go build -o flight-lab flight.go
./flight-lab
go tool trace -http=127.0.0.1:8082 flight.trace
\end{codeblock}

查看请求任务的延迟分布，选择一个实例，检查 backend 的 sleep 与 wait-result 的 channel 等待如何重叠。两个 region 持续时间都接近同一等待窗口，不表示请求延迟等于二者相加。再在后台加入实际 CPU 工作，观察运行和等待状态如何变化。

\section{如何评估 trace 的有效性与扰动}
一次可靠的 trace 实验应回答三个问题：是否覆盖目标事件、是否保留足够的前因后果、是否显著改变被观察系统。先用日志或请求指标确认问题确实落在采集窗口，再检查 task/region 是否跨越窗口边界、子任务是否缺失、数据是否能完整解析。只看到异常的最后一步，不足以还原根因。

采集代价包含事件生成、栈记录、缓冲与写出，随 goroutine 切换频率、标注数量和事件密度变化。官方 2024 年文章报告了 Go 1.21–1.22 两次版本演进带来的开销改进及其特定测试背景；教材将其作为设计进步证据，不把报道比例视作任意服务的上界。\cite{tools-trace2024}

设计 A/B 实验时用相同负载分布交替执行“关闭、开启、关闭”，并比较吞吐、p99、CPU、内存与错误率。如果 trace 使系统从接近饱和推到饱和，尾延迟影响可能远大于平均 CPU 开销。可缩短窗口、降低业务注解密度、抽样选实例或采用异常触发保留。

最终报告应同时交付原始 trace、工具链版本、可复现命令、关键事件时间区间及机制解释；截图只用于导航，不能替代原始证据。对于 flight recorder，还应注明配置与实际覆盖时长，特别是容量优先导致的历史缩短。

当事件量很大时，可以先根据业务日志选定异常时间附近的任务，而不是遍历所有事件。筛选条件应该与待验证假设对应，例如指定请求类型、延迟阈值和实例；不要只挑最符合预期的一张图。保留正常样本作为对照，才能判断某种调度或 GC 形态是否真是异常特征。

\section{实验任务}
\subsection*{从任务延迟回溯运行时因果链}
\begin{enumerate}
\item 使用 Go 1.25+ 运行 flight.go，确认 flight.trace 可解析并能找到 request 类型任务。

\item 选择一个 task，记录 backend、wait-result 与 task 的起止，解释重叠为何不能相加。

\item 导出 sched、sync、net profile，说明哪些有内容、哪些可能为空以及原因。

\item 将 sleep 替换为计算工作，并分别在 GOMAXPROCS=1 与多个 P 下观察 running/runnable 状态。

\item 把 recorder 的 MaxBytes 调小，比较实际覆盖时长和文件大小；验证它不是历史时长与内存双重硬保证。

\end{enumerate}
\section{自测与解析}
\noindent\textbf{1. goroutine 已被唤醒但尚未运行的时间，主要属于哪一类？}
\begin{enumerate}[label=\Alph*.]
\item CPU 执行
\item runnable 调度等待
\item 必定是网络阻塞
\item 必定是 GC 暂停
\end{enumerate}
\noindent{\color{forest}\textbf{答案 B。}}解除阻塞后进入 runnable，获得执行资源后才 running；应结合 P 与系统 CPU 供给分析。

\noindent\textbf{2. 关于 task 与 region，哪项正确？}
\begin{enumerate}[label=\Alph*.]
\item 二者都必须只存在于单个 goroutine
\item task 可跨 goroutine，region 应在同一 goroutine 正确配对
\item region 可任意跨 goroutine End
\item task 会自动等待全部子 goroutine
\end{enumerate}
\noindent{\color{forest}\textbf{答案 B。}}context 关联逻辑任务，region 标记当前 goroutine 区间；任务结束时间需要应用显式维护。

\noindent\textbf{3. Go 1.25 的 MaxBytes=32 MiB 保证什么？}
\begin{enumerate}[label=\Alph*.]
\item 进程总 RSS 不超过 32 MiB
\item 任何快照文件严格小于 32 MiB
\item 作为窗口容量提示参与保留策略，不提供绝对内存和输出硬上限
\item 至少保留 32 秒
\end{enumerate}
\noindent{\color{forest}\textbf{答案 C。}}源码文档明确将其视为 hint，容量配置优先于 MinAge，并不提供输出或内存硬保证。

\noindent\textbf{4. Go 1.25 FlightRecorder 能否与 trace.Start 同时运行？}
\begin{enumerate}[label=\Alph*.]
\item 能，但同一时间仅一个 recorder 和一个普通 trace 消费者
\item 不能，二者完全互斥
\item 能启动无限多个 recorder
\item 只允许在 Windows 上并行
\end{enumerate}
\noindent{\color{forest}\textbf{答案 A。}}标准库明确支持一个 flight recorder 与普通 trace 消费者并存；不要沿用早期实验 API 的限制推断。

\section{分析题与参考答案}
\noindent\textbf{题 1：}假设请求在 0 ms 启动，两个子任务分别在 10–70 ms、20–100 ms 执行，父任务等待两者后在 110 ms 返回。计算请求延迟与子任务累计区间，并说明优化优先级。

\noindent\textbf{参考答案：}请求延迟为 110 ms；两个子任务区间分别为 60 和 80 ms，累计 140 ms，因重叠不能等同延迟。第二个任务在 100 ms 最晚结束，通常处于关键路径；缩短第一个任务而不改变依赖可能不改善总延迟。还需检查 0–20 ms 的调度/启动以及 100–110 ms 的收尾。

\noindent\textbf{题 2：}按与记录器窗口一致的字节口径，假设某服务 trace 数据以恒定的 12 MiB/s 产生，希望保留 15 秒。配置 MinAge=15s、MaxBytes=32 MiB 是否足以规划该历史窗口？应如何设计实验？

\noindent\textbf{参考答案：}仅数据量模型就需要约 180 MiB；32 MiB 相当于约 2.7 秒的数据量，并非保证的覆盖时长，在该模型下不足以规划 15 秒历史。真实开销与代际粒度也会影响结果。应测定事件速率分布与峰值，增加容量预算、减少不必要的业务标注，或选择事件量更适合预算的采集实例，并以快照中实际最早事件时间验证覆盖。MaxBytes 非硬上限，还要同时监控额外内存及导出 I/O。

\medskip\noindent\textbf{本章主要阅读：}\cite{tools-trace-doc}\cite{tools-trace-cmd}\cite{tools-trace-annotation}\cite{tools-go122}\cite{tools-trace2024}\cite{tools-go125}\cite{tools-flight-blog}\cite{tools-flight-source}。
\chapter{Delve、崩溃与现场调试}
\label{ch:09}
\noindent{\large\color{forest}在理解暂停语义和现场边界的前提下检查程序状态}\par\medskip
\noindent\textbf{先修要求：}第 02、07、08 章；掌握 Go 错误、panic/recover、栈帧与并发基础。

\noindent\textbf{完成本章后，应能够：}
\begin{itemize}
\item 选择 debug、exec、attach、test 与 core 的适用场景
\item 使用条件断点、goroutine 和 frame 检查并发程序
\item 理解优化构建、DWARF、内联与变量不可用的关系
\item 安全保存崩溃现场，并区分状态证据与性能证据
\end{itemize}
\section{调试器回答状态问题，而非平均成本问题}
Profile 通过汇总寻找高成本路径，调试器则在具体执行点检查寄存器、局部变量、栈帧与 goroutine。“缓存键为何错误”“哪个条件阻止退出”“panic 时指针是什么”适合 Delve；“哪个函数占 20\% CPU”通常应先用 pprof。断点暂停和单步会改变时间、锁竞争及外部超时，因此调试会话不能直接作为生产性能基准。

Delve 理解 Go 的 goroutine、类型和运行时数据结构，并通过平台后端访问进程状态。以下列出不同的调试入口，各自会启动一个独立会话，应按场景选择执行；路径、包名与进程号需要替换成实际对象：

\begin{codeblock}
dlv version
dlv debug ./cmd/service
dlv exec ./service -- -config=test.json
dlv test . -- -test.run '^TestPipeline$'
dlv attach 12345
dlv core ./service core.12345
\end{codeblock}

\texttt{debug} 会构建关闭优化的程序并调试；\texttt{exec} 使用已有二进制；\texttt{test} 构建测试目标；\texttt{attach} 附加到已有进程；\texttt{core} 读取静态转储。参数分隔符 \texttt{-\allowbreak{}-\allowbreak{}} 后面的内容交给被调试程序，而不是 Delve。本章依据 Delve v1.25.2 文档描述命令语义，安装版本应与目标 Go 工具链兼容，保留默认版本检查。\cite{tools-delve-debug}\cite{tools-delve-exec}\cite{tools-delve-core}

附加到正在运行的进程通常会暂停或干扰服务，应优先调试本地复现程序、隔离副本，或分析已有 core。现场保全包括二进制、构建 ID、源代码提交、Go 与 Delve 版本、命令行和必要配置；配置中的秘密应另外控制访问。

一次调试记录应明确“观察到的状态”与“对过去的推断”。某个字段现在为零，不能独自证明它从未赋值；某个 goroutine 当前等待，也不能说明它此前一直没有进展。条件断点可以在状态将发生变化的位置取证，日志或 trace 可以补足时间顺序，core 则提供某一时刻更完整的静态现场。根据问题组合这些工具，比让单一调试器承担全部诊断更可靠。

版本兼容还包含调试器自身的构建版本。Delve v1.25.2 这一固定版本的常规 Go 版本检查范围是 1.22–1.25，riscv64 还要求目标使用 Go 1.24 或更新版本；遇到 DWARFv5 可执行文件时，其兼容检查还要求 Delve 自身由 Go 1.25 或更新版本构建。只看 dlv 的发布号、或用 \texttt{-\allowbreak{}-\allowbreak{}check-\allowbreak{}go-\allowbreak{}version=false} 压过错误，都不能证明运行时布局和 DWARF 读取已经兼容。版本检查通过也不代表支持所有平台与架构，仍应核对对应后端和 core 的支持矩阵。\cite{tools-delve-compat}

\section{构建信息、优化和可调试性}
\begin{figure}[H]
\centering
\begin{tikzpicture}[x=1cm,y=1cm]
\node[draw=linegray!75!black,fill=linegray!13,rounded corners=3pt,text width=3.33cm,minimum height=1.12cm,align=center,font=\small] (p0-src) at (1.88,0.97) {源码变量 x};
\node[draw=leaf!75!black,fill=leaf!13,rounded corners=3pt,text width=3.33cm,minimum height=1.12cm,align=center,font=\small] (p0-opt) at (5.62,0.00) {优化/\allowbreak{}寄存器};
\node[draw=linegray!75!black,fill=linegray!13,rounded corners=3pt,text width=3.33cm,minimum height=1.12cm,align=center,font=\small] (p0-dwarf) at (9.38,0.00) {DWARF 位置};
\node[draw=linegray!75!black,fill=linegray!13,rounded corners=3pt,text width=3.33cm,minimum height=1.12cm,align=center,font=\small] (p0-pc) at (1.88,-0.97) {当前 PC};
\node[draw=leaf!75!black,fill=leaf!13,rounded corners=3pt,text width=3.33cm,minimum height=1.12cm,align=center,font=\small] (p0-value) at (13.12,0.00) {重建变量值};
\draw[-{Stealth[length=2mm]},forest!70!black] (p0-src.east) -- (p0-opt.west);
\draw[-{Stealth[length=2mm]},forest!70!black] (p0-opt.east) -- (p0-dwarf.west);
\draw[-{Stealth[length=2mm]},forest!70!black] (p0-pc.north) -- (1.88,2.12) -- (9.38,2.12) -- (p0-dwarf.north);
\draw[-{Stealth[length=2mm]},forest!70!black] (p0-dwarf.east) -- (p0-value.west);
\end{tikzpicture}
\caption{优化构建为何未必能读到每个变量}\label{fig:09-debugger}
\begin{minipage}{0.96\linewidth}\footnotesize 教学示意：用于解释机制，不是运行时的实测数据或完整实现。

\textbf{读图顺序：}1. 源码不是运行态变量表\quad 2. 变量位置依赖 PC\quad 3. 调试器尝试重建\quad 4. 不能反推性能。
\textbf{连线语义：}优化/寄存器\ensuremath{\rightarrow}\allowbreak{}DWARF 位置：位置描述；当前 PC\ensuremath{\rightarrow}\allowbreak{}DWARF 位置：有效范围；DWARF 位置\ensuremath{\rightarrow}\allowbreak{}重建变量值：可恢复时。
\textbf{机制依据：}\cite{tools-delve-debug}\cite{tools-delve-cli}。
\end{minipage}
\end{figure}
调试器依赖调试信息将机器地址映射到源码、类型与变量位置。Go 默认优化构建通常仍带有有用调试信息，但内联、寄存器分配、常量折叠和变量生命周期缩短会使某些变量显示为优化掉、不可访问或不再位于直观位置。\texttt{-\allowbreak{}ldflags='-\allowbreak{}s -\allowbreak{}w'} 去除符号/调试信息后，源码级调试能力会进一步受限。

\begin{codeblock}
go build -gcflags='all=-N -l' -o service.debug ./cmd/service
dlv exec ./service.debug
go version -m ./service
go tool buildid ./service
\end{codeblock}

\texttt{-\allowbreak{}N} 关闭优化，\texttt{-\allowbreak{}l} 关闭内联，是构造本地调试副本的常见方式；该副本的性能与调度不再等同生产二进制。重建一个“相同源码”的调试版本，也不能拿它解释既有生产 core 的地址。core 与二进制必须匹配，最好还保存原始构建环境与依赖信息。\cite{tools-delve-exec}

若生产问题只在优化后发生，应先使用原二进制检查可获得的信息，再构造最小复现。关闭优化后不再出现，只说明复现条件改变，不能直接断言编译器有 bug。并发时序、未同步访问或 unsafe 误用都可能受到影响。

源码目录移动时，可以配置路径映射而不是改造二进制。内联函数可能对应多个机器位置，一个源码断点也可能安装到多个位置；设置后检查断点列表及实际命中栈。变量读取失败时先核对当前 goroutine、frame 和变量作用域，再考虑调试信息是否缺失。

\section{断点、条件、栈帧和检查命令}
一个有效会话从明确问题开始。例如“只有 id=2 的 worker 返回错误”，就应在 worker 的关键位置设置条件断点，而非反复暂停全部任务。下面是 Delve 提示符中的命令，不是 shell 命令：

\begin{codeblock}
break workerEntry main.worker
condition workerEntry id == 2
continue
args
locals
print id
stack 10
frame 1
list
breakpoints
\end{codeblock}

\texttt{break} 支持函数、文件行与其他 location spec；\texttt{condition} 支持断点名称或编号。条件只能使用该执行位置可访问的表达式，错误条件或优化后的不可用变量都需要检查。\texttt{args} 与 \texttt{locals} 分别查看参数和局部变量；\texttt{frame} 改变检查上下文，不会让程序倒退执行。\cite{tools-delve-cli}

\texttt{next} 跨过当前源码行的调用，\texttt{step} 进入可单步的调用，\texttt{stepout} 运行到当前函数返回。它们是源码级控制指令，不承诺每次正好执行一个机器指令；编译器移动、内联、defer 与调度都会影响所见位置。继续运行时也可能先命中其他断点。

\texttt{trace} 命令设置 tracepoint，命中后显示信息并继续，而非等待用户交互。它与 \texttt{runtime/\allowbreak{}trace} 执行追踪完全不同；其实现仍可能需要陷入调试器并造成开销。\texttt{set} 修改状态，\texttt{call} 可执行函数且可能阻塞或产生副作用，应明确区分观察与干预。排查初期优先只读变量、栈和反汇编，避免把现场改成另一个状态。

查看复杂对象时应逐步展开。先检查长度、容量、关键字段和指针是否为空，再按需读取部分元素，避免一次打印巨大的映射或缓冲区拖慢调试器并暴露无关数据。界面截断也不代表对象真的只有显示出的元素；需要核对输出配置和实际长度。

对于错误路径，可以在生成错误的位置断住，向上检查调用者是否丢弃、包装或错误地重试。若只在最终 panic 位置检查，原始输入和中间状态可能早已不可用。断点位置应由具体假设驱动，并在每次命中时记录必要字段；大量无条件断点会改变调度，增加噪声。

\section{Goroutine 与线程：选择检查上下文与恢复执行}
Go 的 goroutine 在多个操作系统线程间调度，goroutine ID 与线程 ID 不是同一标识。\texttt{goroutines} 列出逻辑执行单元，\texttt{threads} 列出后端跟踪的系统线程。选择 goroutine 后可以查看它的栈，即使它当前停在同步等待而没有独占某个线程。

\begin{codeblock}
goroutines
threads
goroutine 17
stack 20
goroutine 17 print requestID
thread 1234
\end{codeblock}

这里的 17 与 1234 仅为示例，应使用当前列表中的实际 ID。goroutine ID 只在该进程现场有意义，不能在重启后拿它当稳定业务标识。调试器通常需要协调多个线程的停止状态；选择 goroutine 只是选择观察/操作上下文，并不意味着后续 continue 只让它运行。

若需要针对已知 goroutine 限定断点，Delve 支持条件表达式，例如：

\begin{codeblock}
condition workerEntry runtime.curg.goid == 17
\end{codeblock}

这是 Delve 提供的调试表达式惯例，不是业务程序应依赖的公开 Go goroutine ID API。普通业务代码不应通过解析栈文本获取 ID 建立逻辑。条件求值和断点命中仍会干扰调度，应限制频度。\cite{tools-delve-cli}

调试死锁时先画出等待依赖：A 等哪个锁或 channel，谁能解除该等待，解除方又等什么。检查锁顺序、channel 的生产者与关闭者、WaitGroup 计数以及取消信号。单步可能暂时消除竞态，暂停也可能让外部系统超时；如果现场只在真实并发下出现，转向低侵入 trace、race detector 的独立实验或录制重放方案更合适。

\section{Panic、core 与离线分析}
\begin{figure}[H]
\centering
\begin{tikzpicture}[x=1cm,y=1cm]
\node[draw=linegray!75!black,fill=linegray!13,rounded corners=3pt,text width=4.58cm,minimum height=1.12cm,align=center,font=\small] (p0-panic) at (2.50,0.97) {未恢复 panic /\allowbreak{} fatal};
\node[draw=linegray!75!black,fill=linegray!13,rounded corners=3pt,text width=4.58cm,minimum height=1.12cm,align=center,font=\small] (p0-traceback) at (7.50,0.97) {运行时栈文本};
\node[draw=linegray!75!black,fill=linegray!13,rounded corners=3pt,text width=4.58cm,minimum height=1.12cm,align=center,font=\small] (p0-extra) at (12.50,0.97) {SetCrashOutput};
\node[draw=linegray!75!black,fill=linegray!13,rounded corners=3pt,text width=4.58cm,minimum height=1.12cm,align=center,font=\small] (p0-core) at (7.50,-0.97) {OS core};
\node[draw=leaf!75!black,fill=leaf!13,rounded corners=3pt,text width=4.58cm,minimum height=1.12cm,align=center,font=\small] (p0-binary) at (2.50,-0.97) {匹配二进制};
\node[draw=leaf!75!black,fill=leaf!13,rounded corners=3pt,text width=4.58cm,minimum height=1.12cm,align=center,font=\small] (p0-debugger) at (12.50,-0.97) {离线调试};
\draw[-{Stealth[length=2mm]},forest!70!black] (p0-panic.east) -- (p0-traceback.west);
\draw[-{Stealth[length=2mm]},forest!70!black] (p0-traceback.east) -- (p0-extra.west);
\draw[-{Stealth[length=2mm]},forest!70!black] (p0-panic.east) -- (p0-core.west);
\draw[-{Stealth[length=2mm]},forest!70!black] (p0-core.east) -- (p0-debugger.west);
\draw[-{Stealth[length=2mm]},forest!70!black] (p0-binary.north) -- (2.50,2.12) -- (12.50,2.12) -- (p0-debugger.north);
\end{tikzpicture}
\caption{崩溃日志和 core 是两条取证路径}\label{fig:09-crash}
\begin{minipage}{0.96\linewidth}\footnotesize core 路径要求 GOTRACEBACK=crash 等崩溃模式及 OS 收集配置；默认未恢复 panic 通常打印栈后 exit(2)，不能仅因系统允许 core 就假定有转储。

\textbf{读图顺序：}1. 先保存运行时信息\quad 2. 额外输出不是 core\quad 3. OS 是否生成 core 另行配置\quad 4. 组合匹配产物。
\textbf{连线语义：}运行时栈文本\ensuremath{\rightarrow}\allowbreak{}SetCrashOutput：额外文件；未恢复 panic / fatal\ensuremath{\rightarrow}\allowbreak{}OS core：崩溃模式＋OS配置；匹配二进制\ensuremath{\rightarrow}\allowbreak{}离线调试：符号与映射。
\textbf{机制依据：}\cite{tools-delve-core}\cite{tools-runtime-debug-api}。
\end{minipage}
\end{figure}
未恢复 panic 打印的栈，是获取成本较低的崩溃证据之一，但通常不足以恢复全部堆对象和其他线程状态。\texttt{GOTRACEBACK} 控制崩溃时的栈输出范围；\texttt{crash} 在相应平台上以触发系统崩溃转储的方式退出。实际是否生成 core 还取决于 shell 限额、内核 core\_pattern、容器配置和 systemd-coredump 等收集器。\cite{tools-runtime-debug}

以下 Linux shell 命令仅用于本章的隔离实验程序；先按 9.6 构建 debug-lab，\texttt{-\allowbreak{}crash} 会故意结束该测试进程：

\begin{codeblock}
ulimit -c unlimited
cat /proc/sys/kernel/core_pattern
GOTRACEBACK=crash ./debug-lab -crash
\end{codeblock}

不要假设文件一定名为 \texttt{core} 并出现在当前目录。若系统交给转储服务管理，应先查询本机收集方式；获得实际文件后使用匹配的二进制打开：

\begin{codeblock}
dlv core ./debug-lab ./actual.core
\end{codeblock}

也可在 Delve 暂停的本地实验进程内执行 \texttt{dump snapshot.\allowbreak{}core}。v1.25.2 文档说明原生 core 输入覆盖 Linux amd64/arm64、Windows amd64 minidump 及 Delve 自己生成的 dump；平台和架构支持不是无限的，应查本地 \texttt{dlv help core}。离线 core 没有“继续执行”的能力，它是转储采集时刻的寄存器与内存现场。\cite{tools-delve-core}\cite{tools-delve-cli}

Core 可能包含请求正文、认证材料、密钥及堆中的业务数据，体积也可能很大。保全时记录校验值、访问范围和保留期；共享教材或工单只提供脱敏的栈与必要字段。源码重建、调试符号和 core 的匹配性比截图更重要。

转储分析的结论也有边界。一个线程停在系统调用里，并不能仅凭 core 得到该调用已经持续多久；一个锁处于锁定状态，也不必然说明死锁。静态现场需要与故障前日志、指标、超时记录以及其他 goroutine 的依赖一起解读。对于尚在运行的复现进程，必要时可以间隔采集多次受控转储，比较状态是否持续无进展；已经退出的进程无法补采过去的现场。

现场涉及本地原生库时，还应保存对应共享库与映射信息。只有 Go 主二进制可能不足以解释原生栈；缺少符号时可以先记录地址与模块偏移，明确未解析范围，而不是把未知帧随意归入 Go 运行时。

每份转储应同时标注采集触发原因、时间边界和已知缺失的映射，防止后续审阅者把不完整现场理解为全部执行历史。

\subsection{Go 1.23+：使用 SetCrashOutput 保存额外崩溃文本}
\texttt{debug.\allowbreak{}SetCrashOu\allowbreak{}tput(f, debug.\allowbreak{}CrashOptions\{\})} 从 Go 1.23 起支持把未处理 panic 和其他 fatal error 的文本额外写到一个文件，标准错误输出仍然保留。它配置的是单一额外目标：再次调用会替换旧目标；传入 nil 可关闭。运行时复制文件描述符，因此成功返回后调用方可以关闭原 \texttt{f}，但仍须检查设置错误。若替换目标与崩溃并发，部分正在写出的内容可能继续进入旧目标。\cite{tools-crash-output}

它不是 panic 恢复接口，不生成可供 Delve 使用的 core，也不捕获被应用 recover 吞掉的 panic。SIGKILL、容器 OOM kill 或机器断电未必给 Go 运行时输出机会；不能用“没有崩溃文件”证明程序正常退出。应在启动阶段设置受控输出，结合退出状态、日志和系统转储收集机制分析。崩溃文本同样可能包含敏感栈信息，需要控制文件权限与保留范围。

\section{完整实验与远程调试边界}
以下完整程序保存为 \texttt{debug.\allowbreak{}go}，三个 worker 的参数和局部变量可供观察，\texttt{-\allowbreak{}crash} 是显式的本地崩溃实验开关。

\begin{codeblock}
package main

import (
	"flag"
	"fmt"
)

func worker(id int, out chan<- int) {
	total := 0
	for i := 0; i < 5; i++ {
		total += id + i
	}
	out <- total
}

func main() {
	crash := flag.Bool("crash", false, "panic after local experiment")
	flag.Parse()
	out := make(chan int)
	for id := 1; id <= 3; id++ {
		go worker(id, out)
	}
	for i := 0; i < 3; i++ {
		fmt.Println(<-out)
	}
	if *crash {
		panic("intentional local experiment")
	}
}
\end{codeblock}

\begin{codeblock}
go build -gcflags='all=-N -l' -o debug-lab debug.go
dlv exec ./debug-lab
\end{codeblock}

按 9.3 设置条件断点后检查 id；在循环内部设置行断点观察 total，再检查其余 goroutine 停在何处、是否等待发送，以及是否已经进入 worker。进程暂停时显示的运行时状态不表示它此刻仍在消耗 CPU。程序打印顺序不固定，条件断点不会把并发程序变成具有固定执行顺序的程序。

远程调试服务具有读取和修改目标内存、继续执行乃至调用函数的能力，不是可公开的只读诊断网页。若隔离开发环境确需 headless 模式，应绑定 loopback，通过受控隧道连接并保留默认同用户检查：

\begin{codeblock}
dlv exec ./debug-lab --headless --listen=127.0.0.1:2345 --api-version=2
\end{codeblock}

另一终端通过 \texttt{dlv connect 127.\allowbreak{}0.\allowbreak{}0.\allowbreak{}1:\allowbreak{}2345} 连接。不要把 \texttt{-\allowbreak{}-\allowbreak{}accept-\allowbreak{}multiclient} 理解成认证，也不要为了方便将监听地址改成公网通配地址。容器中的 ptrace 权限需要按平台最小范围配置，不能照搬“完全特权”作为默认安装步骤。\cite{tools-delve-exec}

\section{实验任务}
\subsection*{从条件断点到离线现场}
\begin{enumerate}
\item 以调试构建运行 debug.go，设置 id==2 的 worker 条件断点并检查参数、局部变量和调用栈。

\item 列出 goroutines 和 threads，说明两种 ID 的差别以及哪些 goroutine 正在等待。

\item 保存一个本地 Delve dump，退出后用匹配二进制尝试 dlv core；若平台不支持，记录准确错误与支持矩阵。

\item 用默认优化重新构建另一个二进制并比较单步和变量可见性，不能将新二进制用于旧 core。

\item 完成后退出调试器和监听服务；交付脱敏命令记录、版本及状态解释，不用断点期间耗时宣称性能结论。

\end{enumerate}
\section{自测与解析}
\noindent\textbf{1. 选择 goroutine 17 后执行 continue，意味着什么？}
\begin{enumerate}[label=\Alph*.]
\item 保证只有 goroutine 17 运行
\item 选择检查上下文不等于只恢复一个 goroutine
\item 所有其他 goroutine 永久停止
\item 线程 ID 也自动变为 17
\end{enumerate}
\noindent{\color{forest}\textbf{答案 B。}}goroutine 与系统线程由调度器映射，调试上下文选择不等于独占执行或线程调度策略。

\noindent\textbf{2. 生产 core 应搭配哪一个二进制？}
\begin{enumerate}[label=\Alph*.]
\item 同源码重建且关闭优化的任意文件
\item 实际生成该现场的匹配二进制
\item 只要 main 函数同名即可
\item 最新版服务二进制
\end{enumerate}
\noindent{\color{forest}\textbf{答案 B。}}地址布局、构建信息和调试数据必须匹配；同源码重建也可能改变地址与布局。

\noindent\textbf{3. Delve tracepoint 与 runtime/trace 的关系是什么？}
\begin{enumerate}[label=\Alph*.]
\item 同一种事件文件的两个入口
\item 前者是自动继续的调试断点，后者是运行时事件追踪
\item 二者都不产生开销
\item 前者能导出全部 GC 时序
\end{enumerate}
\noindent{\color{forest}\textbf{答案 B。}}名字相似但机制与数据不同；调试 tracepoint 仍可能造成陷入与停顿。

\noindent\textbf{4. 调试构建消除了问题，能否直接证明优化器错误？}
\begin{enumerate}[label=\Alph*.]
\item 能，关闭优化就是证明
\item 不能，并发时序、未同步访问和 unsafe 误用等也可能受到影响
\item 能，只要某个变量不可见
\item 不能，但可直接忽略生产问题
\end{enumerate}
\noindent{\color{forest}\textbf{答案 B。}}构建改变了性能和调度，需要最小复现、正确性检查与原构建证据，才能讨论编译器缺陷。

\section{分析题与参考答案}
\noindent\textbf{题 1：}goroutine A 持有 mu 并等待 ack 接收；goroutine B 必须取得 mu 后才能向 ack 发送确认。说明死锁依赖并给出检查与修复方案。

\noindent\textbf{参考答案：}A 等待 B 的确认，B 等待 A 释放 mu，形成 A\ensuremath{\rightarrow}\allowbreak{}B\ensuremath{\rightarrow}\allowbreak{}A 的环。用 A、B 的栈、channel 与锁的地址，以及源码中的加锁、解锁和收发路径，确认它们确实引用同一组资源。普通 sync.Mutex 不记录一个可直接查询的所属 goroutine，不能只看锁结构就确定谁应解锁。可在不破坏不变量的前提下先释放锁再等待确认，或重新设计所有权与消息协议，避免持锁等待依赖自己锁的操作；通过超时、取消和并发重复测试验证。

\noindent\textbf{题 2：}某变量在 Delve 显示不可用。给出从低成本到高成本的排查顺序。

\noindent\textbf{参考答案：}先确认当前 goroutine、frame、源码位置与作用域；再检查是否已过变量存活期、函数内联和优化情况；核对二进制与源码及 core 是否匹配、DWARF 是否被移除；最后在隔离环境以 -N -l 构建复现。不可用不等于变量为零，也不允许用调试副本替代生产 core 的匹配文件。

\medskip\noindent\textbf{本章主要阅读：}\cite{tools-delve-debug}\cite{tools-delve-exec}\cite{tools-delve-cli}\cite{tools-delve-core}\cite{tools-delve-compat}\cite{tools-runtime-debug}\cite{tools-crash-output}。
\chapter{运行时诊断与操作系统工具}
\label{ch:10}
\noindent{\large\color{forest}建立从低成本指标到内核观测的分层证据链}\par\medskip
\noindent\textbf{先修要求：}第 02、05–09 章；具备 Linux 进程、线程、系统调用和容器资源的基本知识。

\noindent\textbf{完成本章后，应能够：}
\begin{itemize}
\item 读取 runtime/metrics 并正确区分计数器、瞬时量和直方图
\item 建立受控的 pprof/trace 诊断端点并管理采集开销
\item 解释 GODEBUG、perf、strace 与 eBPF 各自的测量边界
\item 把 Go 调度、内核线程和容器资源约束关联起来
\end{itemize}
\section{分层工具与生态地图}
\begin{figure}[H]
\centering
\begin{tikzpicture}[x=1cm,y=1cm]
\node[draw=linegray!75!black,fill=linegray!13,rounded corners=3pt,text width=3.33cm,minimum height=1.12cm,align=center,font=\small] (p0-question) at (1.88,0.00) {待解释现象};
\node[draw=linegray!75!black,fill=linegray!13,rounded corners=3pt,text width=3.33cm,minimum height=1.12cm,align=center,font=\small] (p0-go) at (5.62,0.97) {Go 语义};
\node[draw=linegray!75!black,fill=linegray!13,rounded corners=3pt,text width=3.33cm,minimum height=1.12cm,align=center,font=\small] (p0-os) at (5.62,-0.97) {OS 线程/\allowbreak{}内核};
\node[draw=linegray!75!black,fill=linegray!13,rounded corners=3pt,text width=3.33cm,minimum height=1.12cm,align=center,font=\small] (p0-profile) at (9.38,0.97) {pprof /\allowbreak{} trace};
\node[draw=linegray!75!black,fill=linegray!13,rounded corners=3pt,text width=3.33cm,minimum height=1.12cm,align=center,font=\small] (p0-kernel) at (9.38,-0.97) {perf /\allowbreak{} strace};
\node[draw=leaf!75!black,fill=leaf!13,rounded corners=3pt,text width=3.33cm,minimum height=1.12cm,align=center,font=\small] (p0-join) at (13.12,0.00) {按边界关联};
\draw[-{Stealth[length=2mm]},forest!70!black] (p0-question.east) -- (p0-go.west);
\draw[-{Stealth[length=2mm]},forest!70!black] (p0-question.east) -- (p0-os.west);
\draw[-{Stealth[length=2mm]},forest!70!black] (p0-go.east) -- (p0-profile.west);
\draw[-{Stealth[length=2mm]},forest!70!black] (p0-os.east) -- (p0-kernel.west);
\draw[-{Stealth[length=2mm]},forest!70!black] (p0-profile.east) -- (p0-join.west);
\draw[-{Stealth[length=2mm]},forest!70!black] (p0-kernel.east) -- (p0-join.west);
\end{tikzpicture}
\caption{先选证据层，再选工具}\label{fig:10-tool-map}
\begin{minipage}{0.96\linewidth}\footnotesize 教学示意：用于解释机制，不是运行时的实测数据或完整实现。

\textbf{读图顺序：}1. 先写缺失的量\quad 2. 语言层\quad 3. 系统层\quad 4. 组合不是混加。
\textbf{机制依据：}\cite{go-diagnostics}\cite{tools-perf}\cite{tools-strace}。
\end{minipage}
\end{figure}
当服务慢时，不应同时开启所有工具并希望自动出现答案。先写出待验证问题：是否 CPU 饱和、是否堆增长、是否同步竞争、是否下游等待、是否被容器节流。然后选择能观察该机制且成本最低的工具，再用更细证据解释异常。

\begingroup\small\setlength{\tabcolsep}{4pt}
\begin{longtable}{>{\raggedright\arraybackslash}p{0.215\linewidth}>{\raggedright\arraybackslash}p{0.215\linewidth}>{\raggedright\arraybackslash}p{0.215\linewidth}>{\raggedright\arraybackslash}p{0.215\linewidth}}
\toprule
\textbf{层次} & \textbf{工具} & \textbf{能回答} & \textbf{主要边界} \\
\midrule
\endfirsthead
\toprule
\textbf{层次} & \textbf{工具} & \textbf{能回答} & \textbf{主要边界} \\
\midrule
\endhead
运行时总量 & runtime/metrics、GODEBUG & GC、分配、调度和锁等待趋势 & 缺少完整业务归因 \\
\addlinespace[3pt]
调用栈聚合 & pprof & CPU、分配、同步热点 & 不保留完整时间顺序 \\
\addlinespace[3pt]
运行时时序 & go tool trace & goroutine 依赖与调度延迟 & 不覆盖所有进程/内核细节 \\
\addlinespace[3pt]
具体状态 & Delve、core & 变量、栈和崩溃现场 & 暂停会扰动执行 \\
\addlinespace[3pt]
系统执行 & perf & 计数器、CPU 栈、内核事件 & 线程不等于 goroutine \\
\addlinespace[3pt]
系统调用 & strace & syscall 参数、结果、时长 & 不能直接观察 Go channel \\
\addlinespace[3pt]
可编程内核观测 & eBPF & 调度、I/O、网络等定制事件 & 受权限、内核与展开能力约束 \\
\addlinespace[3pt]
\bottomrule
\end{longtable}
\endgroup
保持单位和身份可对齐：记录 PID、容器、构建、时间窗口及逻辑请求类别。进程 PID 可以复用，容器内外 PID 命名空间不同；一个 Go goroutine 会跨线程迁移。只有名称类似的事件不足以证明它们属于同一请求。

\subsection{生态工具地图：补足等待、发现、指标与机器码}
\subsubsection{fgprof：同时看计算和等待的聚合栈}
fgprof v0.9.5 是 wall-clock profiler：按名义 99 Hz 调用 goroutine 栈采样接口，覆盖正在执行以及等待 I/O、channel、锁等状态的 goroutine，并输出 pprof 或 folded 格式。它适合“请求很慢但 CPU profile 没有大热点”的混合计算与等待负载。汇总权重近似为 goroutine 时间，多条并发或空闲等待栈可以叠加超过墙钟窗口；它不是进程 CPU 使用量。\cite{tools-fgprof}\cite{tools-fgprof-source}

与内核线程 off-CPU 工具不同，fgprof 观察 Go goroutine 层；但聚合结果仍不保留请求依赖与完整事件顺序，不能直接证明等待因果。长期空闲 worker 可能很宽，应结合业务类别和 trace 判断。开销随 goroutine 数增加，采样 goroutine 自身还可能受到调度延迟影响；旧 Go 版本有明显暂停风险，版本升级和本机评估都必要。把处理器注册到受控诊断 mux 后，可以用熟悉的 pprof 界面阅读。

\subsubsection{google/gops：发现进程，再进入已启用的诊断能力}
gops v0.3.28 可列出 Go 进程的 PID、构建版本和程序路径；启动诊断 agent 的进程会被特别标记。能发现进程并不表示能读取它的 Go 栈或 profile：高级能力要求目标程序显式运行 \texttt{agent.\allowbreak{}Listen}，并能访问其诊断端口。典型查询入口如下，数字仅为示例 PID：

\begin{codeblock}
gops
gops 12345
gops stack 12345
gops memstats 12345
\end{codeblock}

本地 PID 模式要求同用户环境和可用的端口发现配置，远程模式使用明确的地址。v0.3.28 默认监听 loopback 随机 TCP 端口，但源码明确同机其他程序也能连接；同用户发现约定并非 TCP 身份认证。诊断接口还可触发 GC、改变 GC 参数和采集 trace，不能按只读网页对待。进程元数据权限、容器 PID 命名空间以及 listener 可达性都可能限制使用。\cite{tools-gops}\cite{tools-gops-agent}

\subsubsection{expvar：暴露聚合计数，不提供调用栈}
标准库 expvar 将整数、浮点数、映射与回调计算值以 JSON 暴露在 \texttt{/\allowbreak{}debug/\allowbreak{}vars}。它适合请求计数、缓存条目等简单业务状态；使用自定义 mux 时应显式注册 \texttt{expvar.\allowbreak{}Handler()}。导入包不会自动启动服务器，默认还会注册命令行参数与 memstats，诊断边界应与 pprof 一起管理。\cite{tools-expvar}

原子更新一个计数器不代表读取多个变量得到同一时刻的事务快照；回调暴露的数据也需要应用自行保证并发安全。计数器做区间差分，再除以窗口时长，才能得到该窗口的平均速率。但这些计数无法告诉你哪条调用栈消耗 CPU、哪些对象被保留或一次请求如何排队。用它发现趋势，再用 metrics、pprof 和 trace 查机制，职责更清晰。

\subsubsection{nm 与 objdump：验证编译器实际生成了什么}
\texttt{go tool nm} 列出目标文件、归档或可执行文件中的符号、地址和类型，\texttt{-\allowbreak{}size} 显示大小，按 size 排序可检查代码或数据体积；\texttt{go tool objdump} 反汇编匹配符号，\texttt{-\allowbreak{}S} 交错显示可找到的源码。例如可使用第 09 章已有的调试二进制：

\begin{codeblock}
go tool nm -size -sort size ./debug-lab
go tool objdump -S -s 'main[.]worker' ./debug-lab
\end{codeblock}

它们适合验证函数是否保留、调用是否内联以及热循环生成的指令。符号大小不是执行成本，指令条数也不等于运行时间；必须与 profile 的热点权重和真实构建相结合。内联、裁剪和符号剥离会影响结果，分析生产问题时应使用生产匹配产物。\cite{tools-nm}\cite{tools-objdump}

fgprof 的一个具体反例：只考虑主 goroutine 等待 10 秒、四个 worker 同时各活动 10 秒这五条栈，聚合总量约为 50 goroutine 秒，主等待栈约占 20\%。这不表示主 goroutine 消耗了 20\% CPU，也不表示删除其等待就能加速请求；它可能恰好是正确等待子任务的汇合点。v0.9.5 只去除采集器自身栈，并不会自动识别这种业务等待为无用工作。\cite{tools-fgprof-source}

\section{runtime/metrics：稳定名称、明确种类和区间差分}
\begin{figure}[H]
\centering
\begin{tikzpicture}[x=1cm,y=1cm]
\node[anchor=east,font=\small] at (1,0.0) {泳道1};
\draw[linegray] (1.2,-0.35) -- (14.799999999999999,-0.35);
\path[fill=forest!60] (1.200,-0.250) rectangle (1.472,0.250);
\draw[linegray] (1.336,0.250) -- (1.336,0.680);
\node[circle,fill=mist,inner sep=2pt,font=\scriptsize] at (1.336,0.850) {1};
\path[fill=forest!60] (1.472,-0.250) rectangle (14.528,0.250);
\node[align=center,font=\small,text width=12.96cm] at (8.000,0.750) {观察区间 10s};
\path[fill=forest!60] (14.528,-0.250) rectangle (14.800,0.250);
\draw[linegray] (14.664,0.250) -- (14.664,0.680);
\node[circle,fill=mist,inner sep=2pt,font=\scriptsize] at (14.664,0.850) {3};
\draw[-{Stealth[length=2mm]},forest] (1.2,-1.15) -- (14.999999999999998,-1.15);
\draw[forest] (1.20,-1.15) -- (1.20,-1.25);
\node[below,font=\scriptsize] at (1.20,-1.25) {0};
\draw[forest] (3.92,-1.15) -- (3.92,-1.25);
\node[below,font=\scriptsize] at (3.92,-1.25) {2};
\draw[forest] (6.64,-1.15) -- (6.64,-1.25);
\node[below,font=\scriptsize] at (6.64,-1.25) {4};
\draw[forest] (9.36,-1.15) -- (9.36,-1.25);
\node[below,font=\scriptsize] at (9.36,-1.25) {6};
\draw[forest] (12.08,-1.15) -- (12.08,-1.25);
\node[below,font=\scriptsize] at (12.08,-1.25) {8};
\draw[forest] (14.80,-1.15) -- (14.80,-1.25);
\node[below,font=\scriptsize] at (14.80,-1.25) {10};
\node[anchor=east,font=\scriptsize] at (14.799999999999999,-1.9) {时间 / s};
\node[align=left,anchor=north west,text width=14cm,font=\scriptsize] at (1,-2.15) {1=t0 读数=100；3=t1 读数=160};
\end{tikzpicture}
\caption{累计计数器怎样换算为区间增量}\label{fig:10-counters}
\begin{minipage}{0.96\linewidth}\footnotesize 教学示意：用于解释机制，不是运行时的实测数据或完整实现。

\textbf{读图顺序：}1. 读取起点\quad 2. 观察区间\quad 3. 读取终点\quad 4. 检查失效条件。
\textbf{机制依据：}\cite{tools-metrics}\cite{tools-metrics-read}。
\end{minipage}
\end{figure}
\texttt{runtime/\allowbreak{}metrics} 提供带单位的名称与类型，使用 \texttt{metrics.\allowbreak{}All()} 查询当前工具链支持的集合，再读 \texttt{Description.\allowbreak{}Kind}、\texttt{Cumulative} 和说明文字。名称未知时返回 \texttt{KindBad}；不要因开发机存在某个指标就假定所有生产 Go 版本都支持。估计速率时，计数器需要差分后除以窗口时长；瞬时量可直接观察数值与趋势；累计直方图需要逐桶相减得到区间分布。前后样本应来自同一未重启进程，并使用一致的指标定义和桶边界。\cite{tools-metrics}\cite{tools-metrics-read}

常用指标包括 \texttt{/\allowbreak{}gc/\allowbreak{}heap/\allowbreak{}allocs:\allowbreak{}bytes}、\texttt{/\allowbreak{}gc/\allowbreak{}heap/\allowbreak{}live:\allowbreak{}bytes}、\texttt{/\allowbreak{}sched/\allowbreak{}goroutines:\allowbreak{}goroutines}、\texttt{/\allowbreak{}sched/\allowbreak{}latencies:\allowbreak{}seconds} 和 \texttt{/\allowbreak{}sync/\allowbreak{}mutex/\allowbreak{}wait/\allowbreak{}total:\allowbreak{}seconds}。其中 live 反映上一次 GC 标记的存活对象占用量，并不是任意时刻所有已分配堆内存的精确值。最后一项近似累计 goroutine 在 sync.Mutex、sync.RWMutex 和运行时内部锁上的阻塞时间，用于发现整体竞争变化；要归因仍需 mutex/block profile。调度直方图统计 runnable 到实际运行的时间，不包括所有阻塞时间。

一个易错点是 \texttt{/\allowbreak{}cpu/\allowbreak{}classes/\allowbreak{}total:\allowbreak{}cpu-\allowbreak{}seconds}：它表示由 GOMAXPROCS 积分得到的估计可用 CPU 时间，包含该分类体系中的 idle，并不是操作系统记录的实际进程 CPU 消耗。官方明确 \texttt{/\allowbreak{}cpu/\allowbreak{}classes} 的估计值不能直接与系统 CPU 时间比较，宜在同一分类家族内比较比例。\cite{tools-metrics}

下面是完整的指标读取程序，保存为 \texttt{metrics.\allowbreak{}go}。它打印标量指标的当前读数与累计直方图的事件总数：其中分配字节和锁等待是累计量，live 则是上一次 GC 得到的存活堆值。这些输出均不应直接标为每秒速率：

\begin{codeblock}
package main

import (
	"fmt"
	"runtime/metrics"
	"time"
)

func main() {
	samples := []metrics.Sample{
		{Name: "/gc/heap/allocs:bytes"},
		{Name: "/gc/heap/live:bytes"},
		{Name: "/sched/latencies:seconds"},
		{Name: "/sync/mutex/wait/total:seconds"},
	}
	for round := 0; round < 3; round++ {
		metrics.Read(samples)
		for _, s := range samples {
			switch s.Value.Kind() {
			case metrics.KindUint64:
				fmt.Println(s.Name, s.Value.Uint64())
			case metrics.KindFloat64:
				fmt.Println(s.Name, s.Value.Float64())
			case metrics.KindFloat64Histogram:
				h := s.Value.Float64Histogram()
				var n uint64
				for _, count := range h.Counts {
					n += count
				}
				fmt.Println(s.Name, "cumulative events", n)
			default:
				fmt.Println(s.Name, "unsupported")
			}
		}
		time.Sleep(time.Second)
	}
}
\end{codeblock}

这个短程序没有主动制造业务负载或触发 GC，部分读数可能很小或保持为零，应先核对指标定义再解读。重复读取可能复用内部存储，特别是直方图切片；若保存旧窗口或跨 goroutine 传递，必须深拷贝所需数据，避免前一次快照被覆盖或发生数据竞争。同一批读取也不应被想当然地当作所有字段的事务性业务快照。\cite{tools-metrics-read}

\section{GODEBUG：快速观察 GC 与调度趋势}
GODEBUG 适合在受控进程中临时输出运行时诊断信息；不同 Go 版本支持项和输出格式可能改变，解析前阅读对应版本文档。启动时指定仅影响该命令，不要为了一个实验修改全局环境。

\begin{codeblock}
GODEBUG=gctrace=1 ./service
GODEBUG=schedtrace=1000,scheddetail=1 ./service
GODEBUG=scavtrace=1 ./service
\end{codeblock}

\texttt{gctrace} 输出 GC 周期、相关时间与堆规模等信息；\texttt{schedtrace=1000} 大约每 1000 ms 输出调度摘要，\texttt{scheddetail=1} 加入更细的 P、M、G 状态；\texttt{scavtrace} 观察向操作系统归还内存的情况，区分后台清退和主动清退等统计。摘要中的字段不是统一口径的请求延迟，更不应把不同列的 wall time 和 CPU time 随意相加。\cite{tools-runtime-debug}

调度队列持续增长而可用 P 全忙，支持 CPU 供给不足的假设；goroutine 多而队列短、P 空闲，则可能正常等待或依赖外部资源。这里“空闲”需要和容器配额一起解释：进程可见很多逻辑核，但 cgroup 的时间配额较小，会出现宿主机看似空闲、容器却被节流的情况。

Go 1.25 改进了容器环境下默认 GOMAXPROCS 的选择和自动更新，但具体行为还受模块语言版本影响：本书 lab 的 go.mod 声明 go 1.22，因此在没有显式覆盖相关设置时默认保留旧行为，详见第 02 章。对配额实验同时记录 Go 工具链、go.mod、显式 GOMAXPROCS 与相关 GODEBUG，不能仅凭运行 go version 就推断新策略已启用。不要把手工设置一个很大的 GOMAXPROCS 当作普遍优化；更多并发执行者可能加重争用和配额突发。应记录实际值并检查相应版本说明。\cite{tools-go125}

运行时文本诊断可能产生大量 stderr 数据，详细调度输出尤其如此。小范围、短窗口使用并记录启停时间；长期监控优先使用结构化 metrics。

\section{runtime/debug：区分读取、干预和完整堆转储}
\texttt{runtime/\allowbreak{}debug} 既提供诊断数据，也提供改变运行时行为的控制接口；名称中带 debug 不代表只读。\texttt{ReadGCStats} 读取 GC 次数与有限的近期暂停记录，适合快速检查，不是整个进程生命周期的完整暂停事件日志。长期监控应优先使用结构化指标和明确定义的统计窗口。\cite{tools-runtime-debug-api}

\begingroup\small\setlength{\tabcolsep}{4pt}
\begin{longtable}{>{\raggedright\arraybackslash}p{0.298\linewidth}>{\raggedright\arraybackslash}p{0.298\linewidth}>{\raggedright\arraybackslash}p{0.298\linewidth}}
\toprule
\textbf{API} & \textbf{行为} & \textbf{使用边界} \\
\midrule
\endfirsthead
\toprule
\textbf{API} & \textbf{行为} & \textbf{使用边界} \\
\midrule
\endhead
SetGCPercent & 改变 GC 增长目标并返回旧值 & 影响整个进程，不能当作无副作用 getter \\
\addlinespace[3pt]
SetMemoryLimit & 改变软内存限制并返回旧值 & 负参数只读取当前限制 \\
\addlinespace[3pt]
FreeOSMemory & 强制 GC 并尽力归还内存 & 会扰动执行，不适合请求热路径 \\
\addlinespace[3pt]
SetTraceback & 改变崩溃栈详细程度 & 不能降低环境变量要求的详细级别 \\
\addlinespace[3pt]
WriteHeapDump & 输出对象级堆转储 & 写出期间暂停全部 goroutine \\
\addlinespace[3pt]
\bottomrule
\end{longtable}
\endgroup
特别注意，\texttt{debug.\allowbreak{}SetMemoryL\allowbreak{}imit(-\allowbreak{}1)} 是读取当前限制的合法方式，但 \texttt{debug.\allowbreak{}SetGCPercent(-\allowbreak{}1)} 会关闭按普通增长目标触发的 GC；达到软内存限制时仍可能回收，不能因为两者名字类似就推断负值含义相同。动态参数实验应由单一管理者实施，记录原值、变更时刻和恢复方式。

\texttt{FreeOSMemory} 可以帮助区分运行时保留页与不可回收存活集，但调用后 RSS 降低只证明此时能归还部分内存，不能证明此前的分配策略合理。它还改变后续分配和缺页行为，所以不应通过每请求强制回收“修复”内存曲线。

\texttt{WriteHeapD\allowbreak{}ump} 与 \texttt{pprof.\allowbreak{}WriteHeapP\allowbreak{}rofile} 完全不同：前者输出对象与引用层面的特定堆转储格式，后者输出按分配栈采样聚合的数据。完整转储需要专门读取工具，不能直接交给 \texttt{go tool pprof}。官方还要求输出目标不能是由同一个 Go 进程另一端消费的管道或套接字，因为全部 goroutine 被暂停时可能无法排空写入。这种操作必须在已评估暂停和存储预算的隔离场景下使用。

\section{诊断端点：可访问性、隔离与采集控制}
以下是完整的本地诊断服务器，保存为 \texttt{diagnostics.\allowbreak{}go}。实际服务可将相同路由放进独立的诊断 listener。使用显式 mux 避免误以为注册到 DefaultServeMux 的端点会自动出现在自定义服务器里；Go 1.22+ 的 \texttt{GET .\allowbreak{}.\allowbreak{}.\allowbreak{}} 方法路由限制常见非读取方法，但标准 ServeMux 的 GET 模式也匹配 HEAD。自定义 mux 必须自行采用方法模式；处理器函数本身不普遍替调用方检查方法。

\begin{codeblock}
package main

import (
	"log"
	"net/http"
	httppprof "net/http/pprof"
	"time"
)

func main() {
	mux := http.NewServeMux()
	mux.HandleFunc("GET /debug/pprof/", httppprof.Index)
	mux.HandleFunc("GET /debug/pprof/profile", httppprof.Profile)
	mux.HandleFunc("GET /debug/pprof/trace", httppprof.Trace)
	srv := &http.Server{
		Addr:              "127.0.0.1:6060",
		Handler:           mux,
		ReadHeaderTimeout: 3 * time.Second,
	}
	log.Fatal(srv.ListenAndServe())
}
\end{codeblock}

导入 \texttt{net/\allowbreak{}http/\allowbreak{}pprof} 本身不会启动 listener，也不会自动启用 block 和 mutex 采集率。heap 默认以二进制输出；未设置 seconds 时，\texttt{debug>0} 切换可读文本，\texttt{heap?gc=1} 先触发 GC；CPU 的 \texttt{seconds} 是采集时长，而 heap/block/mutex 等的 \texttt{seconds} 是前后快照差分，含义不可互换。heap 同时设置 seconds 与 gc=1 时进入差分分支，不执行该处理器的强制 GC 分支；差分也不支持 debug>0。\cite{tools-http-pprof}

Loopback 绑定降低误暴露机会，但并不等于完整认证体系；同主机、代理转发和容器网络仍需按部署边界处理。生产访问应通过受控管理网络、已有认证代理或授权隧道。profile、trace 与栈可能包含标签和运行细节，应按诊断数据管理。

采集控制至少要约束并发次数、窗口上限和文件保留。CPU 与普通 trace 各有全局活动限制，多个运维自动任务不能无限重试争抢；高频强制 GC 或超长 trace 也会扰动业务。服务端采集失败要检查 HTTP 状态和文件可解析性，不能把错误响应保存成 \texttt{.\allowbreak{}pprof} 后认为采集成功。

\section{perf 与 strace：观察线程和内核边界}
Linux \texttt{perf} 可读取性能计数器并采样执行栈。先用 \texttt{perf stat} 查看任务时钟、上下文切换及硬件计数，再用 \texttt{perf record} 定位执行位置。以下为针对自有实验进程的示例，先将目标 PID 替换为实际值，并确认进程能覆盖整个采集窗口；权限和可用事件取决于本机 perf\_event 配置。每轮应在独立输出目录运行，避免 perf.data 和后续 syscalls.log 覆盖先前结果。

\begin{codeblock}
target_pid=12345
perf stat -p "$target_pid" -- sleep 10
perf record -F 99 --call-graph dwarf -p "$target_pid" -- sleep 10
perf report
\end{codeblock}

\texttt{-\allowbreak{}-\allowbreak{}call-\allowbreak{}graph dwarf} 需要保存用户栈并展开，可能产生明显数据量与开销；frame pointer 模式依赖目标架构和编译产物保留可靠链。cgo、Go 栈增长、内联、内核符号和容器映射都可能影响栈质量。perf 的线程视角可以补足内核执行热点，但不能自动把一个线程的全部成本分给某个 goroutine。\cite{tools-perf}

\texttt{strace} 记录系统调用边界，适合解释重复失败、futex、epoll、打开文件和网络读写等行为。\texttt{-\allowbreak{}T} 显示 syscall 进入到返回的墙钟间隔，不是纯 CPU 时间；阻塞 syscall 可以持续很久却几乎不消耗 CPU。\texttt{-\allowbreak{}c} 的汇总和 \texttt{-\allowbreak{}w} 的墙钟汇总有不同计时口径，解读前确认实际参数。\cite{tools-strace}

\begin{codeblock}
timeout --signal=INT 10s strace -f -T -tt   -e trace=futex,epoll_wait,epoll_pwait   -o syscalls.log -p "$target_pid"
\end{codeblock}

Go 的许多 channel 操作和用户态调度不会触发 syscall，所以 strace 中看不到不等于它们没有发生。附加追踪器会改变 syscall 性能，尤其高频 I/O；应限制事件种类和时间窗口。若扩展到 read/write 参数，需要考虑输出包含业务正文的风险。

\section{eBPF、容器资源与跨层归因}
eBPF 允许在受验证的内核执行环境中处理事件，并通过 map 或事件缓冲区把聚合数据送到用户态。Tracepoint、kprobe、uprobe 等附着点观察的对象不同；成熟工具通常将内核采集程序、用户态符号化、栈展开与聚合结合起来。验证器保障特定程序安全属性，并不保证观测零开销、数据无丢失或业务归因准确。\cite{tools-ebpf}

On-CPU 采样适合找执行热点；调度切换事件可构造线程 off-CPU 时间；网络和块设备事件能定位外部等待。对 Go 服务必须注意两层调度：goroutine 在用户态等待 channel 时，线程可能立即去执行另一个 goroutine；线程 off-CPU 的开始/结束也不自动对应某个请求。因此声称“无侵入精确得到所有 Go 请求等待时间”的方案，需要明确它如何识别 goroutine、处理版本布局、关联上下文及展开用户栈。

Linux 容器排查还应读取 cgroup 配额、节流计数、内存事件与压力指标，并区分容器视图和宿主机视图。cgroup v2 的 \texttt{cpu.\allowbreak{}stat}、\texttt{cpu.\allowbreak{}max}、\texttt{memory.\allowbreak{}events} 等路径与归属由实际部署决定，不能直接假定任意 \texttt{/\allowbreak{}sys/\allowbreak{}fs/\allowbreak{}cgroup} 文件就是目标服务所属组。结合 \texttt{/\allowbreak{}proc/\allowbreak{}<pid>/\allowbreak{}cgroup} 与部署配置定位后再解释。

当 CPU profile 有稳定计算热点且 perf 同样指向相应用户栈，计算解释更有力；当 trace 显示很长 runnable 间隔而 cgroup 节流同步增加，应优先分析 CPU 供给；当 Go 堆稳定而 RSS 增长，应检查 mmap/cgo 和内核映射；当 syscall 延迟高而进程 CPU 低，则需进一步追踪 I/O 或远端依赖。每项结论都应明确各工具覆盖了什么，以及仍有哪些未知项。

\section{建立可重复的分层诊断流程}
第一步固定负载与身份，记录构建、工具链、PID/容器、时间边界、GOMAXPROCS 和资源限额。第二步观察低成本总量：吞吐、错误率、延迟、CPU、内存、GC、调度与锁等待趋势。第三步用对应 profile 找候选调用路径。第四步在需要因果顺序时补 trace，在需要具体状态时补 Delve/core，在运行时解释不到的位置补 perf、strace 或 eBPF。

每一轮只验证一个主要机制。例如“编码路径的分配速率过高导致 GC assist 增加”需要 alloc profile 指向编码栈、GC/CPU 证据支持 assist 成本，以及减少分配后同等负载的改善；仅看到 gctrace 频繁输出不能完成证明。排查阶段改变了采样率或构建参数，也必须写入实验记录。

跨工具对账不要追求不可能的逐字节或逐纳秒相等。它们可能使用不同对象集合、采样机制、时钟、窗口和版本语义。先列出口径，再解释差异：pprof CPU 是执行采样估计，runtime CPU classes 是运行时分类估计，操作系统 CPU 是内核记账；三者可以趋势一致而数值不完全相同。

最后用不依赖诊断采集的重复基准或服务实验验证修改，检查功能、吞吐、延迟与资源是否共同满足目标。原始数据、命令和清晰的未知项，是可复核工程结论的一部分。下一章将以这套证据选择优化与 PGO 策略。

\section{实验任务}
\subsection*{同一现象的分层观测}
\begin{enumerate}
\item 运行 metrics.go，并通过当前工具链的 metrics.All() 核对名称、种类与 Cumulative；把一个累计计数的区间增量除以窗口时长，转换成平均速率，直方图差分时深拷贝旧桶。

\item 启动 diagnostics.go，确认仅 loopback 监听并分别采集 CPU、heap、goroutine；检查 HTTP 状态与文件可解析性。

\item 在第 07 章锁实验中短时开启 schedtrace，并保证实验持续时间覆盖多个输出周期；分别采集 block/mutex 与 trace，说明各工具的观察粒度。必要时延长循环次数，但各对照轮次保持相同工作量。

\item 在支持的 Linux 隔离环境中，对同一实验进程运行有界 perf 或 strace；权限不足时记录实际限制，不修改宿主机全局安全策略。

\item 提交分层对照表，列出每项指标单位、累计/瞬时属性、时间窗口、版本、采样设置与无法解释的剩余部分。

\end{enumerate}
\section{自测与解析}
\noindent\textbf{1. /cpu/classes/total:cpu-seconds 能否当作进程实际 CPU 消耗？}
\begin{enumerate}[label=\Alph*.]
\item 能，它就是内核 task-clock
\item 不能，它是由 GOMAXPROCS 积分得到的估计可用 CPU 总量
\item 只能在单核 Linux 上精确相等
\item 它表示全部请求延迟
\end{enumerate}
\noindent{\color{forest}\textbf{答案 B。}}该家族为运行时估计分类，total 包含可用容量中的 idle；官方要求不要直接与系统 CPU 时间比较。

\noindent\textbf{2. 保存 Float64Histogram 指针作为上次快照后再次 metrics.Read，主要风险是什么？}
\begin{enumerate}[label=\Alph*.]
\item 直方图单位变成字节
\item 底层存储被复用，旧快照被覆盖或并发访问产生数据竞争
\item 自动清空全部 runtime 指标
\item Go 自动深拷贝，因此没有风险
\end{enumerate}
\noindent{\color{forest}\textbf{答案 B。}}Read 可能复用指针型 Value 的存储；差分保留旧样本或跨 goroutine 传递时，需要深拷贝所需数据。

\noindent\textbf{3. strace 中 epoll\_wait 的 -T 显示 3 秒，说明什么？}
\begin{enumerate}[label=\Alph*.]
\item CPU 执行内核代码 3 秒
\item 系统调用从进入到返回经过约 3 秒墙钟时间
\item 所有 goroutine 都等待 3 秒
\item CPU profile 应出现同样 3 秒
\end{enumerate}
\noindent{\color{forest}\textbf{答案 B。}}阻塞调用的墙钟持续时间包含睡眠，不能直接当作计算消耗或所有 goroutine 的等待。

\noindent\textbf{4. eBPF 线程 off-CPU 数据为何不能直接等同 Go 请求等待？}
\begin{enumerate}[label=\Alph*.]
\item eBPF 不支持 Linux
\item Go 存在 goroutine 到线程的用户态调度与迁移
\item Go 程序从不使用线程
\item 所有 Go 等待都在内核完成
\end{enumerate}
\noindent{\color{forest}\textbf{答案 B。}}一个线程可承载多个 goroutine，用户态阻塞后也可能运行其他工作；需要明确关联机制才能做请求级归因。

\section{分析题与参考答案}
\noindent\textbf{题 1：}在同一未重启进程的 5 秒窗口内，累计分配增加 500 MiB，锁等待指标增加 40 秒，GOMAXPROCS 固定为 4，CPU classes total 增加约 20 秒。分别解释这些数值；若估计平均同时等待人数，还需什么前提？哪些数值不能直接相加？

\noindent\textbf{参考答案：}分配速率约 100 MiB/s。40 秒锁等待是多个 goroutine 等待时间的累计估计，覆盖用户同步锁及运行时内部锁；仅当估计误差与跨窗口等待的记账影响可忽略时，40/5 才近似对应平均 8 个 goroutine 同时等待。若有很长的未完成等待或窗口前已开始的等待，应像第 07 章一样检查边界，不能直接解释为精确人数。CPU classes total 的 20 CPU 秒是 4\ensuremath{\times}5 秒的估计可用容量，不是已使用 20 CPU 秒。分配、等待与可用 CPU 属于不同量纲/对象集合，不能相加，也不能据此直接得到请求耗时或 OS CPU 利用率。

\noindent\textbf{题 2：}假设 runtime.MemStats.HeapInuse 稳定、进程 RSS 持续上升，CPU 与 goroutine 数稳定。提出三种待验证解释及对应工具。

\noindent\textbf{参考答案：}可能是运行时保留但尚未归还的空闲堆：比较 /memory/classes/heap/free:bytes、/memory/classes/heap/released:bytes 与 scavtrace；可能是 cgo 库分配：检查库自身统计、原生分配跟踪和系统映射；可能是 mmap 或其他映射驻留增长：检查目标进程 smaps/映射及相关系统调用。先排除口径和采样窗口差异，再定位所有权，不应仅根据 RSS 宣布 Go 堆泄漏。

\medskip\noindent\textbf{本章主要阅读：}\cite{tools-metrics}\cite{tools-metrics-read}\cite{tools-runtime-debug}\cite{tools-go125}\cite{tools-http-pprof}\cite{tools-perf}\cite{tools-strace}\cite{tools-ebpf}\cite{tools-runtime-debug-api}\cite{tools-fgprof}\cite{tools-fgprof-source}\cite{tools-gops}\cite{tools-gops-agent}\cite{tools-expvar}\cite{tools-nm}\cite{tools-objdump}。
\part*{第三篇 · 优化与工程实践}
\addcontentsline{toc}{part}{第三篇 · 优化与工程实践}
\chapter{优化决策与 PGO}
\label{ch:11}
\noindent{\large\color{forest}用证据选择优化，理解编译器反馈怎样转化为收益。}\par\medskip
\noindent\textbf{先修要求：}第 03–08 章；Go 编译与接口调用基础。

\noindent\textbf{完成本章后，应能够：}
\begin{itemize}
\item 根据瓶颈类别选择算法、内存、同步或编译优化
\item 建立代表性 PGO profile 和可重复构建流程
\item 解释内联、逃逸和去虚化之间的联系
\item 评估分布漂移、收益退化与回滚条件
\end{itemize}
\section{先减工作量，再优化执行方式}
\begin{figure}[H]
\centering
\begin{tikzpicture}[x=1cm,y=1cm]
\node[anchor=east,font=\small] at (1,0.0) {泳道1};
\draw[linegray] (1.2,-0.35) -- (14.799999999999999,-0.35);
\node[anchor=east,font=\small] at (1,-1.7) {泳道2};
\draw[linegray] (1.2,-2.05) -- (14.799999999999999,-2.05);
\path[fill=forest!60] (1.200,-0.250) rectangle (9.360,0.250);
\node[align=center,font=\small,text width=8.06cm] at (5.280,0.750) {原其他 60};
\path[fill=leaf!60] (9.360,-0.250) rectangle (14.800,0.250);
\node[align=center,font=\small,text width=5.34cm] at (12.080,0.750) {原热点 40};
\path[fill=forest!60] (1.200,-1.950) rectangle (9.360,-1.450);
\node[align=center,font=\small,text width=8.06cm] at (5.280,-0.950) {新其他 60};
\path[fill=leaf!60] (9.360,-1.950) rectangle (12.080,-1.450);
\node[align=center,font=\small,text width=2.62cm] at (10.720,-0.950) {新热点 20};
\draw[-{Stealth[length=2mm]},forest] (1.2,-2.8499999999999996) -- (14.999999999999998,-2.8499999999999996);
\draw[forest] (1.20,-2.8499999999999996) -- (1.20,-2.9499999999999997);
\node[below,font=\scriptsize] at (1.20,-2.9499999999999997) {0};
\draw[forest] (3.92,-2.8499999999999996) -- (3.92,-2.9499999999999997);
\node[below,font=\scriptsize] at (3.92,-2.9499999999999997) {20};
\draw[forest] (6.64,-2.8499999999999996) -- (6.64,-2.9499999999999997);
\node[below,font=\scriptsize] at (6.64,-2.9499999999999997) {40};
\draw[forest] (9.36,-2.8499999999999996) -- (9.36,-2.9499999999999997);
\node[below,font=\scriptsize] at (9.36,-2.9499999999999997) {60};
\draw[forest] (12.08,-2.8499999999999996) -- (12.08,-2.9499999999999997);
\node[below,font=\scriptsize] at (12.08,-2.9499999999999997) {80};
\draw[forest] (14.80,-2.8499999999999996) -- (14.80,-2.9499999999999997);
\node[below,font=\scriptsize] at (14.80,-2.9499999999999997) {100};
\node[anchor=east,font=\scriptsize] at (14.799999999999999,-3.5999999999999996) {时间 / 归一化时间};
\end{tikzpicture}
\caption{热点速度翻倍，整体为何只加速 1.25 倍}\label{fig:11-amdahl}
\begin{minipage}{0.96\linewidth}\footnotesize 教学示意：用于解释机制，不是运行时的实测数据或完整实现。

\textbf{读图顺序：}1. 划分原始成本\quad 2. 只加速热点\quad 3. 比较整体\quad 4. 前提不能省略。
\textbf{机制依据：}\cite{go-pgo}。
\end{minipage}
\end{figure}
常见优化顺序是：消除不必要工作 \ensuremath{\rightarrow}\allowbreak{} 选择更合适的算法与数据结构 \ensuremath{\rightarrow}\allowbreak{} 减少数据移动与分配 \ensuremath{\rightarrow}\allowbreak{} 处理共享资源竞争 \ensuremath{\rightarrow}\allowbreak{} 微调局部实现和编译选项。这个顺序并非硬性规定；它用于根据潜在收益和维护成本确定排查优先级。

\begingroup\small\setlength{\tabcolsep}{4pt}
\begin{longtable}{>{\raggedright\arraybackslash}p{0.298\linewidth}>{\raggedright\arraybackslash}p{0.298\linewidth}>{\raggedright\arraybackslash}p{0.298\linewidth}}
\toprule
\textbf{证据} & \textbf{候选干预} & \textbf{验证与风险} \\
\midrule
\endfirsthead
\toprule
\textbf{证据} & \textbf{候选干预} & \textbf{验证与风险} \\
\midrule
\endhead
重复编码占大量 CPU & 复用结果、标准库编码器 & 缓存失效与输入变化 \\
\addlinespace[3pt]
相同工作量下 alloc\_objects 增量很高 & 预分配、避免临时字符串 & 对象存活时间与峰值内存 \\
\addlinespace[3pt]
锁等待集中 & 缩小临界区、分片、批处理 & 顺序和一致性语义 \\
\addlinespace[3pt]
runnable 延迟高 & 削减并发、检查配额 & 吞吐与公平性 \\
\addlinespace[3pt]
下游等待高 & 截止时间、并发预算、批请求 & 重试放大和错误率 \\
\addlinespace[3pt]
稳定 CPU 热点 & PGO 及编译器优化 & profile 代表性与代码体积 \\
\addlinespace[3pt]
\bottomrule
\end{longtable}
\endgroup
\texttt{sync.\allowbreak{}Pool}可以复用部分临时对象，减少重复分配，但不提供持久缓存保证：池中的对象可随时被自动移除。采用对象池前，应评估对象大小、复用率、敏感内容的清理方式和并发访问成本。零拷贝可能延长底层大缓冲区的生命周期，因此应同时测量 CPU 收益与堆内存保留量。\cite{go-sync-pool}

\section{Go PGO 的输入、位置与开关}
\begin{figure}[H]
\centering
\begin{tikzpicture}[x=1cm,y=1cm]
\node[draw=linegray!75!black,fill=linegray!13,rounded corners=3pt,text width=2.58cm,minimum height=1.12cm,align=center,font=\small] (p0-traffic) at (1.50,0.00) {代表性运行};
\node[draw=linegray!75!black,fill=linegray!13,rounded corners=3pt,text width=2.58cm,minimum height=1.12cm,align=center,font=\small] (p0-profile) at (4.50,0.00) {CPU profile};
\node[draw=linegray!75!black,fill=linegray!13,rounded corners=3pt,text width=2.58cm,minimum height=1.12cm,align=center,font=\small] (p0-compiler) at (7.50,0.00) {编译器决策};
\node[draw=linegray!75!black,fill=linegray!13,rounded corners=3pt,text width=2.58cm,minimum height=1.12cm,align=center,font=\small] (p0-binary) at (10.50,0.97) {候选二进制};
\node[draw=leaf!75!black,fill=leaf!13,rounded corners=3pt,text width=2.58cm,minimum height=1.12cm,align=center,font=\small] (p0-eval) at (13.50,0.00) {独立负载评价};
\node[draw=linegray!75!black,fill=linegray!13,rounded corners=3pt,text width=2.58cm,minimum height=1.12cm,align=center,font=\small] (p0-fallback) at (10.50,-0.97) {保留通用路径};
\draw[-{Stealth[length=2mm]},forest!70!black] (p0-traffic.east) -- (p0-profile.west);
\draw[-{Stealth[length=2mm]},forest!70!black] (p0-profile.east) -- (p0-compiler.west);
\draw[-{Stealth[length=2mm]},forest!70!black] (p0-compiler.east) -- (p0-binary.west);
\draw[-{Stealth[length=2mm]},forest!70!black] (p0-compiler.east) -- (p0-fallback.west);
\draw[-{Stealth[length=2mm]},forest!70!black] (p0-binary.east) -- (p0-eval.west);
\draw[-{Stealth[length=2mm]},forest!70!black] (p0-eval.north) -- (13.50,2.12) -- (1.50,2.12) -- (p0-traffic.north);
\end{tikzpicture}
\caption{PGO 是有回退的反馈优化闭环}\label{fig:11-pgo}
\begin{minipage}{0.96\linewidth}\footnotesize 教学示意：用于解释机制，不是运行时的实测数据或完整实现。

\textbf{读图顺序：}1. 采集代表性路径\quad 2. 只构建成功不够\quad 3. 常见路径专门化\quad 4. 训练评价分开。
\textbf{连线语义：}编译器决策\ensuremath{\rightarrow}\allowbreak{}保留通用路径：条件去虚化；独立负载评价\ensuremath{\rightarrow}\allowbreak{}代表性运行：更新代表性。
\textbf{机制依据：}\cite{go-pgo}\cite{go-pgo-blog}。
\end{minipage}
\end{figure}
剖析引导优化（profile-guided optimization，PGO）将代表性运行的 CPU profile 作为编译输入，辅助决定哪些调用值得内联、哪些动态调用目标值得生成专门路径。Go 1.21 正式支持 PGO。编译器需要符合要求的 CPU 剖析数据，不能用 heap 或 mutex profile 代替。\cite{go-pgo}\cite{go-pgo-blog}

\begin{codeblock}
# main 包目录中放置代表性的 CPU profile
cp representative.cpu.pb.gz default.pgo
go build -o app.pgo .
go build -pgo=off -o app.base .
# 或显式指定输入，避免依赖自动发现
go build -pgo=representative.cpu.pb.gz -o app.pgo .
go version -m app.pgo
\end{codeblock}

默认 \texttt{-\allowbreak{}pgo=auto} 会在 main 包目录中查找 default.pgo。构建多个 main 包时，应分别核对各主包所选用的 profile。要实现可复现构建，应将 profile 及其采集来源、工具链、工作负载比例和更新策略一同纳入版本管理。不要把含敏感标签的生产 profile 未经检查提交到公开仓库。

PGO 不改变 Go 程序的语义契约，但会改变生成代码、热点归属和内联形态。它不能解决网络下游瓶颈、无限增长缓存或算法复杂度不合理等问题。

输入语义比文件扩展名更重要。Go 的 PGO 要求样本代表 CPU 时间；fgprof 的 wall-clock 聚合含等待，不能当作等价训练数据。v0.9.5 虽导出 samples/count，仍不代表 CPU 样本，且其导出函数记录没有设置 Go PGO 所需的 function.start\_line。不能靠格式能否解析来替代采集语义与源码元数据验证。\cite{tools-fgprof-source}

\section{内联、逃逸与条件去虚化}
内联把被调用函数体放入调用者上下文，可能消除调用成本，也可能使常量传播、死代码消除与逃逸分析看到更多信息。收益有时主要来自这些后续优化，而不仅是省去一次函数调用。过度内联也可能增大代码体积和指令缓存压力。

条件去虚化根据 profile 中的常见动态目标产生快路径，并保留其他目标的通用回退路径。例如接口调用大多数落在某具体实现，编译器可检查该类型后走直接调用；profile 不是“以后永远只会调用这个类型”的保证，所以正确的回退不可省略。

Go 官方 PGO 案例解释了更积极内联促使某些对象不再逃逸，从而减少分配的机制；这是一个具体案例，不应写成所有程序都获得相同收益。\cite{go-pgo-blog} 分析 PGO 结果时可对照编译诊断、alloc profile 和 CPU 差分，寻找跨层证据。

\begin{codeblock}
go build -pgo=default.pgo -gcflags='-m=2' . 2> pgo-compile.txt
go build -pgo=off -gcflags='-m=2' . 2> base-compile.txt
\end{codeblock}

编译器诊断文本不是稳定 API。对关键词做报告时保留完整工具链版本，并避免把单个 inlining 提示当作最终性能证明。

\section{采集与合并：权重就是工作负载假设}
profile 要覆盖真实重要路径、代表性输入、请求类型和时间段。只采到冷启动、单租户、单一热点或错误风暴，可能把编译器引向不代表常态的代码。应将线上请求分布与训练 profile 构成对照，并将高价值但稀少路径纳入回归测试。

\begin{codeblock}
go tool pprof -proto cpu-a.pb.gz cpu-b.pb.gz > merged.cpu.pb.gz
\end{codeblock}

在类型和单位兼容的前提下，pprof 合并会累加样本权重，不会自动按采集时长或请求数求平均。CPU 密集实例或采集更长的窗口会产生更多权重。如果希望不同业务按特定比例影响编译，就必须明确缩放策略和输入来源；不能因为每种业务各提供一个文件，就认为它们对优化的影响相同。\cite{go-pgo}

代表性还涉及空间与时间：不同地区、租户、硬件、特性开关、数据尺寸和版本都可能造成分布漂移。profile 随代码变化可能逐渐变旧；Go 有一定的源码变化容忍能力，但无法免除匹配率、收益和新路径验证。

\section{基准与灰度：把训练和评价分开}
至少设置三组构建：禁用 PGO 的基线、使用候选 profile 的构建，以及使用偏置 profile 的构建。评价负载应与训练采集分离，避免只证明编译器适应了训练样本。保持相同源代码、依赖、工具链和输入批次，使用第 03 章统计协议比较 CPU/请求、延迟、内存和代码大小。

之后在受控灰度中验证：按固定比例或独立实例组运行，记录配置和流量差异；同时监视错误率、p99、GC CPU 和超时。若基线随时间变化，交替窗口或配对实例可以减轻漂移，但不能完全消除不同用户分配造成的混杂。

配套实验包含 PGO 构建冒烟测试，即检查构建流程能否完成的最小测试。它不能仅凭退出码证明 profile 包含可用热点，或证明编译器已据此实施优化。Go 1.25 编译器接受并忽略空文件或无样本的 profile；构建信息记录 -pgo 路径也不证明有优化发生。需另查 profile 非空、CPU sample type、有效源码匹配、编译诊断，并以独立性能实验判断收益。小示例没有显著收益，并不意味着 PGO 对所有程序都无效；同样，构建冒烟测试也不能替代生产收益评估。正式结论必须来自独立工作负载的重复实验。

\section{优化停止规则与维护预算}
一个可交付优化应同时满足：正确性通过；指标达到预设阈值；收益大于噪声且可重复；维护与内存成本可接受；有观测和回滚路径。应在实验开始前定义停止规则，避免不断挑选有利运行直到得到想要的数字。

当热点只占很小比例、干预复杂度很高、profile 样本不足或负载代表性无法保证时，先改善证据质量可能比继续改代码更有价值。持续剖析的作用之一就是让这些决策有长期基线支持，而非让工程师不停地优化排名第一的函数。

\section{实验任务}
\subsection*{构建可审计的 PGO 对照}
\begin{enumerate}
\item 在 lab 中运行 CPU 基准并生成 profile，分别构建禁用 PGO 的基线，以及使用该 profile 的候选版本。

\item 保存构建信息和 profile 校验值；检查 CPU 样本、有效权重与源码匹配，并结合编译诊断判断 profile 是否实际参与了优化。仅构建成功不算通过这一检查。

\item 使用不同于训练样本的输入长度做重复评价；不预设 PGO 一定获益。

\item 故意用偏向某一种工作负载的 profile 再构建，对比另一类输入上的变化。

\item 提交收益、无收益或退化结果，并解释 profile 代表性和统计限制。

\end{enumerate}
\section{自测与解析}
\noindent\textbf{1. Go 编译器 PGO 使用什么主要输入？}
\begin{enumerate}[label=\Alph*.]
\item heap profile
\item CPU pprof profile
\item goroutine 文本栈
\item 任意 trace 文件
\end{enumerate}
\noindent{\color{forest}\textbf{答案 B。}}Go PGO 使用 CPU profile；格式相似不能替代采集语义。

\noindent\textbf{2. 合并两份 profile 默认是否意味着两种业务等权？}
\begin{enumerate}[label=\Alph*.]
\item 是
\item 否，主要按样本权重累加
\item 只有文件名相同时才等权
\item 完全随机
\end{enumerate}
\noindent{\color{forest}\textbf{答案 B。}}采样时间和 CPU 消耗影响总权重。

\noindent\textbf{3. PGO 构建冒烟测试通过，能够证明什么？}
\begin{enumerate}[label=\Alph*.]
\item 生产加速 10\%
\item 该构建命令完成，仍须验证 profile 和优化是否生效
\item 所有路径都优化了
\item 不再需要基准
\end{enumerate}
\noindent{\color{forest}\textbf{答案 B。}}空文件或无有效样本也可能被编译器接受并忽略；成功构建不证明使用了有效热点，更不证明获得性能收益。

\section{分析题与参考答案}
\noindent\textbf{题 1：}业务 A 占请求 90\%但每次 1 ms CPU，B 占 10\%但每次 20 ms CPU。CPU 权重比例是多少？

\noindent\textbf{参考答案：}以 100 请求计算 A=90 ms、B=200 ms，因此 CPU 权重约 31\%和 69\%。按请求比例加权与按 CPU 成本加权回答不同问题，训练目标必须说明。

\noindent\textbf{题 2：}PGO 后 allocs/op 下降而热点函数调用次数未减少，可能是什么机制？

\noindent\textbf{参考答案：}更积极内联可能改善逃逸分析，使对象从堆分配转为栈或被消除。需对照编译诊断和 alloc profile，不能仅凭最终数值认定唯一原因。

\medskip\noindent\textbf{本章主要阅读：}\cite{go-pgo}\cite{go-pgo-blog}\cite{tools-fgprof-source}\cite{go-sync-pool}。
\chapter{生产环境持续剖析}
\label{ch:12}
\noindent{\large\color{forest}让短时诊断升级为可控成本、可检索和可治理的长期证据。}\par\medskip
\noindent\textbf{先修要求：}第 04–10 章；HTTP、容器和基础可观测性知识。

\noindent\textbf{完成本章后，应能够：}
\begin{itemize}
\item 设计从 agent 到符号存储再到查询的完整链路
\item 制定采样、基数、访问控制与开销预算
\item 关联 profile 与 metrics/logs/traces 而不混淆语义
\item 诊断数据丢失、符号缺失与采集偏差
\end{itemize}
\section{持续剖析的体系结构}
\begin{figure}[H]
\centering
\begin{tikzpicture}[x=1cm,y=1cm]
\node[draw=linegray!75!black,fill=linegray!13,rounded corners=3pt,text width=2.08cm,minimum height=1.12cm,align=center,font=\small] (p0-app) at (1.25,0.00) {目标进程};
\node[draw=linegray!75!black,fill=linegray!13,rounded corners=3pt,text width=2.08cm,minimum height=1.12cm,align=center,font=\small] (p0-agent) at (3.75,0.00) {采集/\allowbreak{}展开};
\node[draw=linegray!75!black,fill=linegray!13,rounded corners=3pt,text width=2.08cm,minimum height=1.12cm,align=center,font=\small] (p0-ingest) at (6.25,0.00) {校验/\allowbreak{}摄取};
\node[draw=linegray!75!black,fill=linegray!13,rounded corners=3pt,text width=2.08cm,minimum height=1.12cm,align=center,font=\small] (p0-store) at (8.75,0.00) {列/\allowbreak{}符号存储};
\node[draw=leaf!75!black,fill=leaf!13,rounded corners=3pt,text width=2.08cm,minimum height=1.12cm,align=center,font=\small] (p0-query) at (11.25,0.00) {过滤/\allowbreak{}聚合};
\node[draw=leaf!75!black,fill=leaf!13,rounded corners=3pt,text width=2.08cm,minimum height=1.12cm,align=center,font=\small] (p0-view) at (13.75,0.00) {火焰图};
\draw[-{Stealth[length=2mm]},forest!70!black] (p0-app.east) -- (p0-agent.west);
\draw[-{Stealth[length=2mm]},forest!70!black] (p0-agent.east) -- (p0-ingest.west);
\draw[-{Stealth[length=2mm]},forest!70!black] (p0-ingest.east) -- (p0-store.west);
\draw[-{Stealth[length=2mm]},forest!70!black] (p0-store.east) -- (p0-query.west);
\draw[-{Stealth[length=2mm]},forest!70!black] (p0-query.east) -- (p0-view.west);
\end{tikzpicture}
\caption{持续剖析中的数据流与丢失位置}\label{fig:12-continuous}
\begin{minipage}{0.96\linewidth}\footnotesize 教学示意：用于解释机制，不是运行时的实测数据或完整实现。

\textbf{读图顺序：}1. 采集不是完整观察\quad 2. 校验语义\quad 3. 有界缓冲与存储\quad 4. 查询会再次选样。
\textbf{连线语义：}采集/展开\ensuremath{\rightarrow}\allowbreak{}校验/摄取：样本。
\textbf{机制依据：}\cite{proj-pyroscope}\cite{proj-parca-agent}。
\end{minipage}
\end{figure}
持续剖析周期性或连续地采集程序的资源使用数据，并通过查询按服务、版本、环境和时间窗口聚合调用栈及其权重。典型链路为：进程内或 eBPF 采集器 \ensuremath{\rightarrow}\allowbreak{} agent 批处理 \ensuremath{\rightarrow}\allowbreak{} 摄取与验证 \ensuremath{\rightarrow}\allowbreak{} 栈/符号归并 \ensuremath{\rightarrow}\allowbreak{} 长期存储 \ensuremath{\rightarrow}\allowbreak{} 查询聚合 \ensuremath{\rightarrow}\allowbreak{} 火焰图/差分/告警。第 15 章比较 Pyroscope 与 Parca 的真实实现。

架构首先解决采什么，而不是先选可视化界面。CPU、分配、存活堆、锁等待和 off-CPU 数据的单位与时间边界不同，不能合并成一个没有类型的“性能值”。采集器也不等于符号化器；某些环境在 agent 展开栈，在服务端补符号，两端版本和构建元数据都影响可靠性。

连续并不意味着 100\%捕获每次事件。低开销统计采样牺牲单次请求完整性，换取长期趋势。发现异常后，可再用定向追踪、飞行记录器或调试器获取局部细节，使不同工具提供的证据相互补充。

\section{采集策略：pull、push 与 eBPF}
pull 模式由中心或 agent 定时请求 HTTP profile，便于沿用 net/http/pprof；push 模式由 SDK 控制采集与上传，便于穿越网络边界和补业务标签；eBPF 模式常从主机采样线程并关联容器，减少应用显式埋点，但会受到内核权限、栈展开、运行时识别和平台支持限制。

\begingroup\small\setlength{\tabcolsep}{4pt}
\begin{longtable}{>{\raggedright\arraybackslash}p{0.298\linewidth}>{\raggedright\arraybackslash}p{0.298\linewidth}>{\raggedright\arraybackslash}p{0.298\linewidth}}
\toprule
\textbf{方式} & \textbf{主要优势} & \textbf{主要代价} \\
\midrule
\endfirsthead
\toprule
\textbf{方式} & \textbf{主要优势} & \textbf{主要代价} \\
\midrule
\endhead
进程内 HTTP 采集 & 标准接口、Go 语义直接 & 端点保护、采集调度冲突 \\
\addlinespace[3pt]
应用 SDK push & 生命周期与标签可控 & 依赖和运行时开销 \\
\addlinespace[3pt]
主机 eBPF 采样 & 跨服务覆盖、集中治理 & 内核权限、平台与展开限制 \\
\addlinespace[3pt]
事件追踪 & 提供调度与因果上下文 & 事件量、时间窗与分析复杂度 \\
\addlinespace[3pt]
\bottomrule
\end{longtable}
\endgroup
Go 原生\texttt{runtime/\allowbreak{}pprof}在同一进程中只允许一个活动的 CPU 采集，因此应指定负责启动和停止它的唯一控制器，并协调临时诊断请求。这一限制不等于所有外部采集器都无法同时运行；例如 eBPF 采样可能与 Go 原生采样并存，但仍须评估二者的开销和测量语义。\cite{tools-pprof}多来源生成相同 pprof 格式也不自动等价：采样频率、丢失、栈深、cgo 展开和内核帧覆盖都应单独校验。

\section{标签基数、数据量与预算}
稳定标签建议包括 service、environment、version、region 和有限集合的 workload；不要将 user\_id、完整 URL、SQL 文本或随机 request\_id 直接作为长期聚合维度。高基数不但增加索引和存储成本，也可能泄露敏感信息。

教学容量模型：若 N 个实例每分钟产生 B 字节压缩 profile，保留 D 天，则在忽略副本、索引和压缩率变化时，载荷总量约为 \texttt{N \ensuremath{\times} B \ensuremath{\times} 1440 \ensuremath{\times} D}。假设 1000 实例、每分钟 100 KiB、保留 7 天，已压缩 profile 载荷总量约 961.3 GiB；真实值要用自己的样本测量，不能拿此示例作为产品存储承诺。

预算至少包括 agent CPU、应用吞吐变化、内存缓冲上限、网络带宽、摄取队列、存储增长、查询延迟和损失率。应在相同工作量下比较开启与关闭采集时的开销。不能只统计 agent 进程的 CPU 消耗，因为进程内采集的成本会记在被测应用上。

\section{关联与检索：相关性仍需上下文}
\begin{figure}[H]
\centering
\begin{tikzpicture}[x=1cm,y=1cm]
\path[fill=leaf!25] (0.00,0.00) rectangle (13.97,0.78);
\node[align=center,text width=13.92cm,font=\small] at (7.00,0.40) {总 CPU 样本};
\path[fill=bluegray!38] (0.00,1.05) rectangle (12.57,1.83);
\node[align=center,text width=12.52cm,font=\small] at (6.30,1.45) {已知函数};
\path[fill=bluegray!51] (2.80,2.10) rectangle (12.58,2.88);
\node[align=center,text width=9.72cm,font=\small] at (7.70,2.50) {encode};
\path[fill=leaf!38] (12.60,1.05) rectangle (13.97,1.83);
\node[font=\scriptsize] at (13.30,1.45) {3};
\node[anchor=north west,align=left,text width=14cm,font=\scriptsize] at (0,-.45) {1=总 CPU 样本: flat 0 /\allowbreak{} cum 100；2=已知函数: flat 20 /\allowbreak{} cum 90；3=unknown: flat 10 /\allowbreak{} cum 10；4=encode: flat 70 /\allowbreak{} cum 70};
\end{tikzpicture}
\caption{符号恢复会改变归因，却未必改变总成本}\label{fig:12-unknown}
\begin{minipage}{0.96\linewidth}\footnotesize 构造的归因示例；known 节点自身 flat 为 20，encode 是其子调用，不能将两者 cum 重复相加。

\textbf{读图顺序：}1. 符号不完整\quad 2. 修复符号匹配\quad 3. 比较同等可见性。
\textbf{机制依据：}\cite{proj-pprof}。
\end{minipage}
\end{figure}
指标帮助发现资源或服务状态异常；profile 反映选定窗口内的成本分布；分布式 trace 展示请求路径和各 span 的时间关系，执行追踪还可揭示调度与等待状态；日志补充具体事件。有效关联需要一致的服务身份、版本、时间范围以及明确的时钟和聚合粒度。profile 窗口为 60 秒而 trace 只包含某一次请求时，不能直接宣布窗口内最大栈就是该请求根因。

把 trace ID 附着到每个样本并长期索引可能导致基数爆炸。可以考虑有界的代表性样本关联（exemplar），或只在临时调查中建立关联；采用哪种方式，仍需结合数据模型、覆盖率与隐私要求评估。Go 原生 pprof labels 也不保证所有 profile 类型均支持标签传播。

符号化失败会把成本推入 unknown，从而使已知函数占比看似改变。查询界面应展示未知帧、截断栈、丢样和过滤总量，并能定位到构建 ID。健康度先于美观：一张信息缺失的漂亮火焰图不能作为优化依据。

\section{安全、隔离与故障模式}
诊断端点独立监听在回环或管理网络，通过已有认证代理、网络策略和审计控制访问。profile 与 core 按敏感诊断数据管理，设置最小留存与按角色查询。尤其要避免公开暴露 cmdline、goroutine、trace 等端点。

采集器要有超时、并发上限、退避、队列边界与丢弃策略；后端故障不能反向拖垮业务。短生命周期容器可能在一个采样窗口结束前就退出；如果退出前未能完成缓冲数据上传（flush），这类实例还会进一步缺失。两种情况都可能造成代表性偏差。丢失的数据如果集中在 CPU 高峰或崩溃前，长期查询会系统性低估最需要观察的区间。

常见故障排查顺序是：目标发现 \ensuremath{\rightarrow}\allowbreak{} 权限/连接 \ensuremath{\rightarrow}\allowbreak{} 采集是否生成有效数据 \ensuremath{\rightarrow}\allowbreak{} 解析单位与标签 \ensuremath{\rightarrow}\allowbreak{} 栈展开与符号 \ensuremath{\rightarrow}\allowbreak{} 摄取排队和丢弃 \ensuremath{\rightarrow}\allowbreak{} 存储和查询过滤。不要因为页面为空就立即更改采样率。

\section{上线与质量验证协议}
小规模试点先覆盖不同 Go 版本、架构、容器限制和 cgo 使用。以受控 CPU 热点、内存保留和锁竞争注入验证工具能否看见预期机制；同时设计负例，例如纯等待不应被误判为 CPU 热点。

扩大范围前定义验收条件：可识别帧比例、预期热点排名稳定性、开销上界、未知数据可见性、服务隔离与访问审计。逐步扩大覆盖范围，并保留能够停用采集的开关。在常规发布流水线保留符号、构建信息和 PGO profile 来源，形成性能基线与代码版本之间的长期连接。

诊断系统自身也应设置服务等级目标（SLO）：从发现目标到数据可查询的延迟、查询延迟的 p95、agent 离线比例、数据完整性，以及留存策略的达成情况。没有这些观测，所谓“连续”可能只是间歇性采到幸存实例。

\section{实验任务}
\subsection*{为一个 20 实例服务设计剖析方案}
\begin{enumerate}
\item 选择 CPU 和一种内存指标，定义采样周期、唯一采集所有者与错误处理。

\item 列出有限集合标签和明确禁止的高基数标签，并估算 7 天数据量。

\item 设计 A/B 开销实验与至少三个故障注入：上传失败、符号缺失、实例短生命周期。

\item 说明如何从 CPU 异常跳转到相同版本和时间段的 profile，再用 trace 验证等待假设。

\item 写出诊断系统回滚、数据留存、权限和审计方案。

\end{enumerate}
\section{自测与解析}
\noindent\textbf{1. profile 窗口和某条 trace 重合能证明什么？}
\begin{enumerate}[label=\Alph*.]
\item 最大热点一定造成该请求慢
\item 仅提供关联线索，仍需确认上下文与路径
\item 两种数据完全等价
\item 无需版本标识
\end{enumerate}
\noindent{\color{forest}\textbf{答案 B。}}窗口聚合与单请求时序的粒度不同。

\noindent\textbf{2. 哪个标签最适合长期聚合？}
\begin{enumerate}[label=\Alph*.]
\item 随机 request\_id
\item 完整 SQL
\item 有限集合的 workload 类别
\item 用户手机号
\end{enumerate}
\noindent{\color{forest}\textbf{答案 C。}}稳定低基数标签兼顾可查询性、成本与隐私。

\noindent\textbf{3. eBPF 采集意味着什么？}
\begin{enumerate}[label=\Alph*.]
\item 完全无开销、全语言同等支持
\item 减少某些应用改动，但仍有平台与栈展开限制
\item 一定比进程内 profile 准确
\item 不需要权限
\end{enumerate}
\noindent{\color{forest}\textbf{答案 B。}}eBPF 能力受内核、硬件、运行时和采集实现约束。

\section{分析题与参考答案}
\noindent\textbf{题 1：}某 agent 只在低负载时成功上传，高峰时大量丢样。按可用样本求平均会如何偏差？

\noindent\textbf{参考答案：}缺失与目标变量相关，属于非随机缺失；结果会低估高峰成本。应监视丢弃率和队列，对高峰做独立受控采集，不能靠简单补零修复。

\noindent\textbf{题 2：}某函数占比从 10\%升到 20\%，同时 unknown 帧大幅下降，能直接认定回归吗？

\noindent\textbf{参考答案：}不能。符号化改善可能把原未知样本重新归属该函数。需固定符号状态、比较绝对成本及业务工作量，并检查构建 ID 和展开质量。

\medskip\noindent\textbf{本章主要阅读：}\cite{go-diagnostics}\cite{go-http}\cite{go-pprof}\cite{tools-pprof}。
\chapter{综合案例与证据链}
\label{ch:13}
\noindent{\large\color{forest}把工具组合成可证伪的诊断过程。}\par\medskip
\noindent\textbf{先修要求：}第 01–12 章；能够运行附带实验模块。

\noindent\textbf{完成本章后，应能够：}
\begin{itemize}
\item 为 CPU、分配、保留、锁等待与下游等待建立鉴别诊断
\item 避免根据单一 profile 过早下结论
\item 区分正证据、负证据和未观测项
\item 完成包含正确性与性能验证的调查报告
\end{itemize}
\section{案例 A：编码路径消耗 CPU}
\begin{figure}[H]
\centering
\begin{tikzpicture}[x=1cm,y=1cm]
\node[draw=linegray!75!black,fill=linegray!13,rounded corners=3pt,text width=3.33cm,minimum height=1.12cm,align=center,font=\small] (p0-slow) at (1.88,0.00) {请求变慢};
\node[draw=linegray!75!black,fill=linegray!13,rounded corners=3pt,text width=3.33cm,minimum height=1.12cm,align=center,font=\small] (p0-highcpu) at (5.62,0.97) {CPU/\allowbreak{}工作↑};
\node[draw=linegray!75!black,fill=linegray!13,rounded corners=3pt,text width=3.33cm,minimum height=1.12cm,align=center,font=\small] (p0-lowcpu) at (5.62,-0.97) {CPU 不高};
\node[draw=linegray!75!black,fill=linegray!13,rounded corners=3pt,text width=3.33cm,minimum height=1.12cm,align=center,font=\small] (p0-alloc) at (9.38,0.97) {分配/\allowbreak{}GC 证据};
\node[draw=linegray!75!black,fill=linegray!13,rounded corners=3pt,text width=3.33cm,minimum height=1.12cm,align=center,font=\small] (p0-waiting) at (9.38,-0.97) {等待/\allowbreak{}排队证据};
\node[draw=leaf!75!black,fill=leaf!13,rounded corners=3pt,text width=3.33cm,minimum height=1.12cm,align=center,font=\small] (p0-verify) at (13.12,0.00) {同负载干预};
\draw[-{Stealth[length=2mm]},forest!70!black] (p0-slow.east) -- (p0-highcpu.west);
\draw[-{Stealth[length=2mm]},forest!70!black] (p0-slow.east) -- (p0-lowcpu.west);
\draw[-{Stealth[length=2mm]},forest!70!black] (p0-highcpu.east) -- (p0-alloc.west);
\draw[-{Stealth[length=2mm]},forest!70!black] (p0-lowcpu.east) -- (p0-waiting.west);
\draw[-{Stealth[length=2mm]},forest!70!black] (p0-alloc.east) -- (p0-verify.west);
\draw[-{Stealth[length=2mm]},forest!70!black] (p0-waiting.east) -- (p0-verify.west);
\end{tikzpicture}
\caption{从症状到反例的诊断分岔}\label{fig:13-diagnosis}
\begin{minipage}{0.96\linewidth}\footnotesize 教学示意：用于解释机制，不是运行时的实测数据或完整实现。

\textbf{读图顺序：}1. 先保留业务分母\quad 2. 计算候选\quad 3. 等待候选\quad 4. 用反例验结论。
\textbf{连线语义：}CPU/工作↑\ensuremath{\rightarrow}\allowbreak{}分配/GC 证据：并非唯一解释。
\textbf{机制依据：}\cite{go-diagnostics}。
\end{minipage}
\end{figure}
症状：相同输入下，/cpu 的吞吐低于预期，进程 CPU 较高。候选机制包括重复格式化、字符串复制、GC 或 CPU 配额节流。先固定请求数与输入长度，记录每请求 CPU 时间、每请求分配字节数和错误率；在函数基准中另记录 B/op。随后采集 CPU 和 alloc profile。

lab 的 encodeSlow 逐字节 fmt.Sprintf 并反复拼接字符串。合理预测是格式化、分配与复制相关栈可见，encodeFast 使用 encoding/hex 减少工作。但只有运行后的实际 profile 能证明这次环境中哪些栈占主导，不能把代码直觉当作测量结果。

干预把编码器替换为标准库实现，先验证所有字节值、空输入与输出格式一致。随后在多个输入长度下重复 benchmark，并保留关闭剖析采集时的对照结果，避免把采集开销计入优化收益。若长输入收益更大，可能与重复拼接的累计复制有关；需要规模扫描支撑复杂度解释。

可反驳条件：CPU profile 没有预期热点、分配没有变化、吞吐受发生器限制、错误请求比例变化，或者环境正在被节流。任意一项出现都要求修订解释。

\section{案例 B：分配压力不等于内存泄漏}
症状：/alloc 负载下 alloc\_space 快速增长，inuse\_space 却在 GC 后回落，RSS 未立刻下降。它支持“临时对象制造分配压力”，不支持“内存一定泄漏”。/retain 则把 1 MiB 对象保存在全局切片中，上限 64 MiB；如果统一 GC 及 profile 记账的可见性后，该分配栈的存活量仍随请求增加，就支持“可达对象持续积累”的解释。

\begin{codeblock}
curl -o alloc.pb.gz http://127.0.0.1:6060/debug/pprof/allocs
curl -o live.pb.gz 'http://127.0.0.1:6060/debug/pprof/heap?gc=1'
go tool pprof -sample_index=alloc_space -top alloc.pb.gz
go tool pprof -sample_index=inuse_space -top live.pb.gz
curl -X POST http://127.0.0.1:8080/reset
\end{codeblock}

采样、GC 周期和内存释放时机都会影响数值。\texttt{gc=1}会主动触发 GC，具有观察者效应；释放引用后，仍需允许 GC 完成，且向 OS 归还驻留页可能滞后。RSS 还包含栈、runtime 元数据、代码、共享库与其他映射。

修复判断应以相同操作次数、稳定阶段、多轮 GC 后的存活趋势与业务正确性为基础。不要仅用“调用 reset 后 RSS 没有马上降”来否定引用已释放，也不要仅用 alloc\_space 单调增加来证明泄漏。

\section{案例 C：共享锁把并发变成排队}
lab 的/lock 在共享 mutex 内 sleep，故意模拟慢操作占据临界区。并发增加时，单请求延迟与等待总量可增长，而 CPU 不一定很高。先同时看完成率、在途数和请求延迟，再用 mutex、block 与 trace 区分持锁路径、阻塞位置和调度延迟。

mutex profile 常将成本归于临界区结束的解锁栈；它不是简单告诉你“Lock 调用花了多少 CPU”。block profile 记录等待发生处，trace 帮助看到阻塞与唤醒的时间关系。等待总量可以超过墙钟采集时长，因为多个 goroutine 同时等待。

理想化计算：临界区 10 ms 且完全串行，不考虑其他成本，吞吐上界约 100 req/s。将并发从 10 增到 100 不会让这个临界区并行，只会扩大队列。真实结果还受到调度、定时器精度和负载生成方式的影响，因此这是模型上界，不是实验测得的吞吐。

干预首先判断 sleep 或 I/O 能否移出锁；若不能，考虑分片、批处理或明确串行队列的背压。改变锁粒度可能改变一致性边界，必须检验状态更新顺序与取消行为。

\section{案例 D：等待导致长尾，CPU 图却很平}
/wait 通过 timer 等待；它模拟等待形态而非真实网络栈。对比/cpu 和/wait 可学习一个重要负证据：CPU profile 热点弱或总样本少，不等于没有性能问题。请求可以绝大多数时间处于等待，仍有很高的用户可见延迟。

\begin{codeblock}
curl -o trace.out 'http://127.0.0.1:6060/debug/pprof/trace?seconds=3'
go tool trace trace.out
go tool trace -pprof=sched trace.out > sched.pb.gz
go tool pprof -top sched.pb.gz
\end{codeblock}

时间线区分 timer 等待、runnable 等待和实际运行。若要证明真实下游慢，应加入真实可控下游、网络 trace 或系统调用证据，不能从 timer 示例推广到 DNS、TCP 拥塞或磁盘服务时间。

可考虑减少串行等待、限制并发、设置总截止时间，并将取消信号正确传递到下游操作。重试须有总预算与抖动，避免把下游故障放大为重试风暴。这里的因果链最终要由故障注入和恢复实验检验。

\section{案例 E：混合负载下的错觉}
生产请求通常混合 CPU 密集、低频大对象和等待型业务。总体 CPU 占比变化可能只是请求组合变化，而非代码退化。利用有限集合 workload 标签分析 CPU 分布，同时记录每类请求成功数、输入大小与延迟。标签缺失的 profile 类型应通过独立实验或版本/实例隔离补证，而非假定能按标签过滤。

例如，发布后的同长度窗口中，alloc\_space 增量上升 50\%，而大输入请求比例也从 10\%升到 40\%。应分别比较同类请求的每请求分配字节数，并以相同的输入组合重放负载。若没有保留类别分母，最多能报告总分配上升，不能精确归因到发布。

混合负载还会产生相互干扰：内存型业务引发 GC，影响 CPU 型业务；锁持有者被调度延迟，放大等待者长尾。分析时需要资源共享模型，不能将各类独立 benchmark 简单相加当作混跑结果。

\section{把发现写成可以复查的结论}
建议每条发现用四句结构：现象是什么；原始证据来自何处；当前最合理机制是什么；下一步干预如何区分替代解释。例如：“固定批次下 B/op 上升；alloc\_space 集中于某编码栈；重复格式化可能增加临时对象；替换实现并验证输出后观察分配与 CPU 是否同步下降。”

报告还应说明验证边界：本地结果不能直接代表生产；未采集相应硬件事件，就不能声称缓存未命中率有所改善。race detector 仅检测本次执行路径上的数据竞争，即使检查通过，也不能排除未执行路径的数据竞争、死锁或逻辑并发错误。仅凭单次运行结果，通常也不足以给出可靠的不确定性估计。\cite{go-race}

附带 verification 目录保存本次实际执行的命令、原始 profile、日志与检查结果。教材中的模型数值均标记为示例或假设，实测结果只在验证报告中按环境解释。

\section{实验任务}
\subsection*{完成五种负载的盲诊断}
\begin{enumerate}
\item 先只看每种负载的 CPU、延迟和内存趋势，不阅读对应 handler 实现，写出候选假设。

\item 为每种情况选择两种相互补充的工具；收集 CPU、heap、block、mutex 和 trace。

\item 打开源码验证机制，记录哪些初始假设被否定。

\item 至少完成一个保持语义的优化或一个负结果实验，并进行独立重复评价。

\item 将模型预测、实际观测和生产外推限制分别写入报告。

\end{enumerate}
\section{自测与解析}
\noindent\textbf{1. alloc\_space 持续增大能否单独证明泄漏？}
\begin{enumerate}[label=\Alph*.]
\item 能
\item 不能，累计分配本来就可持续增大
\item 只要超过 RSS 就能
\item 采样率高就能
\end{enumerate}
\noindent{\color{forest}\textbf{答案 B。}}泄漏分析要关注可达存活对象、生命周期与长期趋势。

\noindent\textbf{2. 10 ms 串行临界区的理想吞吐上界约是多少？}
\begin{enumerate}[label=\Alph*.]
\item 100 req/s
\item 10000 req/s
\item 并发多少就多少倍
\item 无法给模型上界
\end{enumerate}
\noindent{\color{forest}\textbf{答案 A。}}1/0.01s=100 req/s，只是忽略其他开销的模型上界。

\noindent\textbf{3. 混合请求比例变化后 CPU 占比变了，首先需要什么？}
\begin{enumerate}[label=\Alph*.]
\item 直接回滚所有代码
\item 按业务类别和相同工作量比较
\item 移除所有标签
\item 只看函数排名
\end{enumerate}
\noindent{\color{forest}\textbf{答案 B。}}工作负载构成变化本身可以改变 profile 分布。

\section{分析题与参考答案}
\noindent\textbf{题 1：}mutex profile 中的累计等待为 20 秒，而你本次采集持续 5 秒。这是否说明数据损坏？

\noindent\textbf{参考答案：}不一定。应先确认 20 秒是累计总量还是本次基线差：累计值可能包含此前的历史。即使它是区间增量，多个 goroutine 并发等待也会使总等待超过墙钟时长；事件记账边界还可能跨越窗口。应核对版本、采样缩放和时间边界，不能把 20/5 解释为 400\%的 CPU 使用率。

\noindent\textbf{题 2：}临时内存下降而 RSS 不变，列出两种解释和补证工具。

\noindent\textbf{参考答案：}运行时保留已空闲堆页、其他内存映射或 cgo 占用均可能。结合 runtime/metrics 内存类别、heap、/proc 映射与后续 GC/归还趋势；不要仅据 RSS 推断 Go 对象仍可达。

\medskip\noindent\textbf{本章主要阅读：}\cite{go-pprof}\cite{go-http}\cite{go-memory}\cite{go-race}。
\chapter{实验课程与研究项目}
\label{ch:14}
\noindent{\large\color{forest}把命令、数据、推理和复现打包为完整学习成果。}\par\medskip
\noindent\textbf{先修要求：}第 01–13 章；至少一台可运行 Go 的开发机器。

\noindent\textbf{完成本章后，应能够：}
\begin{itemize}
\item 从零运行本地诊断服务和自动验证脚本
\item 完成从采集到解释再到干预的六组实验
\item 设计有消融和效度讨论的研究型项目
\item 按明确评分量规审核自己的结业报告
\end{itemize}
\section{实验包布局、环境与运行方式}
配套目录包含 \texttt{lab/\allowbreak{}main.\allowbreak{}go}、\texttt{main\_\allowbreak{}test.\allowbreak{}go}、\texttt{go.\allowbreak{}mod}、FlightRecorder 示例，以及 \texttt{scripts/\allowbreak{}verify\_\allowbreak{}lab.\allowbreak{}py} 和验证报告。普通实验使用 Go 1.22 或更新版本；FlightRecorder 要求 Go 1.25 及以上版本。本文按 Go 1.25.0 核查运行时源码，不声称它是发布时最新补丁版；实际部署应采用受支持的安全补丁版本。

\begin{codeblock}
cd lab
go test ./...
go run .
# 在第二个终端进行采集
curl http://127.0.0.1:8080/health
curl 'http://127.0.0.1:8080/cpu?n=1000'
\end{codeblock}

服务分开监听 127.0.0.1:8080 业务端口和 127.0.0.1:6060 诊断端口。CPU 循环、临时分配、保留对象和等待时间均有教学上限；/reset 要求 POST。实验启用高密度 block/mutex 采样以便短时间内观察，因此不可将其配置直接复制到生产。

用项目根目录下的 \texttt{python3 scripts/\allowbreak{}verify\_\allowbreak{}lab.\allowbreak{}py} 可重跑验证。脚本自动选择空闲回环端口，使用有界负载，只停止自己启动的进程，并把本次结果写入 verification/lab。手动启动可用 -app-port 与 -diag-port 指定其他端口。/health 的 finished\_handlers 只统计通过 annotate 包装且正常返回的业务处理器，包含其中的失败与取消，不统计 /health、/reset，也不能作为成功请求分母；负载脚本单独统计成功响应。

这些上限有范围：/alloc 的 n 和 size 只约束单次请求，多个并发请求仍可叠加显著内存；64 MiB 是 /retain 的教学保留上限，并非进程 RSS 或全部内存的上限。/lock 故意使用不可取消的 Mutex.Lock 与 Sleep，等待中的请求即使取消也不会立即释放这段工作；应把它当成反例，而非完整取消设计。

实验服务为兼容展示而手动注册诊断 handler，没有全部使用第 10 章的 GET 方法模式；因此不能把标准库在 DefaultServeMux 上的默认方法限制当作该手工 mux 的属性。正式部署应沿用受控监听并明确方法和认证策略。

\section{六组实验及提交清单}
\begingroup\small\setlength{\tabcolsep}{4pt}
\begin{longtable}{>{\raggedright\arraybackslash}p{0.298\linewidth}>{\raggedright\arraybackslash}p{0.298\linewidth}>{\raggedright\arraybackslash}p{0.298\linewidth}}
\toprule
\textbf{实验} & \textbf{核心问题} & \textbf{必须提交的证据} \\
\midrule
\endfirsthead
\toprule
\textbf{实验} & \textbf{核心问题} & \textbf{必须提交的证据} \\
\midrule
\endhead
L1 基准 & 优化真的少做了工作吗 & 相同输出测试、10 次原始结果、输入规模扫描 \\
\addlinespace[3pt]
L2 CPU & 哪个调用路径消耗 CPU & CPU profile、\texttt{top}/\texttt{top -\allowbreak{}cum}/\texttt{list}输出、成功请求分母 \\
\addlinespace[3pt]
L3 内存 & 是分配压力还是长期保留 & 四种 sample index、GC 前后趋势、生命周期解释 \\
\addlinespace[3pt]
L4 并发 & 等待发生在哪里、谁持锁 & block、mutex、trace 与吞吐—并发曲线 \\
\addlinespace[3pt]
L5 调试 & 变量状态与优化构建有何差别 & Delve 会话记录、构建参数、已验证边界 \\
\addlinespace[3pt]
L6 综合 & PGO/持续采集如何改变表现 & 对照构建、开销实验、代表性分析、回滚方案 \\
\addlinespace[3pt]
\bottomrule
\end{longtable}
\endgroup
每组实验必须有负例或反证：例如纯等待负载不应被 CPU 图解释为 CPU 热点；alloc 快速增长但 live 稳定时不能称为泄漏；PGO 构建成功不应等同于运行更快。负例能检验学习者是否掌握测量语义，而非只会执行命令。

\section{手动采集操作卡}
先在一个终端持续运行服务，在第二个终端持续产生负载，在第三个终端采集。采集 CPU profile 时如果没有工作，得到少量样本是正常现象。

\begin{codeblock}
# 终端2：有界并发负载，Linux/macOS 环境需有curl和xargs
seq 1 200 | xargs -P 8 -I{} curl -s -o /dev/null   'http://127.0.0.1:8080/cpu?n=10000'
# 终端3：采集时终端2应仍在运行
curl -o cpu.pb.gz 'http://127.0.0.1:6060/debug/pprof/profile?seconds=5'
curl -o heap.pb.gz 'http://127.0.0.1:6060/debug/pprof/heap?gc=1'
curl -o mutex.pb.gz http://127.0.0.1:6060/debug/pprof/mutex
curl -o block.pb.gz http://127.0.0.1:6060/debug/pprof/block
curl -o goroutine.txt 'http://127.0.0.1:6060/debug/pprof/goroutine?debug=2'
curl -o trace.out 'http://127.0.0.1:6060/debug/pprof/trace?seconds=3'
\end{codeblock}

上述\texttt{-\allowbreak{}o}参数使用相对路径，输出文件保存在执行命令时的当前目录；如果要保存在 lab 中，应先切换到该目录。建议为不同实验建立独立子目录，并记录负载和采集时间。alloc 等累计指标、以及启用采样后积累的 block/mutex 数据，可能包含当前负载之前的历史；应使用新进程或语义兼容的基线差，避免将旧负载的成本计入本次实验。CPU 区间采样、goroutine 快照和当前存活堆各有不同时间语义，不能一概按累计总量处理。

本书自动脚本仅进行短时功能验证，不代替以上完整性能研究。若缺少 Linux perf 权限、core 转储配置或外部调试器，报告应注明未执行的实验，不补造结果。

\section{研究型项目：三条可选路线}
\begin{figure}[H]
\centering
\begin{tikzpicture}[x=1cm,y=1cm]
\node[draw=linegray!75!black,fill=linegray!13,rounded corners=3pt,text width=7.08cm,minimum height=1.12cm,align=center,font=\small] (p0-base) at (3.75,2.92) {无 PGO};
\node[draw=linegray!75!black,fill=linegray!13,rounded corners=3pt,text width=7.08cm,minimum height=1.12cm,align=center,font=\small] (p0-rep) at (3.75,0.97) {代表性训练};
\node[draw=linegray!75!black,fill=linegray!13,rounded corners=3pt,text width=7.08cm,minimum height=1.12cm,align=center,font=\small] (p0-bias) at (3.75,-0.98) {偏置训练};
\node[draw=linegray!75!black,fill=linegray!13,rounded corners=3pt,text width=7.08cm,minimum height=1.12cm,align=center,font=\small] (p0-old) at (3.75,-2.92) {过期训练};
\node[draw=leaf!75!black,fill=leaf!13,rounded corners=3pt,text width=7.08cm,minimum height=1.12cm,align=center,font=\small] (p0-test) at (11.25,0.00) {独立评价};
\draw[-{Stealth[length=2mm]},forest!70!black] (p0-base.east) -- (p0-test.west);
\draw[-{Stealth[length=2mm]},forest!70!black] (p0-rep.east) -- (p0-test.west);
\draw[-{Stealth[length=2mm]},forest!70!black] (p0-bias.east) -- (p0-test.west);
\draw[-{Stealth[length=2mm]},forest!70!black] (p0-old.east) -- (p0-test.west);
\end{tikzpicture}
\caption{研究项目需要对照和消融}\label{fig:14-design}
\begin{minipage}{0.96\linewidth}\footnotesize 教学示意：用于解释机制，不是运行时的实测数据或完整实现。

\textbf{读图顺序：}1. 建立基线\quad 2. 检验训练作用\quad 3. 挑战代表性假设\quad 4. 独立评价。
\textbf{机制依据：}\cite{go-pgo}。
\end{minipage}
\end{figure}
路线 A：比较 CPU 采样与执行追踪对低频尾延迟事件的诊断能力。控制事件发生率和持续时间，测量发现率、定位准确率、数据量与扰动，并探索触发式 flight recorder。必须讨论“已知人工注入”的机制与未知生产异常之间的差距。

路线 B：研究 PGO 的训练负载漂移。设计三类请求与多种混合比例，训练和评价分离，测量 CPU/请求、代码尺寸和尾延迟；用 profile 权重分布变化解释收益退化。消融实验包括无 PGO、代表性 profile、偏置 profile 和过期 profile。

路线 C：比较进程内与 eBPF CPU 采样。固定相同工作量、构建与时间窗，评价栈完整性、unknown 比例、热点排序稳定性和开销。须报告 cgo、内联、架构、内核、容器和权限差异；不能仅以火焰图相似就宣布两采集器等价。

第 16 章将这些项目连接到近五年原始工作。优秀项目应提出一个可以通过实验检验并可能被推翻的假设，而不是“部署某工具并截图”。

\section{报告结构、评分量规与复现要求}
\begingroup\small\setlength{\tabcolsep}{4pt}
\begin{longtable}{>{\raggedright\arraybackslash}p{0.298\linewidth}>{\raggedright\arraybackslash}p{0.298\linewidth}>{\raggedright\arraybackslash}p{0.298\linewidth}}
\toprule
\textbf{维度} & \textbf{分值} & \textbf{合格条件} \\
\midrule
\endfirsthead
\toprule
\textbf{维度} & \textbf{分值} & \textbf{合格条件} \\
\midrule
\endhead
问题与模型 & 15 & 明确指标、分母、假设和可推翻条件 \\
\addlinespace[3pt]
实验设计 & 20 & 控制变量、独立重复、负载代表性与对照 \\
\addlinespace[3pt]
工具与数据 & 20 & 原始文件、版本、单位、采集边界正确 \\
\addlinespace[3pt]
机制解释 & 20 & 跨工具证据、替代解释、消融或干预 \\
\addlinespace[3pt]
正确性与限制 & 15 & 输出一致、并发安全、效度边界诚实 \\
\addlinespace[3pt]
复现与表达 & 10 & 脚本可运行、路径清晰、结论可审查 \\
\addlinespace[3pt]
\bottomrule
\end{longtable}
\endgroup
这份 100 分评分量规评价论证质量，不以性能提升的百分比作为评分标准。负结果只要设计严谨、解释充分，同样可获得高分。提交包包含 README、源码、命令、原始数据、处理脚本、图表及来源清单；隐去敏感信息时说明脱敏规则及其对复现的影响。

至少讨论内部效度（是否真由干预引起）、外部效度（是否能推广）、构念效度（指标是否代表目标）和统计结论效度（样本与不确定性是否足够）。这四个角度让工程诊断达到研究生课程的论证要求。

\section{口试与同行复核}
口试应围绕原始证据追问：若 profile 采样窗口减半结论是否稳定？若请求组合变化是否仍获益？为什么这个 cum 百分比不能相加？为何 heap 下降但 RSS 不变？为什么 perf 与 Go CPU profile 可能出现不同栈？答案必须回到时间、单位、采样、生命周期和平台边界。

同行复核者首先运行正确性测试，再复现一项核心结论，然后尝试构造反例。作者根据反例更新结论范围，不应把所有偏差都归为“机器不同”。最终交付的是可维护的知识：哪些条件下有效、哪些条件下无效、怎样快速判断当前系统属于哪一类。

\section{完整实验服务源码}
以下为配套 lab/main.go 的完整源码，使用标准库，无第三方运行依赖。文件 go.mod 声明 module example.com/golab 与 go 1.22。

\begin{codeblock}
// GoLab: bounded, local-only diagnostics laboratory. See README.md.
package main

import (
	"context"
	"encoding/hex"
	"flag"
	"fmt"
	"log"
	"net/http"
	"net/http/pprof"
	"runtime"
	rpprof "runtime/pprof"
	"runtime/trace"
	"strconv"
	"sync"
	"sync/atomic"
	"time"
)

var retained struct {
	sync.Mutex
	blobs [][]byte
}
var serial sync.Mutex

// This counts finished handlers, including canceled or rejected requests.
var finishedHandlers atomic.Uint64

func encodeSlow(src []byte) string {
	out := ""
	for _, b := range src {
		out += fmt.Sprintf("%02x", b)
	}
	return out
}
func encodeFast(src []byte) string { return hex.EncodeToString(src) }
func bounded(r *http.Request, key string, fallback, max int) int {
	v, err := strconv.Atoi(r.URL.Query().Get(key))
	if err != nil || v < 1 {
		return fallback
	}
	if v > max {
		return max
	}
	return v
}
func annotate(kind string, next func(context.Context, http.ResponseWriter, *http.Request)) http.HandlerFunc {
	return func(w http.ResponseWriter, r *http.Request) {
		rpprof.Do(r.Context(), rpprof.Labels("workload", kind), func(ctx context.Context) {
			ctx, task := trace.NewTask(ctx, "request")
			defer task.End()
			trace.Log(ctx, "route", kind)
			region := trace.StartRegion(ctx, kind)
			defer region.End()
			next(ctx, w, r)
			finishedHandlers.Add(1)
		})
	}
}
func main() {
	appPort := flag.Int("app-port", 8080, "loopback application port")
	diagPort := flag.Int("diag-port", 6060, "loopback diagnostics port")
	flag.Parse()
	appAddr := fmt.Sprintf("127.0.0.1:%d", *appPort)
	diagAddr := fmt.Sprintf("127.0.0.1:%d", *diagPort)
	// Aggressive rates make short lab runs visible. Do not copy into production.
	runtime.SetBlockProfileRate(1)
	runtime.SetMutexProfileFraction(1)
	diagnostic := http.NewServeMux()
	diagnostic.HandleFunc("/debug/pprof/", pprof.Index)
	diagnostic.HandleFunc("/debug/pprof/cmdline", pprof.Cmdline)
	diagnostic.HandleFunc("/debug/pprof/profile", pprof.Profile)
	diagnostic.HandleFunc("/debug/pprof/symbol", pprof.Symbol)
	diagnostic.HandleFunc("/debug/pprof/trace", pprof.Trace)
	go func() {
		log.Println("diagnostics http://" + diagAddr + "/debug/pprof/")
		log.Fatal(http.ListenAndServe(diagAddr, diagnostic))
	}()
	app := http.NewServeMux()
	app.HandleFunc("/health", func(w http.ResponseWriter, r *http.Request) {
		fmt.Fprintf(w, "ok finished_handlers=%d\n", finishedHandlers.Load())
	})
	app.HandleFunc("/cpu", annotate("cpu", func(ctx context.Context, w http.ResponseWriter, r *http.Request) {
		src := []byte("profiling-with-go-0123456789abcdef")
		iterations := bounded(r, "n", 1000, 100000)
		fast := r.URL.Query().Get("fast") == "1"
		var value string
		for i := 0; i < iterations; i++ {
			if i%256 == 0 {
				select {
				case <-ctx.Done():
					return
				default:
				}
			}
			if fast {
				value = encodeFast(src)
			} else {
				value = encodeSlow(src)
			}
		}
		fmt.Fprintln(w, value)
	}))
	app.HandleFunc("/alloc", annotate("alloc", func(ctx context.Context, w http.ResponseWriter, r *http.Request) {
		n := bounded(r, "n", 2000, 20000)
		size := bounded(r, "size", 512, 4096)
		// Retain the batch within this request so allocations cannot disappear.
		blobs := make([][]byte, n)
		for i := range blobs {
			blobs[i] = make([]byte, size)
			blobs[i][0] = byte(i)
		}
		fmt.Fprintf(w, "temporary=%d checksum=%d\n", n*size, blobs[n-1][0])
		runtime.KeepAlive(blobs)
	}))
	app.HandleFunc("/retain", annotate("retain", func(ctx context.Context, w http.ResponseWriter, r *http.Request) {
		retained.Lock()
		defer retained.Unlock()
		if len(retained.blobs) >= 64 {
			http.Error(w, "64 MiB lab cap reached; call POST /reset", http.StatusTooManyRequests)
			return
		}
		b := make([]byte, 1<<20)
		for i := 0; i < len(b); i += 4096 {
			b[i] = 1
		}
		retained.blobs = append(retained.blobs, b)
		fmt.Fprintf(w, "retained=%d MiB\n", len(retained.blobs))
	}))
	app.HandleFunc("/reset", func(w http.ResponseWriter, r *http.Request) {
		if r.Method != http.MethodPost {
			http.Error(w, "POST required", http.StatusMethodNotAllowed)
			return
		}
		retained.Lock()
		retained.blobs = nil
		retained.Unlock()
		fmt.Fprintln(w, "references released; next GC determines reclamation")
	})
	app.HandleFunc("/wait", annotate("wait", func(ctx context.Context, w http.ResponseWriter, r *http.Request) {
		timer := time.NewTimer(time.Duration(bounded(r, "ms", 40, 500)) * time.Millisecond)
		defer timer.Stop()
		select {
		case <-ctx.Done():
			return
		case <-timer.C:
			fmt.Fprintln(w, "waited")
		}
	}))
	app.HandleFunc("/lock", annotate("lock", func(ctx context.Context, w http.ResponseWriter, r *http.Request) {
		// Intentionally bad: hold a shared mutex while sleeping.
		serial.Lock()
		time.Sleep(time.Duration(bounded(r, "ms", 10, 100)) * time.Millisecond)
		serial.Unlock()
		fmt.Fprintln(w, "serialized")
	}))
	server := &http.Server{Addr: appAddr, Handler: app, ReadHeaderTimeout: 3 * time.Second}
	log.Println("application http://" + appAddr)
	log.Fatal(server.ListenAndServe())
}
\end{codeblock}

\section{完整基准与正确性测试源码}
以下为同目录中的 main\_test.go。它只覆盖若干代表性输入，并非完整的输入域验证；研究实验还应补充全部 256 种字节值、不同长度和随机输入，并设计独立重复。

\begin{codeblock}
package main

import "testing"

var result string

func TestEncoders(t *testing.T) {
	for _, in := range [][]byte{nil, {}, []byte("hello"), {0, 1, 15, 16, 127, 128, 255}} {
		if a, b := encodeSlow(in), encodeFast(in); a != b {
			t.Fatalf("encode mismatch: %q != %q", a, b)
		}
	}
}
func BenchmarkEncodeSlow(b *testing.B) {
	input := []byte("profiling-with-go-0123456789abcdef")
	b.ReportAllocs()
	b.SetBytes(int64(len(input)))
	b.ResetTimer()
	for i := 0; i < b.N; i++ {
		result = encodeSlow(input)
	}
}
func BenchmarkEncodeFast(b *testing.B) {
	input := []byte("profiling-with-go-0123456789abcdef")
	b.ReportAllocs()
	b.SetBytes(int64(len(input)))
	b.ResetTimer()
	for i := 0; i < b.N; i++ {
		result = encodeFast(input)
	}
}
\end{codeblock}

\section{实验任务}
\subsection*{提交结业证据包}
\begin{enumerate}
\item 完成六组实验中的至少四组，其中必须包含 CPU、内存和并发。

\item 选择一条研究路线，预先写出假设、主指标、重复单元和停止规则。

\item 保留原始数据并执行至少一个消融或反例实验。

\item 邀请另一人按 README 复现；记录失败原因并修复环境说明。

\item 按 100 分量规自评，写出三条最重要结论和三条适用边界。

\end{enumerate}
\section{自测与解析}
\noindent\textbf{1. 结业项目最关键的产物是什么？}
\begin{enumerate}[label=\Alph*.]
\item 一张颜色丰富的火焰图
\item 可复查的证据、机制解释和复现协议
\item 最长的命令清单
\item 一定超过 50\%的加速
\end{enumerate}
\noindent{\color{forest}\textbf{答案 B。}}课程评估科学推理与可复现性，不预设性能收益。

\noindent\textbf{2. 自动烟测可以替代长期性能统计吗？}
\begin{enumerate}[label=\Alph*.]
\item 可以
\item 不能，它主要验证路径可运行
\item 只要日志无 error 即可
\item 有 PDF 就可以
\end{enumerate}
\noindent{\color{forest}\textbf{答案 B。}}运行正确与统计结论是不同层级的验证。

\noindent\textbf{3. 设计消融实验的目的是什么？}
\begin{enumerate}[label=\Alph*.]
\item 让报告页数更多
\item 隔离某机制对结果的贡献
\item 隐藏失败结果
\item 只看最有利样本
\end{enumerate}
\noindent{\color{forest}\textbf{答案 B。}}消融移除或改变特定因素，帮助区分因果机制。

\section{分析题与参考答案}
\noindent\textbf{题 1：}为 PGO 漂移项目写一个可反驳假设。

\noindent\textbf{参考答案：}例如：在固定源码和资源下，训练与评价的 CPU 权重分布距离增大时，CPU/请求收益下降。设计多种混合比例、无 PGO 对照、独立重复，并接受无单调关系的负结果；不能只比较两个任意点。

\noindent\textbf{题 2：}如何区分“采集器栈展开不完整”与“函数没有执行”？

\noindent\textbf{参考答案：}使用已知执行路径的受控负载、匹配构建符号、比较原始样本和 unknown/truncated 计数，必要时以另一采集机制交叉验证。图上没有函数本身不足以做判断。

\medskip\noindent\textbf{本章主要阅读：}\cite{go-diagnostics}\cite{go-pgo}\cite{go-pprof}。
\part*{第四篇 · 开源实现与研究前沿}
\addcontentsline{toc}{part}{第四篇 · 开源实现与研究前沿}
\chapter[成熟开源项目的实现与核心技术]{成熟开源项目的实现\\与核心技术}
\label{ch:15}
\noindent{\large\color{forest}沿一条样本追踪采集、展开、规范化、存储与查询，并用固定源码核对各步的数据转换}\par\medskip
\noindent\textbf{先修要求：}第 04—12 章；理解指针、哈希表、树、CPU 寄存器、采样误差与基本数据库聚合。

\noindent\textbf{完成本章后，应能够：}
\begin{itemize}
\item 逐步计算含 ASLR、局部 ID、标签和多种单位的 profile 合并结果
\item 从栈展开规则、符号引用重写及列式查询推导数据结构与复杂度
\item 定位部分栈、累计计数重置、缓存淘汰和过滤造成的信息损失
\item 比较暂停调试、进程内剖析与持续剖析的适用边界，并设计交叉验证
\end{itemize}
\section{贯穿案例：同一条调用栈为何有多种身份}
\begin{figure}[H]
\centering
\begin{tikzpicture}[x=1cm,y=1cm]
\node[draw=linegray!75!black,fill=linegray!13,rounded corners=3pt,text width=2.58cm,minimum height=1.12cm,align=center,font=\small] (p0-pc) at (1.50,0.00) {寄存器 PC/\allowbreak{}SP\\{\scriptsize 执行现场}};
\node[draw=linegray!75!black,fill=linegray!13,rounded corners=3pt,text width=2.58cm,minimum height=1.12cm,align=center,font=\small] (p0-stack) at (4.50,0.00) {地址栈\\{\scriptsize [0x1040,0x2080]}};
\node[draw=linegray!75!black,fill=linegray!13,rounded corners=3pt,text width=2.58cm,minimum height=1.12cm,align=center,font=\small] (p0-profile) at (7.50,0.00) {pprof\\{\scriptsize location IDs＋权重}};
\node[draw=linegray!75!black,fill=linegray!13,rounded corners=3pt,text width=2.58cm,minimum height=1.12cm,align=center,font=\small] (p0-fact) at (10.50,0.00) {样本事实\\{\scriptsize 时间/\allowbreak{}标签/\allowbreak{}栈/\allowbreak{}值}};
\node[draw=leaf!75!black,fill=leaf!13,rounded corners=3pt,text width=2.58cm,minimum height=1.12cm,align=center,font=\small] (p0-tree) at (13.50,0.00) {查询火焰树\\{\scriptsize 按调用路径聚合}};
\draw[-{Stealth[length=2mm]},forest!70!black] (p0-pc.east) -- (p0-stack.west);
\draw[-{Stealth[length=2mm]},forest!70!black] (p0-stack.east) -- (p0-profile.west);
\draw[-{Stealth[length=2mm]},forest!70!black] (p0-profile.east) -- (p0-fact.west);
\draw[-{Stealth[length=2mm]},forest!70!black] (p0-fact.east) -- (p0-tree.west);
\end{tikzpicture}
\caption{同一条样本跨越五种表示}\label{fig:15-sample-journey}
\begin{minipage}{0.96\linewidth}\footnotesize 教学示意：用于解释机制，不是运行时的实测数据或完整实现。

\textbf{读图顺序：}1. 取得执行现场\quad 2. 逐帧展开\quad 3. 转换本地 ID\quad 4. 保留事实语义\quad 5. 形成查询结果。
\textbf{连线语义：}寄存器 PC/SP\ensuremath{\rightarrow}\allowbreak{}地址栈：展开；地址栈\ensuremath{\rightarrow}\allowbreak{}pprof：规范化；pprof\ensuremath{\rightarrow}\allowbreak{}样本事实：摄取；样本事实\ensuremath{\rightarrow}\allowbreak{}查询火焰树：过滤归并。
\textbf{机制依据：}\cite{proj-pprof-model}\cite{proj-parca-store}。
\end{minipage}
\end{figure}
\subsubsection{阅读任务：沿数据流分析实现}
本章选择 google/pprof、Delve、Grafana Pyroscope、Parca 与 Parca Agent。这些项目分别把“观测程序”中的不同难点做成了可用系统：pprof 定义和操作带权调用栈；Delve 从停止的进程恢复语言现场；Pyroscope 与 Parca 承担多实例、长时间的数据存储和查询；Agent 负责低侵入的系统级采样。沿同一个样本追踪这些系统，才能解释为什么相同的火焰图外观可能来自不同测量语义。

源码固定为 pprof 提交 \texttt{60bf690a93\allowbreak{}02b860b211\allowbreak{}0fbb4dccdf\allowbreak{}bbeb26ab89}、Delve \texttt{v1.\allowbreak{}24.\allowbreak{}2}、Pyroscope \texttt{v1.\allowbreak{}14.\allowbreak{}0}、Parca \texttt{v0.\allowbreak{}23.\allowbreak{}1}、Parca Agent \texttt{v0.\allowbreak{}38.\allowbreak{}1}。Agent 的 Go 模块替换指向 OpenTelemetry（OTel）eBPF profiler 的 Parca 分支 \texttt{edfa6253a8b2}。这是可复查的历史截面，不是“均为最新版”的声明；尤其不能据 Delve 旧版本推断它支持本教材的 Go 1.25。不同系统是可比较的实现，不要求读者把它们串联部署。\cite{proj-pprof} \cite{proj-delve} \cite{proj-pyroscope} \cite{proj-parca} \cite{proj-parca-agent}

\subsubsection{构造一个可手算的输入}
把两份 profile 记为 P1、P2，假设各自进程使用的二进制构建标识（Build ID）都是 B1。函数 encode、hash、serve、main 在同一映射内的相对地址分别为 \texttt{0x1120、0x1800、0x2020、0x3030}。P1 记录的装载基址为 \texttt{0x400000}，P2 为 \texttt{0x700000}；地址空间布局随机化（ASLR）可以造成这种差异。同一 encode 指令因而可能显示为 \texttt{0x401120} 或 \texttt{0x701120}。映射大小和文件偏移相同；本例没有内联帧。

P1 用位置 ID 7、3、1 表示 encode、serve、main，P2 却使用 42、8、9；hash 在 P2 中用 ID 43。ID 只在各自文件中有效，数字接近或相等都不构成函数同一性的证据。下面的栈均按 pprof 的叶到根顺序书写，值向量的类型依次是 \texttt{samples/\allowbreak{}count}、\texttt{cpu/\allowbreak{}nanoseconds}。

\begingroup\small\setlength{\tabcolsep}{4pt}
\begin{longtable}{>{\raggedright\arraybackslash}p{0.165\linewidth}>{\raggedright\arraybackslash}p{0.165\linewidth}>{\raggedright\arraybackslash}p{0.165\linewidth}>{\raggedright\arraybackslash}p{0.165\linewidth}>{\raggedright\arraybackslash}p{0.165\linewidth}}
\toprule
\textbf{输入记录} & \textbf{局部位置 ID} & \textbf{route 标签} & \textbf{计数} & \textbf{CPU 时间} \\
\midrule
\endfirsthead
\toprule
\textbf{输入记录} & \textbf{局部位置 ID} & \textbf{route 标签} & \textbf{计数} & \textbf{CPU 时间} \\
\midrule
\endhead
P1-A & 7 \ensuremath{\rightarrow}\allowbreak{} 3 \ensuremath{\rightarrow}\allowbreak{} 1 & /pay & 3 & 30 ms \\
\addlinespace[3pt]
P2-A & 42 \ensuremath{\rightarrow}\allowbreak{} 8 \ensuremath{\rightarrow}\allowbreak{} 9 & /pay & 2 & 20 ms \\
\addlinespace[3pt]
P2-B & 43 \ensuremath{\rightarrow}\allowbreak{} 8 \ensuremath{\rightarrow}\allowbreak{} 9 & /pay & 1 & 10 ms \\
\addlinespace[3pt]
P2-C & 42 \ensuremath{\rightarrow}\allowbreak{} 8 \ensuremath{\rightarrow}\allowbreak{} 9 & /admin & 1 & 10 ms \\
\addlinespace[3pt]
\bottomrule
\end{longtable}
\endgroup
这是人为构造的 10 ms 周期示例，不是一次真实服务测量，也不意味着可以任意调整 Go 的标准 CPU profiler。本例中，两份 profile 采自先后运行的进程，对应两个连续采集窗口。后文另讨论并发实例的累计采集时长与墙钟时间为何不同。

正确的语义合并产生三条样本：encode 栈的 \texttt{/\allowbreak{}pay} 为 \texttt{[5,50 ms]}，hash 栈的 \texttt{/\allowbreak{}pay} 为 \texttt{[1,10 ms]}，encode 栈的 \texttt{/\allowbreak{}admin} 为 \texttt{[1,10 ms]}。有三条样本，却只有两条不同调用路径。全局总计 7 次、70 ms；筛选 \texttt{/\allowbreak{}pay} 后是 6 次、60 ms，其中 encode 的 flat 为 50 ms，hash 为 10 ms，serve 与 main 的 cum 都为 60 ms。各函数 cum 相加会重复计算，不能得到总 CPU。

\subsubsection{全章要维持的五个不变量}
从内核到界面，必须分别保持二进制身份、调用顺序、单位、标签分组和时间含义。地址去重不能跨 Build ID 误合并；位置表重编号不能倒转栈；将 count 变为时间要显式乘周期；同栈不同标签不能提前合并；累计量和区间量不能相加后解释成同一种时间序列。

接下来的每个实现都在减少成本：内核使用有限规则，传输复用栈 ID，存储复用符号，查询提前过滤。节约的代价是增加假设和状态。源码阅读的主问题因此是：转换是否保持上述不变量，失败时系统如何告诉我们已经失去哪些信息？

\section{采集与栈展开：从一组寄存器恢复可用调用链}
\begin{figure}[H]
\centering
\begin{tikzpicture}[x=1cm,y=1cm]
\node[draw=linegray!75!black,fill=linegray!13,rounded corners=3pt,text width=2.58cm,minimum height=1.12cm,align=center,font=\small] (p0-regs) at (1.50,0.00) {PC /\allowbreak{} SP /\allowbreak{} FP};
\node[draw=linegray!75!black,fill=linegray!13,rounded corners=3pt,text width=2.58cm,minimum height=1.12cm,align=center,font=\small] (p0-rule) at (4.50,0.00) {帧规则\\{\scriptsize DWARF 或运行时规则}};
\node[draw=linegray!75!black,fill=linegray!13,rounded corners=3pt,text width=2.58cm,minimum height=1.12cm,align=center,font=\small] (p0-cfa) at (7.50,0.97) {CFA/\allowbreak{}调用者 SP};
\node[draw=linegray!75!black,fill=linegray!13,rounded corners=3pt,text width=2.58cm,minimum height=1.12cm,align=center,font=\small] (p0-ret) at (10.50,0.00) {返回地址};
\node[draw=linegray!75!black,fill=linegray!13,rounded corners=3pt,text width=2.58cm,minimum height=1.12cm,align=center,font=\small] (p0-next) at (13.50,0.00) {下一帧 PC};
\node[draw=leaf!75!black,fill=leaf!13,rounded corners=3pt,text width=2.58cm,minimum height=1.12cm,align=center,font=\small] (p0-stop) at (7.50,-0.97) {终止/\allowbreak{}失败};
\draw[-{Stealth[length=2mm]},forest!70!black] (p0-regs.east) -- (p0-rule.west);
\draw[-{Stealth[length=2mm]},forest!70!black] (p0-rule.east) -- (p0-cfa.west);
\draw[-{Stealth[length=2mm]},forest!70!black] (p0-cfa.east) -- (p0-ret.west);
\draw[-{Stealth[length=2mm]},forest!70!black] (p0-ret.east) -- (p0-next.west);
\draw[-{Stealth[length=2mm]},forest!70!black] (p0-next.north) -- (13.50,2.12) -- (4.50,2.12) -- (p0-rule.north);
\draw[-{Stealth[length=2mm]},forest!70!black] (p0-rule.east) -- (p0-stop.west);
\end{tikzpicture}
\caption{逐帧展开：先求上一帧，再继续}\label{fig:15-unwind}
\begin{minipage}{0.96\linewidth}\footnotesize 教学示意：用于解释机制，不是运行时的实测数据或完整实现。

\textbf{读图顺序：}1. 按 PC 选择规则\quad 2. 恢复调用者位置\quad 3. 读取返回地址\quad 4. 保留失败原因。
\textbf{连线语义：}CFA/调用者 SP\ensuremath{\rightarrow}\allowbreak{}返回地址：读取保存位置；下一帧 PC\ensuremath{\rightarrow}\allowbreak{}帧规则：继续；帧规则\ensuremath{\rightarrow}\allowbreak{}终止/失败：无规则/越界。
\textbf{机制依据：}\cite{proj-delve-stack}\cite{proj-otel-ebpf}。
\end{minipage}
\end{figure}
\subsubsection{一个 CPU 样本起初并不是函数名}
系统级采样通常先得到时间戳、进程线程标识、指令指针 PC 和部分寄存器，再读用户内存恢复调用者。知道当前 PC 只够识别栈顶；要得到 encode \ensuremath{\rightarrow}\allowbreak{} serve \ensuremath{\rightarrow}\allowbreak{} main，还需回答每一帧保存返回地址的位置，以及调用者的栈指针如何恢复。Go 的栈可增长和移动，且存在运行时栈切换，不能把“沿内存不断读指针”当成普适算法。

所审读的 OTel eBPF 分支将可执行文件解析放在用户态，生成紧凑的 StackDelta 规则供内核查询，用它们描述不同 PC 区间如何恢复调用者状态。这里 SP 表示栈指针。\texttt{elfgopclntab.\allowbreak{}go:\allowbreak{}480–626} 读取 Go 函数表和 PCSP 信息，将 PC 区间映射成栈恢复规则；x86 路径按情况选择帧指针策略或 SP delta 策略。PCSP 表表示“程序计数器走到这里时，SP 相对函数入口变化了多少”，不是采样权重。\cite{proj-go-unwind}

用一个简化的 x86-64 模型说明。CFA（规范帧地址）是恢复调用者状态的参考地址；本例当前 \texttt{SP=0x1000}，规则给出 \texttt{CFA=SP+40}。返回地址保存于 \texttt{CFA\ensuremath{-}8}，该位置的值为 \texttt{0x40202a}。于是下一帧使用 \texttt{SP=0x1028、PC=0x40202a}，再查询 serve 所处 PC 的规则。真实编译器可以使用不同序言、尾声、寄存器和特殊函数规则；这组地址只用来演示递推，不是 Go ABI 的固定布局。

\subsubsection{为什么需要页索引、二分和分段执行}
DWARF 是一种调试信息格式；在 BPF 的受限执行环境中，不宜每次采样都完整解析它。\texttt{native\_\allowbreak{}stack\_\allowbreak{}trace.\allowbreak{}ebpf.\allowbreak{}c:\allowbreak{}137–217} 先按可执行文件标识和地址页找规则表，再在页内寻找覆盖当前 PC 的最后一条规则。页面为 64 KiB，源码使用至多 16 步的有界二分。它先找严格大于目标的边界，再退一项取得不大于目标的规则，必要时使用前一页延伸来的区间规则。\cite{proj-ebpf-unwind}

这个“前驱查询”很重要。假设规则起点为 10、20、35，目标 PC 偏移为 27，应使用起点 20 的规则，下一条起点 35 尚未生效。规则描述半开区间上的状态；即使目标改为 29、数值上更接近 35，也仍应使用起点 20 的规则。按最近距离选择记录可能得到错误 SP。对有序规则表，二分每一步砍掉约一半候选；其工作量随规则数对数增长，而内核验证器需要额外看到确定的最大迭代次数。

一次 BPF 程序只展开受限数量的本地代码帧（native frame）；需要继续时，通过 BPF tail call 跳转到后续展开程序。\texttt{575–630} 行把本轮展开、遇到解释器帧后的分派、错误码与继续执行组合起来。这样拆分的目的在于让每一段的资源消耗可分析，而不是无限扩大单个内核程序。深栈、特殊运行时边界和解释器切换仍然可能触及限制。

\subsubsection{部分成功比完全失败更难发现}
该实现先把已知当前帧压入结果，再尝试恢复调用者。因此，已经保存 encode、serve 两帧后，仍可能因为缺少 serve 当前 PC 区间的展开规则，或读取其返回地址失败而停止，无法继续恢复 main。返回的栈非空，不代表展开完整；“有一张火焰图”也不能证明它的根节点真实。

规则页不存在、嵌套 BPF 映射（map-of-maps）查找失败、没有可用前驱规则和读栈失败，都有各自的错误路径。对生产数据应分别统计采样事件数、成功取得栈顶数、完整展开数、符号化成功数和标签可读数。若只统计非空栈，某些工作负载被截短的比例会被掩盖。这组质量指标由源码控制流推导，具体名称与采集方式仍需在监测方案中实现。

把失败栈删除也不总是正确：假如 cgo 热点更容易展开失败，删掉它们会系统性高估纯 Go 热点占比。保留错误类别、可用前缀和原始总权重，有助于给查询结果标注“剩余部分未知”。但将所有失败样本统一合入一个正常函数名同样会制造虚假的调用关系。

\subsubsection{Go 栈、Go 标签与兼容性是三个问题}
该 fork 的 \texttt{interpreter/\allowbreak{}golang/\allowbreak{}runtime\_\allowbreak{}data.\allowbreak{}go} 中最新标签布局为 \texttt{go1.\allowbreak{}24}；加载器遇到未知版本会告警并尝试最新布局。标签依赖运行时对象的具体内存结构，正确展开 native 栈并不能证明正确读到了 goroutine 的 pprof 标签。Go 1.25 应分别验证两个能力，不能用旧布局回退当成兼容性保证。\cite{proj-otel-ebpf}

对照实验可以让 \texttt{/\allowbreak{}pay} 与 \texttt{/\allowbreak{}admin} 执行相同函数，只给它们不同 pprof 标签。若系统级结果的栈相同但标签缺失，先检查标签布局和继承路径；不要立即解释为请求没有运行。再加入 cgo、深递归和 goroutine 迁移，比较进程内 profile 的总量与完整栈比例。本节只给出实验设计，尚无这组对照的运行结果。

Delve 与这种连续采样的底层材料有重叠：二者都要从寄存器和二进制恢复栈。但采样器要在极短时间内反复执行，Delve 可以在暂停状态做更丰富的内存读取。选择预计算规则而非内核完整解析，正是性能、权限和覆盖之间的工程折中。

\section{pprof 规范化与合并：建立语义身份后才能相加}
\begin{figure}[H]
\centering
\begin{tikzpicture}[x=1cm,y=1cm]
\node[draw=linegray!75!black,fill=linegray!13,rounded corners=3pt,text width=3.33cm,minimum height=1.12cm,align=center,font=\small] (p0-a) at (1.88,0.97) {文件 A loc=7\\{\scriptsize encode /\allowbreak{} +0x40}};
\node[draw=linegray!75!black,fill=linegray!13,rounded corners=3pt,text width=3.33cm,minimum height=1.12cm,align=center,font=\small] (p0-b) at (1.88,-0.97) {文件 B loc=7\\{\scriptsize hash /\allowbreak{} +0x90}};
\node[draw=linegray!75!black,fill=linegray!13,rounded corners=3pt,text width=3.33cm,minimum height=1.12cm,align=center,font=\small] (p0-key) at (5.62,0.00) {语义键\\{\scriptsize mapping＋地址＋行}};
\node[draw=linegray!75!black,fill=linegray!13,rounded corners=3pt,text width=3.33cm,minimum height=1.12cm,align=center,font=\small] (p0-dest) at (9.38,0.00) {新 ID 空间\\{\scriptsize encode=1 /\allowbreak{} hash=2}};
\node[draw=leaf!75!black,fill=leaf!13,rounded corners=3pt,text width=3.33cm,minimum height=1.12cm,align=center,font=\small] (p0-sum) at (13.12,0.00) {同键才求和\\{\scriptsize 单位类型必须兼容}};
\draw[-{Stealth[length=2mm]},forest!70!black] (p0-a.east) -- (p0-key.west);
\draw[-{Stealth[length=2mm]},forest!70!black] (p0-b.east) -- (p0-key.west);
\draw[-{Stealth[length=2mm]},forest!70!black] (p0-key.east) -- (p0-dest.west);
\draw[-{Stealth[length=2mm]},forest!70!black] (p0-dest.east) -- (p0-sum.west);
\end{tikzpicture}
\caption{局部 ID 相同，不意味着函数相同}\label{fig:15-id-merge}
\begin{minipage}{0.96\linewidth}\footnotesize 独立于正文 P1/P2 的教学反例：两个文件都用 loc=7 却代表不同位置。归并先建立语义身份，再重写局部 ID，绝不能仅按整数 ID 相加。

\textbf{读图顺序：}1. 读取两份 profile\quad 2. 建立规范化键\quad 3. 重写引用\quad 4. 相同语义才累加。
\textbf{机制依据：}\cite{proj-pprof-model}\cite{proj-pprof-merge}。
\end{minipage}
\end{figure}
\subsubsection{先重写对象引用，再计算样本身份}
pprof 的 profile 不是一张平铺函数表，而是一张有依赖的对象图：Sample 引用 Location，Location 引用 Mapping 和带行号的 Function，函数名与路径等另有字符串表示。输入文件里的 ID 是局部索引。\texttt{profile/\allowbreak{}merge.\allowbreak{}go:\allowbreak{}135–147} 因而同时维护“本次输入的 ID 缓存”和“所有输入共享的语义表”。处理新文件时前者重建，后者继续累积。\cite{proj-pprof-merge}

回到案例，P1 的 location 7 与 P2 的 location 42 只有在映射和函数信息完成归一化后才能判断等价。映射匹配使用按页对齐的大小、文件偏移和身份；有 Build ID 优先用它，否则退化到文件名，缺失时还存在较弱的占位路径。这里的风险非常实际：同名二进制被重新编译后，其地址含义可能已经改变；文件名相同远不如 Build ID 相同可靠。

\texttt{mapMapping} 为匹配映射计算“已有起始地址减去新输入起始地址”的重定位差。P2 的 \texttt{0x701120} 加上 \texttt{0x400000\ensuremath{-}0x700000} 后成为 \texttt{0x401120}。后续 location 的语义键使用相对映射起点的地址，同时包含映射 ID、折叠状态，以及函数、行、列组成的序列。于是 ASLR 被消除，内联信息仍参与身份。只用函数名或绝对 PC 合并都会丢失不同层次的信息。

\subsubsection{源码精读一：为什么值不在去重键里}
下面摘录 \texttt{profile/\allowbreak{}merge.\allowbreak{}go:\allowbreak{}155–164} 的核心，省去函数后半段创建新样本的代码。这是源码片段，不是可独立编译的程序。

\begin{codeblock}
func (pm *profileMerger) mapSample(src *Sample) *Sample {
    // Check memoization table
    k := pm.sampleKey(src)
    if ss, ok := pm.samples[k]; ok {
        for i, v := range src.Value {
            ss.Value[i] += v
        }
        return ss
    }
\end{codeblock}

参数 \texttt{src} 是待合并的输入样本。\texttt{sampleKey} 先把局部位置变为合并后的语义位置，再编码标签。查表成功说明两条记录代表同一调用路径与同一组标签。循环按列相加，保留 count 与 nanoseconds 两个量纲；最后返回已存在的样本对象。若值本身也进入键，\texttt{[3,30 ms]} 与 \texttt{[2,20 ms]} 会得到不同键，从而永远不能汇总。

源码 \texttt{196–245} 行使用变长整数及带长度的字符串编码，避免简单拼接导致边界歧义。字符串标签的键排序后写入，消除 Go map 遍历顺序的不稳定；数值标签连同单位进入键。注意它排序的是标签键，不是任意重排同一键下的全部值。读取这类代码时，不能把“看到了 sort”扩大成所有标签表示都具有任意顺序不变性。

现在可以逐行执行案例：P1-A 首次生成 encode 栈 \texttt{/\allowbreak{}pay} 的条目；P2-A 映射后命中并加为 \texttt{[5,50 ms]}；P2-B 的栈不同，创建第二条；P2-C 的标签不同，创建第三条。先删标签再合并虽然可以得到 encode 共 60 ms，却无法以后恢复 \texttt{/\allowbreak{}pay} 的 50 ms。按完整键合并能够保持各分组的总权重，但不再保留每个输入各自贡献多少，因此不是对原始记录的可逆压缩。删除标签还会进一步丢失分组信息。

\subsubsection{合并头部不是重新估计采样过程}
\texttt{combineHea\allowbreak{}ders:\allowbreak{}468–518} 先检查 period 类型和有序的 sample 类型/单位是否兼容，再累加 Duration，取最早 TimeNanos，并使用最大的 Period。它没有因为选了最大 Period 就重新缩放所有值。已按纳秒给出的 CPU 值应作为纳秒相加；不同频率来源的 count 要换算时间，则必须保留或利用各自周期，不能事后对合并 count 乘一个任意周期。

例如两个来源各 1 次，周期分别为 10 ms 和 20 ms，正确估计时间为 30 ms。若仅剩总 count 2 与最大周期 20 ms，再相乘就变成 40 ms。这个反例不是说 pprof 的所有报告都会这样算，而是说明不能把头部 Period 选择误解为已经完成了重新加权。第 05 章讨论的采样误差仍然存在；格式兼容不等于估计量可直接比较。

Duration 相加也有明确边界。连续窗口的 10 s 与 20 s 合并可解释为累计采集 30 s；两个实例在同一个 10 s 墙钟区间采集，合并 Duration 为 20 s，并不表示服务经历了 20 s。计算核占用或每请求成本时，分母应来自所选估计目标，不能盲从某个文件头字段。

\subsubsection{差分与报告为何还能改变结论}
\texttt{internal/\allowbreak{}driver/\allowbreak{}fetch.\allowbreak{}go:\allowbreak{}63–89} 的差分路径可先调用 \texttt{p.\allowbreak{}Normalize(pbase)}：按各 sample type 的总量，把当前 profile 缩放到基线总量。例如当前总量 200、基线总量 100，当前值会统一乘 0.5。随后才给基线值乘以 \ensuremath{-}1，与当前 profile 合并，再进行符号化。

差分是带符号权重，不能当作普通正权重样本；“百分比上升”也可能来自总量变化。该文件还允许部分输入获取失败后合并成功取得的输入，因此必须检查警告，不能只看到输出文件就认定全部实例都在分母中。\cite{proj-pprof}

复杂度也可从对象图推导。设总栈引用数为 R，标签编码字节为 B，每条样本有 k 个标签键；在常见哈希表平均假设下，主要工作约为 \texttt{O(R + B + \ensuremath{\Sigma} k log k)}，另加映射、函数和位置的规范化成本。它不是最坏情况严格线性证明；深栈、大量独特标签与无效输入都能增加常数和内存。源码的稠密切片与稀疏 map 两级 ID 缓存减少常见连续 ID 的查表成本，但不改变语义身份必须先建立这一前提。

\section{Pyroscope：共享栈树、符号重写与累计量差分}
\begin{figure}[H]
\centering
\begin{tikzpicture}[x=1cm,y=1cm]
\path[fill=leaf!25] (0.00,0.00) rectangle (13.97,0.78);
\node[align=center,text width=13.92cm,font=\small] at (7.00,0.40) {root};
\path[fill=bluegray!38] (0.00,1.05) rectangle (11.64,1.83);
\node[align=center,text width=11.59cm,font=\small] at (5.83,1.45) {serve};
\path[fill=bluegray!51] (0.00,2.10) rectangle (6.97,2.88);
\node[align=center,text width=6.92cm,font=\small] at (3.50,2.50) {encode};
\path[fill=bluegray!51] (7.00,2.10) rectangle (11.64,2.88);
\node[align=center,text width=4.59cm,font=\small] at (9.33,2.50) {hash};
\path[fill=bluegray!38] (11.67,1.05) rectangle (13.97,1.83);
\node[align=center,text width=2.25cm,font=\small] at (12.83,1.45) {cleanup};
\node[anchor=north west,align=left,text width=14cm,font=\scriptsize] at (0,-.45) {1=root: flat 0 /\allowbreak{} cum 12；2=serve: flat 0 /\allowbreak{} cum 10；3=encode: flat 6 /\allowbreak{} cum 6；4=hash: flat 4 /\allowbreak{} cum 4；5=cleanup: flat 2 /\allowbreak{} cum 2};
\end{tikzpicture}
\caption{共享调用前缀怎样避免重复存储}\label{fig:15-trie}
\begin{minipage}{0.96\linewidth}\footnotesize 前缀树教学重建；不同项目、版本可采用不同方向和编号，不能据此认定内部布局完全相同。

\textbf{读图顺序：}1. 第一条调用路径\quad 2. 复用相同前缀\quad 3. 独立分支\quad 4. 词典与压缩是不同层。
\textbf{机制依据：}\cite{proj-pyro-tree}\cite{proj-pyroscope-store}。
\end{minipage}
\end{figure}
\begin{figure}[H]
\centering
\begin{tikzpicture}[x=1cm,y=1cm]
\node[align=center,text width=7cm,font=\small\bfseries] at (3.55,2.05) {正常累计增长};
\node[draw=linegray!75!black,fill=linegray!13,rounded corners=3pt,text width=1.95cm,minimum height=1.12cm,align=center,font=\scriptsize] (p0-input) at (1.18,0.12) {本次累计值\\{\scriptsize 值：160}};
\node[draw=linegray!75!black,fill=linegray!13,rounded corners=3pt,text width=1.95cm,minimum height=1.12cm,align=center,font=\scriptsize] (p0-base) at (1.18,-1.52) {前一基线\\{\scriptsize 值：100}};
\node[draw=leaf!75!black,fill=leaf!13,rounded corners=3pt,text width=1.95cm,minimum height=1.12cm,align=center,font=\scriptsize] (p0-check) at (3.55,-0.70) {是否单调};
\node[draw=leaf!75!black,fill=leaf!13,rounded corners=3pt,text width=1.95cm,minimum height=1.12cm,align=center,font=\scriptsize] (p0-delta) at (5.92,0.12) {区间输出\\{\scriptsize 值：60}};
\node[draw=linegray!75!black,fill=linegray!13,rounded corners=3pt,text width=1.95cm,minimum height=1.12cm,align=center,font=\scriptsize] (p0-next) at (5.92,-1.52) {下次基线\\{\scriptsize 值：160}};
\draw[-{Stealth[length=2mm]},forest!70!black] (p0-input.east) -- (p0-check.west);
\draw[-{Stealth[length=2mm]},forest!70!black] (p0-base.east) -- (p0-check.west);
\draw[-{Stealth[length=2mm]},forest!70!black] (p0-check.east) -- (p0-delta.west);
\draw[-{Stealth[length=2mm]},forest!70!black] (p0-check.east) -- (p0-next.west);
\node[align=center,text width=7cm,font=\small\bfseries] at (11.25,2.05) {检测下降并重置};
\node[draw=linegray!75!black,fill=linegray!13,rounded corners=3pt,text width=1.95cm,minimum height=1.12cm,align=center,font=\scriptsize] (p77-input) at (8.88,0.12) {本次累计值\\{\scriptsize 值：20}};
\node[draw=linegray!75!black,fill=linegray!13,rounded corners=3pt,text width=1.95cm,minimum height=1.12cm,align=center,font=\scriptsize] (p77-base) at (8.88,-1.52) {前一基线\\{\scriptsize 值：160}};
\node[draw=leaf!75!black,fill=leaf!13,rounded corners=3pt,text width=1.95cm,minimum height=1.12cm,align=center,font=\scriptsize] (p77-check) at (11.25,-0.70) {是否单调};
\node[dashed,draw=linegray!75!black,fill=linegray!13,rounded corners=3pt,text width=1.95cm,minimum height=1.12cm,align=center,font=\scriptsize] (p77-delta) at (13.62,0.12) {区间输出\\{\scriptsize 不纳入}};
\node[draw=leaf!75!black,fill=leaf!13,rounded corners=3pt,text width=1.95cm,minimum height=1.12cm,align=center,font=\scriptsize] (p77-next) at (13.62,-1.52) {下次基线\\{\scriptsize 值：20}};
\draw[-{Stealth[length=2mm]},forest!70!black] (p77-input.east) -- (p77-check.west);
\draw[-{Stealth[length=2mm]},forest!70!black] (p77-base.east) -- (p77-check.west);
\draw[-{Stealth[length=2mm]},linegray,dashed] (p77-check.east) -- (p77-delta.west);
\draw[-{Stealth[length=2mm]},forest!70!black] (p77-check.east) -- (p77-next.west);
\end{tikzpicture}
\caption{累计计数重置时，空区间不等于零分配}\label{fig:15-delta-reset}
\begin{minipage}{0.96\linewidth}\footnotesize 固定版本单栈教学重建；“不纳入”的输出节点表示空样本，不是数值 0。多栈原地更新与重置交互另见正文及复现记录，不能由单栈例证明整个算法正确。

\textbf{读图顺序：}1. 首次只建立基线\quad 2. 正常累计增长\quad 3. 检测下降并重置\quad 4. 在新基线上继续。
\textbf{连线语义：}是否单调\ensuremath{\rightarrow}\allowbreak{}区间输出：可信时相减；是否单调\ensuremath{\rightarrow}\allowbreak{}下次基线：更新基线。
\textbf{机制依据：}\cite{proj-pyro-delta}。
\end{minipage}
\end{figure}
\subsubsection{样本值和调用栈应分开保存}
假设服务每秒产生许多 encode \ensuremath{\rightarrow}\allowbreak{} serve \ensuremath{\rightarrow}\allowbreak{} main 样本。若每条都携带三组函数名、文件名和行号，存储和传输成本会主要花在重复符号上。Pyroscope 的符号分区把字符串、映射、函数、位置分别去重，再为调用栈分配可复用 ID；profile 记录保存时间、类型及“栈 ID—权重”。这种设计主要减少跨样本重复存储的符号和调用路径。\cite{proj-pyroscope-store}

\texttt{dedup\_\allowbreak{}slice.\allowbreak{}go:\allowbreak{}32–104} 的处理次序是字符串、映射、函数、位置、样本。这个顺序来自依赖关系：位置引用映射和函数，后两者又引用字符串；只有被引用对象的新 ID 已知，引用者才能安全改写。跨分区压缩的 \texttt{rewriter.\allowbreak{}go:\allowbreak{}128–217} 进一步先找需要的栈和符号，再按依赖顺序追加和改号。把旧 ID 原样复制到目标分区会把“42”错误指向另一个位置，文件格式可以合法而语义已经损坏。

\subsubsection{一棵从根到叶插入、从叶到根恢复的树}
\texttt{stacktrace\_\allowbreak{}tree.\allowbreak{}go} 的可变节点保存四个整数：父节点 p、位置引用 r、第一个子节点 fc、下一个兄弟节点 ns。输入栈是叶到根，插入时反向遍历，因而共享调用前缀。对于本章两条栈，可以得到下面的教学树；数字是本例假设的节点号，不承诺与任意运行实例一致。

\begin{codeblock}
0 sentinel
+- 1 main
   +- 2 serve
      +- 3 encode    stack ID = 3
      +- 4 hash      stack ID = 4
/pay 与 /admin 可以引用同一个 stack ID 3
标签区分样本，树只表示调用路径
\end{codeblock}

根方向的 main、serve 只保存一次，叶节点号即可作为完整栈的入口。这里称“共享调用前缀”是以根到叶阅读而言；若按 pprof 的叶到根数组阅读，它共享的是数组后缀。没有先约定方向，讨论前缀压缩很容易互相矛盾。\cite{proj-pyro-tree}

\subsubsection{源码精读二：恢复路径只需不可变字段}
下面摘自 \texttt{stacktrace\_\allowbreak{}tree.\allowbreak{}go:\allowbreak{}89–101}，仅省略源码中的注释。它是片段，接收已经存在的树和栈 ID。

\begin{codeblock}
func (t *stacktraceTree) resolve(dst []int32, id uint32) []int32 {
    dst = dst[:0]
    if id >= uint32(len(t.nodes)) {
        return dst
    }
    for i := int32(id); i > 0; i = t.nodes[i].p {
        dst = append(dst, t.nodes[i].r)
    }
    return dst
}
\end{codeblock}

第一条赋值清空逻辑长度，并在容量足够时复用缓冲区，减少重复分配；遇到越界 ID 则直接返回空结果。循环从叶节点开始，读取当前位置引用 r，再跳父节点 p；当遇到保留的节点 0 时停止。以 ID 3 输入，输出 encode、serve、main，恰好恢复 pprof 所需方向。恢复时不读取 fc 和 ns，所以新增兄弟节点不改变已有路径。

并发安全仍需进一步看切片。\texttt{partition\_\allowbreak{}memory.\allowbreak{}go:\allowbreak{}131–149} 在读锁下复制节点切片头，再释放锁逐栈解析；注释解释 p/r 插入后不变，但切片扩容会改切片头，必须处理这一层竞争。持久化的 parentPointerTree 只需 p/r，省去插入专用的 child/sibling 字段。这里的空间节约来自读写阶段分工，不能简单归纳为“树天然线程安全”。

恢复长度 d 的栈需要 O(d) 次父指针访问。插入却不保证 O(d)：源实现沿兄弟链寻找匹配子节点，某一层分支越多，线性扫描越长；代价是各层经过兄弟数之和。给每个节点增加 child 哈希表可以改善宽树查找，却引入额外分配和哈希表空间。该实现另用栈哈希到 ID 的缓存加速重复栈，先读锁查命中，未命中才写锁插入。源码的 size 方法还明确没有计入 map 占用，因此不能直接把它当作完整内存测量。

\subsubsection{累计分配 profile 必须先变成区间增量}
CPU profile 通常已经描述某段时间的样本；\texttt{alloc\_\allowbreak{}space} 和 \texttt{alloc\_\allowbreak{}objects} 则具有累计语义。Pyroscope \texttt{delta.\allowbreak{}go} 的这条实现路径仅对 memory 的这两种 sample 类型启用差分，不应外推到所有 CPU、live heap 或 block 数据。\cite{proj-pyro-delta}

该模块维护 \texttt{fingerprint \ensuremath{\rightarrow}\allowbreak{} stack ID \ensuremath{\rightarrow}\allowbreak{} 最高累计值}。fingerprint 是根据标签集合计算的系列标识。第一份 profile 只建立基线，返回空样本；后续对每条栈做当前值减历史值。如果发现计数下降，代码会重建基线并丢弃这一可疑区间。可能的原因包括进程重启、计数重置或旧数据迟到，需结合输入记录判断。假设单栈累计分配字节依次为 100、160、20、45，其区间输出为“空、60、空、25”。第一个空并不表示第一段零分配，也不是上传丢失。

为什么不把 20\ensuremath{-}160 得到的 \ensuremath{-}140 当成释放？因为 alloc\_space 是累计分配量，释放要由其他指标表达，负增长违背该计数的模型。另一种方案是遇到重置就输出当前 20；这样可能保留重启后的活动，但需要明确知道新起点时间，否则区间速率的分母不确定。当前实现选择保守丢弃，可比性与覆盖率之间存在取舍。

还需审查原地更新顺序。该版本 \texttt{deltaSamples:\allowbreak{}85–99} 一边遍历一边修改输入 Values；若较后的栈出现下降，\texttt{computeDelta:\allowbreak{}63–68} 会从已被部分修改的输入重建基线。

考虑两条栈：旧值为 \texttt{[100,100]}，新值为 \texttt{[160,20]}。第一条栈先被减成 60，第二条栈随后触发重置，重建所见便是 \texttt{[60,20]}。下一份 \texttt{[180,30]} 因而得到 \texttt{[120,10]}；若使用原始新值重建基线，则应为 \texttt{[20,10]}。

本书用逐字复制的原始 \texttt{delta.\allowbreak{}go} 和 Samples 方法完成了最小隔离 Go 测试，复现了这个序列，并用正常增长和首栈重置作对照。输入保持同一系列、同一 sample type 和排序后的稳定栈 ID，符合真实调用链允许单份 profile 含多条栈的条件。外围标签与类型使用最小桩；本次没有运行完整服务或上游测试套件，结论也不外推到后续版本。

可验证的替代设计有两种：先检查全部栈是否需要重置，再统一相减；或保存原始输入，专门用于重建基线。\cite{proj-pyro-delta}

系列身份因此成为算法的一部分。若两个进程共享不合适的 fingerprint，交错到达的累计值会彼此污染；若旧数据迟到，也可能伪装成计数重置。排查“刚重启后曲线空一段”时，应先核对基线与到达顺序，再决定是否需要重试。评估应同时报告差分后覆盖区间与丢弃原因，不能仅展示一条平滑曲线。

\subsubsection{查询和写入的近似必须分别审计}
该版本的 distributor 在较昂贵的规范化之前执行租户限流。它可将数据发送到 ingester（写入节点）、segment writer（段写入组件），或同时发给两者。可选聚合路径标注的 best-effort（尽力处理）和 at-most-once（至多一次）只约束该路径，不代表全系统的交付保证。查询则按时间拆分内存 ingester 与历史 store gateway 的范围，分别取得树后归并，最终才受最大节点数等展示限制影响。\cite{proj-pyroscope} \cite{proj-pyroscope-store}

“写入前丢弃一个 profile”和“报告省略小节点”都让图变小，但前者改变可观测总量，后者可能只改变展示粒度。诊断时应沿输入数、入库数、查询总权重和渲染节点数分别检查。实际租户限流、对象存储和查询延迟仍需部署后测量。

\section{Parca：按需补栈、列式执行与延后符号化}
\begin{figure}[H]
\centering
\begin{tikzpicture}[x=1cm,y=1cm]
\node[draw=linegray!75!black,fill=linegray!13,rounded corners=3pt,text width=2.08cm,minimum height=1.12cm,align=center,font=\small] (p0-rows) at (1.25,0.00) {样本列\\{\scriptsize 时间/\allowbreak{}服务/\allowbreak{}栈/\allowbreak{}值}};
\node[draw=linegray!75!black,fill=linegray!13,rounded corners=3pt,text width=2.08cm,minimum height=1.12cm,align=center,font=\small] (p0-filter) at (3.75,0.00) {时间＋标签\\{\scriptsize 保留满足条件的行}};
\node[draw=linegray!75!black,fill=linegray!13,rounded corners=3pt,text width=2.08cm,minimum height=1.12cm,align=center,font=\small] (p0-convert) at (6.25,0.00) {逐行换算\\{\scriptsize count \ensuremath{\times} 本行 period}};
\node[draw=linegray!75!black,fill=linegray!13,rounded corners=3pt,text width=2.08cm,minimum height=1.12cm,align=center,font=\small] (p0-group) at (8.75,0.00) {栈键分组\\{\scriptsize 兼容权重求和}};
\node[draw=linegray!75!black,fill=linegray!13,rounded corners=3pt,text width=2.08cm,minimum height=1.12cm,align=center,font=\small] (p0-symbols) at (11.25,0.00) {符号解析\\{\scriptsize 字典/\allowbreak{}构建身份}};
\node[draw=leaf!75!black,fill=leaf!13,rounded corners=3pt,text width=2.08cm,minimum height=1.12cm,align=center,font=\small] (p0-flame) at (13.75,0.00) {火焰树\\{\scriptsize 共享前缀归并}};
\draw[-{Stealth[length=2mm]},forest!70!black] (p0-rows.east) -- (p0-filter.west);
\draw[-{Stealth[length=2mm]},forest!70!black] (p0-filter.east) -- (p0-convert.west);
\draw[-{Stealth[length=2mm]},forest!70!black] (p0-convert.east) -- (p0-group.west);
\draw[-{Stealth[length=2mm]},forest!70!black] (p0-group.east) -- (p0-symbols.west);
\draw[-{Stealth[length=2mm]},forest!70!black] (p0-symbols.east) -- (p0-flame.west);
\end{tikzpicture}
\caption{查询先选择事实，再重建调用树}\label{fig:15-column-query}
\begin{minipage}{0.96\linewidth}\footnotesize 重建 Parca 查询算子顺序。samples/count 先逐行乘本行 period 并更改结果单位，再聚合；已是时间权重的数据不再重复乘周期。数值为教学反例。

\textbf{读图顺序：}1. 定义查询样本范围\quad 2. 先按各行周期换算\quad 3. 聚合同类权重\quad 4. 恢复可读栈\quad 5. 呈现成本分布。
\textbf{机制依据：}\cite{proj-parca-query}\cite{proj-parca-store}。
\end{minipage}
\end{figure}
\begin{figure}[H]
\centering
\begin{tikzpicture}[x=1cm,y=1cm]
\node[draw=linegray!75!black,fill=linegray!13,rounded corners=3pt,text width=2.58cm,minimum height=1.12cm,align=center,font=\small] (p0-sample) at (1.50,0.00) {样本批次\\{\scriptsize 时间/\allowbreak{}标签/\allowbreak{}stack ID}};
\node[draw=linegray!75!black,fill=linegray!13,rounded corners=3pt,text width=2.58cm,minimum height=1.12cm,align=center,font=\small] (p0-backend) at (4.50,0.00) {后端字典\\{\scriptsize 哪些 ID 未知？}};
\node[draw=linegray!75!black,fill=linegray!13,rounded corners=3pt,text width=2.58cm,minimum height=1.12cm,align=center,font=\small] (p0-request) at (7.50,0.00) {缺失 ID 列表};
\node[draw=leaf!75!black,fill=leaf!13,rounded corners=3pt,text width=2.58cm,minimum height=1.12cm,align=center,font=\small] (p0-cache) at (10.50,0.00) {Agent 栈缓存\\{\scriptsize LRU 可能淘汰}};
\node[dashed,draw=linegray!75!black,fill=linegray!13,rounded corners=3pt,text width=2.58cm,minimum height=1.12cm,align=center,font=\small] (p0-full) at (13.50,-0.97) {发送完整栈\\{\scriptsize 不纳入}};
\node[draw=leaf!75!black,fill=leaf!13,rounded corners=3pt,text width=2.58cm,minimum height=1.12cm,align=center,font=\small] (p0-missing) at (13.50,0.97) {标记不完整\\{\scriptsize is\_complete=false}};
\draw[-{Stealth[length=2mm]},forest!70!black] (p0-request.east) -- (p0-cache.west);
\draw[-{Stealth[length=2mm]},forest!70!black] (p0-cache.east) -- (p0-missing.west);
\end{tikzpicture}
\caption{先发栈 ID，后端只索要未知栈}\label{fig:15-stack-exchange}
\begin{minipage}{0.96\linewidth}\footnotesize 重建 Parca Agent 双阶段传输；箭头按步骤显示所走分支。权重与结构分开减少重复，但新栈有额外往返，缓存失配需显式保留完整性。

\textbf{读图顺序：}1. 先发送轻量样本\quad 2. 后端只返回缺失 ID\quad 3. 查 Agent 缓存并发送\quad 4. 覆盖缓存淘汰分支。
\textbf{连线语义：}后端字典\ensuremath{\rightarrow}\allowbreak{}缺失 ID 列表：有未知ID；Agent 栈缓存\ensuremath{\rightarrow}\allowbreak{}发送完整栈：命中；Agent 栈缓存\ensuremath{\rightarrow}\allowbreak{}标记不完整：已淘汰。
\textbf{机制依据：}\cite{proj-parca-agent}。
\end{minipage}
\end{figure}
\subsubsection{Agent 先发送什么，后端还需要什么}
Parca Agent 的 \texttt{ReportTrac\allowbreak{}eEvent} 先缓存 TraceHash 对应的栈，再将标签、时间戳和权重写入 Arrow builder。每次事件都检查栈缓存，因为按最近最少使用策略（LRU）管理的条目可能已被淘汰。标签提取、重标记过滤、无效 UTF-8 跳过也发生在这条路径；“采到了事件”不自动等于“后端得到带完整标签的样本”。\cite{proj-parca-agent}

采样来源决定值和单位。普通 CPU 采样写入 value=1、sample 类型 samples/count、period 类型 cpu/nanoseconds；off-CPU 来源用 OffTime 写入 wallclock/nanoseconds。后者表达离开 CPU 的时间，不能加入 CPU count 后画成统一时间。\texttt{buildSampl\allowbreak{}eRecord:\allowbreak{}1096–1120} 对该采样批次写入配置频率对应的周期，并标记 delta 语义。

这解释了本章案例与线上传输的差异：表中的 \texttt{[3,30 ms]} 是教学用聚合记录，Agent 可以先发三条 value=1 的事件。只有在明确周期为 10 ms 时，它们才合计 30 ms。不能假设每个采集器的传输记录都直接携带 count 和 cpu 两列，更不能因列名都叫 value 就跨来源相加。

\texttt{reportData\allowbreak{}ToBackend:\allowbreak{}969–1081} 首先发送样本批次，其中保留权重、时间和标签，但调用栈仅用 ID 表示。若后端已知这些 ID，返回空记录即可结束；否则会在 \texttt{stacktrace\_\allowbreak{}id} 列中返回尚未解析的 ID 列表，Agent 查询缓存后再发送完整栈。响应字段数、列名和二进制类型都有检查。这个双阶段协议把“经常变化的权重/时间”和“多数时候不变的结构”分开。\cite{proj-parca-agent}

对大量重复栈，它减少重复传输；对全新栈，它增加一次请求往返及状态管理。若发样本后栈被淘汰，\texttt{buildStack\allowbreak{}traceRecor\allowbreak{}d:\allowbreak{}1133–1137} 会输出缺失 locations 和 \texttt{is\_\allowbreak{}complete=false}；缺少可执行文件信息也使记录不完整。不能将未知栈静默补成空调用链并参与正常占比。较大的缓存、主动随批携带首次栈、后端持久字典各有内存、带宽和一致性代价，选择要看栈复用率与批次延迟。

\subsubsection{列存保存的不是一棵预先画好的火焰图}
Parca 的 pprof WriteRaw 入口通过 normalizer 转换为 Arrow record，再交给 ingester 插入表。normalizer 逐 sample type 拆分值，保留类型、单位、周期、时间和标签。因而同一 profile 的 count 与 cpu 两列不会未经说明就变成同一数值序列。\cite{proj-parca-store}

\texttt{pkg/\allowbreak{}profile/\allowbreak{}schema.\allowbreak{}go:\allowbreak{}38–158} 给出了实际物理布局：labels 是可空动态列，类型与单位适合字典编码；stacktrace 是重复的可空字符串列，配合字典编码和 LZ4；时间和值使用 delta binary packed 编码。将同列值放在一起可以让重复标签和相近时间更容易压缩，也使查询只读取所需维度。这里的 delta 是整数编码方式，和上一节累计分配差分是两个概念。

Arrow 则描述内存中批量数据的布局及生命周期；不能仅凭使用 Arrow 就推断落盘格式、持久化保证或查询一定零拷贝。该版本 ingester 在插入前显式初始化字典，注释说明这是避免惰性初始化并发竞争的处理。优秀架构不只在于列式大方向，也在于引用计数、字典初始化和错误路径这些细节。

\subsubsection{源码精读三：乘周期应发生在聚合之前}
\texttt{pkg/\allowbreak{}parcacol/\allowbreak{}querier.\allowbreak{}go:\allowbreak{}1431–1479} 先选择 stacktrace 和请求的分组列，再把 count 转成 period 的量纲，最后聚合。核心片段如下：\cite{proj-parca-query}

\noindent\begin{minipage}{\linewidth}
\begin{codeblock}
var valueCol logicalplan.Expr = logicalplan.Col(profile.ColumnValue)
if isSamplesCount(queryParts.Meta.SampleType) {
    valueCol = logicalplan.Mul(
        logicalplan.Col(profile.ColumnValue),
        logicalplan.Col(profile.ColumnPeriod),
    ).Alias(profile.ColumnValue)
    resultType = queryParts.Meta.PeriodType
}
\end{codeblock}
\end{minipage}

初始表达式读取原始值；只有 sample type 确认是 samples/count 才进入分支。乘法按行结合各自周期，随后把结果类型改成周期类型，以保证数值与元数据一致。若先 \texttt{sum(count)} 再乘单一周期，不同频率的来源会产生前节 30 ms 误算为 40 ms 的问题。这里的写法展示了估计量语义如何直接约束查询算子顺序。

完整执行链是扫描、过滤、投影、聚合、最终投影，之后再符号化。\texttt{selectMerge} 对时间使用包含两端的比较，而 \texttt{QueryRange} 的部分范围路径使用严格比较；读者实现相邻窗口拼接时应检查具体接口，不能凭经验假定一律半开区间。还应明确 \texttt{start/\allowbreak{}end} 端点是否包含在查询范围内。

把教学输入代入查询：先限定 \texttt{/\allowbreak{}pay} 与时间，再分别算 3\ensuremath{\times}10、2\ensuremath{\times}10、1\ensuremath{\times}10 ms，按栈汇总为 50 与 10。若不保留 route 分组且不过滤，它会得到 encode 60、hash 10；这不是引擎算错，而是查询已经要求丢掉 route 维度。为了回答两个接口分别消耗多少，必须在过滤或分组阶段保留该维度，而不是等火焰图生成后猜测。

\subsubsection{延后符号化节省了什么，又依赖什么}
\texttt{QueryMerge:\allowbreak{}1367–1401} 先获取聚合结果，再调用 SymbolizeArrowRecord。若许多事件共享同一栈，先聚合可减少重复符号处理。符号化路径取决于配置和文件中的可用信息，可以使用外部地址解析器、DWARF、Go pclntab 或符号表；缓存和二进制定位以 Build ID 等信息为关键。\cite{proj-parca}

代价是查询依赖符号保留策略。部署时只保留正在运行的二进制，可能使一个月前的剖析只能显示地址。文件路径不能替代构建身份。源码行归因需要可用的文件与行号映射，但未必需要完整 DWARF：该版本 GoLiner 可通过 Go 符号表的 \texttt{PCToLine} 获得文件和行号。恢复内联调用则需要额外信息；\texttt{pkg/\allowbreak{}symbol/\allowbreak{}addr2line/\allowbreak{}go.\allowbreak{}go:\allowbreak{}92–105} 的注释明确指出这条 Go 路径尚未展开内联函数。应分别检查地址、函数、行号和内联帧的解析程度。

对于过滤后 N 行、唯一栈 U、平均深度 d，可以粗略推导：扫描受读取列和谓词选择性影响，哈希聚合通常约 O(N)，还原栈/符号的工作与 U 和 d 更相关。没有实测就不能声称某后端一定快于另一个；高基数标签使 N 与 U 的压缩关系变差，复杂正则或符号缺失也可能主导查询。这个成本分解比产品口号更能指导实验。

\section{Delve 交叉验证：暂停现场如何补齐采样解释}
\subsubsection{一个 goroutine 不总对应一条正在运行的线程}
当采样结果提示 serve 热点时，Delve 能回答“当前局部变量是什么、goroutine 停在哪里、谁持有状态”。它不能仅凭一次暂停给出 serve 消耗了多少 CPU。时间权重来自重复采样或计时，调试快照来自一个具体执行点；混淆两者会把等待很久的 goroutine 错认成 CPU 热点。

\texttt{pkg/\allowbreak{}proc/\allowbreak{}stack.\allowbreak{}go:\allowbreak{}130–146} 明确区分寄存器来源：G 正在某线程运行时从线程读取；不在线程运行时，从保存的 PC、SP、BP、LR 构造起始寄存器，再结合映像静态基址。选择错误来源会把 OS 线程当前执行的另一个 goroutine 当作目标 G。\cite{proj-delve-stack}

随后 stackIterator 逐帧前进，\texttt{advanceRegs:\allowbreak{}490–553} 查询当前 PC 对应的帧描述条目（FDE），建立 DWARF 帧上下文，再执行 CFA 与寄存器恢复规则。CFA 可以看作对调用者帧进行定位的参考位置，具体由规则计算；不是所有架构都等于当前 SP 加固定字节数。缺少 FDE 时还有架构相关修正路径，不能把 DWARF 作为唯一信息来源。

\subsubsection{物理栈帧与逻辑内联调用不是一回事}
编译器把某函数内联到 encode 后，机器栈中未必多一帧，但源码层仍希望看到该函数。\texttt{appendInli\allowbreak{}neCalls:\allowbreak{}426–487} 读取 DWARF 内联树，生成标记 \texttt{Inlined=true} 的逻辑帧；非顶层帧还对 call PC 做减一处理，避免把返回后的指令错误归到下一行或函数边界。

因此 Delve 多出的某一逻辑帧不一定证明采样器漏采；需要先比较是否启用内联展开、二者所用符号是否一致、栈中的 PC 是调用点还是返回点。相反，源程序优化导致变量无法恢复，也不自动表示程序数据被破坏。调试信息、寄存器分配和生命周期共同决定可观测性。

\subsubsection{部分栈和缓存都带时间条件}
\texttt{stack.\allowbreak{}go:\allowbreak{}380–412} 在已经有帧时，把后续展开错误附为带 Err 的 Stackframe；若第一帧就失败则直接返回错误。这和 eBPF 的部分成功问题相呼应：必须保留失败边界，不能只数正常显示的行。

\texttt{target.\allowbreak{}go:\allowbreak{}65–68} 注释要求每次恢复执行都清空 goroutine 缓存。一次暂停时得到的 G 地址、栈和状态不能无条件沿用到下一次暂停。\texttt{native/\allowbreak{}proc.\allowbreak{}go:\allowbreak{}395–417} 还把 ptrace 操作固定到同一个 OS 线程，利用 channel 串行提交。应用层 goroutine 调度自由，不代表底层调试协议也允许操作任意换线程。\cite{proj-delve}

一套有解释力的验证流程是：先从 profile 定位 serve 热路径，保留对应构建产物；在可控环境用兼容版本的 Delve 暂停相同代码，检查具体分支与状态；再恢复运行，用基准测试验证修改是否降低对应时间权重。Delve 本身改变调度和时序，所以不能拿“单步时延”替代服务延迟测量。本章所读旧 tag 用于展示机制，真实 Go 1.25 实验应另核对支持版本。

\section{把实现比较变成可检验的系统设计}
\subsubsection{一个查询结果应能追溯回哪些中间量}
现在反向追踪 \texttt{/\allowbreak{}pay} 中 encode 的 50 ms：界面宽度来自聚合权重；权重来自过滤后两个兼容来源；二者依 Build ID 和相对地址归为同一路径；局部 ID 已全部重写。栈顶 PC 来自采样时的寄存器，其余调用路径则依赖正确的栈展开。任一层不满足条件，50 ms 都可能变成另一个含义。

为系统建立检查点比单独验证火焰图更有效。教学案例的预期如下：规范化前四条聚合记录，语义合并后三条，唯一调用路径两条，全局 70 ms，\texttt{/\allowbreak{}pay} 60 ms，encode 在 \texttt{/\allowbreak{}pay} 为 50 ms。树去重可以减少结构节点，却不能改变总权重；标签过滤会改变所选样本的权重；百分比分母可以是过滤前的全局总量，也可以是过滤后的所选总量，必须明确采用哪一种；渲染裁剪不应伪装成完整证据。

将这些性质写成变形测试：重编号所有局部 ID，结果不变；给同一映射的地址与基址加同一偏移，结果不变；把一条样本拆成两个相同键且权重和不变的样本，结果不变；仅改变 route，应改变分组结果而不改变全局总量；改变 Build ID，不能继续无条件按原地址解释。这些测试针对语义不变量，比对着某个内部函数重复写相同实现更有价值。

\subsubsection{四类设计选择及其可验证代价}
\begingroup\small\setlength{\tabcolsep}{4pt}
\begin{longtable}{>{\raggedright\arraybackslash}p{0.215\linewidth}>{\raggedright\arraybackslash}p{0.215\linewidth}>{\raggedright\arraybackslash}p{0.215\linewidth}>{\raggedright\arraybackslash}p{0.215\linewidth}}
\toprule
\textbf{设计选择} & \textbf{节约的资源} & \textbf{增加的约束} & \textbf{适合如何验证} \\
\midrule
\endfirsthead
\toprule
\textbf{设计选择} & \textbf{节约的资源} & \textbf{增加的约束} & \textbf{适合如何验证} \\
\midrule
\endhead
内核预计算展开规则 & 每次采样的解析工作 & 规则覆盖、版本布局、有限深度 & 完整率与权重偏差随栈类型变化 \\
\addlinespace[3pt]
栈字典与共享父指针树 & 重复字符串及公共路径 & ID 重写、字典生命周期、宽树插入 & 固定样本下验证去重率和峰值内存 \\
\addlinespace[3pt]
列存先过滤聚合再符号化 & 无关行扫描、重复符号工作 & 标签基数、历史二进制可得性 & 相同查询语义下测试延迟和缺符号比例 \\
\addlinespace[3pt]
暂停调试恢复完整状态 & 运行时持续观测成本 & 进程扰动、单点时序、版本支持 & 以状态解释为目标，另测运行性能 \\
\addlinespace[3pt]
\bottomrule
\end{longtable}
\endgroup
这张表不提供产品排名。单进程问题通常可从 pprof 和 Delve 入手，必要时补充 trace 与操作系统观测；跨实例回归需要持续采集、构建身份、时间与标签查询；系统级观测还需要 eBPF 权限和不同语言栈的质量评估。前一个场景的最小配置，不应被下一个场景的全部复杂性淹没。

\subsubsection{故障注入比只测正常吞吐更能暴露设计问题}
建议至少注入四种异常：删除一个二进制的调试信息；缩小栈缓存造成按需补栈未命中；让累计分配序列乱序到达；让两个构建使用相同文件名。分别观察符号缺失、完整性标记、delta 基线重置和 Build ID 隔离。异常是否可见，比系统是否始终返回一张漂亮图更重要。

实验报告应记录原始输入与查询约定，明确样本量、单位、窗口端点、失败计数、输出总量与符号完整率。若比较 Pyroscope 与 Parca，要同时控制 profile 类型、标签、采样率、保留时间和查询范围；只比较两个默认界面的延迟没有可解释性。

本章完成了源码审读、算法推导、手算例证，以及 Pyroscope 原函数的隔离复现。尚未部署这些平台，也没有平台开销排名或生产规模实验结果。附带的研究记录给出文件与行号，可沿同一版本重复检查。下一章将用相同标准阅读研究论文：先辨认估计目标，再检查方法如何改变数据，最后判断实验到底支持了多大范围的结论。

\section{实验任务}
\subsection*{实验 15：追踪同一组样本的身份、权重与信息损失}
\begin{enumerate}
\item 使用正文四条合成记录建立表格，记录 Build ID、映射起点、局部位置 ID、route、count 和 cpu 单位。先手算三条语义样本及 70/60/50 ms 三个检查点。

\item 用自选语言实现教学版语义键：先以映射内相对地址和 Build ID 建立位置身份，再编码边界明确的有序调用路径与标签；对同键权重逐列求和。写明简化范围，避免将其当作完整 pprof 解析器。

\item 进行五种变形：局部 ID 重编号、ASLR 整体平移、样本等量拆分、route 修改、Build ID 修改。前三种应保持查询结果，第四种改变标签分组但保留全局总量，第五种要求重新建立符号身份。

\item 手工建立父指针树，验证栈 ID 3 恢复 encode/serve/main，且 /pay 与 /admin 复用树路径时仍保留不同样本。对宽分支输入，比较兄弟链扫描和按孩子建立哈希索引的查找次数与内存。

\item 先输入单栈累计分配 100\ensuremath{\rightarrow}\allowbreak{}160\ensuremath{\rightarrow}\allowbreak{}20\ensuremath{\rightarrow}\allowbreak{}45；再从头执行，并在 160 之后插入一份时间戳早于它、累计值为 150 的迟到记录，再继续输入 20、45。分别记录基线、区间输出及被丢弃的区间，说明按实例隔离和按时间排序各解决什么问题。

\item 在隔离环境运行实际项目时，先核对 Go、内核和项目的兼容版本，再用真实源码入口验证这些不变量。报告应区分教学实现、原函数隔离测试与完整平台运行结果。

\end{enumerate}
\section{自测与解析}
\noindent\textbf{1. P1 与 P2 的 encode 地址不同、位置 ID 也不同，什么条件最能支持它们合并？}
\begin{enumerate}[label=\Alph*.]
\item 局部 ID 都是正整数
\item 归一化后的映射身份、相对地址、函数行信息及样本标签兼容
\item 函数名包含相同子串
\item 两张火焰图的颜色相同
\end{enumerate}
\noindent{\color{forest}\textbf{答案 B。}}局部 ID 与 ASLR 地址不是跨文件身份。应先规范化对象依赖，再依据栈与标签形成语义键。

\noindent\textbf{2. 两条 CPU count 样本各为 1，周期分别 10 ms 和 20 ms，如何汇总 CPU 时间？}
\begin{enumerate}[label=\Alph*.]
\item count 求和后乘最大周期，得到 40 ms
\item count 求和后乘最小周期，得到 20 ms
\item 分别乘周期再求和，得到 30 ms
\item 取平均 count，得到 1 ms
\end{enumerate}
\noindent{\color{forest}\textbf{答案 C。}}不同来源要按行保留权重。pprof 合并头部选择最大 Period，并不自动为全部 count 重新加权。

\noindent\textbf{3. Pyroscope 该路径收到第一份 alloc\_space profile 后输出空样本，最合理的解释是？}
\begin{enumerate}[label=\Alph*.]
\item 应用完全没有分配
\item 首次建立累计计数基线，尚不能可靠估计该区间增量
\item 树为空所以忽略一切数据
\item 网络一定丢包
\end{enumerate}
\noindent{\color{forest}\textbf{答案 B。}}所读 delta 实现首次只建立基线；空输出不能当作零分配。计数重置时也存在保守丢弃。

\noindent\textbf{4. Agent 已返回两帧且 goroutine 标签存在，能否断言调用栈完整？}
\begin{enumerate}[label=\Alph*.]
\item 能，两帧足够
\item 能，标签能证明全部 PC 正确
\item 不能；栈展开和标签读取需分别验证，部分栈仍可能带错误
\item 不能，所有 eBPF 栈均不可靠
\end{enumerate}
\noindent{\color{forest}\textbf{答案 C。}}当前帧可以在下一帧恢复失败前已被保存。应检查错误和完整性信息，避免把非空与完整混为一谈。

\section{分析题与参考答案}
\noindent\textbf{题 1：}在贯穿案例中先删除 route 标签再合并，得到几条样本？仍能回答 /pay 的 encode CPU 时间吗？

\noindent\textbf{参考答案：}得到两条调用栈样本：encode 为 60 ms，hash 为 10 ms，总量仍为 70 ms。无法从这两条记录恢复 /pay 的 encode 50 ms，因为 /admin 的 10 ms 已混入。保留原始记录或在聚合键中保留 route 才可回答。

\noindent\textbf{题 2：}把共享栈树的孩子查找从兄弟链扫描改成哈希表，是否一定更快、更省内存？给出反例与测量方案。

\noindent\textbf{参考答案：}不一定。每层只有一个孩子的深链几乎只需一次比较；哈希表会增加对象、桶和哈希计算的成本。宽树可能受益。应分别构造深窄、浅宽、重复率高和持续出现新栈的输入，控制路径总数与权重，测插入时间、查找次数、分配和峰值内存；恢复出的栈必须一致。

\noindent\textbf{题 3：}Delve 暂停时发现许多 goroutine 的栈都经过 serve，但同一负载的 CPU profile 中 serve 的 cum 很低，这一定矛盾吗？

\noindent\textbf{参考答案：}不一定。暂停快照中的这些 goroutine 可能大多在等待，而 CPU profile 统计的是采集期间实际运行时的样本权重。应再核对构建、采集窗口、内联展开、PC 归属和栈完整性，不能用暂停时 goroutine 数乘一个假定时间来替代 CPU 总量。

\medskip\noindent\textbf{本章主要阅读：}\cite{proj-pprof}\cite{proj-pprof-model}\cite{proj-pprof-merge}\cite{proj-delve}\cite{proj-delve-stack}\cite{proj-pyroscope}\cite{proj-pyroscope-store}\cite{proj-pyro-tree}\cite{proj-pyro-delta}\cite{proj-parca}\cite{proj-parca-store}\cite{proj-parca-query}\cite{proj-parca-agent}\cite{proj-otel-ebpf}\cite{proj-ebpf-unwind}\cite{proj-go-unwind}。
\chapter{近五年前沿工作与研究方向}
\label{ch:16}
\noindent{\large\color{forest}从三篇原始论文的算法、图表与失效条件，走向 Go 的持续观测和反馈优化研究}\par\medskip
\noindent\textbf{先修要求：}第 01、03、05、06、08、11、12、15 章；会区分累计量和区间量，理解采样、实验基线和基本概率。

\noindent\textbf{完成本章后，应能够：}
\begin{itemize}
\item 复述三篇代表论文的估计目标、关键算法和原始实验分母
\item 计算阈值采样、测量窗口偏移和 TCP 时间戳匹配的教学例子
\item 识别论文结果的适用边界、对照混杂、有效性威胁及公式一致性问题
\item 把 PGO、Flight Recorder 与开放 profile 格式转化为有基线、可证伪的 Go 研究方案
\end{itemize}
\section{研究地图与证据方法：近五年意味着什么}
\subsubsection{先固定时间，再区分证据类型}
本章检索窗口为 2021-10-02 至 2026-10-02，按论文正式发表时间或官方工程文章发布日期判断是否入选。选择三篇已取得原始全文的学术工作，分别研究多资源归因、细粒度硬件计量和网络延迟观测；再以 Go PGO、Trace、Flight Recorder 和 OpenTelemetry Profiles 讨论工程演进。它们不是全部相关文献的穷尽目录，也不是按引用数自动筛出的排行榜。

\begingroup\small\setlength{\tabcolsep}{4pt}
\begin{longtable}{>{\raggedright\arraybackslash}p{0.215\linewidth}>{\raggedright\arraybackslash}p{0.215\linewidth}>{\raggedright\arraybackslash}p{0.215\linewidth}>{\raggedright\arraybackslash}p{0.215\linewidth}}
\toprule
\textbf{工作} & \textbf{时间与载体} & \textbf{研究的问题} & \textbf{主要原始证据} \\
\midrule
\endfirsthead
\toprule
\textbf{工作} & \textbf{时间与载体} & \textbf{研究的问题} & \textbf{主要原始证据} \\
\midrule
\endhead
Triangulating Python Performance Issues with Scalene & OSDI 2023 & CPU、内存、复制如何共同解释源码性能 & 算法章节，图 3—8，表 1—3 \\
\addlinespace[3pt]
Fine-Grained Hardware Profiling – Are You Using the Right Tools? & SIGMOD Record 53(2)，2024 年 6 月 & 短区间硬件计数的边界与自扰动 & Listings 2—4，图 1—5 \\
\addlinespace[3pt]
Efficient Continuous Latency Monitoring with eBPF & PAM 2023 & 高流量下持续 RTT 观测的覆盖与成本 & 图 1—9，表 1，附录算法 1 \\
\addlinespace[3pt]
Go PGO、执行追踪、Flight Recorder & 2023—2025 官方工程资料 & 将观测接入优化和异常回溯 & 发布说明、实现说明、Go 1.25.0 源码 \\
\addlinespace[3pt]
OpenTelemetry Profiles & 2024 官方介绍及 2026-10-02 规范快照 & 不同剖析来源与其他信号如何互操作 & 数据模型、状态与兼容性约束 \\
\addlinespace[3pt]
\bottomrule
\end{longtable}
\endgroup
SIGMOD Record 的期刊文章不能写成 SIGMOD 主会长文；官方技术文章也不能写成经过同样同行评审的论文。OTel 在检索日的状态快照不是一篇发表于 2026 年的新研究。入选工程资料中，明确发布日期最晚的是 2025-09-26 的 Flight Recorder 文章；这并不证明 2026 年没有其他工作。

\subsubsection{用同一套问题阅读每篇论文}
先写出估计目标：是执行时间、分配活动、存活内存，还是某段网络的往返时间（RTT）？再找观测机制：信号、分配器钩子、硬件性能监控单元（PMU）计数或协议字段。随后分析从观测到估计的规则、基线为何失败、实验如何构造真值，以及算法在哪些输入上失去前提。

本章的论文数字均标为作者原实验结果；手算例子明确标为教学构造；Go 迁移方案属于本教材推导，尚未实施。逐图阅读比转述摘要更重要：表中中位数、图中单个案例与某应用最高加速不能混写成同一个普遍结论。源码与文献的实际阅读范围及证据条目保存在配套研究记录中。

\section{Scalene：改变采样对象，比增加采样频率更关键}
\begin{figure}[H]
\centering
\begin{tikzpicture}[x=1cm,y=1cm]
\node[draw=linegray!75!black,fill=linegray!13,rounded corners=3pt,text width=4.58cm,minimum height=1.12cm,align=center,font=\small] (p0-alloc) at (2.50,0.97) {malloc 事件};
\node[draw=linegray!75!black,fill=linegray!13,rounded corners=3pt,text width=4.58cm,minimum height=1.12cm,align=center,font=\small] (p0-free) at (2.50,-0.97) {free 事件};
\node[draw=linegray!75!black,fill=linegray!13,rounded corners=3pt,text width=4.58cm,minimum height=1.12cm,align=center,font=\small] (p0-rate) at (7.50,0.97) {分配流量信号};
\node[draw=linegray!75!black,fill=linegray!13,rounded corners=3pt,text width=4.58cm,minimum height=1.12cm,align=center,font=\small] (p0-net) at (7.50,-0.97) {净增长信号};
\node[draw=leaf!75!black,fill=leaf!13,rounded corners=3pt,text width=4.58cm,minimum height=1.12cm,align=center,font=\small] (p0-line) at (12.50,0.00) {源码行归因};
\draw[-{Stealth[length=2mm]},forest!70!black] (p0-alloc.east) -- (p0-rate.west);
\draw[-{Stealth[length=2mm]},forest!70!black] (p0-alloc.east) -- (p0-net.west);
\draw[-{Stealth[length=2mm]},forest!70!black] (p0-free.east) -- (p0-net.west);
\draw[-{Stealth[length=2mm]},forest!70!black] (p0-rate.east) -- (p0-line.west);
\draw[-{Stealth[length=2mm]},forest!70!black] (p0-net.east) -- (p0-line.west);
\end{tikzpicture}
\caption{分配流量和净保留量需要分开观察}\label{fig:16-scalene}
\begin{minipage}{0.96\linewidth}\footnotesize 仅示意分配流量与净变化的区别：假设每次分配和释放的对象等大，且释放对象来自所示分配。100、20 等为教学事件数，不是 Scalene 的字节阈值、公式或原实验结果。

\textbf{读图顺序：}1. 临时对象很多\quad 2. 对象持续保留\quad 3. 两种信号分开\quad 4. 迁移到 Go 需重设计。
\textbf{连线语义：}malloc 事件\ensuremath{\rightarrow}\allowbreak{}净增长信号：增加；free 事件\ensuremath{\rightarrow}\allowbreak{}净增长信号：减少。
\textbf{机制依据：}\cite{paper-scalene}。
\end{minipage}
\end{figure}
\begin{figure}[H]
\centering
\begin{tikzpicture}[x=1cm,y=1cm]
\node[draw=linegray!75!black,fill=linegray!13,rounded corners=3pt,text width=3.33cm,minimum height=1.12cm,align=center,font=\small] (p0-a) at (1.88,0.97) {累计分配 A\\{\scriptsize 值：9}\\{\scriptsize 自上次采样以来}};
\node[draw=linegray!75!black,fill=linegray!13,rounded corners=3pt,text width=3.33cm,minimum height=1.12cm,align=center,font=\small] (p0-f) at (1.88,-0.97) {累计释放 F\\{\scriptsize 值：4}\\{\scriptsize 自上次采样以来}};
\node[draw=leaf!75!black,fill=leaf!13,rounded corners=3pt,text width=3.33cm,minimum height=1.12cm,align=center,font=\small] (p0-net) at (5.62,0.00) {净变化 |A\ensuremath{-}F|\\{\scriptsize 值：5}};
\node[draw=leaf!75!black,fill=leaf!13,rounded corners=3pt,text width=3.33cm,minimum height=1.12cm,align=center,font=\small] (p0-threshold) at (9.38,0.00) {阈值 T=4 MiB\\{\scriptsize 值：4}\\{\scriptsize 教学值}};
\node[draw=leaf!75!black,fill=leaf!13,rounded corners=3pt,text width=3.33cm,minimum height=1.12cm,align=center,font=\small] (p0-emit) at (13.12,0.00) {输出并清零\\{\scriptsize 重新累计 A/\allowbreak{}F}};
\draw[-{Stealth[length=2mm]},forest!70!black] (p0-a.east) -- (p0-net.west);
\draw[-{Stealth[length=2mm]},forest!70!black] (p0-f.east) -- (p0-net.west);
\draw[-{Stealth[length=2mm]},forest!70!black] (p0-net.east) -- (p0-threshold.west);
\draw[-{Stealth[length=2mm]},forest!70!black] (p0-threshold.east) -- (p0-emit.west);
\end{tikzpicture}
\caption{净变化阈值：抵消活动不触发，达到阈值才记录}\label{fig:16-scalene-threshold}
\begin{minipage}{0.96\linewidth}\footnotesize 按 Scalene 第 3.2 节净变化机制重建，T=4 MiB 是教学值；原论文阈值为略大于 10 MB 的素数。适合显著存量变化，不替代分配流量观测。

\textbf{读图顺序：}1. 分配 2 MiB\quad 2. 释放后抵消\quad 3. 再分配释放一次\quad 4. 连续增长跨过阈值\quad 5. 重置累计窗口。
\textbf{连线语义：}阈值 T=4 MiB\ensuremath{\rightarrow}\allowbreak{}输出并清零：达到阈值。
\textbf{机制依据：}\cite{paper-scalene}。
\end{minipage}
\end{figure}
\subsubsection{问题与失败基线：一行 Python 不等于一种资源}
Berger、Stern 和 Pizzorno 的 OSDI 2023 工作指出，仅给每行分配一个 CPU 百分比不足以指导优化。一行 Python 可能在解释器中执行，也可能等待本地扩展计算，还可能不断分配和复制大型对象。把本地扩展成本归到 Python 解释器，会诱导读者优化错误层次；把驻留集大小（RSS）当作程序已分配内存，又会混入触页和物理驻留行为。\cite{paper-scalene}

论文第 2 节利用 CPython 的信号交付特征做 CPU 归因。信号处理通常要等主线程回到可处理 Python 信号的位置；设按 CPU 时间计的目标采样间隔为 q，两次信号处理之间实际流逝的进程虚拟时间为 T。Scalene 用 q 估计 Python 部分，将延后的 \texttt{T\ensuremath{-}q} 归于本地代码（native）部分。图 3 展示这种区别。它依赖解释器行为，不能拿同一公式直接处理 Go goroutine。

多线程更复杂。第 2.2 节为阻塞操作设置可间断的等待与线程状态，检查线程当前栈及字节码；停留在 CALL 可作为正在本地调用中的线索。这是结合运行时语义的推断，不是精确记录每条 CPU 指令。异步机制、解释器版本或扩展对全局解释器锁（GIL）的使用方式变化，都要求重新验证前提。

\subsubsection{算法一：净变化阈值采样}
传统按分配字节数触发的采样擅长估计分配活动。如果一个程序反复分配又立即释放 1 MiB，它可能需要处理大量样本，即使存活规模几乎不变。Scalene 第 3.2 节转而维护自上次采样以来分配 A 和释放 F 的字节计数，当 \texttt{abs(A\ensuremath{-}F) \ensuremath{\geq} T} 时输出记录并重置计数。论文实现中的 T 为略大于 10 MB 的素数，以降低固定步长与阈值产生周期性耦合的风险。

用 T=4 MiB 做教学例子：依次分配 2、释放 2、分配 2、释放 2 MiB，分配加释放活动已累计 8 MiB，但净变化始终没有达到阈值，不触发净变化样本；随后连续分配 3 MiB 和 2 MiB，净增长 5 MiB，触发一次记录。输出的是显著存量变化，不是对所有短命分配逐一计数。

若目标是画内存规模趋势，重复的分配释放可能对趋势贡献很小，却让按活动采样产生大量记录。反过来，若目标是找导致 GC 压力的高频临时对象，这些被抵消的活动恰恰重要。因此阈值法减少样本不等于保留了所有性能信息。迁移到 Go 时应并列考察 \texttt{alloc\_\allowbreak{}space}、\texttt{inuse\_\allowbreak{}space} 和 GC 时间，不能用一张净增长曲线替代全部堆剖析。

\subsubsection{算法二：钩子重入、归因与泄漏候选}
第 3.1 节在系统分配器与 Python 分配器两层拦截。上层分配器可能继续调用下层，如果两层都计数，分配量会重复。Scalene 用线程局部的“已处在分配器中”标记，使嵌套钩子只转发而不重复记录，也避免 profiler 自身分配触发无限递归。这是多层观测必须回答的所有权问题。

样本触发时，C++ 扩展读取 Python 栈并找到受分析的源码行；后台线程读取采样文件更新统计。短时启用 \texttt{PyEval\_\allowbreak{}SetTrace}，直到执行离开当前行，用来处理行级平均内存归属。该设计将高频轻量计数与低频昂贵归因分开；与第 15 章预计算展开规则的思想相通，但优化的事件和运行时不同。

第 3.4 节在达到新内存高点时跟踪一个对象是否最终释放，将持续不释放的位置列为候选，再以概率阈值及整体增长趋势过滤报告。这里必须保留“候选”措辞：业务缓存也可以长期不释放，泄漏涉及程序意图，不能只由统计值定义。

原文还提供一个值得研究生练习的公式一致性问题。其定义说 mallocs 是跟踪分配总次数，frees 是其中回收次数，却将泄漏概率写为 \texttt{1\ensuremath{-}(frees+1)/\allowbreak{}(mallocs\ensuremath{-}frees+2)}。按这一定义代入 mallocs=frees=3 得 \ensuremath{-}1，不在概率范围内；这已核对 PDF 印刷页 56，而非 OCR 推测。

如果采用 Beta(1,1) 先验，并将是否回收建模为独立同分布的伯努利观测，下一对象不回收的后验预测概率为 \texttt{1\ensuremath{-}(frees+1)/\allowbreak{}(mallocs+2)}。后一个式子只是通用统计推导，不能擅自替换作者实现；还需核查对应发布源码的变量含义，才能判断是排版、定义还是实现问题。

\subsubsection{原始实验：哪些结果可以准确引用}
第 6.1 节使用 2022-12-08 发布的原型、CPython 3.10.9、Linux 5.13.0-35、8 核 AMD Ryzen 7、32 GB 内存和 NVIDIA RTX 2070。选取 pyperformance 中耗时最长的十个基准，重复到各自超过 10 s；每个基准运行十次，开销使用四分位间均值，即根据中间 50\% 的观测值计算均值。它是特定软硬件和工作负载下的测量，不是今天所有 Python 或 Go 服务的成本上界。

\begingroup\small\setlength{\tabcolsep}{4pt}
\begin{longtable}{>{\raggedright\arraybackslash}p{0.298\linewidth}>{\raggedright\arraybackslash}p{0.298\linewidth}>{\raggedright\arraybackslash}p{0.298\linewidth}}
\toprule
\textbf{原图表} & \textbf{作者报告的数据} & \textbf{可支持的解释} \\
\midrule
\endfirsthead
\toprule
\textbf{原图表} & \textbf{作者报告的数据} & \textbf{可支持的解释} \\
\midrule
\endhead
图 6，512 MB 数组、改变触页比例 & Scalene 与 Fil 的误差在真实分配大小的 1\% 内，Memray 在 6\% 内 & 该受控输入中，直接分配观测比 RSS 代理更贴近已分配大小 \\
\addlinespace[3pt]
表 2，pprint & 按活动 7976 个样本，阈值 23 个，约 347 倍差异 & 大量抵消活动使净变化采样明显减量 \\
\addlinespace[3pt]
表 2，sympy 与整体 & 6757 对 10，约 676 倍；十个基准中，两种方法样本数之比的中位数为 18 & 减少的是样本数量，不是程序得到同倍数加速 \\
\addlinespace[3pt]
表 3 & CPU 模式中位相对时间 1.02\ensuremath{\times}，完整模式 1.32\ensuremath{\times} & 完整模式中位开销约 32\% \\
\addlinespace[3pt]
表 3，pprint 完整模式 & 4.03\ensuremath{\times} & 中位开销不能代表每一种负载 \\
\addlinespace[3pt]
\bottomrule
\end{longtable}
\endgroup
论文叙述处出现完整模式约 30\% 的概括，本章引用表 3 可核对的 1.32\ensuremath{\times}，不把近似文字与表中数值混合。图 5 比较两种语义等价的写法：在循环中调用另一个函数，以及把相同逻辑直接写入循环；据此检查归因精度。这里的“内联”指源码展开，不是 Go 编译器的自动内联。不能把任一图上的最差工具表现泛化为所有采样器。\cite{paper-scalene}

\subsubsection{有效性威胁与 Go 迁移}
首先，净变化采样可能忽略大量迅速分配又释放的对象，这是估计目标选择的后果；其次，比较工具与版本是历史截面，不能用于今天的产品排名；第三，源码行、native 调用和 GPU 指标属于不同时间边界，联合呈现仍可能产生关联而非因果。UI 还进行曲线简化和约 100 点下采样，显示轨迹本身有信息损失。

可迁移的研究问题是：在固定 CPU/内存预算下，能否同时用低频存活堆快照、累计分配量差分及复制热点，区分 Go 服务的“高分配但平稳”和“低分配但持续保留”？基线应包含现有 Go pprof，评估目标分别设为 GC 成本定位、存活增长定位和优化后的实际收益。不能用 Python 的信号规则替代 Go runtime 事件，也不能因净变化样本更少就宣布检测能力更强。

\section{硬件计量：时间边界错误怎样污染优化决策}
\begin{figure}[H]
\centering
\begin{tikzpicture}[x=1cm,y=1cm]
\node[anchor=east,font=\small] at (1,0.0) {泳道1};
\draw[linegray] (1.2,-0.35) -- (14.799999999999999,-0.35);
\node[anchor=east,font=\small] at (1,-1.7) {泳道2};
\draw[linegray] (1.2,-2.05) -- (14.799999999999999,-2.05);
\node[anchor=east,font=\small] at (1,-3.4) {泳道3};
\draw[linegray] (1.2,-3.75) -- (14.799999999999999,-3.75);
\path[fill=forest!60] (1.200,-0.250) rectangle (5.733,0.250);
\node[align=center,font=\small,text width=4.43cm] at (3.467,0.750) {工具启动};
\path[fill=forest!60] (5.733,-0.250) rectangle (14.800,0.250);
\node[align=center,font=\small,text width=8.97cm] at (10.267,0.750) {外部计数窗口};
\path[fill=leaf!60] (2.711,-1.950) rectangle (7.244,-1.450);
\node[align=center,font=\small,text width=4.43cm] at (4.978,-0.950) {短工作负载};
\path[fill=forest!60] (2.711,-3.650) rectangle (7.244,-3.150);
\node[align=center,font=\small,text width=4.43cm] at (4.978,-2.650) {同步读数区间};
\draw[-{Stealth[length=2mm]},forest] (1.2,-4.55) -- (14.999999999999998,-4.55);
\draw[forest] (1.20,-4.55) -- (1.20,-4.6499999999999995);
\node[below,font=\scriptsize] at (1.20,-4.6499999999999995) {0};
\draw[forest] (3.92,-4.55) -- (3.92,-4.6499999999999995);
\node[below,font=\scriptsize] at (3.92,-4.6499999999999995) {18};
\draw[forest] (6.64,-4.55) -- (6.64,-4.6499999999999995);
\node[below,font=\scriptsize] at (6.64,-4.6499999999999995) {36};
\draw[forest] (9.36,-4.55) -- (9.36,-4.6499999999999995);
\node[below,font=\scriptsize] at (9.36,-4.6499999999999995) {54};
\draw[forest] (12.08,-4.55) -- (12.08,-4.6499999999999995);
\node[below,font=\scriptsize] at (12.08,-4.6499999999999995) {72};
\draw[forest] (14.80,-4.55) -- (14.80,-4.6499999999999995);
\node[below,font=\scriptsize] at (14.80,-4.6499999999999995) {90};
\node[anchor=east,font=\scriptsize] at (14.799999999999999,-5.3) {时间 / 教学时间单位};
\end{tikzpicture}
\caption{短任务测量最怕计数窗口和执行窗口错位}\label{fig:16-pmu-window}
\begin{minipage}{0.96\linewidth}\footnotesize 重建测量窗口错位机制，非原论文实测时间线；不能据此宣称 perf 整体无效。

\textbf{读图顺序：}1. 先看真实工作区间\quad 2. 外部工具准备有延迟\quad 3. 窗口还可能包含无关活动\quad 4. 同步到被测区域。
\textbf{机制依据：}\cite{paper-hardware}。
\end{minipage}
\end{figure}
\subsubsection{问题不是有没有计数器，而是计了哪一段}
Kakaraparthy 与 Patel 的 2024 年文章研究短区间 PMU 测量。目标是在每批哈希表查询前后读取 CPU 周期数（cycles）、已提交指令数（retired instructions）和缓存未命中数（cache misses）。如果启动工具、目标区间和停止工具没有对齐，读数可以看似合理却覆盖错误窗口；测量本身还可能改变被测工作。\cite{paper-hardware}

第 2 节的 Listings 2—4 对照三条路径。第一种为每批工作 fork/exec 一个 perf stat 并附加到父进程；第二种在进程内部调用 PAPI，借助 perf\_event 接口配置与读取；第三种 PMU-metrics 直接配置 PMU 寄存器并用 rdpmc 读取。第三种低层路径依赖硬件、权限和计数器所有权，不能直接当作多租户生产环境的默认建议。

\subsubsection{误差模型：漏掉工作和改变工作可以同时发生}
设目标函数真实开始于 a、结束于 b，计数器从 a+\ensuremath{\delta} 才开始，读数近似只包含 \texttt{[a+\ensuremath{\delta},b]} 的事件，漏掉前面的工作。即使采样或计数器本身完全准确，观测窗口仍有偏差。若启动附加子进程还使目标因写时复制产生缺页和地址转换缓存（TLB）扰动，那么被测期间的实际执行也已从原基线改变。

图 3 正是在区分这两件事：原始单批约 330 ms，fork 后 processRequests 本体约 429 ms，完整批次约 452 ms。工具开始计数较晚，导致记录偏小；同时程序被扰动变慢。看到“计数少、时间长”不能立即推断每周期指令数（IPC）提高或缓存行为改善，二者可能来自边界和自扰动的组合。

一个教学例子是：无工具任务耗时 1 ms，启动后延迟 0.4 ms 才开始计数，而工具使工作延长到 1.2 ms。计数覆盖的只是后 0.8 ms，但用户测得总时间却多了 20\%。用 0.8 ms 对应 cycles 去解释原始 1 ms 的行为，已经把两个不同执行过程拼在一起。正确实验需要同步开始、尽量减少被测路径外扰，并单独测量空窗口成本。

\subsubsection{原始图 1—3：准确度数字的适用范围}
第 3 节在独占服务器上运行：Intel Xeon Silver 4114、10 核，通过 taskset 固定核心，CPU 调频策略设为 performance，对应所用硬件的 3 GHz 频率；每种实验五次取中位数。总计 1 亿次哈希查询，每批 1000 万次，未剖析时每批约 330 ms。perf 情况使用两个核心以容纳其进程。这个细节也意味着配置条件并非仅有工具名称不同。

图 1 用 \texttt{估计 cycles = 已测时间 t \ensuremath{\times} 频率 f} 作交叉参照。作者报告 perf、PAPI、PMU-metrics 的 cycles 准确度分别 23.9\%、99.3\%、99.7\%。图 2 中整体执行时间相对无剖析基线分别增加 40\%、3\%、2\%。这些数字针对论文所用的按批 fork/exec 方法，不能读成“perf 无法准确计量”。常驻工具、同步控制或直接 perf\_event 读数是不同实验。

\texttt{t\ensuremath{\times}f} 也不是无条件真值：它依赖频率估计和执行边界，若发生频率变化、抢占或虚拟化计数语义变化，需另外处理。论文还说明不同库预设事件可能映射到不同 native events，例如 retired-instruction 的事件变体；名称相似不能证明测的是同一硬件事件。比较 cache misses 时，更不能把 cycles 的校准自动推广为各类 cache 的独立真值验证。

\subsubsection{原始图 4—5：测量越来越细会改变最优选择}
第 4.1 节保持总查询数 1 亿不变，将每批 1000 万逐步降到 100，使单批从约 0.33 s 降至约 3.3 \ensuremath{\mu}s。每个工具以自身粗粒度结果为基线，五次取中位数。图 4 在最细粒度报告 PAPI 指令数偏离约 +12\%、运行时间 +268\%；PMU-metrics 则约 +2\%、+8\%。这里的 +268\% 不是图 2 无工具基线的同一个比较，应保留“相对各自粗粒度”的分母。

直观模型是 \texttt{总成本 \ensuremath{\approx} 有效工作 + 测量次数 \ensuremath{\times} 每次固定开销}。批次越小，测量次数越多，即使每次仅有少量系统调用，也可能超过实际任务。直接读取硬件计数器降低固定成本，但必须解决计数器复用、任务切换和访问权限，不能只看指令执行时延。

第 4.2 节把计数器放进 BLARE 的正则策略选择：多臂老虎机依据各策略的 cycles 作为奖励信号，选择后续策略。实验使用 100 MB 数据集，十次取中位数。图 5 相对“最佳正则策略”的额外运行时间为 perf\_event 20\%、PAPI 25\%、PMU-metrics 10\%。这个基线不是“同一自适应程序关闭计量”，因此同时包含学习选择等成本，不能把全部差异都称为计数器读取开销。

这揭示了更深的问题：测量不仅用于事后报告，还会反馈到算法。相同的固定开销加在各策略上，并不会颠倒它们的单次耗时排序；但若策略的尝试次数、测量频率或计量路径不同，观测成本就可能改变奖励信号。反馈系统应分别测量目标工作与观测成本，检查策略选择是否受后者驱动。

\subsubsection{迁移到 Go：goroutine 与核心不是同一个归属单位}
Go 函数可能在不同线程执行，线程又可能迁核；进程级 PMU 统计还可能包含 GC、调度和其他 goroutine。单纯在函数两侧读“当前核心计数器”，两次读数不保证属于同一核心和同一工作实体。固定 OS 线程或核心只是受控实验工具，还要明确系统调用、异步工作和 cgo 是否计入。

可设计三组 Go 实验：较长批次用稳定 PMU 观测校准；逐步缩短批次测量固定开销占比；再让服务开启正常调度与 GC，检查计量从受控环境迁移后的误差。报告原始事件配置、核心亲和性、频率、空窗口成本和是否多路复用。对 PGO 等反馈优化还应增加最终吞吐和延迟验证，避免“profile 变好”只是观测对象变了。

\section{ePPing：协议语义、观测覆盖与低开销聚合}
\begin{figure}[H]
\centering
\begin{tikzpicture}[x=1cm,y=1cm]
\node[draw=linegray!75!black,fill=linegray!13,rounded corners=3pt,text width=3.33cm,minimum height=1.12cm,align=center,font=\small] (p0-out) at (1.88,0.97) {观察发出包\\{\scriptsize 0 us: TSval=100}};
\node[draw=linegray!75!black,fill=linegray!13,rounded corners=3pt,text width=3.33cm,minimum height=1.12cm,align=center,font=\small] (p0-key) at (5.62,0.00) {建立待匹配项\\{\scriptsize 流＋方向＋100; t0=0}};
\node[draw=linegray!75!black,fill=linegray!13,rounded corners=3pt,text width=3.33cm,minimum height=1.12cm,align=center,font=\small] (p0-echo) at (1.88,-0.97) {观察回程包\\{\scriptsize 1500 us: TSecr=100}};
\node[draw=linegray!75!black,fill=linegray!13,rounded corners=3pt,text width=3.33cm,minimum height=1.12cm,align=center,font=\small] (p0-match) at (9.38,0.97) {按键匹配\\{\scriptsize 取出发送时刻}};
\node[draw=linegray!75!black,fill=linegray!13,rounded corners=3pt,text width=3.33cm,minimum height=1.12cm,align=center,font=\small] (p0-delta) at (13.12,0.00) {时间差估计\\{\scriptsize 1500 \ensuremath{-} 0 = 1500 us}};
\node[draw=leaf!75!black,fill=leaf!13,rounded corners=3pt,text width=3.33cm,minimum height=1.12cm,align=center,font=\small] (p0-expire) at (9.38,-0.97) {超时/\allowbreak{}容量边界};
\draw[-{Stealth[length=2mm]},forest!70!black] (p0-out.east) -- (p0-key.west);
\draw[-{Stealth[length=2mm]},forest!70!black] (p0-key.east) -- (p0-match.west);
\draw[-{Stealth[length=2mm]},forest!70!black] (p0-echo.north) -- (1.88,2.12) -- (9.38,2.12) -- (p0-match.north);
\draw[-{Stealth[length=2mm]},forest!70!black] (p0-match.east) -- (p0-delta.west);
\draw[-{Stealth[length=2mm]},forest!70!black] (p0-key.east) -- (p0-expire.west);
\end{tikzpicture}
\caption{ePPing 用往返可匹配字段估计网络延迟}\label{fig:16-epping}
\begin{minipage}{0.96\linewidth}\footnotesize 按论文首 TSval 配对机制重建，0/200/1500 us 均为教学数值。匹配键包含流与方向；RTT 边界由观察点及 TCP 回显规则决定。

\textbf{读图顺序：}1. 保存首次出现的时间戳\quad 2. 重复值不覆盖首时间\quad 3. 反向 ACK 回显\quad 4. 计算部分往返时间\quad 5. 保留观测限制。
\textbf{连线语义：}建立待匹配项\ensuremath{\rightarrow}\allowbreak{}超时/容量边界：未匹配。
\textbf{机制依据：}\cite{paper-epping}。
\end{minipage}
\end{figure}
\subsubsection{问题与基线：看见更多报文才能正确配对}
Sundberg 等人的 PAM 2023 论文将被动 TCP RTT 观测放进 eBPF。传统 PPing 在用户态捕获报文并匹配，流量升高时抓包路径可能丢失大量报文。若算法需要“第一次出现的某个时间戳”，丢失第一次而看到第二次不仅减少样本，还会改变 RTT 数值。这是覆盖率与准确性耦合的典型案例。\cite{paper-epping}

ePPing 在 XDP（较早的收包处理点）和 tc（流量控制钩子）等内核路径解析必要头部，以流标识及协议时间戳建立状态，再在用户态加载程序、清理过期映射和输出结果。高频配对在内核进行，低频控制与展示在用户态。它减少报文复制和每包用户态处理，但仍占用 CPU 与 map 内存；“运行在内核”本身不是低开销证明。

\subsubsection{核心算法：保存首个 TSval，匹配回显 TSecr}
TCP 时间戳选项携带发送端 TSval；接收端在响应里用 TSecr 回显所观察的值。简化的 RTT 配对过程如下：对流的某个方向，首次见到时间戳 v 时存下本机观察时间 t0；反向 ACK 首次回显 v 时取当前时间 t1，得到 \texttt{RTT\_\allowbreak{}est=t1\ensuremath{-}t0}，并按规则处理状态。匹配键需包含流和方向，不能只拿时间戳整数全局查表。

教学例子：观察点在 0 \ensuremath{\mu}s 看到 \texttt{TSval=100}，在 200 \ensuremath{\mu}s 再看到同值，1500 \ensuremath{\mu}s 看到反向 \texttt{TSecr=100}。保留同一 TSval 的首次观测时间，得到 1500 \ensuremath{\mu}s；如果第二包覆盖 t0，会误成 1300 \ensuremath{\mu}s。时间戳粒度使同一值对应多个报文，简单“最后写入覆盖”破坏了被动估计。

这个 RTT 是相对观察点的部分往返时间，包含该段路径、远端 TCP 协议处理及确认等待。它不是 Go handler 从接收请求到完成响应的耗时。放在中间盒观察，既不能恢复请求到达服务器前的全部路径，也不能直接拆出应用排队。它适合给第 13 章的跨层证据链增加网络信号，而不是替代 trace。

\subsubsection{附录算法为何关系到测量定义}
论文附录 A、图 8 与算法 1 解释 TCP 时间戳回显规则：当 \texttt{SEG.\allowbreak{}TSval \ensuremath{\geq} TS.\allowbreak{}Recent} 且 \texttt{SEG.\allowbreak{}SEQ \ensuremath{\leq} Last.\allowbreak{}ACK.\allowbreak{}sent} 时，更新待回显的时间戳缓存 TS.Recent。这里 SEG 表示当前报文，SEQ 是其起始序号，Last.ACK.sent 是最近已发送 ACK 中的确认序号。结果是时间戳方法常包含 delayed ACK，而常见基于序号的匹配可能选择更晚的报文起点。两者出现小差异，不一定是谁计时错误，首先要比较测量定义。

重传也不能一概消除歧义：如果重传有新的 TSval，匹配更容易区分；若快速重传仍共享同一个时间戳值，估计仍可能偏大，通常受时间戳周期约束。附录用受控实验展示异常样本，目的不是承诺所有网络栈都满足固定误差上界，而是暴露协议粒度和 ACK 策略的后果。

\subsubsection{原始实验分三步揭示不同限制}
第 3 节的端主机为 i7-7700、16 GB、Linux 5.16，中间盒为 Xeon E5-1650、32 GB、Linux 5.19，使用 100 Gbps 链路；在中间盒关闭 GRO/GSO/TSO 等接收聚合与发送分段卸载功能，以明确观测报文粒度。实验设置会影响可迁移性，尤其不能直接外推到开启合并卸载的云网络。

第一步验证 RTT。图 3 用 netem 将延迟从 0 到 100 ms、每 10 ms 一档、每档 10 s，流量为单条限速 100 Mbps 流。图 4 比较 ePPing、PPing 与 tshark；PPing/ePPing 获得的 RTT 数比 tshark 约少 13\%，因此作者对匹配 ACK 的共同子集比较。ePPing 相对 tshark 常高约 1—5 \ensuremath{\mu}s，与 delayed ACK 定义有关；共同子集的一致性不等于所有报文的完整覆盖。

第二步测多核转发。图 5 在十条并发流、六核条件下，无监测平均 CPU 利用率 47\%，ePPing 57\%，PPing 77\%。增加的是 10 与 30 个百分点，不能写成 10\% 与 30\% 的相对增幅。PPing 只处理略高于 6\% 的报文，ePPing 看到约 16 倍数量；吞吐相近还受端主机瓶颈限制，不能据此宣称观测没有成本。

第三步把中间盒限制为单核，使观测开销直接体现在转发能力。表 1 的基线总包率约 1.90—1.95 Mpps；图 6 中，流数从 1 增至 1000，ePPing 从 1.53 降到 1.13 Mpps。PPing 转发率看似从 0.97 升至 1.18 Mpps，但它实际处理的包率从约 170 kpps 降到 60 kpps。后者“更快”是漏看更多，不能算成更高效地完成同一观测任务。

开销实验每组运行十次、每次 120 s、丢弃最初 20 s。包率比字节吞吐更贴近每包算法成本：流数改变 ACK 数量，字节数相近时需要处理的报文数仍可不同。论文把瓶颈位置和工作量分母显式改动，值得在 Go 服务实验中照做。

\subsubsection{两种降本算法：控制采样对象和压缩输出}
为了降低开销，作者保留“看见每包”的能力，但限制同一流写入新待匹配时间戳的频率：保存一次后，至少等待 t 才为该流保存下一次。与先随机丢包不同，它更容易保留首边缘配对语义，也避免大流独占全部样本预算。不过限流后的观测仍会遗漏短时变化，不能说稀疏流天然获得完整 RTT。

图 7a 使用 t=0、10、100、1000 ms。在论文的时间戳配置下，对应每流每秒最多约 1000、100、10、1 个 RTT。t=0 的约 1000 来自时间戳更新粒度，并非数学上 1/0。1000 流时，t=1000 ms 的转发能力为 1.54 Mpps，未启用采样限流时约 1.14 Mpps；图 6 邻近情形写成 1.13，不应人为合并成一个无误差数字。

另一方案在 BPF 内维护 RTT 直方图及最小最大值，用户态每秒拉一次，减少逐样本上报。图 7b 的 1000 流情况下，未启用采样限流时，转发包率约从 1.14 提升到 1.45 Mpps。代价是丢失单个 RTT 与精确时刻的关联，桶边界还影响分位数分辨率。它与第 15 章只发送栈 ID 的压缩不同：直方图本身是有损统计，不能日后重建原始事件时间线。

\subsubsection{适用边界与 Go 研究方案}
论文主要测试实验室中的持续大流量 iperf 传输，未覆盖运营商网络中的各种混合流量，流数最高 1000；协议依赖时间戳，不能直接推广到全部 QUIC、DNS。作者讨论每包额外处理很短仍可能通过队列累积放大端到端时延，这说明局部纳秒成本与业务尾延迟之间不是简单加法关系。

迁移到 Go 服务时，可联合 handler trace、调度延迟、CPU profile 与 RTT 分布，分别注入网络延迟、CPU 饱和和锁竞争，制造表现相似但原因不同的慢请求。按连接和时间窗对齐网络与应用证据；若连接复用了多个请求，还需避免把某次 RTT 直接归给其中一个请求。验证系统能否把“网络变慢”和“本地排队”分开。若只在 CPU 低时提高网络采样，会改变观测选择，应把采样策略一并记录，防止将策略变化误报为业务改善。

\section{Go 工程前沿：PGO 与异常回溯形成反馈环}
\begin{figure}[H]
\centering
\begin{tikzpicture}[x=1cm,y=1cm]
\path[fill=bluegray!25] (0.00,0.00) rectangle (13.97,0.78);
\node[align=center,text width=13.92cm,font=\small] at (7.00,0.40) {CPU 权重};
\path[fill=leaf!38] (0.00,1.05) rectangle (1.38,1.83);
\node[font=\scriptsize] at (0.70,1.45) {2};
\path[fill=leaf!38] (1.40,1.05) rectangle (4.17,1.83);
\node[align=center,text width=2.72cm,font=\small] at (2.80,1.45) {路径 B};
\path[fill=leaf!38] (4.20,1.05) rectangle (13.97,1.83);
\node[align=center,text width=9.72cm,font=\small] at (9.10,1.45) {路径 C};
\node[anchor=north west,align=left,text width=14cm,font=\scriptsize] at (0,-.45) {1=CPU 权重: flat 0 /\allowbreak{} cum 100；2=路径 A: flat 10 /\allowbreak{} cum 10；3=路径 B: flat 20 /\allowbreak{} cum 20；4=路径 C: flat 70 /\allowbreak{} cum 70};
\end{tikzpicture}
\caption{训练热点和评价热点可能不一致}\label{fig:16-pgo-drift}
\begin{minipage}{0.96\linewidth}\footnotesize 人为构造的 CPU 权重分布，不是原论文曲线；用于设计分布漂移实验。

\textbf{读图顺序：}1. 训练分布\quad 2. 分布发生漂移\quad 3. 观察收益是否迁移\quad 4. 保留对照与回退。
\textbf{机制依据：}\cite{go-pgo}\cite{front-go121}。
\end{minipage}
\end{figure}
\subsubsection{PGO 把样本变成编译决策，偏差也会随之放大}
Go 1.21 将 PGO 作为正式功能，构建可通过 \texttt{-\allowbreak{}pgo=auto} 使用主包目录的 \texttt{default.\allowbreak{}pgo}。CPU profile 影响内联与带保护的去虚化等优化：热点接口调用可在运行时检查实际类型，命中时走更易优化的直接调用路径，未命中仍保留通用路径。因此性能收益取决于代表性、热点稳定性和代码变化，而非 profile 文件是否能解析。\cite{front-go121}

这与硬件论文的 BLARE 存在共同结构：观测决定下一轮执行策略。若训练 profile 只含 \texttt{/\allowbreak{}pay}，却把结果部署到 \texttt{/\allowbreak{}admin} 占主导的负载，优化目标已经漂移。假设教学流量中 encode 50 ms、hash 10 ms，后来流量反过来，旧热点排序不能代表新收益。不能仅用同一训练负载自证 PGO 有效。

研究级评估应把负载分为采集、调参和未见测试集合，同时报告总吞吐、尾延迟、二进制体积与构建变化。修改源代码后还需检查 profile 匹配覆盖，而不仅是编译器没有报错。比较均匀采样与自适应采样时，最终优化收益比火焰图相似度更接近任务目标；但必须保留 CPU 总量与采样偏差，防止优化只惠及被过度观测的路径。

\subsubsection{Trace 结构重写解决了持续事件数据的工程障碍}
2024-03-14 的 Go 官方文章介绍 Go 1.22 执行追踪重写，通过按代组织事件、改进时间信息等方式，降低采集开销并改善数据处理。分段、自包含的数据有利于增量处理和保留窗口；但文章发表时 \texttt{go tool trace} 仍需要整体加载，不能因格式支持分段便宣称当时所有工具都能无限流式分析。\cite{front-trace2024}

采样 profile 擅长估计一段时间中“哪里消耗资源”，trace 保存调度、阻塞和事件顺序。一次短时出现的尾延迟尖峰可能在总 CPU 样本里权重很小，却需要精确的先后关系解释。二者联合不是多画一张图，而是要把相同时间窗、实例和请求的证据对齐，并确认时钟与 ID 的边界。

\subsubsection{Flight Recorder 把问题改成何时保留证据}
2025-09-26 的官方文章和 Go 1.25.0 源码提供 Flight Recorder：持续记录有限窗口，异常触发后写出快照，从而保留“知道故障之前”的事件。设慢请求在时刻 t 被检测到。如果此时才开启 trace，就可能错过此前的排队过程；记录器可保留 t 之前的一段历史。\cite{front-flight2025} \cite{front-flight-source}

Go 1.25.0 的接口有明确约束：同一进程只能启用一个 FlightRecorder，但可与 \texttt{trace.\allowbreak{}Start} 并存；同一时刻只允许一个 \texttt{WriteTo}，重叠调用返回错误。窗口大小目标 \texttt{MaxBytes} 优先于历史保留时长目标 \texttt{MinAge}，但仍只是提示，不保证快照文件或进程内存始终低于该数值。高事件率下为了满足空间目标，保留的有效历史长度可能缩短。

因此可研究的变量不只是缓冲区大小，还包括触发延迟、事件生成速率、故障前因跨度和重复触发策略。假设教学窗口实际保留 5 s，但报警经过 8 s 聚合才触发，最有价值的前因仍然可能被覆盖。正确指标是“异常所需因果片段的保留率”，再结合内存与 CPU 成本，而不是配置里写了几秒。

联合设计可以先以低成本指标识别候选异常，再保存 Flight Recorder 快照，必要时补充 CPU 或 heap profile。快照可能同时涵盖多个请求和异步后台工作；请求关联需结合显式的 trace task/region/log 注解，不能假定 pprof 标签会自动成为 trace 注解。应按第 08 章的方法重建时间线，辨认与异常有关的活动。研究方案应加入假阳性触发、重复异常和缓存高压，才能评估持续观测的实际价值。

\section{开放 Profiles：互操作问题本质上是语义守恒}
\begin{figure}[H]
\centering
\begin{tikzpicture}[x=1cm,y=1cm]
\node[draw=linegray!75!black,fill=linegray!13,rounded corners=3pt,text width=3.33cm,minimum height=1.12cm,align=center,font=\small] (p0-native) at (1.88,0.97) {Go 原生采样};
\node[draw=linegray!75!black,fill=linegray!13,rounded corners=3pt,text width=3.33cm,minimum height=1.12cm,align=center,font=\small] (p0-ebpf) at (1.88,-0.97) {eBPF 采样};
\node[draw=linegray!75!black,fill=linegray!13,rounded corners=3pt,text width=3.33cm,minimum height=1.12cm,align=center,font=\small] (p0-format) at (5.62,0.00) {交换格式};
\node[draw=linegray!75!black,fill=linegray!13,rounded corners=3pt,text width=3.33cm,minimum height=1.12cm,align=center,font=\small] (p0-semantics) at (9.38,0.97) {单位/\allowbreak{}周期/\allowbreak{}时窗};
\node[draw=linegray!75!black,fill=linegray!13,rounded corners=3pt,text width=3.33cm,minimum height=1.12cm,align=center,font=\small] (p0-quality) at (9.38,-0.97) {展开/\allowbreak{}丢样/\allowbreak{}身份};
\node[draw=leaf!75!black,fill=leaf!13,rounded corners=3pt,text width=3.33cm,minimum height=1.12cm,align=center,font=\small] (p0-compare) at (13.12,0.00) {可比较的结论};
\draw[-{Stealth[length=2mm]},forest!70!black] (p0-native.east) -- (p0-format.west);
\draw[-{Stealth[length=2mm]},forest!70!black] (p0-ebpf.east) -- (p0-format.west);
\draw[-{Stealth[length=2mm]},forest!70!black] (p0-format.east) -- (p0-semantics.west);
\draw[-{Stealth[length=2mm]},forest!70!black] (p0-format.east) -- (p0-quality.west);
\draw[-{Stealth[length=2mm]},forest!70!black] (p0-semantics.east) -- (p0-compare.west);
\draw[-{Stealth[length=2mm]},forest!70!black] (p0-quality.east) -- (p0-compare.west);
\end{tikzpicture}
\caption{同样的格式不等于同样的测量语义}\label{fig:16-interoperability}
\begin{minipage}{0.96\linewidth}\footnotesize 教学示意：用于解释机制，不是运行时的实测数据或完整实现。

\textbf{读图顺序：}1. 先能交换数据\quad 2. 再检查目标量\quad 3. 再检查观测质量\quad 4. 最后才做可比性实验。
\textbf{机制依据：}\cite{front-otel-spec}\cite{proj-otel-ebpf}。
\end{minipage}
\end{figure}
\subsubsection{从格式兼容到可比较的观测}
OpenTelemetry 在 2024 年介绍 Profiles 方向，意图把剖析纳入日志、指标和 trace 之外的共同观测体系。2026-10-02 阅读的规范快照仍标为 Alpha。原始公告说明发展方向，当前规范定义接口要求，二者都不是“所有后端已兼容”的实验结果。\cite{front-otel2024} \cite{front-otel-spec}

模型包含资源属性，以及采集库名称、版本等作用域信息（InstrumentationScope）、去重字典、样本与 trace/span 的关联，以及保留原始载荷的机制。它们各解决一个具体问题：资源属性用于区分服务及其实例；字典减少重复符号；关联帮助把 CPU 权重与请求时间线连接；原始载荷用于保留尚未映射的格式细节。可扩展字段并不自动保证分析端理解它们。

第 15 章的 count 与 period 示例就是互操作测试的最小反例。一个来源上报 count=1、period=10 ms，另一个来源直接上报 cpu=10 ms；转换器可以让二者统一为时间，但必须显式改变数值和单位。若只把字段名映射到 value，数据传输成功而估计量错误。累计分配量与区间增量的时间语义（temporality）、标签缺失与空标签、内联行序列和部分栈标记，也都需要明确的转换规则。

\subsubsection{关联 ID 不提供因果证明}
假设后台 goroutine 继承请求标签，却在响应结束后继续执行，profile 与 span 关联可能仍然成立，但这段 CPU 是否计入“请求延迟”要由服务语义定义。多个请求共享批处理任务时，若把全部 CPU 重复记到每个 span，总量就不守恒；若只记到其中一个 span，总量虽守恒，分摊仍可能失真。关联 ID 帮助定位，如何分摊资源成本还需要单独定义。

可用守恒检查发现转换错误：同一合成输入经过 pprof\ensuremath{\rightarrow}\allowbreak{}通用格式\ensuremath{\rightarrow}\allowbreak{}后端查询后，换算到相同单位后，合法样本的总权重应保持；按互斥且覆盖全部样本的标签组拆分后，求和应等于未拆分总量；当后端只支持部分帧或标签时，应显式标注信息丢失。非法单位与未知版本应按明确规则拒绝或保留。合法负差分必须保留符号；后端不支持时应明确报错，不能悄悄改为零。

\subsubsection{eBPF 与开放格式各自负责哪一层}
eBPF profiler 可减少语言 SDK 接入，却仍面对 Go 栈变化、cgo、内联符号和标签布局；开放格式可携带完整性信息，却无法修复采集时没有观察到的栈。把它们结合的研究价值在于让采集质量成为跨后端可比较的数据，而不是统一输出一张图就宣布可观测性统一。

一个有意义的原型是“语义保真转换测试集”：覆盖不同采样率、累计重置、乱序、ASLR、Build ID 变更、标签基数、深栈与部分符号。先确定每个转换应保持的量，再测转换吞吐和体积。按格式能解析多少样本只是最低层的成功标准；真正的终点是相同查询是否回答同一个问题。该测试集仍是研究计划，尚无跨后端验证结果。

\section{从论文到研究项目：统一评估目标、对照与失败条件}
\subsubsection{三篇论文共享的技术命题}
Scalene 改变要记录的内存事件，减少与目标无关的记录；硬件论文缩短每次观测路径，并校准边界；ePPing 把配对放在更接近数据的位置，再对输出限速或聚合。它们表明，除了降低采样率，还可以通过选择记录对象、缩短观测路径和聚合输出减少开销。每次改变机制，都必须重新确认估计对象有没有变化。

对 Go 工具，需要分别评价四个维度。它们可能相互影响，不能相互替代。准确性衡量数值或归因接近什么参考；覆盖率衡量哪些事件、栈、请求和时间段被观察；开销衡量 CPU、内存、吞吐与尾延迟；行动收益衡量据此实施的优化是否改善真实服务。这四个轴不能压成一个“准确率”，也不能靠某一平台内部统计完全自证。

\subsubsection{三项可以形成硕士研究课题的设计}
第一项是保留估计量的自适应采样。根据负载和异常程度调整 Go CPU、heap 与 trace 预算，但保留每个来源实际采样概率或周期。以固定预算均匀采样为基线，测试稳定热点、短突发和低频高成本路径；评价热点排序、总量误差和尾事件召回。若只提高告警期频率而不重加权，平均成本会偏向异常期。

第二项是运行时感知的跨层归因。将 Go 调度状态、profile 栈与 eBPF 网络 RTT 对齐，区分 CPU 饱和、锁等待、GC 与网络延迟。先用可控注入建立原因标签，再在混合原因中评估。基线不能只是随机分类器，应包含现有工具组合及熟练人工流程；报告未知类别与误判，而不仅是挑几张成功图。

第三项是观测驱动优化的闭环稳定性。对 PGO 或动态批处理，比较不同 profile 新鲜度、采样策略与训练负载的优化效果。把观测开销和策略收益分开测量，保留未见负载作为外部验证。若优化对训练负载有效却损害另一接口尾延迟，应算作边界发现，而不是删除失败样本。

\subsubsection{如何设置真值、消融与复现材料}
对合成计算，可用已知调用次数、受控持续时间和总分配字节作局部参考；对真实服务，不存在一个无干扰的完美 profiler，应使用多个独立证据及误差范围。硬件 \texttt{t\ensuremath{\times}f}、RSS、trace 时间戳和应用计数都有前提，引用参考值时要同时说明它能证明什么。

消融实验一次去掉一项机制。例如阈值采样去掉净变化而按活动触发；ePPing 去掉内核聚合而逐样本输出；将自适应采样决策替换为固定频率，并使总体观测预算尽量相近。只有这样才能分清收益来自算法、实现语言、批量程度还是默认配置。不同工具功能不等价时，至少提供共同功能子集比较，并单独展示额外信息的成本。

统计报告保留每次重复结果、运行顺序和异常退出，区分中位数、均值、置信区间与最差负载。不要把 ePPing 的百分点改成相对百分数，也不要把 Scalene 的样本减少倍数写成加速倍数。硬件研究的粗粒度自身基线与无工具基线必须分表表达，才能让读者重算。

\subsubsection{完成研究的最低交付条件}
交付应包括固定工具与源码版本、公开或可生成的输入、实验脚本、采样单位与窗口约定、原始测量、图表生成方式、失败案例和环境约束。若数据涉及真实请求，只保留必要且脱敏的标识，并说明哪些脱敏影响可复现性。图表应能追溯到数字，而不是仅保留截图。

本章提供的是原始论文审读和研究设计，没有重新运行三篇论文的实验，没有完成 PGO 新算法、eBPF 部署或互操作基准。作者结果与本书手算、推论保持分离，才能让读者继续验证、反驳和扩展。研究生教材的终点应是能提出有证据支撑、也可能被实验推翻的问题，而不只是记住更多工具名称。

\section{实验任务}
\subsection*{实验 16：复核一个论文结论，再设计 Go 对照实验}
\begin{enumerate}
\item 在三篇论文中选一篇，打开原始全文，逐项记录估计目标、采样机制、软件硬件、基线、重复次数、图表编号与单位。至少选一张支持结论和一张暴露局限的图表。

\item 从原表重算一个数值：Scalene pprint 的 7976/23；ePPing CPU 从 47\% 到 57\% 的百分点与相对增幅；或硬件图 3 的 452/330 与 429/330。不要将手算写成实验复现。

\item 构造三个最小反例：净变化为零但累计分配很大、计数窗口晚启动、相同 TSval 的首次观测时间被后包覆盖。分别说明目标不匹配、时间边界错误或状态更新错误如何影响结论。

\item 为 Go 迁移制定一个可证伪假设，确定现有工具基线、工作负载、控制变量、消融、误差与开销上限。加入至少一个预期失败的条件。

\item 若仅完成设计，交付协议和预期判定，不填写虚构结果。若执行实验，固定运行时及内核，保存每次原始数据，并把本人测量、论文结果和理论推导分开成表。

\item 最后写一页有效性审查：内部混杂、测量真值的前提、外部可迁移性、符号与标签覆盖、尚未完成的验证。论证应能由另一位读者依图表和数据重算。

\end{enumerate}
\section{自测与解析}
\noindent\textbf{1. Scalene 表 2 中两种方法的样本数之比，中位数为 18；表 3 完整模式的相对运行时间中位数为 1.32\ensuremath{\times}。应如何解释？}
\begin{enumerate}[label=\Alph*.]
\item 程序加速了 18 倍
\item 净变化采样减少记录；完整模式相对运行时间的中位数表明约 32\% 的额外开销
\item 所有内存活动都被无损压缩 18 倍
\item 任意 Go 服务开销至多 32\%
\end{enumerate}
\noindent{\color{forest}\textbf{答案 B。}}样本数量与运行时间是不同指标，净变化还改变了目标事件；该结果仅属于论文特定实验。

\noindent\textbf{2. 硬件论文按批 fork/exec perf 的读数偏小，最合适的结论是什么？}
\begin{enumerate}[label=\Alph*.]
\item 所有 perf 功能无效
\item PMU 永远不会产生测量误差
\item 该测量方式存在窗口错位和自扰动，应校准具体使用路径
\item 只要增加重复次数就消除系统性偏差
\end{enumerate}
\noindent{\color{forest}\textbf{答案 C。}}论文研究特定细粒度使用方式。重复可以减小随机波动，不能修复固定的启动延迟或错误边界。

\noindent\textbf{3. ePPing 图 5 的 CPU 从 47\% 到 57\%，应报告什么？}
\begin{enumerate}[label=\Alph*.]
\item 增加 10 个百分点；相对增幅约 21.3\%
\item 增加 10\% 相对开销
\item 增加 57 个百分点
\item 吞吐没下降所以开销为零
\end{enumerate}
\noindent{\color{forest}\textbf{答案 A。}}百分点差是 57\ensuremath{-}47=10；相对增幅为 10/47。端主机可能是吞吐瓶颈，不能由吞吐不变否认 CPU 成本。

\noindent\textbf{4. 采样器支持开放 Profile 格式，是否足以证明 Go 1.25 标签归因正确？}
\begin{enumerate}[label=\Alph*.]
\item 足够，格式会自动修复运行时布局
\item 不足，采集布局、转换语义与查询合同需要分别验证
\item 足够，只要 HTTP 成功
\item 不足，所以开放格式没有价值
\end{enumerate}
\noindent{\color{forest}\textbf{答案 B。}}互操作携带信息，无法补回采集端未正确读取的信息。须独立核查运行时版本和语义守恒。

\section{分析题与参考答案}
\noindent\textbf{题 1：}在采用净变化阈值采样的教学模型中，T=4 MiB。若 Go 程序反复分配并立即释放 1 MiB 一千次，净变化样本可能很少。这能说明程序无需优化吗？

\noindent\textbf{参考答案：}不能。存量可保持平稳，但累计分配达到 1000 MiB，还可能造成 GC 和分配器成本。应将估计目标从“存量趋势”扩展到 alloc\_space、对象数量、GC CPU 和业务吞吐。净变化法在自己的目标上高效，不代表覆盖全部性能问题。

\noindent\textbf{题 2：}ePPing 记录 TSval=8 的首包在 100 \ensuremath{\mu}s、第二包在 400 \ensuremath{\mu}s，ACK 回显在 2100 \ensuremath{\mu}s。覆盖首时间戳会造成什么误差？

\noindent\textbf{参考答案：}首边缘估计为 2000 \ensuremath{\mu}s；覆盖后变成 1700 \ensuremath{\mu}s，低估 300 \ensuremath{\mu}s。这是协议时间戳粒度与状态更新造成的误差，增加计时器精度也不能修复。实际算法还需处理流方向、重传和 delayed ACK。

\noindent\textbf{题 3：}一个 PGO 实验用同一批流量采集 profile 并验证性能提高，如何补成研究级证据？

\noindent\textbf{参考答案：}保留原测量但限定其为训练分布内收益；另划分未见负载，加入请求结构漂移、构建变化和不同比例接口，比较无 PGO 基线，重复并报告吞吐、尾延迟和体积。记录 profile 匹配覆盖及采集开销，失败负载不得事后删除。

\noindent\textbf{题 4：}原文公式按声明变量产生负概率时，教材应怎样处理？

\noindent\textbf{参考答案：}引用原始公式与变量定义，给出可重算反例，明确这是印刷表达的一致性疑问。可另给通用统计模型推导，但不能据此宣称源码有同样 bug，也不能静默替作者修正。需要对应发布源码和实现行为验证后再作实现层结论。

\medskip\noindent\textbf{本章主要阅读：}\cite{paper-scalene}\cite{paper-hardware}\cite{paper-epping}\cite{front-go121}\cite{front-trace2024}\cite{front-flight2025}\cite{front-flight-source}\cite{front-otel2024}\cite{front-otel-spec}。
% BEGIN APPENDED MECHANISM CHAPTERS

\begingroup
% Scope Unicode-capable Latin/monospace fonts to the added chapters only.
\setmainfont{DejaVu Serif}
\setsansfont{DejaVu Sans}
\setmonofont[Scale=.85]{DejaVu Sans Mono}
\part*{第五篇 · eBPF、pprof 与 Perfetto 机制深研}
\addcontentsline{toc}{part}{第五篇 · eBPF、pprof 与 Perfetto 机制深研}
本篇接续前四篇的测量、工具、工程与研究基础，从固定版本源码追踪数据如何产生、传输与解释。正文和图示说明算法及其边界；共享教学输入用于验证格式、聚合和查询，不冒充运行时采样或生产收益。

源码基线为 Linux v6.12、libbpf v1.5.0、Go 1.25.0、固定提交的 google/pprof 与 Perfetto v58.2。实际分析客户端为 Go 1.22.12 内置 pprof 与官方 Perfetto v58.2；eBPF 实验完成源码与 API 校验，未进行特权加载。这些边界与原教材各工具实验的验证范围分别保留。

\chapter{工具深研导论：测量模型与证据边界}
\label{ch:17}

本专题围绕“程序为什么慢”建立一条可检验的证据链。eBPF 提供在事件发生处运行受限程序的机制；pprof 把带权调用栈组织成成本分布；Perfetto 组织时间事件、线程调度和应用区间，并通过查询解释它们的关系。三者处于不同层次，可以组成一套诊断系统。Perfetto 与 Linux perf 也不同：perf 是 Linux 性能事件基础设施及工具集；Perfetto 是另一套追踪、处理与可视化体系。名字相近不意味着数据格式或命令兼容。


本专题的阅读顺序是：测量模型 →\allowbreak{} eBPF 如何获得数据 →\allowbreak{} pprof 如何归因与聚合 →\allowbreak{} Perfetto 如何保留并查询时间关系 →\allowbreak{} 同一请求的对照实验 →\allowbreak{} 组合诊断与研读路线。读者需要操作系统、并发编程与基本概率知识；重要公式均解释其假设。图和动画用于展示机制，完整推导与源码索引用于复核。本篇作为第17–22章追加在原第1–16章之后，承接运行时、剖析工具、开源实现与研究方法的基础。


\section{用七个问题定义一次观察}

\Needspace{7.34cm}
\begin{center}\begin{minipage}{\linewidth}
\centering
\begin{tikzpicture}[x=0.016cm,y=-0.016cm]
\path[use as bounding box] (0,0) rectangle (960,340);
\draw[rounded corners=5pt,draw={rgb,255:red,25;green,118;blue,210},fill={rgb,255:red,25;green,118;blue,210},fill opacity=.075] (40.0,75) rectangle (310.0,187);
\node[anchor=center,text={rgb,255:red,35;green,52;blue,72},font=\fontsize{9.1}{11.0}\selectfont] at (175.0,117.5) {eBPF};
\node[anchor=center,text={rgb,255:red,35;green,52;blue,72},font=\fontsize{8.2}{9.9}\selectfont] at (175.0,144.5) {事件触发与受限执行};
\draw[draw={rgb,255:red,99;green,115;blue,134},-{Stealth[length=4pt]}] (313.0,130) -- (341.0,130);
\draw[rounded corners=5pt,draw={rgb,255:red,129;green,80;blue,181},fill={rgb,255:red,129;green,80;blue,181},fill opacity=.075] (345.0,75) rectangle (615.0,187);
\node[anchor=center,text={rgb,255:red,35;green,52;blue,72},font=\fontsize{9.1}{11.0}\selectfont] at (480.0,117.5) {pprof};
\node[anchor=center,text={rgb,255:red,35;green,52;blue,72},font=\fontsize{8.2}{9.9}\selectfont] at (480.0,144.5) {栈 + 类型 + 权重};
\draw[draw={rgb,255:red,99;green,115;blue,134},-{Stealth[length=4pt]}] (618.0,130) -- (646.0,130);
\draw[rounded corners=5pt,draw={rgb,255:red,88;green,167;blue,0},fill={rgb,255:red,88;green,167;blue,0},fill opacity=.075] (650.0,75) rectangle (920.0,187);
\node[anchor=center,text={rgb,255:red,35;green,52;blue,72},font=\fontsize{9.1}{11.0}\selectfont] at (785.0,117.5) {Perfetto};
\node[anchor=center,text={rgb,255:red,35;green,52;blue,72},font=\fontsize{8.2}{9.9}\selectfont] at (785.0,144.5) {事件 + 身份 + 时间};
\node[anchor=center,text={rgb,255:red,35;green,52;blue,72},font=\fontsize{8.6}{10.4}\selectfont] at (480,235) {可分别使用，也可通过明确的数据转换组合};
\node[anchor=center,text={rgb,255:red,35;green,52;blue,72},font=\fontsize{8.6}{10.4}\selectfont] at (480,264) {采集能力、聚合语义与时间语义必须逐层说明};
\end{tikzpicture}
\captionof{figure}{三种能力，三个层次。eBPF 是可编程采集机制；pprof 组织栈权重；Perfetto 组织时间事件。箭头表示一种可能的数据通路，不是内置的自动转换。}
\end{minipage}\end{center}

每次观察都应明确：观察对象是谁，触发条件是什么，记录什么数值，权重怎样产生，归因到哪里，使用什么时钟，哪些事件不在覆盖范围。相同的矩形图可以代表 CPU 纳秒、分配字节、等待时间或对象数量。如果没有这个契约，“最宽的函数就是最慢的函数”只是猜测。\cite{pp-proto}\cite{pp-profile-kinds}


把原始观察写成 \(e_i=(t_i,o_i,k_i,x_i,s_i)\)：时间戳、对象身份、事件类型、数值和栈。采样与过滤先选出一部分观察；编码器可能再把多个观察合并。令 \(w_i\) 为选定数值维度最终入账的贡献，已包含生产者进行的采样校正（如有）；它既不必等于事件次数，也不等于包含概率本身。普通栈 profile 可写为 \(P(s,k)=\sum_i w_i\,1\{s_i=s,\,k_i=k\}\)。这个求和保留同一归因桶的总量，却通常丢掉桶内事件的先后。trace 则保留一组有时间戳的事件或区间，仍需知道事件来源、丢失情况和时钟关系，才能解释其顺序。\cite{pp-proto}\cite{pp-cpu-builder}


不能把“profile 有抽样、trace 没有抽样”当成定义。profile 可以由全量事件聚合而来；trace 可以因配置筛选、缓冲区覆盖或主动采样而不完整。真正的区别在于输出保留哪些维度，而不是文件后缀。甚至同一采集器也可以只在内核 map 中增加栈计数，或把每个事件送到用户态：前者节约传输，后者保留更多关联机会，代价随事件速率增长。


\section{CPU、墙钟、等待：先看单位再做加法}

请求墙钟时间是开始到结束的跨度 \(T=t_{end}-t_{start}\)。请求归属集合为 \(G\) 时，其 CPU 消耗可写成 \(C=\sum_{g\in G}\int 1\{g\text{ 正在 CPU 上执行}\}\,dt\)。这里的 \(g\) 是明确归属的执行活动，不能仅凭与请求时间重叠来选取。多个活动并行时，\(C\) 可以大于 \(T\)；\(C/T\) 是容量尺度，不是请求的完成百分比。


对一个给定的 OS 线程，在完整且无重叠的分类条件下，有 \(T_{thread}=T_{running}+T_{runnable}+T_{blocked}\)。running 是得到 CPU；runnable 是可以运行但尚未得到 CPU；blocked 是还不满足继续执行的条件。这是教学三分法；实际内核状态更丰富，trace 开头、结尾和丢事件处还会有未知区间。对一个跨多个 goroutine 的 Go 请求，不能直接把这三个 OS 线程时间相加当作请求分解。G 会在 M 上切换，同一 M 也会先后执行不同 G。\cite{pp-cpu-linux}\cite{pp-sema}


例如一个请求主线程执行 8 ms、等待 12 ms，另有工作线程执行 15 ms，墙钟却只有 20 ms。CPU 是 23 ms；主线程 12 ms 非运行时间又可能由 2 ms runnable 和 10 ms sleeping 构成。CPU profile 可以发现 15 ms 的计算来源，却不会自动显示 10 ms 的等待区间。线程调度 trace 能重建状态变化，却不能仅凭 Sleeping 就证明它在等待磁盘。第21章将让实际工具读取这一合成模型。


\section{观察关系不等于因果关系}

\Needspace{7.34cm}
\begin{center}\begin{minipage}{\linewidth}
\centering
\begin{tikzpicture}[x=0.016cm,y=-0.016cm]
\path[use as bounding box] (0,0) rectangle (960,340);
\draw[rounded corners=5pt,draw={rgb,255:red,88;green,167;blue,0},fill={rgb,255:red,88;green,167;blue,0},fill opacity=.075] (30,60) rectangle (270,160);
\node[anchor=center,text={rgb,255:red,35;green,52;blue,72},font=\fontsize{9.1}{11.0}\selectfont] at (150.0,96.5) {业务请求};
\node[anchor=center,text={rgb,255:red,35;green,52;blue,72},font=\fontsize{8.2}{9.9}\selectfont] at (150.0,123.5) {request=demo};
\draw[rounded corners=5pt,draw={rgb,255:red,129;green,80;blue,181},fill={rgb,255:red,129;green,80;blue,181},fill opacity=.075] (360,60) rectangle (600,160);
\node[anchor=center,text={rgb,255:red,35;green,52;blue,72},font=\fontsize{9.1}{11.0}\selectfont] at (480.0,96.5) {运行时任务};
\node[anchor=center,text={rgb,255:red,35;green,52;blue,72},font=\fontsize{8.2}{9.9}\selectfont] at (480.0,123.5) {G / task};
\draw[rounded corners=5pt,draw={rgb,255:red,25;green,118;blue,210},fill={rgb,255:red,25;green,118;blue,210},fill opacity=.075] (690,60) rectangle (930,160);
\node[anchor=center,text={rgb,255:red,35;green,52;blue,72},font=\fontsize{9.1}{11.0}\selectfont] at (810.0,96.5) {OS 执行实体};
\node[anchor=center,text={rgb,255:red,35;green,52;blue,72},font=\fontsize{8.2}{9.9}\selectfont] at (810.0,123.5) {PID / TID};
\draw[draw={rgb,255:red,99;green,115;blue,134},-{Stealth[length=4pt]}] (275,110) -- (355,110);
\draw[draw={rgb,255:red,99;green,115;blue,134},-{Stealth[length=4pt]}] (605,110) -- (685,110);
\node[anchor=center,text={rgb,255:red,35;green,52;blue,72},font=\fontsize{9.1}{11.0}\selectfont] at (480,220) {一个请求可涉及多个 G；一个 M 会执行不同 G};
\node[anchor=center,text={rgb,255:red,35;green,52;blue,72},font=\fontsize{8.6}{10.4}\selectfont] at (480,258) {观察到时间重叠，只能得到候选关联，不能直接确认归属};
\end{tikzpicture}
\captionof{figure}{三种身份不能混为一谈。请求、Go goroutine、OS 线程分别有生命周期；箭头表示需显式建立的关联，不是ID相等。}
\end{minipage}\end{center}

调用关系、调度关系和业务依赖关系是三种图。调用栈描述同一执行上下文的嵌套；线程调度描述 CPU 上的先后占用；请求依赖描述发起、排队、处理、唤醒和结果交付。worker 完成后 main 才继续，是可能的依赖；把 worker 画成 main 当前栈的子函数，则会捏造调用关系。相反，保留两个独立栈并加正确的任务身份，才保留了二者关系的性质。\cite{pp-labels}


即使两个事件显示在同一时间线上，“A 先于 B”也不充分证明 A 导致 B。需要共享请求 ID、显式 flow、唤醒链或协议证据。跨时钟的差值还受转换误差影响；若两事件的相对误差上界大于观测间隔，严格先后便不能确定。时间戳精度为纳秒，也不表示真实因果排序具备纳秒精度。


三种身份尤其值得单独保存：业务对象身份、运行时对象身份、OS 对象身份。请求 ID 不等于 goroutine ID；goroutine ID 不等于 Linux tid；tid 还可能在进程生命周期外复用。把它们关联起来需要明确的生命周期和映射，不应把一个数值相等当作同一对象。pprof 标签适合分组；Perfetto 的 track/utid 与生命周期适合时序组织；eBPF 可在可观察的边界采集关联数据，任何一方都不会自动知道全部业务语义。


\section{从“热点”到可复核的优化结论}

\Needspace{7.34cm}
\begin{center}\begin{minipage}{\linewidth}
\centering
\begin{tikzpicture}[x=0.016cm,y=-0.016cm]
\path[use as bounding box] (0,0) rectangle (960,340);
\draw[rounded corners=5pt,draw={rgb,255:red,25;green,118;blue,210},fill={rgb,255:red,25;green,118;blue,210},fill opacity=.075] (40.0,75) rectangle (310.0,187);
\node[anchor=center,text={rgb,255:red,35;green,52;blue,72},font=\fontsize{9.1}{11.0}\selectfont] at (175.0,117.5) {控制条件};
\node[anchor=center,text={rgb,255:red,35;green,52;blue,72},font=\fontsize{8.2}{9.9}\selectfont] at (175.0,144.5) {工作量 / 构建 / 环境};
\draw[draw={rgb,255:red,99;green,115;blue,134},-{Stealth[length=4pt]}] (313.0,130) -- (341.0,130);
\draw[rounded corners=5pt,draw={rgb,255:red,129;green,80;blue,181},fill={rgb,255:red,129;green,80;blue,181},fill opacity=.075] (345.0,75) rectangle (615.0,187);
\node[anchor=center,text={rgb,255:red,35;green,52;blue,72},font=\fontsize{9.1}{11.0}\selectfont] at (480.0,117.5) {提出机制};
\node[anchor=center,text={rgb,255:red,35;green,52;blue,72},font=\fontsize{8.2}{9.9}\selectfont] at (480.0,144.5) {热点 / 等待 / 队列};
\draw[draw={rgb,255:red,99;green,115;blue,134},-{Stealth[length=4pt]}] (618.0,130) -- (646.0,130);
\draw[rounded corners=5pt,draw={rgb,255:red,88;green,167;blue,0},fill={rgb,255:red,88;green,167;blue,0},fill opacity=.075] (650.0,75) rectangle (920.0,187);
\node[anchor=center,text={rgb,255:red,35;green,52;blue,72},font=\fontsize{9.1}{11.0}\selectfont] at (785.0,117.5) {干预与复测};
\node[anchor=center,text={rgb,255:red,35;green,52;blue,72},font=\fontsize{8.2}{9.9}\selectfont] at (785.0,144.5) {延迟 / CPU / 错误率};
\node[anchor=center,text={rgb,255:red,35;green,52;blue,72},font=\fontsize{8.6}{10.4}\selectfont] at (480,235) {同时保存原始数据、配置、版本与缺失指标};
\node[anchor=center,text={rgb,255:red,35;green,52;blue,72},font=\fontsize{8.6}{10.4}\selectfont] at (480,264) {图形是证据的呈现，不替代对照实验};
\end{tikzpicture}
\captionof{figure}{观察到的热点怎样变成可检验的结论。优化需要控制工作量并进行干预。源码解释、合成输入验证与真实负载实验支持不同层次的结论。}
\end{minipage}\end{center}

一个可信的结论通常经历四步：用稳定工作量复现症状；用测量定位资源消耗或等待的位置；提出能够改变该机制的干预；在相同工作量和条件下重复验证吞吐、延迟和资源变化。剖析得到的是观察证据，干预前后重复实验才进一步支持因果解释。只比较两张火焰图的形状，会混入请求量、流量组成和采样误差。\cite{pp-diff}\cite{pp-total}


测量开销也要分层。事件触发次数为 \(N\)、每次执行和传输的平均成本为 \(c\)，则一阶模型是额外工作约 \(Nc\)；真实成本还包括缓存扰动、锁竞争、内存分配、调度和下游处理。降低采样频率通常降低部分成本，却增加估计方差；先在内核聚合减少传输，却放弃逐事件时序。这些是可量化的设计权衡，不能用“eBPF 天生低开销”或“追踪一定很重”代替测量。


本专题把证据分为三类：固定版本源码及官方文档，合成数据在真实工具中的执行结果，以及尚未运行的实验建议。合成实验验证格式和算法，不验证生产负载中的采样精度；阅读 verifier 源码不等于在此机器上成功加载了 BPF 程序。这一分层贯穿正文、图注和验证报告。


\chapter{eBPF 性能观测：从可验证程序到可解释证据}
\label{ch:18}


\begin{quote}本章是固定版本的原理分析样本，不把分析版本称为最新版本。Linux 基线为正式标签 \texttt{v6.\allowbreak{}12}，完整提交 \texttt{adc21867\allowbreak{}6eef2557\allowbreak{}54692347\allowbreak{}09c2d871\allowbreak{}85ca223a}；libbpf 为 \texttt{v1.\allowbreak{}5.\allowbreak{}0}，提交 \texttt{09b9e831\allowbreak{}02eb8ab9\allowbreak{}e540d36b\allowbreak{}4559c55f\allowbreak{}3bcdb95d}；Go runtime 补充样本为 \texttt{go1.\allowbreak{}25.\allowbreak{}0}，提交 \texttt{6e676ab2\allowbreak{}b809d466\allowbreak{}23acb598\allowbreak{}8248d95d\allowbreak{}1eb7939c}。源码按不可变提交下载；来源与实际阅读区间见 \texttt{research/\allowbreak{}sources-\allowbreak{}ebpf.\allowbreak{}json}，缓存见 \texttt{research/\allowbreak{}ebpf/\allowbreak{}}。本章的成本数字和布局偏移示例均为教学假设，不能作为实测性能结论。

\end{quote}

\section{先定义观测问题，再选择 eBPF 机制}

\Needspace{7.34cm}
\begin{center}\begin{minipage}{\linewidth}
\centering
\begin{tikzpicture}[x=0.016cm,y=-0.016cm]
\path[use as bounding box] (0,0) rectangle (960,340);
\draw[rounded corners=5pt,draw={rgb,255:red,25;green,118;blue,210},fill={rgb,255:red,25;green,118;blue,210},fill opacity=.075] (40.0,75) rectangle (310.0,187);
\node[anchor=center,text={rgb,255:red,35;green,52;blue,72},font=\fontsize{9.1}{11.0}\selectfont] at (175.0,117.5) {用户态 loader};
\node[anchor=center,text={rgb,255:red,35;green,52;blue,72},font=\fontsize{8.2}{9.9}\selectfont] at (175.0,144.5) {ELF / BTF / 重定位};
\draw[draw={rgb,255:red,99;green,115;blue,134},-{Stealth[length=4pt]}] (313.0,130) -- (341.0,130);
\draw[rounded corners=5pt,draw={rgb,255:red,129;green,80;blue,181},fill={rgb,255:red,129;green,80;blue,181},fill opacity=.075] (345.0,75) rectangle (615.0,187);
\node[anchor=center,text={rgb,255:red,35;green,52;blue,72},font=\fontsize{9.1}{11.0}\selectfont] at (480.0,117.5) {内核加载};
\node[anchor=center,text={rgb,255:red,35;green,52;blue,72},font=\fontsize{8.2}{9.9}\selectfont] at (480.0,144.5) {verifier → JIT 等路径};
\draw[draw={rgb,255:red,99;green,115;blue,134},-{Stealth[length=4pt]}] (618.0,130) -- (646.0,130);
\draw[rounded corners=5pt,draw={rgb,255:red,88;green,167;blue,0},fill={rgb,255:red,88;green,167;blue,0},fill opacity=.075] (650.0,75) rectangle (920.0,187);
\node[anchor=center,text={rgb,255:red,35;green,52;blue,72},font=\fontsize{9.1}{11.0}\selectfont] at (785.0,117.5) {事件触发};
\node[anchor=center,text={rgb,255:red,35;green,52;blue,72},font=\fontsize{8.2}{9.9}\selectfont] at (785.0,144.5) {执行 → map / 输出};
\node[anchor=center,text={rgb,255:red,35;green,52;blue,72},font=\fontsize{8.6}{10.4}\selectfont] at (480,235) {加载成功与测量正确是两个判断};
\node[anchor=center,text={rgb,255:red,35;green,52;blue,72},font=\fontsize{8.6}{10.4}\selectfont] at (480,264) {程序类型、attach 点、helper 和权限共同约束能力};
\end{tikzpicture}
\captionof{figure}{eBPF 的加载与运行生命周期。加载时验证程序，事件发生时执行程序；CO-RE 是加载前后的重定位环节之一，attach 决定何时触发。}
\end{minipage}\end{center}

eBPF 是 Linux 内核中的受约束程序执行机制及其生态，不是某个 profiler 的名字。一个可运行的性能观测系统至少包括程序生成、加载与验证、事件挂载、内核数据结构、用户态消费、符号化和结果解释。libbpf 主要承担 ELF 处理与加载、挂载和部分消费工作；它不会自动决定何种采样代表 CPU 成本，也不会替用户恢复请求因果链。把所有这些层都简称为“eBPF”会掩盖错误发生的位置。例如程序通过验证却没有挂载，得到的不是零事件工作负载，而是没有建立观测入口。\cite{eb-libbpf-overview}\cite{eb-libbpf-load}


应当先写出一份测量契约：观测对象是进程、线程、容器还是请求；事件是 CPU 时钟采样、硬件事件溢出、函数调用还是调度切换；值是次数、纳秒、字节还是按周期加权的计数；时间窗口从哪一步开始；哪些分支会过滤或丢弃；调用栈失败如何保留。比如“统计目标进程十秒内的调度切出次数”可以由 \texttt{sched\_\allowbreak{}switch} 完成；“解释某个 Go 请求的十秒延迟”则还需要 goroutine 状态、应用请求身份与跨线程传播证据。二者可使用同一内核机制，但测量对象和结论强度不同。\cite{eb-sched-trace}\cite{eb-sched-core}\cite{eb-go-runtime}


控制平面负责发现内核能力、选择程序变体、解析对象、创建 map、验证和挂载，并管理链接与对象的生命周期。数据平面在事件发生时执行过滤、读字段、更新聚合值或保留记录；用户态随后读出、符号化、合并并上传。控制平面“启动成功”不能替代数据平面的健康检查：加载日志说明程序被接受，样本计数和失败计数才说明它实际观察到了什么。相反，消费速度慢也不必然意味着内核程序慢，可能是用户态对每条记录进行格式化或同步网络写入。\cite{eb-libbpf-overview}\cite{eb-ring-consume}


eBPF 的价值之一是把筛选和聚合前移。例如每秒百万次事件中只关心一个进程，先比较进程标识，再更新固定数量的桶，可以显著减少跨边界输出。但前移操作仍在被观测系统的执行路径上付费，开销包括挂载机制、读取用户内存、查表、原子操作和通知。能加载的程序并不等于适合以任意事件率长期运行的程序；验证器证明的是特定安全约束，而不是业务开销上限，更不是“生产零开销”。\cite{eb-verifier-doc}\cite{eb-map-array-code}\cite{eb-ring-code}


学习本章后应能从一次输出反向追踪到事件入口、抽象状态、共享数据和消费顺序；给出“哪些内容被观测、哪些没有”的严格说明；并能发现三个常见错误：把线程当 goroutine，把缺失栈当缺失事件，把 CO-RE 通过当语义正确。后续所有图和练习围绕这一证据链展开。\cite{eb-go-runtime}\cite{eb-stack-helper}\cite{eb-core-relo}


\section{编译、ELF、BTF、CO-RE 与加载}

\Needspace{7.34cm}
\begin{center}\begin{minipage}{\linewidth}
\centering
\begin{tikzpicture}[x=0.016cm,y=-0.016cm]
\path[use as bounding box] (0,0) rectangle (960,340);
\draw[rounded corners=5pt,draw={rgb,255:red,25;green,118;blue,210},fill={rgb,255:red,25;green,118;blue,210},fill opacity=.075] (25,60) rectangle (290,180);
\node[anchor=center,text={rgb,255:red,35;green,52;blue,72},font=\fontsize{9.1}{11.0}\selectfont] at (157.5,106.5) {编译时布局（示意）};
\node[anchor=center,text={rgb,255:red,35;green,52;blue,72},font=\fontsize{8.2}{9.9}\selectfont] at (157.5,133.5) {字段 state：偏移 16};
\draw[rounded corners=5pt,draw={rgb,255:red,129;green,80;blue,181},fill={rgb,255:red,129;green,80;blue,181},fill opacity=.075] (345,60) rectangle (610,180);
\node[anchor=center,text={rgb,255:red,35;green,52;blue,72},font=\fontsize{9.1}{11.0}\selectfont] at (477.5,106.5) {BTF + relocation};
\node[anchor=center,text={rgb,255:red,35;green,52;blue,72},font=\fontsize{8.2}{9.9}\selectfont] at (477.5,133.5) {定位同名结构成员};
\draw[rounded corners=5pt,draw={rgb,255:red,88;green,167;blue,0},fill={rgb,255:red,88;green,167;blue,0},fill opacity=.075] (665,60) rectangle (930,180);
\node[anchor=center,text={rgb,255:red,35;green,52;blue,72},font=\fontsize{9.1}{11.0}\selectfont] at (797.5,106.5) {目标布局（示意）};
\node[anchor=center,text={rgb,255:red,35;green,52;blue,72},font=\fontsize{8.2}{9.9}\selectfont] at (797.5,133.5) {字段 state：偏移 24};
\draw[draw={rgb,255:red,99;green,115;blue,134},-{Stealth[length=4pt]}] (295,120) -- (340,120);
\draw[draw={rgb,255:red,99;green,115;blue,134},-{Stealth[length=4pt]}] (615,120) -- (660,120);
\node[anchor=center,text={rgb,255:red,35;green,52;blue,72},font=\fontsize{9.1}{11.0}\selectfont] at (480,240) {迁移的是结构访问意图；功能探测和语义兼容仍需另行处理};
\end{tikzpicture}
\captionof{figure}{CO-RE：把字段意图迁移到目标布局。图中字段偏移为示意数值。CO-RE 用 BTF 类型与 relocation 信息重写相应访问，不保证 helper、内核功能和语义都兼容。}
\end{minipage}\end{center}

\subsection{指令机器：寄存器、调用与受限栈}

BPF 的寄存器编号为 \texttt{R0} 至 \texttt{R10}，均为六十四位：\texttt{R0} 放返回值，\texttt{R1–R5} 传参数，\texttt{R6–R9} 按约定由被调用方保存；\texttt{R10} 是只读的栈帧指针，程序不能像普通本机代码那样自由修改栈指针。基本指令占八字节，编码操作码、源/目标寄存器、十六位偏移与三十二位立即数；\texttt{ldimm64} 使用两个连续槽、共十六字节，不能把第二槽当独立指令入口。\cite{eb-stack-helper}\cite{eb-verifier-doc}\cite{eb-design-qa}\cite{eb-isa}


普通 BPF 子程序调用有返回点，成功的 tail call 则转移执行而不返回。固定 v6.12 的调用帧上限为八；验证器按 \texttt{MAX\_\allowbreak{}BPF\_\allowbreak{}STACK=512} 检查普通调用链的累计栈用量，不是允许每层任意使用五百一十二字节。子程序调用混合 tail call 时还检查更严格条件：目标子程序含 tail call 且前序帧累计深度达到二百五十六字节会被拒绝。须同时核验版本、调用形态与架构支持，不能把单个常量当作无限递归许可。\cite{eb-isa}\cite{eb-verifier-code}\cite{eb-verifier-state}\cite{eb-design-qa}


传统 helper 由静态编号调用，属于内核对 BPF 承诺的接口；kfunc 通过 BTF 描述与解析暴露的内核函数，接口不具有相同的稳定性承诺。二者都不是任意内核函数调用权限。kfunc 按程序类型注册，验证器还检查参数类型、可空性、引用获取/释放和可睡眠上下文；仅有函数名或 BTF 记录不足以证明当前程序可调用。增加 kfunc 依赖时应探测目标能力并记录版本，不能期待 CO-RE 自动补出缺失实现。\cite{eb-design-qa}\cite{eb-kfuncs}\cite{eb-isa}


\subsection{目标代码包含程序，也包含解释程序所需的信息}

典型构建由 Clang 把受约束的 C 编译为 BPF 指令，输出 ELF 对象。程序段、map 定义、数据段、符号表和重定位信息在同一对象中协作。\texttt{SEC("tracepoint/\allowbreak{}sched/\allowbreak{}sched\_\allowbreak{}switch")} 并不是 C 语言赋予的内核能力；它把函数放入带约定名称的 ELF section，由 libbpf 解释程序类型与自动挂载线索。\texttt{SEC}、\texttt{\_\allowbreak{}\_\allowbreak{}uint} 和 \texttt{\_\allowbreak{}\_\allowbreak{}type} 等宏帮助编码元数据，最终仍须经过加载器和内核检查。\cite{eb-helpers}\cite{eb-libbpf-overview}


普通 ELF 重定位与 CO-RE 重定位必须分开理解。前者解决符号、数据或函数在对象中的关联。例如固定版本 LLVM 重定位文档中，\texttt{R\_\allowbreak{}BPF\_\allowbreak{}64\_\allowbreak{}64} 用于 \texttt{ld\_\allowbreak{}imm64}，\texttt{R\_\allowbreak{}BPF\_\allowbreak{}64\_\allowbreak{}32} 用于调用；调用的重定位计算包含按八字节指令单位换算。CO-RE 则回答“编译时对某种类型的访问，在目标内核类型中应如何实现”。两类记录可能同处一个 ELF，却不解决同一问题；把 ELF 所有 relocation 都叫“内核字段偏移适配”会遗漏全局变量和 BPF 子程序链接。\cite{eb-llvm-reloc}\cite{eb-core-relo}


BTF 用紧凑的类型记录和字符串表描述类型关系，类型 ID 在相应 BTF 对象中引用其他记录，并不是跨内核永恒不变的类型编号。\texttt{.\allowbreak{}BTF.\allowbreak{}ext} 在 ELF 中承载函数、行号和 CO-RE 信息，头部有真实字段 \texttt{core\_\allowbreak{}relo\_\allowbreak{}off} 与 \texttt{core\_\allowbreak{}relo\_\allowbreak{}len}。还要注意单位：某些内核 API 中指令位置以指令条数表示，ELF 扩展信息使用字节偏移；把二者混用会把正确的指令指向错误位置。调试信息让故障可定位，但不会替代指令验证。\cite{eb-btf}


\subsection{CO-RE 是结构适配，不是完整兼容性证明}

CO-RE 的合作方是编译器保留的访问意图、编译时类型信息、目标 BTF 和加载器。常见流程读取 \texttt{/\allowbreak{}sys/\allowbreak{}kernel/\allowbreak{}btf/\allowbreak{}vmlinux}，也可以针对目标类型信息生成 \texttt{vmlinux.\allowbreak{}h}。加载时 libbpf 匹配目标类型与访问路径，计算字段偏移、尺寸、存在性以及类型或枚举相关值，再修补指令的立即数或内存访问偏移。固定版本 \texttt{relo\_\allowbreak{}core.\allowbreak{}c} 中可以直接读到计算与修补函数；这不是运行中每次读字段都查询 BTF 的动态反射。\cite{eb-libbpf-overview}\cite{eb-core-relo}


设教学结构 \texttt{task\_\allowbreak{}demo} 的 \texttt{latency} 字段在构建环境偏移为十六字节，在目标环境偏移为二十四字节。结构重定位可以把访问位置修正到二十四；假如字段在目标内核中的单位从纳秒变为微秒，即使名字、宽度和位置都能成功匹配，乘除比例仍可能完全错误。类似地，字段何时有效、在何种锁保护下读取、函数是否覆盖全部请求，以及读取值是否需要配对事件，都不由布局适配保证。这一例子中的偏移和字段为虚构教学值，不是任何真实 \texttt{task\_\allowbreak{}struct} 布局。\cite{eb-core-relo}\cite{eb-btf}


字段不存在也不意味着所有对象必须立刻失败。存在性检测可以引导程序走其他分支；无法完成的访问可能被加载器改写为会触发拒绝的“毒化”指令，只有实际可达时才阻止加载。因此，正确模式是能力检测与有意义的替代语义配合，而不是把失败吞掉后输出零。不要把未执行的分支解释为“该功能被完整支持”；应该把功能不可用与观测值为零分别呈现。\cite{eb-core-relo}\cite{eb-verifier-code}


\subsection{按阶段定位启动失败}

libbpf 的生命周期可以概括为打开、加载、挂载与释放。打开阶段解析 ELF，收集外部引用、BTF、map、程序和重定位；此时可以修改配置。加载阶段准备目标 BTF、解析外部依赖和重定位，创建内核对象，向内核提交程序。固定的 v1.5.0 代码中，顶层加载顺序包含目标 BTF、外部引用、map 清理、重定位、BTF 加载、map 创建与程序加载；不能机械地宣称“所有重定位都发生在 map 创建之后”。不同重定位类别的最终处理散布在相应阶段。\cite{eb-libbpf-load}


内核程序加载路径检查程序类型、权限与策略，调用安全钩子和验证器，再选择执行引擎，成功后返回可引用该程序的文件描述符。权限能力受到程序类型、内核配置、BPF token 和系统策略影响，不能用一句“只要 root”概括全部部署条件。文件描述符表示已加载对象，不表示某个事件已经引用它；必须完成相应挂载，并在运行期间保留需要的对象引用。关闭、解除链接和清理 map 也属于实验设计，否则下一次实验可能混入旧状态。\cite{eb-load-syscall}\cite{eb-libbpf-overview}


故障报告应保存失败阶段、内核和加载器版本、选中的程序变体与验证日志。例如找不到 section 是对象问题，目标类型不匹配是重定位问题，非法内存访问是验证问题，目标 tracepoint 不存在是挂载问题。按这个顺序缩小问题，比反复调整采样频率更有效。对教学实验而言，先验证“无业务字段读取的计数程序”，再增加读字段和栈收集，可以使首次失败对应的机制变化可解释。\cite{eb-libbpf-load}\cite{eb-attach-api}\cite{eb-verifier-doc}


\section{验证器：抽象状态、路径收敛与资源义务}

\Needspace{7.98cm}
\begin{center}\begin{minipage}{\linewidth}
\centering
\begin{tikzpicture}[x=0.016cm,y=-0.016cm]
\path[use as bounding box] (0,0) rectangle (960,380);
\draw[rounded corners=5pt,draw={rgb,255:red,25;green,118;blue,210},fill={rgb,255:red,25;green,118;blue,210},fill opacity=.075] (20,105) rectangle (255,215);
\node[anchor=center,text={rgb,255:red,35;green,52;blue,72},font=\fontsize{9.1}{11.0}\selectfont] at (137.5,146.5) {map\_lookup\_elem};
\node[anchor=center,text={rgb,255:red,35;green,52;blue,72},font=\fontsize{8.2}{9.9}\selectfont] at (137.5,173.5) {PTR\_OR\_NULL};
\draw[rounded corners=5pt,draw={rgb,255:red,189;green,101;blue,0},fill={rgb,255:red,189;green,101;blue,0},fill opacity=.075] (350,20) rectangle (610,125);
\node[anchor=center,text={rgb,255:red,35;green,52;blue,72},font=\fontsize{9.1}{11.0}\selectfont] at (480.0,59.0) {空指针分支};
\node[anchor=center,text={rgb,255:red,35;green,52;blue,72},font=\fontsize{8.2}{9.9}\selectfont] at (480.0,86.0) {结束 / 不解引用};
\draw[rounded corners=5pt,draw={rgb,255:red,88;green,167;blue,0},fill={rgb,255:red,88;green,167;blue,0},fill opacity=.075] (350,220) rectangle (610,325);
\node[anchor=center,text={rgb,255:red,35;green,52;blue,72},font=\fontsize{9.1}{11.0}\selectfont] at (480.0,259.0) {非空分支};
\node[anchor=center,text={rgb,255:red,35;green,52;blue,72},font=\fontsize{8.2}{9.9}\selectfont] at (480.0,286.0) {有效 map value 指针};
\draw[draw={rgb,255:red,99;green,115;blue,134},-{Stealth[length=4pt]}] (260,150) -- (345,75);
\node[anchor=center,text={rgb,255:red,35;green,52;blue,72},font=\fontsize{7.7}{9.3}\selectfont] at (302.5,97.5) {== NULL};
\draw[draw={rgb,255:red,99;green,115;blue,134},-{Stealth[length=4pt]}] (260,170) -- (345,265);
\node[anchor=center,text={rgb,255:red,35;green,52;blue,72},font=\fontsize{7.7}{9.3}\selectfont] at (302.5,202.5) {!= NULL};
\draw[rounded corners=5pt,draw={rgb,255:red,129;green,80;blue,181},fill={rgb,255:red,129;green,80;blue,181},fill opacity=.075] (715,220) rectangle (935,325);
\node[anchor=center,text={rgb,255:red,35;green,52;blue,72},font=\fontsize{9.1}{11.0}\selectfont] at (825.0,259.0) {再检查偏移与边界};
\node[anchor=center,text={rgb,255:red,35;green,52;blue,72},font=\fontsize{8.2}{9.9}\selectfont] at (825.0,286.0) {允许受限访问};
\draw[draw={rgb,255:red,99;green,115;blue,134},-{Stealth[length=4pt]}] (615,273) -- (710,273);
\node[anchor=center,text={rgb,255:red,35;green,52;blue,72},font=\fontsize{9.6}{11.6}\selectfont] at (785,80) {证明模型内安全};
\node[anchor=center,text={rgb,255:red,35;green,52;blue,72},font=\fontsize{8.2}{9.9}\selectfont] at (785,117) {不证明业务语义正确};
\end{tikzpicture}
\captionof{figure}{验证器跟踪的是抽象状态。示例表达空指针检查如何收窄指针状态。实际 verifier 还追踪标量范围、已知位、栈初始化、引用及路径等；图不覆盖全部规则。}
\end{minipage}\end{center}

\subsection{验证的对象不是源代码文本}

验证器分析加载时的 BPF 指令及其环境。它维护寄存器、栈槽、指针来源、数值范围、调用帧以及资源引用等抽象状态。入口寄存器中含受限的上下文指针；未初始化寄存器不能直接读取。调用 helper 后，返回值进入 \texttt{R0}，\texttt{R1} 至 \texttt{R5} 按调用约定失效，\texttt{R6} 至 \texttt{R9} 可保存值。理解这一点才能解释：源代码中一个变量看似没变化，编译成指令后却因为 helper 调用跨越了寄存器生存期而不能直接沿用。\cite{eb-verifier-doc}


抽象状态不是每个输入各跑一次。对标量，它可同时记录有符号和无符号上下界，以及跟踪已知位与未知位的 tnum；对指针，它还记录对象类别、固定偏移、可变偏移和相关身份。例如一个读入的八位数起初范围为零到二百五十五；经过“小于八”的条件分支，某条路径可以缩小为零到七，从而证明以它访问八元素数组安全。若条件只约束了有符号比较，而后续计算按无符号解释，则不能仅凭源代码直觉认定范围足够。\cite{eb-verifier-doc}\cite{eb-verifier-code}


tnum 可写成 \texttt{(value,\allowbreak{} mask)}：\texttt{mask} 为一的位未知，为零的位由 \texttt{value} 指定。官方文档的字节例子从 \texttt{(0,\allowbreak{}0xff)} 出发，按位或 \texttt{0x40} 后为 \texttt{(0x40,\allowbreak{}0xbf)}。这不是统计概率，未知位之间也不承诺独立；加法的进位会扩大不确定性，因此按位信息与数值区间需要协同。对性能观测程序，位掩码、直方图索引和长度计算常直接影响是否能证明 map 访问有界。\cite{eb-verifier-doc}


\subsection{安全证明要求沿每条可达路径成立}

map 查询返回的值可能为空。必须先排除空指针，再按 map 的 value 大小读写；指针的身份关系还能让验证器把同一来源的空值检查传播到相关副本。栈内存则必须在被读取或传给 helper 前完成初始化。只初始化事件结构中的几个可见字段而漏掉填充字节，可能导致拒绝或数据泄露风险；教学代码应显式初始化所有将输出的字节。helper 可用性和参数类型由程序类型与上下文约束，不能因为另一个示例能调用就直接复制。\cite{eb-verifier-doc}\cite{eb-verifier-code}


验证器还跟踪资源义务。ring buffer 的成功 reserve 获得一项待释放引用，每条后续可达退出路径都要 submit 或 discard；保留成功后提前 \texttt{return} 会违反这一义务。它与“指针是否落在合法范围内”是两类约束：前者要求生命周期闭合，后者要求访问位置安全。把失败分支画成状态机比只盯着 happy path 更容易发现遗漏，特别是在增加多个条件过滤之后。\cite{eb-ring-doc}\cite{eb-ring-code}\cite{eb-verifier-code}


\subsection{剪枝不是看到相同指令地址就停止}

路径分支会使状态数量增长。若某个程序点的既有状态已经被证明安全，新状态又被其安全条件覆盖，则无需重走后续路径。v6.12 的 \texttt{regsafe} 会考虑类型、范围、tnum、身份关联及存活性；\texttt{states\_\allowbreak{}equal} 还比较调用帧、调用点、锁、RCU 和可睡眠状态。仅仅两个寄存器当前数值相等，不足以说明状态可互换；一个是 map 指针、另一个是普通整数，就不能以同样方式继续访问内存。\cite{eb-verifier-code}


存活性允许忽略将来不会再被读取的差异。设两条路径在一个无后续用途的临时寄存器中分别留下三和五，而用于数组索引的寄存器都满足同样的安全范围，前一个差异可能无需阻止剪枝。反之，两个值之间共享身份所表达的相关性若影响后续边界检查，就必须保留。状态收敛因此是一种保守抽象与证明复用，不能简单归结为哈希去重，也不能随意把所有相似路径折叠。\cite{eb-verifier-code}


\subsection{循环与复杂度：要读当前源码，不能复制旧口号}

v6.12 的 \texttt{verifier.\allowbreak{}rst} 开头仍含“DAG 检查禁止循环”的旧文字，但同一版本源码有循环检查、回调状态和开放编码迭代器收敛逻辑。因此本章不把“eBPF 不支持循环”当作现行规则。有界循环并不等于验证器必然接受：边界能否证明、分支组合、状态精度和处理指令复杂度都可能成为限制。源码里的已处理指令计数反映探索成本，不等于最终字节码条数；一百条指令的多分支循环可能比更长的直线程序更难验证。\cite{eb-verifier-doc}\cite{eb-verifier-code}


遇到“无限循环”或“程序过复杂”应先读验证日志，明确是循环没有证明终止，还是状态探索达到限制。一个尚未完成证明的循环里重现相似状态，不可冒充“已验证安全的状态缓存命中”。常见改进是缩短可证明的循环上界、把过滤提前、减少循环体内分支，并让边界检查紧邻危险访问。验证通过仍不证明结果公式正确，也不证明并发聚合没有竞态；安全性、终止性、统计有效性与性能预算是分别需要回答的问题。\cite{eb-verifier-code}


\section{执行引擎与挂载入口决定观测语义}

\subsection{JIT 与解释执行}

通过验证的程序由内核选择执行方式。通用代码有解释器分派表，也会调用架构 JIT，把 BPF 指令转换为本机指令。不能简单说 JIT 总是可选优化：\texttt{CONFIG\_\allowbreak{}BPF\_\allowbreak{}JIT\_\allowbreak{}ALWAYS\_\allowbreak{}ON}，以及涉及某些 kfunc 的程序，会要求 JIT；相关失败可能返回不支持，而非总能回退。常量混淆等准备步骤也可能改写指令，故原始 ELF 指令、加载后指令和最终机器码不是同一个观察层次。\cite{eb-runtime-select}


x86 的 JIT 会生成函数序言、算术与访存指令、跳转和调用，并用多轮计算稳定指令地址与生成长度。v6.12 中可见先估计每条指令位置，再反复 \texttt{do\_\allowbreak{}jit}，待程序长度收敛后分配并生成最终镜像。不同分支编码长度影响跳转距离，这是多轮过程的原因之一。JIT 降低解释分派成本，但不能消除 helper 自身、用户内存访问、缓存竞争或事件触发机制的费用。\cite{eb-x86-jit}\cite{eb-ring-code}


\subsection{四类常用入口}


{\small
\begin{longtable}{@{}p{3.665cm}@{\hspace{0.38cm}}p{3.665cm}@{\hspace{0.38cm}}p{3.665cm}@{\hspace{0.38cm}}p{3.665cm}@{}}
\toprule
入口 & 触发条件与数据含义 & 主要适用问题 & 需要核验的边界 \\
\midrule\endhead
\texttt{perf\_\allowbreak{}event} & 时间或硬件事件达到采样条件 & CPU 热点、事件加权热点 & 频率与周期、节流、复用、特权范围、栈质量 \\
tracepoint & 内核显式埋点执行 & 调度、系统调用、块设备等状态变化 & 是否真的覆盖所需状态，字段与版本语义 \\
fentry/fexit & 可挂载函数入口或退出，经 trampoline 调用 & 内核函数调用次数、时延与参数分析 & BTF、架构和内核支持、目标函数生命周期 \\
uprobe/uretprobe & 用户程序中所选文件位置触发 & 用户态库或应用函数观测 & ELF 偏移、二进制身份、优化、ABI、返回路径 \\
\bottomrule\end{longtable}}
该表描述入口机制，不给出跨环境统一开销排名。\texttt{perf\_\allowbreak{}event} 连接 BPF 程序时，libbpf 可使用 BPF link，或走 \texttt{PERF\_\allowbreak{}EVENT\_\allowbreak{}IOC\_\allowbreak{}SET\_\allowbreak{}BPF} 等兼容路径，然后启用事件。tracepoint 的 libbpf 路径也会建立相应 perf event。相同的 \texttt{SEC} 写法隐藏了加载器兼容分支，所以研究代码时应跟踪到实际系统调用与内核路径，而不止看外层函数名。\cite{eb-libbpf-load}\cite{eb-perf-abi}


fentry/fexit 的核心不是把函数名当字符串反复查询，而是依据目标信息建立 trampoline：它按调用模型保存必要状态、组织附加程序，并在需要时调用原函数。固定版本源码可见架构生成 trampoline、保护代码镜像，再通过 ftrace direct 或架构 text poke 注册入口。只有入口程序与同时存在退出程序时，trampoline 的标志与调用组织也不同。因此“所有 BPF 观测都靠断点陷入”不成立，反过来“所有挂载都是直接调用”也不成立。\cite{eb-trampoline}


uprobe 的标识关联可执行文件 inode 与文件偏移；进程虚拟地址还要结合映射的起点和文件页偏移换算。固定源码包含软件断点写入以及原始指令保存、异位执行相关结构，所以不能用 JIT 很快推导出每次 uprobe 触发没有陷入成本。libbpf 可以从符号名定位文件偏移，但符号存在不保证它代表所有逻辑调用：内联函数可能没有独立入口，二进制替换会改变身份，Go ABI 与寄存器参数布局也不能套用普通 C 约定。\cite{eb-uprobe}\cite{eb-libbpf-load}


函数入口与退出配对需要区分递归和并发，还要处理缺失事件与任务退出。用线程 ID 保存一个开始时间，会在递归调用时覆盖外层时间；若观测目标是 goroutine，运行时迁移还会破坏跨入口、退出的线程关联。选择入口本身是一项建模工作：已有语义明确的 tracepoint 通常比私有函数名更容易解释，但其覆盖范围仍要从调用点与事件定义证明。\cite{eb-go-runtime}\cite{eb-sched-core}\cite{eb-sched-trace}


\section{Maps 与 ring buffer：共享状态如何变为用户态证据}

\Needspace{7.34cm}
\begin{center}\begin{minipage}{\linewidth}
\centering
\begin{tikzpicture}[x=0.016cm,y=-0.016cm]
\path[use as bounding box] (0,0) rectangle (960,340);
\draw[rounded corners=5pt,draw={rgb,255:red,88;green,167;blue,0},fill={rgb,255:red,88;green,167;blue,0},fill opacity=.075] (35,90) rectangle (245,170);
\node[anchor=center,text={rgb,255:red,35;green,52;blue,72},font=\fontsize{9.1}{11.0}\selectfont] at (140.0,130.0) {A 已提交};
\draw[rounded corners=5pt,draw={rgb,255:red,189;green,101;blue,0},fill={rgb,255:red,189;green,101;blue,0},fill opacity=.075] (260,90) rectangle (470,170);
\node[anchor=center,text={rgb,255:red,35;green,52;blue,72},font=\fontsize{9.1}{11.0}\selectfont] at (365.0,130.0) {B busy};
\draw[rounded corners=5pt,draw={rgb,255:red,88;green,167;blue,0},fill={rgb,255:red,88;green,167;blue,0},fill opacity=.075] (485,90) rectangle (695,170);
\node[anchor=center,text={rgb,255:red,35;green,52;blue,72},font=\fontsize{9.1}{11.0}\selectfont] at (590.0,130.0) {C 已提交};
\draw[rounded corners=5pt,draw={rgb,255:red,99;green,115;blue,134},fill={rgb,255:red,99;green,115;blue,134},fill opacity=.075] (710,90) rectangle (920,170);
\node[anchor=center,text={rgb,255:red,35;green,52;blue,72},font=\fontsize{9.1}{11.0}\selectfont] at (815.0,130.0) {空闲容量};
\draw[draw={rgb,255:red,99;green,115;blue,134},-{Stealth[length=4pt]}] (140,240) -- (140,175);
\node[anchor=center,text={rgb,255:red,35;green,52;blue,72},font=\fontsize{8.6}{10.4}\selectfont] at (140,275) {消费者};
\node[anchor=center,text={rgb,255:red,189;green,101;blue,0},font=\fontsize{9.1}{11.0}\selectfont] at (490,215) {读完 A 后在 B 等待};
\node[anchor=center,text={rgb,255:red,35;green,52;blue,72},font=\fontsize{8.6}{10.4}\selectfont] at (490,255) {C 已提交也不能越过 B};
\end{tikzpicture}
\captionof{figure}{ring buffer 的保留顺序与可见性。多生产者共享 ring buffer。较早记录仍 busy 时，消费者不能跳过它消费后面的已提交记录；reserve 失败必须计数。}
\end{minipage}\end{center}

\subsection{聚合优先，但必须保留质量信息}

map 是有类型、大小与具体实现的数据结构，而不是无限容量的字典。ARRAY 的元素预分配并初始化为零，键界定在固定范围内；数组 value 按八字节对齐。按 CPU 分片的 PERCPU\_\allowbreak{}ARRAY 为每个可能 CPU 配置独立 value，BPF 端查询获得当前 CPU 分片，用户态查询需要为所有 possible CPU 的对齐值准备缓冲区。这里不能用容器配额 CPU 数或当前 online CPU 数替代内核期望的 possible CPU 数。\cite{eb-map-array-doc}\cite{eb-map-array-code}\cite{eb-attach-api}


按 CPU 分片减少跨 CPU 写同一缓存行的竞争，却不是普适的“无并发”。还需要考虑程序是否可能嵌套执行、用户态是否同时读取、多个字段是否要求联合一致。全局计数器的读改写可能需要原子操作或适用的同步机制；用户态依次读多个 CPU 也不是全机同一瞬间的原子快照。对于停止后对账，应先解除入口并让在途执行和消费达到稳定，再读取最终值；持续观测时要明确快照误差或使用版本化、分段窗口等协议。\cite{eb-map-array-doc}\cite{eb-map-array-code}\cite{eb-ring-consume}


内存预算必须包含分片倍数。假设一个直方图有六十四个元素，每个 value 八字节，possible CPU 为一百二十八，则纯 value 数据为 \texttt{64×8×128=65536} 字节，还没算指针、map 元数据和分配开销。全局 map 的较小内存与 per-CPU 的较少共享写之间存在权衡。把全部原始事件都输出也不是免费替代：它把压力转移到 ring、消费回调和后续存储。实际工具常用内核聚合，同时保留少量诊断事件和各阶段失败计数。\cite{eb-map-array-code}\cite{eb-ring-doc}


\subsection{reserve、submit 与 consume 是三个独立动作}

BPF ring buffer 是多生产者、单消费推进模型的环形记录通道，跨 CPU 共享容量。它的 map key/value 大小为零，不能像普通 map 那样 lookup/update 数据；容量须满足相应的幂次和页对齐约束。\texttt{bpf\_\allowbreak{}ringbuf\_\allowbreak{}output} 把已有数据复制进去；\texttt{bpf\_\allowbreak{}ringbuf\_\allowbreak{}reserve} 返回 ring 中可填写的区域，随后必须调用公开 helper \texttt{bpf\_\allowbreak{}ringbuf\_\allowbreak{}submit} 或 \texttt{bpf\_\allowbreak{}ringbuf\_\allowbreak{}discard}。文档中的“commit”描述发布动作，内核内部函数也叫 commit，但可复制的 BPF 示例应使用真实公开名称 submit。\cite{eb-ring-doc}\cite{eb-ring-code}


普通 reserve helper 的大小需要验证器能确定，适合固定格式事件；不能据此断言 eBPF 所有 ring API 都不支持动态大小，因为还有 dynptr 等不同接口与约束。本章用固定事件避免引入另一套资源证明。reserve 失败立即返回，不会等消费者腾出空间；失败原因既可能是容量不足，也可能是特定执行上下文中锁不能取得。尤其 NMI 路径使用 trylock，即使有空间，也不承诺 reserve 一定成功。\cite{eb-ring-doc}\cite{eb-ring-code}


每条记录前面有八字节头部，记录总长度再按八字节对齐。头部长度字段含 busy 和 discard 标志，另有供提交定位 ring 的信息。生产者串行确定保留位置，然后可以并行填充；提交清除 busy，消费者通过 acquire 语义看见已发布的数据，并通过 release 推进消费位置。双重虚拟映射把绕回的物理页在地址空间中连续呈现，简化跨尾记录访问；它不意味着物理数据容量翻倍。\cite{eb-ring-doc}\cite{eb-ring-code}\cite{eb-ring-consume}


跨 CPU 的输出顺序是保留位置的顺序，而不是谁先调用 submit。设 A 先 reserve 后被延迟，B 后 reserve 并立即 submit；消费者看到 A 仍 busy 时会停下，不能先跳过 A 消费 B。只有 A submit 或 discard 后才能向前推进。这样提供跨 CPU 的保留顺序一致性，也引入头部阻塞。事件在 reserve 之前取时间戳时，时间戳顺序还可能与保留顺序不同；二者均不能自动升级为业务因果顺序。\cite{eb-ring-doc}\cite{eb-ring-code}\cite{eb-ring-consume}


容量计算要看实现边界。对容量 \texttt{R=4096}、payload 五十六字节的教学例子，每条占 \texttt{align8(56+8)=64} 字节。v6.12 判断新生产位置与消费位置、最早 pending 位置的距离不能超过 \texttt{mask=R−1}，所以完全等长、未消费情形最多保留 \texttt{floor(4095/\allowbreak{}64)=63} 条，而不是六十四条。这个推导只计算数据环可容纳的记录，不包含单独的元数据页，也不排除 NMI 锁竞争先使保留失败。\cite{eb-ring-code}


\subsection{用户态消费也影响观测质量}

libbpf 的消费者先检查 busy/discard，再调用用户回调，随后才推进消费位置。每条记录都打印完整 JSON 会延迟空间回收，从而增加生产端的保留失败。即使回调返回负值，固定版本实现也会推进越过当前记录后返回错误，所以不能假设下一次 poll 自动重放该条。应为解码失败、回调失败和上传失败设置独立指标，不要让它们与内核 ring 丢失共用一个含糊的 \texttt{lost} 数。\cite{eb-ring-consume}


自适应唤醒、强制唤醒和禁止唤醒控制通知策略，不改变记录发布与消费的基本正确性要求。批处理可降低唤醒成本，但增加延迟与突发队列占用；共享 ring 节省 per-CPU 空闲容量，却也让所有生产者共同依赖消费者吞吐。扩容能吸收短突发，不能修复长期生产率大于消费率。若必须做低延迟告警，应同时约束最大回调时间、消费调度间隔和允许积压量，而不只调大 map 大小。\cite{eb-ring-doc}\cite{eb-ring-consume}


\section{采样、栈展开与 Go 运行时边界}

\Needspace{7.34cm}
\begin{center}\begin{minipage}{\linewidth}
\centering
\begin{tikzpicture}[x=0.016cm,y=-0.016cm]
\path[use as bounding box] (0,0) rectangle (960,340);
\draw[rounded corners=5pt,draw={rgb,255:red,25;green,118;blue,210},fill={rgb,255:red,25;green,118;blue,210},fill opacity=.075] (40.0,75) rectangle (310.0,187);
\node[anchor=center,text={rgb,255:red,35;green,52;blue,72},font=\fontsize{9.1}{11.0}\selectfont] at (175.0,117.5) {采样现场};
\node[anchor=center,text={rgb,255:red,35;green,52;blue,72},font=\fontsize{8.2}{9.9}\selectfont] at (175.0,144.5) {PC / SP / 寄存器};
\draw[draw={rgb,255:red,99;green,115;blue,134},-{Stealth[length=4pt]}] (313.0,130) -- (341.0,130);
\draw[rounded corners=5pt,draw={rgb,255:red,129;green,80;blue,181},fill={rgb,255:red,129;green,80;blue,181},fill opacity=.075] (345.0,75) rectangle (615.0,187);
\node[anchor=center,text={rgb,255:red,35;green,52;blue,72},font=\fontsize{9.1}{11.0}\selectfont] at (480.0,117.5) {展开与身份};
\node[anchor=center,text={rgb,255:red,35;green,52;blue,72},font=\fontsize{8.2}{9.9}\selectfont] at (480.0,144.5) {FP / DWARF / runtime};
\draw[draw={rgb,255:red,99;green,115;blue,134},-{Stealth[length=4pt]}] (618.0,130) -- (646.0,130);
\draw[rounded corners=5pt,draw={rgb,255:red,88;green,167;blue,0},fill={rgb,255:red,88;green,167;blue,0},fill opacity=.075] (650.0,75) rectangle (920.0,187);
\node[anchor=center,text={rgb,255:red,35;green,52;blue,72},font=\fontsize{9.1}{11.0}\selectfont] at (785.0,117.5) {符号化与归因};
\node[anchor=center,text={rgb,255:red,35;green,52;blue,72},font=\fontsize{8.2}{9.9}\selectfont] at (785.0,144.5) {构建 ID / 内联 / 标签};
\node[anchor=center,text={rgb,255:red,35;green,52;blue,72},font=\fontsize{8.6}{10.4}\selectfont] at (480,235) {CO-RE 不会自动修复 Go 的用户态栈展开};
\node[anchor=center,text={rgb,255:red,35;green,52;blue,72},font=\fontsize{8.6}{10.4}\selectfont] at (480,264) {截断、缺失映射、原生代码和 G/M 切换都可能影响结果};
\end{tikzpicture}
\captionof{figure}{取得用户栈并不是读取几个地址。PC、寄存器、映射和运行时元数据决定可展开范围。FP、DWARF 及语言运行时展开策略不同；失败应进入未知归因。}
\end{minipage}\end{center}

\subsection{频率参数不是样本数量保证}

\texttt{perf\_\allowbreak{}event\_\allowbreak{}attr} 以同一字段表达采样 period 或 frequency，由 \texttt{freq} 位区分。按事件周期采样时，一个样本可以代表该周期内的一组事件；频率模式则尝试调整周期接近目标频率。内核会依据中断开销和频率反馈调整、节流，因此设置九十九赫兹不保证十秒恰好九百九十个有效栈。被排除的内核或用户态、目标线程未运行的时间、事件复用与处理失败都会影响数据覆盖。\cite{eb-perf-abi}\cite{eb-perf-core}


PMU 事件发生与采样记录中的指令位置也可能存在 skid；\texttt{precise\_\allowbreak{}ip} 表达相关要求和能力，不能把任意采样 IP 都当作触发硬件事件的精确指令。对复用计数器，\texttt{time\_\allowbreak{}enabled} 与 \texttt{time\_\allowbreak{}running} 有助于解释实际运行覆盖并缩放计数值；它们不能恢复在未运行期间从未采集的调用栈，也不能保证热路径在各阶段均匀出现。计数修正与分布修复必须分开。\cite{eb-perf-abi}


\subsection{收集地址、展开栈、符号化是三件事}

栈展开从采样寄存器与可读内存恢复调用地址序列；符号化再把地址映射为函数、源文件、行号和内联层次。用户态 ASLR、映射变化和二进制版本会影响符号化，因此保存进程映射与 build ID 很重要。\texttt{bpf\_\allowbreak{}get\_\allowbreak{}stack} 的 build-ID 模式还规定了文件偏移的含义：它相对于文件起点，不是相对于对象装载基址；映射到 ELF 符号时需按节地址与文件偏移调整。把所有偏移一律加基址会产生看似合理的错误符号。\cite{eb-stack-helper}


\texttt{bpf\_\allowbreak{}get\_\allowbreak{}stackid} 把地址序列放入 STACK\_\allowbreak{}TRACE map，并返回可作为聚合键的标识。它并不保证无限多不同栈都有唯一 ID。v6.12 实现对栈做哈希，映射到桶；默认会比较内容并在冲突等情况返回错误。\texttt{BPF\_\allowbreak{}F\_\allowbreak{}FAST\_\allowbreak{}STACK\_\allowbreak{}CMP} 允许仅按哈希比较，\texttt{BPF\_\allowbreak{}F\_\allowbreak{}REUSE\_\allowbreak{}STACKID} 允许冲突时丢弃旧栈。若计数 map 仍按旧 ID 累积，而栈表已被覆盖，后续可能把旧计数解释到新栈，因此窗口轮转和栈表生命周期属于正确性条件。\cite{eb-stack-helper}\cite{eb-stackmap}


固定 Linux x86 路径中，helper 经 \texttt{get\_\allowbreak{}perf\_\allowbreak{}callchain} 到架构的 \texttt{perf\_\allowbreak{}callchain\_\allowbreak{}user}，后者从采样寄存器的 BP 出发，读取下一帧与返回地址，直到深度上限或读取失败。此路径依赖可用的 frame pointer 链，不是通用 DWARF 解释器。内核代码还针对函数入口处 BP 尚未建立、以及 uretprobe trampoline 返回地址等情况做修正，这些细节说明“有地址就能完整展开”过于乐观。栈深度限制、不可访问内存与体系结构差异均应作为质量信息报告。\cite{eb-stackmap}\cite{eb-callchain}\cite{eb-x86-callchain}


\texttt{perf record -\allowbreak{}-\allowbreak{}call-\allowbreak{}graph dwarf} 走另一种策略：记录一定范围的用户栈字节，供后续依据展开信息处理，固定 v6.12 文档默认栈转储为八千一百九十二字节；内核栈使用其自身的展开机制。该版本 \texttt{perf record -\allowbreak{}-\allowbreak{}off-\allowbreak{}cpu} 文档明确其 BPF off-CPU 路径只用 frame pointer，但这不能推广为“所有 eBPF profiler 不可能支持 DWARF”。其他工具可以在用户态预处理展开规则、向 BPF 提供受限表或实现专用 unwinder，其能力与开销必须逐项目、逐架构核验。\cite{eb-perf-record}\cite{eb-x86-callchain}


\subsection{Go 的物理执行栈不等于请求历史}

Go 的 G 表示 goroutine，M 是 OS 线程，P 提供执行用户 Go 代码所需的调度和分配资源。内核采样首先命中 M 当前所执行的代码；一个等待 channel 的 G 没有因此占据某个正在运行的用户栈。不同 G 可先后在同一个 M 上运行，同一个 G 也可在不同 M 上恢复，所以 \texttt{tid} 既不是稳定 goroutine ID，也不是请求 ID。看到一个进程 CPU 栈，不能推出该进程全部 goroutine 的等待分布。\cite{eb-go-runtime}


Go 用户栈动态增长或收缩；runtime 的 \texttt{copystack} 会分配新栈、复制数据，调整调度上下文、defer、panic 和栈内指针，再更新 G 的栈边界。M 还拥有系统栈和信号栈，runtime 会在需要时切换过去执行。只按照某一时刻线程寄存器读取普通 FP 链的工具，必须面对这些边界；更专门的 Go unwinder 可以利用 runtime 元数据，但其布局与读取时机不能假定跨所有 Go 版本稳定。\cite{eb-go-runtime}\cite{eb-go-stack}


同样，普通物理栈不会自然包含“创建这个 goroutine 的另一个 goroutine 的整个调用链”，更不包含远端请求的上下游关系。cgo、汇编、优化与内联还会影响物理帧和显示帧的对应；对 Go 程序，应先用带已知调用结构的纯 Go、小型 cgo 和调度等待样本验证目标工具实际保留哪些帧，再评估它能否回答问题。若目标是 goroutine 调度和等待原因，runtime trace 与应用任务标记通常提供不同于 CPU 采样的证据维度；不能把一张火焰图解释为所有时间的总账。\cite{eb-go-runtime}\cite{eb-go-stack}\cite{eb-perf-record}


\section{端到端成本、丢数对账与因果边界}

\subsection{可复算的开销模型}

令每秒入口事件数为 \texttt{λ}，每个入口的挂载及过滤成本为 \texttt{c₀}；命中比例为 \texttt{q}，命中后读状态和聚合成本为 \texttt{c₁}；命中事件中执行栈收集的比例为 \texttt{s}、栈成本为 \texttt{c₂}；输出记录平均内核成本为 \texttt{c₃}、用户态消费成本为 \texttt{c₄}。在没有丢失且未加入额外共享竞争的简化模型中，CPU 秒每秒为 \texttt{λc₀+λq(c₁+sc₂+c₃+c₄)}，成本以秒计。这个式子描述预算，不声称事件相互独立，也不包含非线性缓存竞争、调度与上传费用。\cite{eb-map-array-code}\cite{eb-ring-code}\cite{eb-ring-consume}


教学取值：\texttt{λ=100000/\allowbreak{}s}，\texttt{q=0.\allowbreak{}2}，\texttt{s=0.\allowbreak{}25}；\texttt{c₀=80ns}，\texttt{c₁=120ns}，\texttt{c₂=2000ns}，\texttt{c₃=180ns}，\texttt{c₄=400ns}。入口成本为 \texttt{0.\allowbreak{}008} CPU 秒每秒，命中后为 \texttt{20000×1200ns=0.\allowbreak{}024}，合计 \texttt{0.\allowbreak{}032}，即一个核的百分之三点二。若在过滤前对全部事件展开栈，仅栈项就变成 \texttt{0.\allowbreak{}2} CPU 秒每秒，而原栈项为 \texttt{0.\allowbreak{}01}，增加零点一九个核。数字是虚构可复算示例，不能用于承诺特定机器上的开销。\cite{eb-stackmap}\cite{eb-ring-code}


生产评估还应分离编译加载成本、稳定运行成本、突发期成本、消费与符号化成本，以及被观测服务自身延迟变化。重复实验要固定工作负载、机器拓扑、事件率与保留策略，报告观察器开启前后的吞吐和延迟分布。低平均 CPU 不代表尾延迟没有变化；一个在极热同步路径触发的探针可能以很少总 CPU 改变关键队列。为了测量观察器成本而加入的观测同样可能扰动系统，应明确仪表层级与误差来源。\cite{eb-uprobe}\cite{eb-ring-code}\cite{eb-perf-core}


\subsection{队列积压与容量推导}

设实际输出每秒两万条、每条占六十四字节，则生产速率为 \texttt{1280000 B/\allowbreak{}s}；消费者每秒处理一万五千条，回收速率为 \texttt{960000 B/\allowbreak{}s}，积压按 \texttt{320000 B/\allowbreak{}s} 增长。取 \texttt{R=1048576}，初始已占 \texttt{262144} 字节，可用记录边界预算近似为 \texttt{R−1−262144}，达到满环约需 \texttt{2.\allowbreak{}4576s}。离散记录按六十四字节对齐，精确首次失败还由初始位置和到达时序决定。把 ring 加倍会延后失败，不会消除持续的五千条每秒吞吐缺口。\cite{eb-ring-code}\cite{eb-ring-consume}


ring reserve 失败、perf 的 \texttt{PERF\_\allowbreak{}RECORD\_\allowbreak{}LOST}、\texttt{PERF\_\allowbreak{}RECORD\_\allowbreak{}LOST\_\allowbreak{}SAMPLES}、栈 helper 返回错误、用户态解码失败和上传重试耗尽是不同阶段的质量事件。不能依赖 perf 的丢失通知自动覆盖自建 ring 的 reserve 失败；必须在程序中显式计数。失败计数本身也有数据类型、并发更新和读取窗口问题，例如全局非原子加一会把“丢失计数器”也写丢。质量指标需要与主指标同样认真地设计。\cite{eb-perf-abi}\cite{eb-ring-code}\cite{eb-map-array-doc}


\subsection{一份能够对齐的账本}

令 \texttt{M} 为通过目标过滤的事件；\texttt{F} 为 reserve 失败；\texttt{S} 为成功 submit；\texttt{D} 为 reserve 后主动 discard；\texttt{C} 为用户回调已接收记录；\texttt{Q} 为仍排队记录。若每个命中只尝试一次 reserve，且没有在途未决引用，生产端有 \texttt{M=F+S+D}；在无其他消费错误、窗口一致的条件下有 \texttt{S=C+Q}。例如 \texttt{M=100000,\allowbreak{}F=2000,\allowbreak{}D=0,\allowbreak{}S=98000,\allowbreak{}C=97900,\allowbreak{}Q=100} 自洽。停挂并排空后，才可期望 \texttt{Q=0} 并直接对齐提交与消费；运行中不可把尚未消费的一百条称为已经丢失。\cite{eb-ring-code}\cite{eb-ring-consume}


假设已消费的九万七千九百条中有三千条栈未知，这三千条仍是收到的事件，应留在 \texttt{unknown\_\allowbreak{}stack} 桶中，而不能再次加到 reserve 失败数。若为了画漂亮火焰图把未知栈删除，图中的分母会从全部事件变成可展开事件。报告至少给出事件覆盖率、栈成功率、输出失败率及明确的统计窗口；比较两个版本时还要比较这些质量指标，否则“热点占比下降”可能只是新版本更难展开。\cite{eb-stack-helper}\cite{eb-stackmap}\cite{eb-ring-consume}


丢数还可能引入选择偏差。设正常阶段与高负载阶段各发生一万条事件；正常阶段全保留，高负载阶段只保留两成。真实两个阶段各占一半，收到的样本却显示高负载阶段仅占 \texttt{2000/\allowbreak{}12000≈16.\allowbreak{}7\%}。把总数乘以统一补偿系数无法恢复两种阶段的正确比例。若丢失恰好发生在某个长栈、热锁或突发 I/O 路径，偏差会直接扭曲诊断。需要降低事件量、改进聚合或分阶段报告，不能只给火焰图增加“已归一化”的标签。\cite{eb-ring-code}\cite{eb-perf-core}


\subsection{相关性到因果结论仍有距离}

某次调度切出与某次磁盘完成时间接近，不足以证明这个线程正在等待那笔 I/O。证明需要可关联的任务或请求身份、正确的状态转换和配对规则，并处理线程复用、跨 CPU 迁移、事件漏失与重试。一次切出与再次切入之间的间隔含阻塞等待与可运行排队等不同组成；若只读两个时间戳，不能自动拆分。ring 的有序消费只维护自己的记录顺序，不生成缺失的关联边。\cite{eb-sched-trace}\cite{eb-sched-core}\cite{eb-ring-doc}


可用的结论层级应从弱到强：看到了某事件；在给定窗口和覆盖率下统计到某分布；多源证据支持某条等待链；改变被怀疑因素后，受控复测确认预期变化。eBPF 擅长补充系统边界处的证据，但因果定位仍依赖测量设计和干预验证。本章后续实验因此只承诺它实际建立的第一层、第二层事实，不借助工具名称暗示更多能力。\cite{eb-go-runtime}\cite{eb-sched-core}\cite{eb-perf-abi}


\section{完整实验与研究生练习}

\subsection{最小 C/libbpf 实验的契约}

配套文件为 \texttt{research/\allowbreak{}ebpf/\allowbreak{}example/\allowbreak{}switches.\allowbreak{}bpf.\allowbreak{}c}、\texttt{switches.\allowbreak{}c}、共享头和 README。入口为 \texttt{sched:sched\_\allowbreak{}switch}，通过当前 TGID 筛选目标进程，输出时间戳、TGID、TID 和 CPU，使用 per-CPU 计数器记录命中数与 reserve 失败数。Linux v6.12 的实际调用点位于 \texttt{context\_\allowbreak{}switch} 之前，因此当前任务对应即将切出的任务。程序不读取 tracepoint 上下文字段，也不读取内核私有结构，因此这个最小实验不演示 CO-RE 字段重定位。输入应为主机视角 TGID，不能直接套用容器 PID namespace 内的号码。\cite{eb-sched-core}\cite{eb-sched-trace}\cite{eb-stack-helper}\cite{eb-attach-api}


它测得的是挂载有效期间被命中的调度切出观察，不是 CPU 时间、完整 off-CPU 时长或 Go 请求延迟。没有读取 \texttt{prev\_\allowbreak{}state}，不能区分主动阻塞与被抢占；没有配对切入，不能计算等待区间。用户态先打开并加载对象、设置目标 TGID、创建消费者，之后才挂载，避免在配置未完成时产生无意义输出。运行十秒后解除链接、排空消费者并汇总所有 possible CPU 的计数，示范完整的配置、运行和收尾链路。\cite{eb-sched-trace}\cite{eb-sched-core}\cite{eb-attach-api}


环境要求是 Linux、具备 BPF 后端的 Clang、与 v1.5.0 API 相容的 libbpf 开发文件、libelf/zlib，以及允许目标 tracepoint 和 BPF/perf 挂载的本地策略。具体权限由目标环境决定，实验不建议关闭安全策略。当前编写环境只有 C 编译器与 pkg-config，未发现 Clang 与 libbpf 开发包；没有安装依赖、加载 BPF、修改内核参数或进行特权执行。因此验证状态为源码/API 审查与离线算术校验，不能标注为已编译或已在内核运行。\cite{eb-load-syscall}\cite{eb-attach-api}


\subsection{练习与评分要点}

\textbf{练习一：结构与语义。} 给出一个字段宽度不变、偏移改变、单位也改变的假想版本升级，说明 CO-RE 能解决哪一部分，再设计独立的语义兼容测试。评分看是否把“加载成功”“字段可读”“结果单位正确”分成三个条件，并给出已知输入与期望输出，而不是只检查内核版本字符串。\cite{eb-core-relo}


\textbf{练习二：验证状态。} 为一个先查询 map、再 reserve 事件、最后按条件退出的函数画控制流。找出空指针访问与未释放引用两种独立错误，给出每个出口的状态。增加循环后解释为什么字节码数量不变也可能增加验证成本。评分看是否使用类型、范围和资源义务，而不是背诵“BPF 不允许循环”。\cite{eb-verifier-doc}\cite{eb-verifier-code}\cite{eb-ring-doc}


\textbf{练习三：队列。} 将环容量改为八千一百九十二字节，payload 改为七十三字节，求单条占用与最大等长未消费记录数。应得 \texttt{align8(73+8)=88}，\texttt{floor(8191/\allowbreak{}88)=93}。再说明为什么真实运行可能在未达到九十三条时 reserve 失败：NMI 锁竞争、已有占用以及未决保留均须纳入条件。\cite{eb-ring-code}


\textbf{练习四：Go 诊断。} 某请求延迟两秒，CPU 样本只显示十毫秒解析工作。列出至少三种 CPU 采样没有充分覆盖的解释，并提出分别需要的证据。合格答案可以涉及 goroutine 等待、可运行排队、跨进程依赖与 cgo 阻塞，但不能仅凭这组数字认定其中任何一种；应使用 runtime 状态或可关联事件缩小假设。\cite{eb-go-runtime}\cite{eb-sched-trace}


\textbf{练习五：质量对账。} 给定命中、reserve 失败、submit、消费和未知栈计数，判断窗口是否已排空、图中分母是什么、是否重复计入损失。进一步构造“丢失集中在最慢阶段”的反例，证明统一补偿比例不能保证恢复热点分布。评分看是否写出账本恒等式的前提并保留未知类别。\cite{eb-ring-code}\cite{eb-ring-consume}\cite{eb-stackmap}


\textbf{研究任务。} 把最小实验扩展为线程 off-CPU 区间记录：显式处理切出状态、切入配对、任务退出、重复或缺失事件、TID 生命周期与容量限制；再解释它距离 goroutine 或请求等待分析还缺什么。提交物必须包含测量契约、状态图、失败账本、受控工作负载、版本矩阵与实际运行日志。仅给一张火焰图，不足以证明系统正确。\cite{eb-sched-trace}\cite{eb-sched-core}\cite{eb-go-runtime}


\chapter{pprof 的工作原理：从 Go 运行时测量到聚合报告}
\label{ch:19}

本章围绕一个问题展开：屏幕上的“180 ms、37.5\%”究竟怎样产生，它允许我们推出什么，又丢失了什么。源码基线为 Go 1.25.0，完整提交 \texttt{6e676ab2\allowbreak{}b809d466\allowbreak{}23acb598\allowbreak{}8248d95d\allowbreak{}1eb7939c}；分析器源码固定为 google/pprof 提交 \texttt{ebaad5f3\allowbreak{}1b4d25e6\allowbreak{}f51abf4d\allowbreak{}91b9611c\allowbreak{}ff92db7b}。本章的可执行教学验证使用本机 Go 1.22.12 内置 \texttt{go tool pprof}，两者明确区分。检索与验证日期为 2026-10-04；没有把合成样本冒充运行时实测。


\section{先建立测量契约：pprof 分析的是带语义的权重}

\Needspace{7.34cm}
\begin{center}\begin{minipage}{\linewidth}
\centering
\begin{tikzpicture}[x=0.016cm,y=-0.016cm]
\path[use as bounding box] (0,0) rectangle (960,340);
\draw[rounded corners=5pt,draw={rgb,255:red,25;green,118;blue,210},fill={rgb,255:red,25;green,118;blue,210},fill opacity=.075] (40.0,75) rectangle (310.0,187);
\node[anchor=center,text={rgb,255:red,35;green,52;blue,72},font=\fontsize{9.1}{11.0}\selectfont] at (175.0,117.5) {底层观察};
\node[anchor=center,text={rgb,255:red,35;green,52;blue,72},font=\fontsize{8.2}{9.9}\selectfont] at (175.0,144.5) {时间戳 / PCs / 标签};
\draw[draw={rgb,255:red,99;green,115;blue,134},-{Stealth[length=4pt]}] (313.0,130) -- (341.0,130);
\draw[rounded corners=5pt,draw={rgb,255:red,129;green,80;blue,181},fill={rgb,255:red,129;green,80;blue,181},fill opacity=.075] (345.0,75) rectangle (615.0,187);
\node[anchor=center,text={rgb,255:red,35;green,52;blue,72},font=\fontsize{9.1}{11.0}\selectfont] at (480.0,117.5) {profile.proto};
\node[anchor=center,text={rgb,255:red,35;green,52;blue,72},font=\fontsize{8.2}{9.9}\selectfont] at (480.0,144.5) {同栈同标签累加权重};
\draw[draw={rgb,255:red,99;green,115;blue,134},-{Stealth[length=4pt]}] (618.0,130) -- (646.0,130);
\draw[rounded corners=5pt,draw={rgb,255:red,88;green,167;blue,0},fill={rgb,255:red,88;green,167;blue,0},fill opacity=.075] (650.0,75) rectangle (920.0,187);
\node[anchor=center,text={rgb,255:red,35;green,52;blue,72},font=\fontsize{9.1}{11.0}\selectfont] at (785.0,117.5) {分析报告};
\node[anchor=center,text={rgb,255:red,35;green,52;blue,72},font=\fontsize{8.2}{9.9}\selectfont] at (785.0,144.5) {flat / cum / 火焰图};
\node[anchor=center,text={rgb,255:red,35;green,52;blue,72},font=\fontsize{8.6}{10.4}\selectfont] at (480,235) {逐事件先后在聚合时丢失};
\node[anchor=center,text={rgb,255:red,35;green,52;blue,72},font=\fontsize{8.6}{10.4}\selectfont] at (480,264) {重新排列已聚合的 Sample，不能恢复请求时间线};
\end{tikzpicture}
\captionof{figure}{时间记录如何压缩为栈权重。Go CPU builder 忽略逐事件时间戳，按栈及标签聚合。编码字段只能保留被写入的信息。}
\end{minipage}\end{center}

pprof 并不是 Go 专属采集器。google/pprof 项目自述的职责是读取 profiling samples，解析 profile.proto，再生成文本、调用图和火焰图等报告。Go 的 runtime/pprof 是生产者之一；其他语言运行时、原生采样器或转换程序也可以生产这种格式。一个程序用 Go 编写，并不意味着它读入的数据来自 Go；一份文件扩展名是 \texttt{.\allowbreak{}pprof}，也不意味着其中记录的是 CPU 时间。\cite{pp-scope}\cite{pp-proto}


应把链路分成四层。测量层决定在哪里安装观察点、何时观察以及抽样概率；编码层决定将观察到的栈、标签和估计值怎样表示；分析层进行身份归一、合并、筛选和粒度转换；展示层把指定数值维度映射成行、边与矩形宽度。这四层的边界很重要：把图形画得更精细不会恢复采集时漏掉的事件；支持某种标签字段也不表示生产者实际采集了标签。\cite{pp-proto}\cite{pp-cpu-builder}


阅读任何 profile 前，至少写清七个要素：被观察的对象集合、数值代表的物理量、抽样规则、栈的归因位置、计时来源、时间窗口，以及已知覆盖缺口。例如 CPU 的对象是实际执行的线程栈，heap 的对象是被抽中的堆分配块，block 观察特定同步操作中的阻塞，mutex 把竞争成本归于结束临界区的释放路径。它们都能生成“热点函数”，但热点所回答的问题完全不同。\cite{pp-cpu-runtime}\cite{pp-heap-scale}\cite{pp-mprof}\cite{pp-sema}


在数学上，一份 profile 可写成记录集合 \(P=\{(L_i,V_i,K_i)\}\)。\(L_i\) 是调用栈位置序列，\(V_i\) 是定长有符号整数向量，\(K_i\) 是标签。选定第 \(j\) 个 sample type 后，分析器使用权重 \(w_i=V_{i,j}\)。两个可合并记录的值向量逐分量相加，而非先把各列换算成某个共同的“耗时”。对象数与字节数不能相加，CPU 时间与同步等待也不能相加。\cite{pp-proto}


这里“样本”有两种容易混淆的用法：一次底层采样事件，以及一个 protobuf Sample 消息。后者可能代表一条事件、数千条已聚合事件，或经过逆概率校正的估计量。数值第一列叫 count，也不保证它就是实际读取了多少个原始事件。区分“消息条数、事件计数、校正后权重”，是理解随后所有报告的前提。\cite{pp-cpu-builder}\cite{pp-heap-scale}\cite{pp-mprof}


普通 profile 也不是调用时间线。协议没有给每个 Sample 规定开始、结束和持续时间三个字段；Profile.time\_\allowbreak{}nanos 和 duration\_\allowbreak{}nanos 是文件级信息。尤其 Go 的 CPU 构建器在读取底层记录时，明确把事件时间戳标为 ignored，然后按栈与标签聚合。因此可以回答“这个窗口内哪些执行路径消耗较多”，不能回答“第几个请求先执行哪条路径”或“这一段 GC 是否先于那次超时”。这些问题需要另有时序证据。\cite{pp-proto}\cite{pp-cpu-builder}


\section{CPU producer：信号、受限栈展开、环缓冲和信息损失}

\Needspace{7.34cm}
\begin{center}\begin{minipage}{\linewidth}
\centering
\begin{tikzpicture}[x=0.016cm,y=-0.016cm]
\path[use as bounding box] (0,0) rectangle (960,340);
\draw[rounded corners=5pt,draw={rgb,255:red,25;green,118;blue,210},fill={rgb,255:red,25;green,118;blue,210},fill opacity=.075] (40.0,75) rectangle (310.0,187);
\node[anchor=center,text={rgb,255:red,35;green,52;blue,72},font=\fontsize{9.1}{11.0}\selectfont] at (175.0,117.5) {CPU 时钟 → SIGPROF};
\node[anchor=center,text={rgb,255:red,35;green,52;blue,72},font=\fontsize{8.2}{9.9}\selectfont] at (175.0,144.5) {名义 100 Hz};
\draw[draw={rgb,255:red,99;green,115;blue,134},-{Stealth[length=4pt]}] (313.0,130) -- (341.0,130);
\draw[rounded corners=5pt,draw={rgb,255:red,129;green,80;blue,181},fill={rgb,255:red,129;green,80;blue,181},fill opacity=.075] (345.0,75) rectangle (615.0,187);
\node[anchor=center,text={rgb,255:red,35;green,52;blue,72},font=\fontsize{9.1}{11.0}\selectfont] at (480.0,117.5) {受限现场展开栈};
\node[anchor=center,text={rgb,255:red,35;green,52;blue,72},font=\fontsize{8.2}{9.9}\selectfont] at (480.0,144.5) {有界日志 / 溢出计数};
\draw[draw={rgb,255:red,99;green,115;blue,134},-{Stealth[length=4pt]}] (618.0,130) -- (646.0,130);
\draw[rounded corners=5pt,draw={rgb,255:red,88;green,167;blue,0},fill={rgb,255:red,88;green,167;blue,0},fill opacity=.075] (650.0,75) rectangle (920.0,187);
\node[anchor=center,text={rgb,255:red,35;green,52;blue,72},font=\fontsize{9.1}{11.0}\selectfont] at (785.0,117.5) {普通 goroutine};
\node[anchor=center,text={rgb,255:red,35;green,52;blue,72},font=\fontsize{8.2}{9.9}\selectfont] at (785.0,144.5) {count × period};
\node[anchor=center,text={rgb,255:red,35;green,52;blue,72},font=\fontsize{8.6}{10.4}\selectfont] at (480,235) {睡眠时间不会自动成为 CPU 样本};
\node[anchor=center,text={rgb,255:red,35;green,52;blue,72},font=\fontsize{8.6}{10.4}\selectfont] at (480,264) {已知 lost 数量可保留部分总量，却无法恢复原业务栈};
\end{tikzpicture}
\captionof{figure}{Go CPU 采样的两种执行环境。Go 1.25 Linux 路径的简化图：线程 CPU 时钟触发信号，受限信号处理器写有界日志，普通 goroutine 异步构建 profile。}
\end{minipage}\end{center}

Go 1.25 的 \texttt{runtime/\allowbreak{}pprof.\allowbreak{}StartCPUProfile} 固定使用名义 100 Hz。这个频率不是 HTTP \texttt{seconds} 参数决定的，也不是 \texttt{go tool pprof} 在读文件时选择的。\texttt{seconds} 改变采集窗口；读取端选择 \texttt{sample\_\allowbreak{}index} 改变展示哪一列；两者都不重写运行时当时的采样规律。底层 \texttt{runtime.\allowbreak{}SetCPUProfileRate} 另有生命周期约束，活动采集期间不能直接改频率，不应和高层 API 随意混用。\cite{pp-cpu-api}\cite{pp-cpu-runtime}


“100 Hz”必须结合 CPU 时钟理解。以 Go 1.25 的 Linux 路径为例，运行时为线程创建 \texttt{CLOCK\_\allowbreak{}THREAD\_\allowbreak{}CPUTIME\_\allowbreak{}ID} 计时器，通过 \texttt{SIGEV\_\allowbreak{}THREAD\_\allowbreak{}ID} 把 SIGPROF 发到目标线程。线程累计执行约一个周期才产生相应采样机会，线程睡眠的十毫秒不会自动变成一个 CPU 样本。创建线程计时器失败时，可以退到进程级计时机制；\texttt{validSIGPROF} 根据来源和线程状态避免同一 CPU 消耗被两种机制重复记账。不能将这一 Linux 实现推广为所有操作系统的实现方式。\cite{pp-cpu-linux}


线程计时器的首次到期还做了随机化：初始延迟取一个周期内的随机值，随后使用固定周期。若新线程只执行了十分之一个周期，理想期望应当是约 0.1 个样本，而不是永远零个。这尤其关系到短暂活跃的 GC 线程。这个设计减少特定偏差，并不使全部样本成为严格独立同分布；周期性工作、信号合并、调度和机器环境仍会影响观察。\cite{pp-cpu-linux}


一次信号到达后，\texttt{sigprof} 取得被中断现场的 PC、SP 等信息，通过 runtime unwinder 恢复 Go 栈。该路径可能在 GC 或其他受限现场中执行，因此不能普通地分配内存，也不能任意获得被中断代码可能持有的锁。源码使用有界临时栈；本版本 CPU 栈上限为 64 个记录槽。cgo、VDSO、Windows libcall 等还有专门分支。展开失败时，样本可能落到 \texttt{\_\allowbreak{}System}、\texttt{\_\allowbreak{}ExternalCode} 或其他抽象位置，不能把未知部分随意分摊给相邻业务函数。\cite{pp-sigprof}\cite{pp-cpu-runtime}


采样端和编码端没有在信号处理器里同步完成全部 protobuf 工作。采样端把头部、栈及标签指针写入 profBuf，由普通 goroutine 异步读取。profBuf 的数据环和标签环并行推进：一个事件的头部包含长度与时间等信息，栈占可变数量的槽，标签另占一个槽。该缓冲组件按一个读者和一个写者设计无锁协议，但 CPU 入口还用 signalLock 序列化相关写入，因此“内部使用无锁环”不等于整条采样链路完全没有同步。\cite{pp-buffer}\cite{pp-cpu-runtime}


有界缓冲会溢出。Go 1.25 的 CPU 数据环配置为 \(2^{17}\) 个 64 位字，约 1 MiB；在最大栈长度下，源码估算约能容纳 1900 个样本。标签环的容量是另一道限制。如果编码端赶不上生产端，后续写入不会无限增加内存，而是被丢弃，并累计 overflow 数。溢出记录保留丢弃数量及首次丢弃时刻，已丢失的栈和标签则无法恢复；读者或写者都可以把 pending overflow 转成正式记录，序列位用于避免清除状态时的 ABA 问题。\cite{pp-cpu-runtime}\cite{pp-buffer}


编码器遇到这种溢出记录，会构造 \texttt{runtime/\allowbreak{}pprof.\allowbreak{}lostProfileEvent} 对应的特殊栈，并把丢弃数量作为 count。因此部分已知丢样仍可能进入总权重，只是归因缺失。外部线程的暂存区满还有 \texttt{\_\allowbreak{}LostExternalCode} 等路径。应把“总量可能得到部分保全”与“业务归因完整”分开判断；反过来，报告里没有特殊 lost 节点，也不能证明操作系统投递、原生栈展开等所有环节都零损失。\cite{pp-cpu-builder}\cite{pp-cpu-runtime}


真正写文件时，\texttt{addCPUData} 先从首条记录得到周期 \(\tau=10^9/hz\) 纳秒，再以栈和 tag 为查找条件累加 count。\texttt{build} 写出的 CPU 向量是 \texttt{[count,\allowbreak{}count×τ]}，类型分别为 samples/count 和 cpu/nanoseconds。默认 100 Hz 下，一份聚合记录 \texttt{[18,\allowbreak{}180000000]} 表示十八份名义十毫秒权重，而不是某个函数恰好被调用十八次，也不是一次调用精确持续 180 ms。\cite{pp-cpu-builder}


多核下，总 CPU 权重可大于墙钟窗口。若原始采集窗口为 \(T\)，总权重估计为 \(C\)，则 \(C/T\) 可近似表示平均占用的 CPU 核数。例如本章教学文件写有 400 ms 窗口、480 ms CPU 权重，读取端显示 120\%；这表示约 1.2 个核的容量尺度，不表示单个请求用时超过自身的百分之百。此解释只适用于口径适当的原始 CPU 窗口；合并、差分后文件的 duration 可能只是头部运算结果，不能继续机械套用。\cite{pp-cpu-builder}\cite{pp-merge}\cite{pp-teaching}


对热点比例的误差分析，也应保持相同谨慎。在理想独立抽样下，由二项模型可推得 \(\hat p=k/n\) 的标准误差近似 \(\sqrt{p(1-p)/n}\)；这是理想模型推导，runtime 采样并不承诺独立性。只观察一次“10\% 降到 8\%”，不能从这两个百分数直接推出优化成立。应同时检验样本量、窗口稳定性、未归因比例及单位业务工作量的 CPU，并用关闭诊断采集的重复实验评价实际收益。\cite{pp-cpu-linux}\cite{pp-cpu-runtime}


\section{Heap producer：为什么一个小对象会代表很多对象}

\Needspace{8.14cm}
\begin{center}\begin{minipage}{\linewidth}
\centering
\begin{tikzpicture}[x=0.016cm,y=-0.016cm]
\path[use as bounding box] (0,0) rectangle (960,390);
\draw[draw={rgb,255:red,99;green,115;blue,134},-{Stealth[length=4pt]}] (90,300) -- (610,300);
\draw[draw={rgb,255:red,99;green,115;blue,134},-{Stealth[length=4pt]}] (90,300) -- (90,35);
\node[anchor=center,text={rgb,255:red,35;green,52;blue,72},font=\fontsize{9.6}{11.6}\selectfont] at (110,23) {q};
\node[anchor=center,text={rgb,255:red,35;green,52;blue,72},font=\fontsize{8.6}{10.4}\selectfont] at (350,343) {分配大小 s / R};
\node[anchor=center,text={rgb,255:red,35;green,52;blue,72},font=\fontsize{7.7}{9.3}\selectfont] at (90.0,323) {0};
\node[anchor=center,text={rgb,255:red,35;green,52;blue,72},font=\fontsize{7.7}{9.3}\selectfont] at (250.0,323) {1};
\node[anchor=center,text={rgb,255:red,35;green,52;blue,72},font=\fontsize{7.7}{9.3}\selectfont] at (410.0,323) {2};
\node[anchor=center,text={rgb,255:red,35;green,52;blue,72},font=\fontsize{7.7}{9.3}\selectfont] at (570.0,323) {3};
\node[anchor=center,text={rgb,255:red,35;green,52;blue,72},font=\fontsize{7.7}{9.3}\selectfont] at (58,300) {0.0};
\node[anchor=center,text={rgb,255:red,35;green,52;blue,72},font=\fontsize{7.7}{9.3}\selectfont] at (58,175) {0.5};
\node[anchor=center,text={rgb,255:red,35;green,52;blue,72},font=\fontsize{7.7}{9.3}\selectfont] at (58,50) {1.0};
\draw[draw={rgb,255:red,88;green,167;blue,0}] (90.0,300.0) -- (98.0,287.8073561251785);
\draw[draw={rgb,255:red,88;green,167;blue,0}] (98.0,287.8073561251785) -- (106.0,276.20935450898986);
\draw[draw={rgb,255:red,88;green,167;blue,0}] (106.0,276.20935450898986) -- (114.0,265.17699410626443);
\draw[draw={rgb,255:red,88;green,167;blue,0}] (114.0,265.17699410626443) -- (122.0,254.68268826949546);
\draw[draw={rgb,255:red,88;green,167;blue,0}] (122.0,254.68268826949546) -- (130.0,244.70019576785123);
\draw[draw={rgb,255:red,88;green,167;blue,0}] (130.0,244.70019576785123) -- (138.0,235.20455517042947);
\draw[draw={rgb,255:red,88;green,167;blue,0}] (138.0,235.20455517042947) -- (146.0,226.17202242967835);
\draw[draw={rgb,255:red,88;green,167;blue,0}] (146.0,226.17202242967835) -- (154.0,217.58001150890982);
\draw[draw={rgb,255:red,88;green,167;blue,0}] (154.0,217.58001150890982) -- (162.0,209.40703790544333);
\draw[draw={rgb,255:red,88;green,167;blue,0}] (162.0,209.40703790544333) -- (170.0,201.63266492815836);
\draw[draw={rgb,255:red,88;green,167;blue,0}] (170.0,201.63266492815836) -- (178.0,194.23745259512168);
\draw[draw={rgb,255:red,88;green,167;blue,0}] (178.0,194.23745259512168) -- (186.0,187.2029090235066);
\draw[draw={rgb,255:red,88;green,167;blue,0}] (186.0,187.2029090235066) -- (194.0,180.511444190254);
\draw[draw={rgb,255:red,88;green,167;blue,0}] (194.0,180.511444190254) -- (202.0,174.14632594785238);
\draw[draw={rgb,255:red,88;green,167;blue,0}] (202.0,174.14632594785238) -- (210.0,168.0916381852537);
\draw[draw={rgb,255:red,88;green,167;blue,0}] (210.0,168.0916381852537) -- (218.0,162.33224102930538);
\draw[draw={rgb,255:red,88;green,167;blue,0}] (218.0,162.33224102930538) -- (226.0,156.85373298718167);
\draw[draw={rgb,255:red,88;green,167;blue,0}] (226.0,156.85373298718167) -- (234.0,151.6424149351498);
\draw[draw={rgb,255:red,88;green,167;blue,0}] (234.0,151.6424149351498) -- (242.0,146.68525586362531);
\draw[draw={rgb,255:red,88;green,167;blue,0}] (242.0,146.68525586362531) -- (250.0,141.96986029286057);
\draw[draw={rgb,255:red,88;green,167;blue,0}] (250.0,141.96986029286057) -- (258.0,137.48443727778883);
\draw[draw={rgb,255:red,88;green,167;blue,0}] (258.0,137.48443727778883) -- (266.0,133.21777092451987);
\draw[draw={rgb,255:red,88;green,167;blue,0}] (266.0,133.21777092451987) -- (274.0,129.1591923447633);
\draw[draw={rgb,255:red,88;green,167;blue,0}] (274.0,129.1591923447633) -- (282.0,125.29855297805054);
\draw[draw={rgb,255:red,88;green,167;blue,0}] (282.0,125.29855297805054) -- (290.0,121.62619921504754);
\draw[draw={rgb,255:red,88;green,167;blue,0}] (290.0,121.62619921504754) -- (298.0,118.13294825850312);
\draw[draw={rgb,255:red,88;green,167;blue,0}] (298.0,118.13294825850312) -- (306.0,114.81006516147286);
\draw[draw={rgb,255:red,88;green,167;blue,0}] (306.0,114.81006516147286) -- (314.0,111.6492409854016);
\draw[draw={rgb,255:red,88;green,167;blue,0}] (314.0,111.6492409854016) -- (322.0,108.64257202344942);
\draw[draw={rgb,255:red,88;green,167;blue,0}] (322.0,108.64257202344942) -- (330.0,105.78254003710745);
\draw[draw={rgb,255:red,88;green,167;blue,0}] (330.0,105.78254003710745) -- (338.0,103.06199345668577);
\draw[draw={rgb,255:red,88;green,167;blue,0}] (338.0,103.06199345668577) -- (346.0,100.47412949866384);
\draw[draw={rgb,255:red,88;green,167;blue,0}] (346.0,100.47412949866384) -- (354.0,98.01247715518852);
\draw[draw={rgb,255:red,88;green,167;blue,0}] (354.0,98.01247715518852) -- (362.0,95.67088101318367);
\draw[draw={rgb,255:red,88;green,167;blue,0}] (362.0,95.67088101318367) -- (370.0,93.44348586261128);
\draw[draw={rgb,255:red,88;green,167;blue,0}] (370.0,93.44348586261128) -- (378.0,91.32472205539665);
\draw[draw={rgb,255:red,88;green,167;blue,0}] (378.0,91.32472205539665) -- (386.0,89.30929157840691);
\draw[draw={rgb,255:red,88;green,167;blue,0}] (386.0,89.30929157840691) -- (394.0,87.39215480565878);
\draw[draw={rgb,255:red,88;green,167;blue,0}] (394.0,87.39215480565878) -- (402.0,85.5685178966284);
\draw[draw={rgb,255:red,88;green,167;blue,0}] (402.0,85.5685178966284) -- (410.0,83.83382080915317);
\draw[draw={rgb,255:red,88;green,167;blue,0}] (410.0,83.83382080915317) -- (418.0,82.18372589695107);
\draw[draw={rgb,255:red,88;green,167;blue,0}] (418.0,82.18372589695107) -- (426.0,80.61410706324546);
\draw[draw={rgb,255:red,88;green,167;blue,0}] (426.0,80.61410706324546) -- (434.0,79.12103944337423);
\draw[draw={rgb,255:red,88;green,167;blue,0}] (434.0,79.12103944337423) -- (442.0,77.70078959058347);
\draw[draw={rgb,255:red,88;green,167;blue,0}] (442.0,77.70078959058347) -- (450.0,76.34980614046609);
\draw[draw={rgb,255:red,88;green,167;blue,0}] (450.0,76.34980614046609) -- (458.0,75.06471093070095);
\draw[draw={rgb,255:red,88;green,167;blue,0}] (458.0,75.06471093070095) -- (466.0,73.84229055388741);
\draw[draw={rgb,255:red,88;green,167;blue,0}] (466.0,73.84229055388741) -- (474.0,72.67948832235311);
\draw[draw={rgb,255:red,88;green,167;blue,0}] (474.0,72.67948832235311) -- (482.0,71.57339662484262);
\draw[draw={rgb,255:red,88;green,167;blue,0}] (482.0,71.57339662484262) -- (490.0,70.52124965597471);
\draw[draw={rgb,255:red,88;green,167;blue,0}] (490.0,70.52124965597471) -- (498.0,69.52041650028829);
\draw[draw={rgb,255:red,88;green,167;blue,0}] (498.0,69.52041650028829) -- (506.0,68.56839455358346);
\draw[draw={rgb,255:red,88;green,167;blue,0}] (506.0,68.56839455358346) -- (514.0,67.66280326510739);
\draw[draw={rgb,255:red,88;green,167;blue,0}] (514.0,67.66280326510739) -- (522.0,66.80137818493742);
\draw[draw={rgb,255:red,88;green,167;blue,0}] (522.0,66.80137818493742) -- (530.0,65.9819653016769);
\draw[draw={rgb,255:red,88;green,167;blue,0}] (530.0,65.9819653016769) -- (538.0,65.20251565630448);
\draw[draw={rgb,255:red,88;green,167;blue,0}] (538.0,65.20251565630448) -- (546.0,64.46108021870961);
\draw[draw={rgb,255:red,88;green,167;blue,0}] (546.0,64.46108021870961) -- (554.0,63.75580501410181);
\draw[draw={rgb,255:red,88;green,167;blue,0}] (554.0,63.75580501410181) -- (562.0,63.084926487108106);
\draw[draw={rgb,255:red,88;green,167;blue,0}] (562.0,63.084926487108106) -- (570.0,62.446767091965995);
\draw[rounded corners=5pt,draw={rgb,255:red,25;green,118;blue,210},fill={rgb,255:red,25;green,118;blue,210},fill opacity=.075] (655,65) rectangle (930,165);
\node[anchor=center,text={rgb,255:red,35;green,52;blue,72},font=\fontsize{9.1}{11.0}\selectfont] at (792.5,101.5) {64 KiB：q ≈ 0.1175};
\node[anchor=center,text={rgb,255:red,35;green,52;blue,72},font=\fontsize{8.2}{9.9}\selectfont] at (792.5,128.5) {权重倍数 ≈ 8.51};
\draw[rounded corners=5pt,draw={rgb,255:red,129;green,80;blue,181},fill={rgb,255:red,129;green,80;blue,181},fill opacity=.075] (655,200) rectangle (930,300);
\node[anchor=center,text={rgb,255:red,35;green,52;blue,72},font=\fontsize{9.1}{11.0}\selectfont] at (792.5,236.5) {512 KiB：q ≈ 0.6321};
\node[anchor=center,text={rgb,255:red,35;green,52;blue,72},font=\fontsize{8.2}{9.9}\selectfont] at (792.5,263.5) {权重倍数 ≈ 1.58};
\end{tikzpicture}
\captionof{figure}{按字节采样的包含概率。理想 Poisson 模型 q(s)=1-exp(-s/R)，R=512KiB。小对象更难入样，命中后携带更大的逆概率权重；实际实现有近似。}
\end{minipage}\end{center}

内存剖析观察的是分配，不是在固定时刻扫描整个堆并逐对象列账。Go 的默认 \texttt{Mem\allowbreak{}Profile\allowbreak{}Rate} 为 512 KiB；运行时在分配字节轴上安排随机采样点。\texttt{nextSample} 希望相邻采样点间距服从均值为 \(R\) 的指数分布，这对应理想 Poisson 过程。实际 \texttt{fastexprand} 使用有限位随机数、快速对数和整数结果，不能把数学理想化描述当成实现没有近似。\cite{pp-heap-poisson}\cite{pp-mprof}


尺寸为 \(s\) 的分配至少覆盖一个采样点的概率为 \(q(s)=1-e^{-s/R}\)。推导很直接：长度 \(s\) 的区间没有 Poisson 点的概率为 \(e^{-s/R}\)，取补集就是至少一次命中。抽样是按分配字节轴进行，所以大块分配比小块分配更容易出现；“每五十二万个对象抽一个”是错误解释。一次大对象跨过多个理想采样点，也不要求向 profile 写出同一个对象的多份原始记录；校正关心其被纳入的概率。\cite{pp-heap-poisson}\cite{pp-heap-scale}


令第 \(i\) 个分配是否命中为 \(I_i\)。理想模型中，对象数可用 \(\widehat N=\sum_i I_i/q(s_i)\) 估计，分配字节可用 \(\widehat B=\sum_i s_iI_i/q(s_i)\) 估计。逆概率权重解释了为什么文件中一个记录的 count 可以远大于实际采到的事件数。这些估计在模型条件下具有相应统计性质，单次观测仍会有方差，尤其稀少的小对象样本可能携带很大的权重。\cite{pp-heap-scale}


举例取 \(R=512\) KiB。一个 64 KiB 分配的包含概率约为 0.1175，权重倍数约为 8.51；一个 512 KiB 分配的包含概率约为 0.6321，倍数约为 1.58。不能对两种尺寸都简单乘八，也不能拿 sample 消息数量比较它们的真实对象数量。如果一条很少出现的小对象路径在一次采集中碰巧被命中，外推值可能明显跳动；应通过重复观测和更合适的采样预算评估，而不是凭一张图宣布泄漏。\cite{pp-heap-scale}


Go 1.25 的具体实现是 \texttt{scaleHeapSample(count,\allowbreak{}size,\allowbreak{}rate)}。当 count 或 size 为零时返回零；rate 不大于一时不做外推；否则以 \texttt{avgSize=size/\allowbreak{}count} 计算 \texttt{1/\allowbreak{}(1-\allowbreak{}exp(-\allowbreak{}avgSize/\allowbreak{}rate))}，分别乘对象数和字节数，最后转换为 int64。运行时内存桶的身份包含栈和分配尺寸，不能将这里的平均尺寸规则理解为“任意混合大小的对象先聚合，再统一校正也完全等价”。将已经校正的 profile 再乘一次采样率，会重复放大数据。\cite{pp-heap-scale}\cite{pp-mprof}


写出时的四维向量依次为 alloc\_\allowbreak{}objects、alloc\_\allowbreak{}space、inuse\_\allowbreak{}objects、inuse\_\allowbreak{}space。前两列描述累计分配，后两列描述采样可见的尚未释放部分。\texttt{heap} 和 \texttt{allocs} 主要改变默认选择的维度，底层并不是两套互不相干的分配记录。“累计已分配两吉字节”完全可能来自一直复用生命周期很短的临时对象，而“仍在使用两百兆字节”也不直接告诉我们哪个容器正在持有它们。栈是分配来源，不是完整引用保留图。\cite{pp-heap-scale}\cite{pp-profile-kinds}


内存账目还存在发布时间。分配发生在实时路径，释放则跟随 GC 的清扫；若立刻把新分配与尚未登记的释放比较，就会系统性偏向“仍然存活”。mprof 通过 active/future 周期桶协调分配与释放，mark termination 推进周期，后续 flush 和 post-sweep 发布可见状态。\texttt{runtime.\allowbreak{}MemProfile} 明确结果可能落后最多两个 GC 周期。因此 profile 的 inuse 既不是读取瞬间的精确活对象清单，也不是操作系统的 RSS。\cite{pp-mprof}


\texttt{MemProfileRate=1} 的含义是记录采集器覆盖的每个堆分配块，不是记录所有语言级对象：栈上对象不在此范围，tiny allocator 可把多个小对象装入一个块，分配槽还可能包含内部碎片。采样率也应在进程早期至多调整一次，分析工具假设整段记录具有恒定采样率。把一个进程先设成全量采样、再恢复默认后混出一份文件，无法靠读取端知道每个旧记录应当使用哪个历史参数。\cite{pp-mprof}\cite{pp-heap-coverage}


正确的泄漏分析因此需要一个更强的链条：相同稳态工作量下，GC 对齐后仍有持续增长；profile 找到增长的分配来源；代码或对象检查解释持有路径；切断该路径后，重复批次的存活集重新有界。仅有高 alloc\_\allowbreak{}space 说明流量，单次高 inuse\_\allowbreak{}space 说明当时可见保留，两者都没有独立完成因果证明。\cite{pp-profile-kinds}\cite{pp-mprof}


\section{同步和 goroutine producer：相似的栈，完全不同的账}

\Needspace{7.34cm}
\begin{center}\begin{minipage}{\linewidth}
\centering
\begin{tikzpicture}[x=0.016cm,y=-0.016cm]
\path[use as bounding box] (0,0) rectangle (960,340);
\draw[rounded corners=5pt,draw={rgb,255:red,25;green,118;blue,210},fill={rgb,255:red,25;green,118;blue,210},fill opacity=.075] (40.0,75) rectangle (310.0,187);
\node[anchor=center,text={rgb,255:red,35;green,52;blue,72},font=\fontsize{9.1}{11.0}\selectfont] at (175.0,117.5) {等待者};
\node[anchor=center,text={rgb,255:red,35;green,52;blue,72},font=\fontsize{8.2}{9.9}\selectfont] at (175.0,144.5) {block：阻塞位置};
\draw[draw={rgb,255:red,99;green,115;blue,134},-{Stealth[length=4pt]}] (313.0,130) -- (341.0,130);
\draw[rounded corners=5pt,draw={rgb,255:red,129;green,80;blue,181},fill={rgb,255:red,129;green,80;blue,181},fill opacity=.075] (345.0,75) rectangle (615.0,187);
\node[anchor=center,text={rgb,255:red,35;green,52;blue,72},font=\fontsize{9.1}{11.0}\selectfont] at (480.0,117.5) {锁持有者};
\node[anchor=center,text={rgb,255:red,35;green,52;blue,72},font=\fontsize{8.2}{9.9}\selectfont] at (480.0,144.5) {mutex：释放路径};
\draw[draw={rgb,255:red,99;green,115;blue,134},-{Stealth[length=4pt]}] (618.0,130) -- (646.0,130);
\draw[rounded corners=5pt,draw={rgb,255:red,88;green,167;blue,0},fill={rgb,255:red,88;green,167;blue,0},fill opacity=.075] (650.0,75) rectangle (920.0,187);
\node[anchor=center,text={rgb,255:red,35;green,52;blue,72},font=\fontsize{9.1}{11.0}\selectfont] at (785.0,117.5) {调度状态};
\node[anchor=center,text={rgb,255:red,35;green,52;blue,72},font=\fontsize{8.2}{9.9}\selectfont] at (785.0,144.5) {runnable → running};
\node[anchor=center,text={rgb,255:red,35;green,52;blue,72},font=\fontsize{8.6}{10.4}\selectfont] at (480,235) {等待权重是并发需求的累计，不是一次持锁时长};
\node[anchor=center,text={rgb,255:red,35;green,52;blue,72},font=\fontsize{8.6}{10.4}\selectfont] at (480,264) {事件完成记账的窗口，不一定等于等待区间与窗口的交集};
\end{tikzpicture}
\captionof{figure}{相似的等待，不同的归因。block 将特定同步阻塞归到等待位置；mutex 通常归到释放路径。二者会描述同一争用的不同视角，不能直接相加。}
\end{minipage}\end{center}

block profile 的观察点嵌入特定同步原语，包括 channel、Mutex/RWMutex、WaitGroup 和 Cond 的相关路径。它将等待归到阻塞位置；网络轮询、任意系统调用、睡眠和调度队列不能全部期待在这里出现。\texttt{SetBlockProfileRate} 的单位是纳秒，其目标是按阻塞时间抽样：对短于阈值的事件按时长提高命中机会，长事件直接采集。它不是每隔一个墙钟周期扫描所有 goroutine。\cite{pp-mprof}\cite{pp-profile-kinds}


用统一时间单位写，阻塞时长为 \(d\)、阈值为 \(r\) 时，理想包含概率为 \(\min(1,d/r)\)。若一个短事件被采中，源码把计数加 \(r/d\)，把等待量加 \(r\)；若事件足够长，则计数加一、等待量加真实观测时长。二者都是逆概率校正。因而 delay/contentions 是经过抽样校正后的粗略均值，既不保留等待分布，也不能计算 p99。改变 rate 还会改变方差与开销，不能只比较文件里的消息行数。\cite{pp-mprof}


mutex profile 的抽样则按竞争事件进行，配置 rate 为十时，平均抽取约十分之一，再把计数与时间各乘十。更容易误读的是它的归因位置：相关栈通常是释放锁、结束临界区的路径，而非锁指令消耗了多少 CPU。若多个 goroutine 同时等待，同一墙钟秒可以贡献多个 goroutine 秒的竞争量。这个总量反映被阻塞的并发需求，不能作为单次持锁时长。\cite{pp-mprof}\cite{pp-sema}


Go 1.25 的 semrelease 实现还说明了“近似”的具体来源。为避免解锁时扫描所有等待者，它用队列头、尾等待时长的平均值乘其他等待者数量，再加即将被唤醒者的等待量。解除队列关系时更新有关起点，避免下一次释放重复收费。这个设计把归因成本保持在适当范围，但不是逐锁实例、逐等待者的精确积分数据库；相同解锁栈上的多个锁还会在 profile 中聚合。\cite{pp-sema}


发布时间同样决定窗口含义。block 的完成路径才调用 blockevent；同步等待可能在窗口开始前已经发生，直到窗口内结束才把整段费用发布。两个累计快照相减得到的是“窗口内新增记账”，不一定是“每条等待区间与窗口交集”。设窗口长 \(T\)、新增等待权重 \(W\)，只有边界影响与抽样误差足够小时，\(W/T\) 才近似平均等待并发。仍未结束的长等待则可能尚未完整计入，应该补当前栈或时序追踪。\cite{pp-sema}\cite{pp-mprof}


对于一次简单唤醒，\texttt{readyWithTime} 先记录释放时刻再将等待者设为 runnable，之后获得 CPU 才实际 running。因此图解应区分“仍不能继续的阻塞”和“已经可运行但尚未调度”。复杂锁路径可能重复睡眠、重试或执行自旋，不能反过来声称 block profile 与 trace 的每一段状态都精确一一对应。block 与 mutex 也常是同一争用的两种归因视角，不能把二者总 delay 相加算请求延迟。\cite{pp-sema}


goroutine profile 又是第三种测量。其数值主要表示采样集合中的栈数量，不是这些栈累计运行了多久。Go 1.25 采用短暂停顿确立参与集合，再让协调者与调度器协作采栈；每个 goroutine 在 Absent、InProgress、Satisfied 之间转换，被读取栈时不能同时运行；清理时再短暂停顿。新创建的 G 和系统 G 有专门规则，不能简单宣称它是“对整个进程所有线程同时拍的一张照片”，也不能说整个遍历始终停世界。\cite{pp-goroutine}


这里还要区分输出模式：上述协作式采栈适用于 goroutine profile 的 debug=0/1 路径；debug>=2 则调用 \texttt{runtime.\allowbreak{}Stack(buf,\allowbreak{}true)}，在格式化全部 traceback 期间持有停世界，缓冲区不足还会扩大后重试。因此 \texttt{/\allowbreak{}debug/\allowbreak{}pprof/\allowbreak{}goroutine?debug=2} 不能直接套用前述短暂停顿结论。\cite{pp-goroutine}


由此可得到一个实用判读例子：一千个 worker 停在 channel receive，可能仅仅是空闲池。一次快照给出了数量，却没有给出每个 worker 的等待时长、工作频度或泄漏原因。只有比较稳定负载下的重复快照、对象生命周期与完成条件，才能把“很多相似栈”转化为可检验的问题。工具之间应互补观察，而不是因为都出现了 channel 函数就被视为同一种数据。\cite{pp-goroutine}\cite{pp-profile-kinds}


\section{编码与符号化：从一个地址到多层逻辑函数}

\Needspace{7.98cm}
\begin{center}\begin{minipage}{\linewidth}
\centering
\begin{tikzpicture}[x=0.016cm,y=-0.016cm]
\path[use as bounding box] (0,0) rectangle (960,380);
\draw[rounded corners=5pt,draw={rgb,255:red,25;green,118;blue,210},fill={rgb,255:red,25;green,118;blue,210},fill opacity=.075] (30,80) rectangle (290,240);
\node[anchor=center,text={rgb,255:red,35;green,52;blue,72},font=\fontsize{9.1}{11.0}\selectfont] at (160.0,133.0) {Sample};
\node[anchor=center,text={rgb,255:red,35;green,52;blue,72},font=\fontsize{8.2}{9.9}\selectfont] at (160.0,160.0) {location\_id[]};
\node[anchor=center,text={rgb,255:red,35;green,52;blue,72},font=\fontsize{8.2}{9.9}\selectfont] at (160.0,187.0) {value[] / label[]};
\draw[rounded corners=5pt,draw={rgb,255:red,129;green,80;blue,181},fill={rgb,255:red,129;green,80;blue,181},fill opacity=.075] (360,50) rectangle (610,155);
\node[anchor=center,text={rgb,255:red,35;green,52;blue,72},font=\fontsize{9.1}{11.0}\selectfont] at (485.0,89.0) {Location};
\node[anchor=center,text={rgb,255:red,35;green,52;blue,72},font=\fontsize{8.2}{9.9}\selectfont] at (485.0,116.0) {address + Line[]};
\draw[rounded corners=5pt,draw={rgb,255:red,189;green,101;blue,0},fill={rgb,255:red,189;green,101;blue,0},fill opacity=.075] (360,220) rectangle (610,315);
\node[anchor=center,text={rgb,255:red,35;green,52;blue,72},font=\fontsize{9.1}{11.0}\selectfont] at (485.0,254.0) {Mapping};
\node[anchor=center,text={rgb,255:red,35;green,52;blue,72},font=\fontsize{8.2}{9.9}\selectfont] at (485.0,281.0) {构建 / 地址映射};
\draw[rounded corners=5pt,draw={rgb,255:red,88;green,167;blue,0},fill={rgb,255:red,88;green,167;blue,0},fill opacity=.075] (690,50) rectangle (930,155);
\node[anchor=center,text={rgb,255:red,35;green,52;blue,72},font=\fontsize{9.1}{11.0}\selectfont] at (810.0,89.0) {Function};
\node[anchor=center,text={rgb,255:red,35;green,52;blue,72},font=\fontsize{8.2}{9.9}\selectfont] at (810.0,116.0) {名字 / 文件 / 行};
\draw[draw={rgb,255:red,99;green,115;blue,134},-{Stealth[length=4pt]}] (295,130) -- (355,100);
\draw[draw={rgb,255:red,99;green,115;blue,134},-{Stealth[length=4pt]}] (615,100) -- (685,100);
\draw[draw={rgb,255:red,99;green,115;blue,134},-{Stealth[length=4pt]}] (485,160) -- (485,215);
\node[anchor=center,text={rgb,255:red,35;green,52;blue,72},font=\fontsize{8.6}{10.4}\selectfont] at (800,260) {所有 value 列按};
\node[anchor=center,text={rgb,255:red,35;green,52;blue,72},font=\fontsize{8.6}{10.4}\selectfont] at (800,290) {sample\_type 解释};
\end{tikzpicture}
\captionof{figure}{profile.proto 的共享表与引用。Sample 的栈叶在前；一个 Location 可以包含多条内联 Line。ID 是文件内引用，不是地址。}
\end{minipage}\end{center}

profile.proto 将重复内容拆成几类表：Sample 保存 location\_\allowbreak{}id 序列、值向量和标签；Location 保存地址、可选的 mapping 引用及若干 Line；Line 引用 Function；Function 保存名字、文件和起始行；Mapping 描述加载地址区间、文件偏移与构建身份；大量字符串通过共享 string\_\allowbreak{}table 索引引用。string\_\allowbreak{}table 的零号项必须为空字符串，各类非零 ID 则是文件内部的引用编号。Location 的 ID 可以来自整数序列，不能把它当成机器地址。\cite{pp-proto}


一个 CPU Sample 的 value 若为 \texttt{[12,\allowbreak{}120000000]}，必须同时读取 sample\_\allowbreak{}type 才能解释为十二个名义样本和一亿两千万纳秒。heap 可以有四列，其他 producer 可以定义更多类型；所有 Sample 的向量长度必须和 sample\_\allowbreak{}type 相同。格式允许有符号 int64，因此差分负值是有效表达，不是文件损坏。period 描述采样间隔的类型与数值，但 value 列可能已经完成了周期换算；不应看到 period=10000000 就把 cpu/nanoseconds 再乘一次。\cite{pp-proto}\cite{pp-cpu-builder}


栈的方向也必须固定：\texttt{Sample.\allowbreak{}location\_\allowbreak{}id[0]} 是叶子，随后走向调用者。本章教学记录 \texttt{[5,\allowbreak{}4,\allowbreak{}2,\allowbreak{}1]} 可表示 Hash 的位置、Encode、Handle、Root。Location 5 内又有两条 Line：HashRound 在前，Hash 在后，表示 HashRound 内联到 Hash 中。因此四个物理位置可以展开为五层逻辑函数；这不表示多做了一次真实函数调用，更不要求栈上存在一个 HashRound 的独立返回地址。\cite{pp-proto}\cite{pp-cpu-builder}


PC 则需要区分采样现场的指令地址、栈展开返回的 return PC、文件格式中的 Location.Address，以及对象文件里的符号地址。Go 的 CallersFrames 对符合其约定的返回地址适当减一，以回到有关调用指令的范围，并展开内联标记。减一只是定位到调用指令内，不承诺指向它的首字节。pprof 的 Go builder 再把恢复的 Frame.PC 写到 Location。对已经完成这种归一化的地址又统一减一，可能产生新的行号偏差。\cite{pp-pc}\cite{pp-cpu-builder}


符号化可以发生在生产端，也可以发生在读取端。Go 文件常已包含函数、源码行及内联信息，读取这样的文件不一定还要拿一个二进制才能列出函数热点。需要反汇编、补全缺失符号或对应精确源码时，匹配构建仍然重要。pprof 本地符号化按 Mapping 归组，打开相应文件，检查双方已知 Build ID 是否相符，再用地址查询源码帧。若已经符号化且未请求强制处理，通常不会重复解析该映射。\cite{pp-symbolizer}


对于 ELF，进程虚拟地址到对象文件地址的转换并不总是简单的 \texttt{pc-\allowbreak{}memory\_\allowbreak{}start}。程序头、映射偏移、可执行文件类型和加载基址都参与计算；pprof 的 ObjAddr 返回 \texttt{addr-\allowbreak{}base}，其中 base 由对应映射和 ELF 信息算出。\texttt{pc-\allowbreak{}memory\_\allowbreak{}start} 可以用于特定映射内的相对位置比较，但不能作为适用于一切二进制的 addr2line 输入公式。快速 nm 路径又可能只有符号名而没有完整文件和行号，需要区分符号化精度。\cite{pp-binutils}


标签扩展了栈的上下文，但没有补回时序。格式支持字符串标签和带单位的数值标签；\texttt{input\_\allowbreak{}bytes=64 bytes} 与 \texttt{input\_\allowbreak{}bytes=64 requests} 不是相同语义，单位也参与合并身份。新协议说明不鼓励同键多值，更不应在同一键下混用字符串和数字。Go 标准库的 \texttt{pprof.\allowbreak{}Labels} 则提供字符串标签，Go 1.25 明确只有 CPU 与 goroutine profiles 使用这些 labels，不能因为协议有 label 就承诺 heap 或 mutex 自动按租户归因。\cite{pp-proto}\cite{pp-merge}\cite{pp-label-support}


\texttt{pprof.\allowbreak{}Do} 把增强后的 context 标签安装到当前 G，结束时恢复传入 context 所带标签；新建 G 继承创建者当时的标签。传入一个有标签的 context，并不会让一个早已存在的 worker 自动换标签。应在真正执行任务的 goroutine 上安装对应上下文，并正确结束请求级生命周期。无限后台 G 继承请求标签后长期运行，会把不属于该请求的后续工作归入同一类别；这属于上下文建模错误，不是查询界面能修复的问题。\cite{pp-labels}


本节的关键结果是“地址、位置、函数和请求类别”具有不同身份。相同函数可以有多个指令位置、多个内联上下文、多个调用方和多类请求；同一逻辑样本也可能在不同进程中拥有不同绝对地址。下一步合并如果只依据函数名或只依据位置编号，就会制造错误的等价关系。\cite{pp-proto}\cite{pp-merge}


\section{合并、过滤、归一化和差分是不同的数学操作}

\Needspace{7.34cm}
\begin{center}\begin{minipage}{\linewidth}
\centering
\begin{tikzpicture}[x=0.016cm,y=-0.016cm]
\path[use as bounding box] (0,0) rectangle (960,340);
\draw[rounded corners=5pt,draw={rgb,255:red,25;green,118;blue,210},fill={rgb,255:red,25;green,118;blue,210},fill opacity=.075] (50,70) rectangle (440,220);
\node[anchor=center,text={rgb,255:red,35;green,52;blue,72},font=\fontsize{9.1}{11.0}\selectfont] at (245.0,131.5) {默认总量：480 ms};
\node[anchor=center,text={rgb,255:red,35;green,52;blue,72},font=\fontsize{8.2}{9.9}\selectfont] at (245.0,158.5) {80 / 480 = 16.67\%};
\draw[rounded corners=5pt,draw={rgb,255:red,88;green,167;blue,0},fill={rgb,255:red,88;green,167;blue,0},fill opacity=.075] (520,70) rectangle (910,220);
\node[anchor=center,text={rgb,255:red,35;green,52;blue,72},font=\fontsize{9.1}{11.0}\selectfont] at (715.0,131.5) {条件总量：320 ms};
\node[anchor=center,text={rgb,255:red,35;green,52;blue,72},font=\fontsize{8.2}{9.9}\selectfont] at (715.0,158.5) {80 / 320 = 25\%};
\node[anchor=center,text={rgb,255:red,35;green,52;blue,72},font=\fontsize{9.1}{11.0}\selectfont] at (480,280) {变化的是分母，80 ms 的绝对成本没有变化};
\end{tikzpicture}
\captionof{figure}{同一成本，两种正确的百分比。合成教学基线：全局480ms、route=A共320ms、A中的HashRound为80ms。默认分母和relative\_percentages回答不同条件集合。}
\end{minipage}\end{center}

Merge 首先检查基本类型兼容性，然后建立输出的 Mapping、Function、Location 与 Sample 表。每读取一个输入 profile，都重建该文件的局部 ID 到输出对象的映射。因此两个输入的 Location 7 不会因为编号相同就自动合并，两个编号不同的等价位置也可以归一。sample key 使用重新映射后的有序位置序列、字符串标签、数值标签及单位；命中同一键后才逐维累加值向量。\cite{pp-merge}


处理 ASLR 时，Mapping 的键考虑对齐后的映射大小、文件偏移和 Build ID；没有 Build ID 时退到文件身份。Location 的身份结合映射、相对地址和源码行/函数信息。若相同构建在两个进程分别加载到 0x1000 和 0x2000，同一位置位于 0x1050 和 0x2050，则有机会归一为同一位置。相反，同名可执行文件发生重建、文件身份缺失或构建 ID 被错误复用，算法能执行合并并不证明物理身份正确。\cite{pp-merge}


合并还会修改头部：本版本的 combineHeaders 取较早的 time\_\allowbreak{}nanos，duration\_\allowbreak{}nanos 求和，而不是对所有采集区间做并集。两个同时采集的实例，每个十秒，合并后 duration 可以为二十秒；这不代表现实中经过了二十秒。协议把这些字段列为信息字段，但界面会展示它们，因此分析者必须检查元数据的来源。特别是差分里出现的 duration，不应被拿来计算一个并不存在的业务吞吐。\cite{pp-proto}\cite{pp-merge}


Aggregate 与 Merge 也不同。按函数聚合时，分析器可以丢掉地址、行号或列号；选择不保留内联时，Location.Line 只留下最外面的实体函数。这些是投影操作：将更细的身份空间映射到更粗的空间，使后续节点聚合更适合问题，但信息已经不可逆。生成函数级报告后，不能声称那些数值仍能精确恢复到原来的某条机器指令。\cite{pp-scale}


过滤至少有两类。focus 和 ignore 决定保留哪些完整样本：至少一帧匹配 focus 且没有任何帧匹配 ignore，样本才保留。hide 和 show 则主要改变保留样本里的可见帧；所有帧都被删光时，样本也会被移除。隐藏中间函数会把图上的祖先和后代拉近，因此一条可见边未必就是源码里的直接调用。标签过滤又是以样本上下文作条件，而非把该标签对应的所有函数全局删除。\cite{pp-filter}


百分比分母是另一个独立选择。默认流程先创建 Report 并保存总量，再应用过滤，所以“只看 route=A”不保证百分比重新加到一百。启用 \texttt{relative\_\allowbreak{}percentages} 时，顺序才变成先过滤再建 Report。本章中 HashRound 属于 A 的权重始终是 80 ms；默认除以全局 480 ms 得 16.67\%，相对百分比除以 A 的 320 ms 得 25\%。同一个成本出现两个正确百分比，差别来自问题的条件集合。\cite{pp-denominator}\cite{pp-teaching}


Normalize 解决的则是体量尺度。对当前源 profile 第 \(j\) 维总量 \(S_j\) 和基线总量 \(B_j\)，它把当前每个值乘 \(B_j/S_j\)；当当前总量为零时，该维缩放为零。操作对象是当前 profile，不是把基线调成当前体量。ScaleN 会对结果取整再转成 int64，因此归一化后不承诺每个原始向量关系都精确保留。它适合观察相对结构，但不会自动使用成功请求数，也不控制输入类别和错误率。\cite{pp-merge}\cite{pp-scale}


固定分析器还存在一个需要与理想公式区分的实现边界：\texttt{ScaleN} 只在比例不等于一的列上更新保留样本的判断。独立实验使用该固定提交未经修改的 profile 子包，将 \texttt{[2,\allowbreak{}0]} 按 \texttt{[1,\allowbreak{}0.\allowbreak{}5]} 缩放，实际样本被删除，而逐维数学结果应为 \texttt{[2,\allowbreak{}0]}。因此混合比例、部分列为零时的变化不应一概归于取整；本章 count 与 CPU 同比例缩放的教学案例不受该反例影响。此处真实执行的是固定提交的子包，不是以它构建的完整 CLI。\cite{pp-scale}\cite{pp-scale-repro}


差分可以写成 \(D=\alpha P_{after}-P_{before}\)，其中没有归一化时 \(\alpha=1\)。\texttt{-\allowbreak{}diff\_\allowbreak{}base} 会先给基线样本加内部标签 \texttt{pprof::base=true}，然后可选归一化当前 profile，再把基线乘负一并合并。这个标签让两类样本在中间结构中仍可区分，用于报告分母；绘图时同一函数的正负贡献仍会相抵。因而不能只看合并消息数量猜测净差分的结构。\cite{pp-diff}\cite{pp-merge}


一般情况下，Report 总量取所选样本值的绝对值之和；若存在非零的 diff-base 总量，则改用基线绝对总量。这解释了一个常见误区：\texttt{-\allowbreak{}base} 不是在任何输入下都以“源总量减基线总量”作分母。只有真正的累计增量等合适条件下，各项非负，绝对值之和才等于总量之差；当既有增长又有下降时，净差为零仍可能具有很大的变化总量。\cite{pp-total}


本章变体的六组净变化为 −40、−40、−20、+40、+40、+20 ms，绝对值合计 200 ms。\texttt{-\allowbreak{}base} 报告的分母就是 200 ms；\texttt{-\allowbreak{}diff\_\allowbreak{}base} 保留基线 480 ms 分母。同一 HashRound 在两个路由上分别减少和增加 40 ms，函数级聚合后抵消为零，不能据此断言业务结构未变。Parse 的净值为 −40 ms，但 flat\% 显示正的幅度，因为 Percentage 使用绝对比例；增减方向必须阅读带符号的数值。\cite{pp-total}\cite{pp-percent}\cite{pp-teaching}


由此可区分三个问题：合并问“这些可比观察总共发生了多少”，条件过滤问“所关心的子集合是什么”，归一化和差分问“在某个尺度下前后结构怎样变化”。把这三步顺序写清，比只保存最后一张差分火焰图更能支持复核。操作可执行只证明数据结构满足要求，不证明两个程序阶段、负载分布或运行时版本具有研究上的可比性。\cite{pp-merge}\cite{pp-filter}\cite{pp-diff}


\section{报告算法：flat、cum、调用图与火焰图为何不同}

\Needspace{7.34cm}
\begin{center}\begin{minipage}{\linewidth}
\centering
\begin{tikzpicture}[x=0.016cm,y=-0.016cm]
\path[use as bounding box] (0,0) rectangle (960,340);
\draw[rounded corners=5pt,draw={rgb,255:red,25;green,118;blue,210},fill={rgb,255:red,25;green,118;blue,210},fill opacity=.075] (40,45) rectangle (230,130);
\node[anchor=center,text={rgb,255:red,35;green,52;blue,72},font=\fontsize{9.1}{11.0}\selectfont] at (135.0,87.5) {A};
\draw[rounded corners=5pt,draw={rgb,255:red,25;green,118;blue,210},fill={rgb,255:red,25;green,118;blue,210},fill opacity=.075] (40,190) rectangle (230,275);
\node[anchor=center,text={rgb,255:red,35;green,52;blue,72},font=\fontsize{9.1}{11.0}\selectfont] at (135.0,232.5) {B};
\draw[rounded corners=5pt,draw={rgb,255:red,129;green,80;blue,181},fill={rgb,255:red,129;green,80;blue,181},fill opacity=.075] (385,110) rectangle (575,205);
\node[anchor=center,text={rgb,255:red,35;green,52;blue,72},font=\fontsize{9.1}{11.0}\selectfont] at (480.0,157.5) {F};
\draw[rounded corners=5pt,draw={rgb,255:red,88;green,167;blue,0},fill={rgb,255:red,88;green,167;blue,0},fill opacity=.075] (730,45) rectangle (920,130);
\node[anchor=center,text={rgb,255:red,35;green,52;blue,72},font=\fontsize{9.1}{11.0}\selectfont] at (825.0,87.5) {X};
\draw[rounded corners=5pt,draw={rgb,255:red,88;green,167;blue,0},fill={rgb,255:red,88;green,167;blue,0},fill opacity=.075] (730,190) rectangle (920,275);
\node[anchor=center,text={rgb,255:red,35;green,52;blue,72},font=\fontsize{9.1}{11.0}\selectfont] at (825.0,232.5) {Y};
\draw[draw={rgb,255:red,99;green,115;blue,134},-{Stealth[length=4pt]}] (235,87) -- (380,150);
\draw[draw={rgb,255:red,99;green,115;blue,134},-{Stealth[length=4pt]}] (235,233) -- (380,170);
\draw[draw={rgb,255:red,99;green,115;blue,134},-{Stealth[length=4pt]}] (580,150) -- (725,87);
\draw[draw={rgb,255:red,99;green,115;blue,134},-{Stealth[length=4pt]}] (580,170) -- (725,233);
\node[anchor=center,text={rgb,255:red,35;green,52;blue,72},font=\fontsize{8.6}{10.4}\selectfont] at (480,305) {调用树 / 栈视图保留更多路径上下文};
\end{tikzpicture}
\captionof{figure}{局部调用边不等于完整路径。两条原始栈 A→F→X 与 B→F→Y 的合并调用图会共享F；图上的连接不能证明 A→F→Y 曾出现。}
\end{minipage}\end{center}

先固定可见栈展开和节点粒度，再定义函数级权重。对于正权重 profile，\(flat(f)=\sum_i w_i\,1\{f\text{ 是样本 }i\text{ 的叶节点}\}\)；\(cum(f)=\sum_i w_i\,1\{f\text{ 至少出现在样本 }i\text{ 中一次}\}\)。源码的普通图构建与这个解释一致：逐样本从根向叶走，用 seenNode 防止同一节点重复累计，用 seenEdge 防止重复边，最后把 flat 加到叶节点。\cite{pp-graph}


因此入口的 cum 很大而 flat 很小，说明它承载下游成本，不说明入口处某条赋值特别昂贵。所有函数 cum 之和通常远大于 profile 总量，因为同一权重归入多个祖先；这不是重复计费的 bug，而是 inclusive 统计定义。对于递归 \texttt{Root→Handle→Walk→Walk}，同一样本的 Walk cum 只加一次。它不是“每深一层就再增加同一份 CPU”，更不能用 cum 除以递归深度计算每次调用耗时。\cite{pp-graph}


“函数”本身也经过粒度选择。按地址时，同一个源码函数中的不同位置可以是不同节点；按行时保留行信息；按函数时许多位置会归并。内联会让逻辑叶子和实体机器函数不同：本例 HashRound 内联到 Hash，默认 flat 180 ms 归到 HashRound，Hash 的 cum 为 180 ms；忽略内联后则把这部分 flat 归入 Hash。两种视图都来自同一批观察，不能把差异解释为优化前后发生了资源变化。\cite{pp-scale}\cite{pp-graph}\cite{pp-teaching}


普通调用图把所有样本中可见的节点和邻接关系汇总；边的权重是流经该关系的样本权重，不是函数调用次数。图构建排除节点到自身的普通边，因此递归深度还会进一步丢失。截断和剪枝可能生成跨过隐藏节点的 residual edge。若真实样本分别是 \texttt{A→F→X} 与 \texttt{B→F→Y}，合并图具有这些局部边，却不能证明 \texttt{A→F→Y} 这条完整路径曾经发生；局部连接关系不等于原始路径集合。\cite{pp-graph}


调用树保留路径上下文，代价是同一函数会在多个调用路径下出现；按完整栈组织的火焰图也保留这类上下文。pprof 的栈数据把样本位置及内联链反转成根到叶的 Sources，附上所选维度的 Value，并记录来源函数供交互聚焦。递归可以在栈中重复出现；用于聚焦的 Places 索引对同一栈的重复来源只记录最外层，避免把同一权重重复计入交互集合。\cite{pp-flame}


火焰图的宽度取决于权重，纵向表达栈层次，横向是布局位置。CPU 火焰图中的宽块可能表示许多短调用，也可能表示少数较长计算；heap 火焰图中的宽块可以是保留字节，block 火焰图可以是等待量。颜色和排序还取决于具体界面，不能脱离图例把颜色当成“更慢”。尤其并排的两个块不具有“先发生左边再发生右边”的时间含义。\cite{pp-flame}\cite{pp-proto}


还要区别“整份 profile 的总量”和“当前展示节点所能归属的总量”。无法符号化、无可见栈、隐藏规则、负值抵消、阈值剪枝等都可能让显示和总量不同。降低 nodefraction 或增加 nodecount 只改变展示覆盖，不会改善采样覆盖；把 unknown 隐藏掉也不会解决归因。分析报告应保留总量、过滤条件和被省略范围，避免把局部画面误作完整事实。\cite{pp-filter}\cite{pp-total}\cite{pp-graph}


从实现复杂度看，核心工作量跟样本的总栈长度、内联展开数量及标签内容有关，而不只是文件中的函数个数。合并利用哈希表避免反复扫描全部已见结构；每个样本的节点去重集合限制递归重复累计；可视化再把这些结构投影成更易读的形状。高基数标签会使本来相同的栈分裂成更多键，增加状态规模，也降低每个类别的统计稳定性。这是为何标签设计属于测量设计，而非纯界面装饰。\cite{pp-merge}\cite{pp-graph}


\section{可复核的 48 条教学样本与配套图解}

为了让上述规则可以手工核算，本章构造一个只有 48 条 protobuf Sample 消息的基线。每条值向量是 \texttt{[1,\allowbreak{}10000000]}，即一份名义十毫秒 CPU 权重；文件 duration 为 400 ms。地址、构建 ID、源码路径和采集时间均为虚构，数据没有宣称来自真实 Go 程序。基线按下表生成，表中用根到叶方向方便阅读，实际编码方向相反；方括号表示同一个 Location 内的内联链。\cite{pp-teaching}



{\small
\begin{longtable}{@{}p{1.750cm}@{\hspace{0.38cm}}p{6.000cm}@{\hspace{0.38cm}}p{1.150cm}@{\hspace{0.38cm}}p{1.600cm}@{\hspace{0.38cm}}p{2.600cm}@{}}
\toprule
组 & 根到叶的逻辑路径 & route & 记录数 & CPU 权重 \\
\midrule\endhead
A/parse & Root→\allowbreak{}Handle→\allowbreak{}Parse & A & 12 & 120 ms \\
A/hash & Root→\allowbreak{}Handle→\allowbreak{}Encode→\allowbreak{}[Hash→\allowbreak{}HashRound] & A & 8 & 80 ms \\
A/recurse & Root→\allowbreak{}Handle→\allowbreak{}Walk→\allowbreak{}Walk & A & 6 & 60 ms \\
B/hash & Root→\allowbreak{}Batch→\allowbreak{}[Hash→\allowbreak{}HashRound] & B & 10 & 100 ms \\
B/alloc & Root→\allowbreak{}Batch→\allowbreak{}Alloc & B & 6 & 60 ms \\
A/encode & Root→\allowbreak{}Handle→\allowbreak{}Encode & A & 6 & 60 ms \\
\bottomrule\end{longtable}}
A 共 32 份、320 ms；B 共 16 份、160 ms；总数 48、总权重 480 ms。还有数字标签 input\_\allowbreak{}bytes，A 为 64 bytes、B 为 256 bytes，用于展示单位明确的上下文，不参与 CPU 总量换算。HashRound 和 Hash 共用一个地址且具有两条 Line，Walk 的位置 ID 在递归栈里出现两次。这样一个小文件同时覆盖分支、内联、递归、标签和值向量，避免靠巨大生产数据隐藏规则。\cite{pp-teaching}


生成器为 \texttt{research/\allowbreak{}pprof/\allowbreak{}make\_\allowbreak{}teaching\_\allowbreak{}profile.\allowbreak{}py}，只使用 Python 标准库手工编码 protobuf，便于逐字段对照协议。\texttt{teaching/\allowbreak{}manifest.\allowbreak{}json} 保存分组、每条记录、权重和文件 SHA-256；\texttt{verify\_\allowbreak{}teaching.\allowbreak{}py} 调用真实 \texttt{go tool pprof}，不是把预期结果自己画成假终端输出。执行时关闭符号化，因为这些地址没有对应的真实 ELF；因此它验证的是格式解析、聚合和报告规则，不是实际程序的符号解析能力。\cite{pp-proto}\cite{pp-teaching}


基线查询的关键结果如下，绝对权重与手算完全相符：



{\small
\begin{longtable}{@{}p{3.665cm}@{\hspace{0.38cm}}p{3.665cm}@{\hspace{0.38cm}}p{3.665cm}@{\hspace{0.38cm}}p{3.665cm}@{}}
\toprule
节点 & flat & cum & 理由 \\
\midrule\endhead
Root & 0 & 480 ms & 所有样本都经过根 \\
Handle & 0 & 320 ms & A 的四组路径 \\
Batch & 0 & 160 ms & B 的两组路径 \\
HashRound & 180 ms & 180 ms & 两条分支的内联叶子之和 \\
Hash & 0 & 180 ms & 包含内联叶子的逻辑外层 \\
Parse & 120 ms & 120 ms & 十二份叶子权重 \\
Encode & 60 ms & 140 ms & 自身六份，加下游哈希八份 \\
Walk & 60 ms & 60 ms & 六份递归样本，每条只计一次 cum \\
Alloc & 60 ms & 60 ms & 六份叶子权重 \\
\bottomrule\end{longtable}}
可重放的最小入口如下，实际的十五条完整命令保存在 verification.json。命令默认在当前任务根目录执行；这只是验证入口，不替代前文的机制推导。



\begin{codeblock}
python3 research/pprof/make_teaching_profile.py
python3 research/pprof/verify_teaching.py
go tool pprof -symbolize=none -top research/pprof/teaching/baseline.pprof
\end{codeblock}

第一组反例是过滤。筛选 A 后，HashRound 为 80 ms，在全局分母下显示 16.67\%；相对分母下显示 25\%。focus=Handle 会删除 B 的完整样本，hide=Handle 则隐藏一个帧但保留 B 的 Batch 分支。可见“只看这个函数路径”和“把这个函数从图上拿掉”是不同操作，不能用其中一个替代另一个。对应输出为 route-a、route-a-relative、focus-handle、hide-handle 四组文本。\cite{pp-teaching}


第二组反例是位置身份。relocated 变体把 mapping 起点和所有指令地址增加 0x1000，同时给位置和函数 ID 增加 100，保持构建身份、相对位置、标签和权重不变。与基线合并后，pprof 报告总权重 960 ms，各函数权重恰好翻倍。这支持“在该合成输入条件下，重定位和局部 ID 不会错误拆分等价位置”的结论，但没有验证真实共享库搜索、Build ID 提取或缺失符号恢复。\cite{pp-merge}\cite{pp-teaching}


第三组反例是变化抵消。changed 变体六组事件权重依次改为 8、4、4、14、10、8，总体仍为 48 份。HashRound 的总量仍为 180 ms，但 A 的哈希从 80 降为 40 ms，B 的哈希从 100 升为 140 ms。Encode 的 flat 从 60 升到 80 ms，cum 却从 140 降到 120 ms，因为下游下降更多。只盯 flat、只盯 cum、只盯 HashRound 总量，都会遗漏一部分真实结构变化。\cite{pp-teaching}


该变体的 Parse 差值为 −40 ms，Alloc 为 +40 ms，Encode flat 为 +20 ms、cum 为 −20 ms，Walk 为 −20 ms；Handle cum 为 −80 ms，Batch cum 为 +80 ms。\texttt{diff.\allowbreak{}txt} 使用 480 ms 基线分母，\texttt{base.\allowbreak{}txt} 使用 200 ms 绝对变化分母。二者标题都可能写“Showing nodes accounting for 0”，但下方仍有正负节点：零是净和，不是没有变化。差分文件里 duration 变为 800 ms，也只是两个输入头部相加，不是新的真实采集窗口。\cite{pp-total}\cite{pp-merge}\cite{pp-teaching}


最后，double-volume 变体把六组权重全部加倍，使用六条聚合 Sample 表示 96 份事件。未经归一化时它有 960 ms；以基线为 diff\_\allowbreak{}base 并归一化后，所有函数净变化为零，但报告仍保留 480 ms 的基线分母。三个变体各六条消息，连同基线共 66 条 protobuf Sample 消息，方便完整检查。“六条消息”与“九十六份事件权重”的差别正好再次证明 Sample 不等于原始事件。\cite{pp-teaching}


读者应能完成三项检验：交换同一文件中 Sample 的排列，判断报告为何基本不变；把一条已聚合记录拆成多个相同键的小记录，判断守恒的是哪些值；增加一个新标签，使部分相同栈不再合并，再说明总量为何仍应守恒。进一步的研究问题则是：这种不变性在哪些环节失效，例如整数缩放取整、采样率变化、丢样、缺少构建身份或不同总体被错误混合。提出这些条件，比记住更多命令更接近可靠的性能研究。\cite{pp-proto}\cite{pp-merge}\cite{pp-scale}


本轮验证结果是十五种查询成功、十四项语义断言通过，原始标准输出、标准错误和版本记录均保留。未运行真实 CPU 信号压力实验、heap 随机采样收敛实验或生产负载，教学 CLI 也未使用所读 google/pprof 提交构建；独立边界验证仅编译执行该提交的 profile 子包；这些方面仍属于源码解释和理想模型。由此得到的是一条可复核的解释链：知道记录如何产生、哪些转换保持权重、哪些操作损失信息，最终才能把一个热点转化为有条件、可反驳的性能结论。\cite{pp-teaching}


\chapter{Perfetto：从共享内存采集到可解释的时序查询}
\label{ch:20}

本章固定研究 Perfetto \textbf{v58.2}，提交 \texttt{add693d8\allowbreak{}b338ba95\allowbreak{}99dbcbc3\allowbreak{}e300b1ab\allowbreak{}8c000897}，查证日期为 2026-10-04。源码位置均对应该提交；运行验证使用同版本官方 Linux amd64 Trace Processor。Go 格式与转换链路的源码分析固定在 Go 1.25.0，实际原生 trace 负例由本机 Go 1.22.12 生成。版本、教学构造与真实执行分别记录，避免把兼容性判断推广到所有历史或未来版本。\cite{pf-release}\cite{pf-go-main}\cite{pf-validation}


\section{测量契约：Perfetto 能保留什么，不能凭空得到什么}

Perfetto 是由采集协议、数据源、服务、文件格式、Trace Processor 和界面组成的系统。只打开界面并不意味着已经采集了调度事件；成功读取一个文件也不意味着文件里存在所需的因果证据。学习时应首先确定观察对象、事件来源、时间域、身份、完整性和归因规则，然后才选择图表。其核心价值是把多个来源的事件放到可查询的时间关系中，而不是把任何性能文件自动变成“根因”。\cite{pf-service}\cite{pf-architecture}


Linux perf、pprof 与 Perfetto 所处的层次有交集。Linux perf 通常接入内核性能事件机制，处理计数、采样或跟踪；pprof 主要分析带栈、标签和类型的聚合权重；Perfetto 擅长接入多种时序数据并进行关联分析。Perfetto v58.2 的内置导入器确实包括 perf.data 和 pprof，但“支持输入格式”不等于“重新获得生产者没有记录的事实”。同一套界面可以显示聚合调用栈与线程时间线，这两者的证据强度仍不同。\cite{pf-importers}\cite{pf-pprof}


设一个请求起止为 \([a,b)\)，墙钟延迟为 \(W=b-a\)。若已经知道各线程哪些执行区间属于该请求，则总 CPU 成本为这些区间长度的和 \(C\)。多核并行时可以有 \(C>W\)；某个请求长时间等待时也可以有 \(C\ll W\)。CPU 成本回答资源消费，墙钟延迟回答完成时间，关键路径回答哪些前后依赖限制完成。这三个问题不能通过换一种可视化布局互相替代。


本章贯穿案例使用合成的、符合 ftrace 文本格式的 \texttt{sched\_\allowbreak{}switch}、\texttt{sched\_\allowbreak{}waking} 与 \texttt{tracing\_\allowbreak{}mark\_\allowbreak{}write} 事件，再交给官方程序真实解析。请求从 0 到 20 ms；主线程运行 0—4 和 16—20 ms，4—6 ms 可运行但未获调度，6—16 ms 睡眠；工作线程在另一 CPU 上运行 1—16 ms。因此主线程 CPU 为 8 ms、工作线程为 15 ms、请求总 CPU 为 23 ms。这个值与 20 ms 墙钟时间没有矛盾。\cite{pf-request}


此例中的线程归属是教学模型显式给定的。真实线程池可同时处理多个请求，不能把“请求期间所有工作线程的 CPU”全部算给这个请求；OS 线程也可能运行多个 Go goroutine。应用阶段、请求标识与调度事实之间必须有可解释的连接。后文将通过区间相交查询计算已归属阶段的成本，而不会把时间重叠当作自动归因。


把关键路径形式化，可以将已证实的阶段与依赖构成有向无环图，节点权重是该阶段持续时间，完成时间受最长依赖链约束。并行分支的资源成本相加，汇合点的完成时间则取受依赖约束的最晚完成者。若依赖边遗漏，计算出的最长路径只是已有证据下的结果，不能声称覆盖真实请求。时间先后是必要线索，仍不能单独证明因果关系。


\section{Producer、service 与 consumer：把热路径和会话控制分开}

\Needspace{7.34cm}
\begin{center}\begin{minipage}{\linewidth}
\centering
\begin{tikzpicture}[x=0.016cm,y=-0.016cm]
\path[use as bounding box] (0,0) rectangle (960,340);
\draw[rounded corners=5pt,draw={rgb,255:red,25;green,118;blue,210},fill={rgb,255:red,25;green,118;blue,210},fill opacity=.075] (40.0,75) rectangle (310.0,187);
\node[anchor=center,text={rgb,255:red,35;green,52;blue,72},font=\fontsize{9.1}{11.0}\selectfont] at (175.0,117.5) {Producer / DataSource};
\node[anchor=center,text={rgb,255:red,35;green,52;blue,72},font=\fontsize{8.2}{9.9}\selectfont] at (175.0,144.5) {TraceWriter / 共享内存};
\draw[draw={rgb,255:red,99;green,115;blue,134},-{Stealth[length=4pt]}] (313.0,130) -- (341.0,130);
\draw[rounded corners=5pt,draw={rgb,255:red,129;green,80;blue,181},fill={rgb,255:red,129;green,80;blue,181},fill opacity=.075] (345.0,75) rectangle (615.0,187);
\node[anchor=center,text={rgb,255:red,35;green,52;blue,72},font=\fontsize{9.1}{11.0}\selectfont] at (480.0,117.5) {Tracing service};
\node[anchor=center,text={rgb,255:red,35;green,52;blue,72},font=\fontsize{8.2}{9.9}\selectfont] at (480.0,144.5) {会话缓冲 / 提交协议};
\draw[draw={rgb,255:red,99;green,115;blue,134},-{Stealth[length=4pt]}] (618.0,130) -- (646.0,130);
\draw[rounded corners=5pt,draw={rgb,255:red,88;green,167;blue,0},fill={rgb,255:red,88;green,167;blue,0},fill opacity=.075] (650.0,75) rectangle (920.0,187);
\node[anchor=center,text={rgb,255:red,35;green,52;blue,72},font=\fontsize{9.1}{11.0}\selectfont] at (785.0,117.5) {Consumer};
\node[anchor=center,text={rgb,255:red,35;green,52;blue,72},font=\fontsize{8.2}{9.9}\selectfont] at (785.0,144.5) {配置 / 读取 / 保存};
\node[anchor=center,text={rgb,255:red,35;green,52;blue,72},font=\fontsize{8.6}{10.4}\selectfont] at (480,235) {写入、commit、flush、stop 是不同阶段};
\node[anchor=center,text={rgb,255:red,35;green,52;blue,72},font=\fontsize{8.6}{10.4}\selectfont] at (480,264) {缓冲策略决定覆盖旧数据还是丢弃新数据等行为};
\end{tikzpicture}
\captionof{figure}{Perfetto 的生产者、服务与消费者。Producer 通过共享内存提交数据；service 管理会话缓冲；consumer 配置和读取。共享内存减少传输开销，不代表全链路零拷贝。}
\end{minipage}\end{center}

Producer 注册数据源并生成事件；consumer 配置与控制采集会话、读取数据；tracing service 居中管理连接、数据源生命周期、会话与中央缓冲区。一个 producer 通常对应一个进程，但同一进程可以有多个独立 producer；每个 producer 有自己的共享内存缓冲区 SMB，由它与 service 共享，同一 producer 内的数据源共同使用它。不能画成所有进程任意访问同一块可写内存。\cite{pf-service}


典型 Unix 连接把控制面和数据面分开。启停、注册、配置、提交通知等走 IPC；事件序列化主要直接写入 SMB，避免每条事件都进行一次普通消息往返。service 收到 commit 后将完成的 chunk 拷入其私有中央缓冲区。这里减少的是 producer 侧的额外临时序列化与消息拷贝，\textbf{不是端到端零拷贝}。使用不支持共享内存的远程传输时，协议允许在 commit 中直接携带 chunk 字节；那是另一条数据路径。\cite{pf-service}\cite{pf-commit}


这种边界首先是所有权与可信度的安排。应用 producer 可能退出、迟迟不交付，甚至给出非法标识和长度；service 不能把共享内存的内容天然当作可信。\texttt{CopyChunkUntrusted} 路径检查提交者可写入的目标缓冲区，并在中央记录中保留由服务侧确定的 producer 身份。数据里的一个数值 ID 不是访问其他会话缓冲区的授权凭证。私有中央存储也使服务能在 producer 回收 SMB 后继续保存数据。\cite{pf-service-impl}\cite{pf-buffer-v1}


一次正常会话可以理解为：consumer 给出配置；service 找到匹配 producer 的数据源并下发设置；producer 创建写入器并生成数据；服务收取 chunk；consumer 停止、等待有关完成通知并读取结果。实际实现还涉及延迟启动、并发会话、超时、断连以及退出前 scraping，不能假定所有应用在开始与结束时刻同时停下。诊断结论需要检查真正覆盖的时间边界，而非只相信请求采集的时长。\cite{pf-service}\cite{pf-service-impl}


读取中央缓冲区通常会推进读游标，不是对同一状态进行无限次幂等查询。后续 SQL 查询针对的是 Trace Processor 已导入的数据，与 consumer 对服务缓冲区的消费操作属于两层。前者可反复运行；后者会影响后续保留、读出与缺口行为。这一区别尤其影响长期流式写文件时的容量计算。\cite{pf-buffer-v1}


\section{共享内存协议：页、chunk、分片与延迟补丁}

\Needspace{7.34cm}
\begin{center}\begin{minipage}{\linewidth}
\centering
\begin{tikzpicture}[x=0.016cm,y=-0.016cm]
\path[use as bounding box] (0,0) rectangle (960,340);
\draw[rounded corners=5pt,draw={rgb,255:red,99;green,115;blue,134},fill={rgb,255:red,99;green,115;blue,134},fill opacity=.075] (30,90) rectangle (234,180);
\node[anchor=center,text={rgb,255:red,35;green,52;blue,72},font=\fontsize{9.1}{11.0}\selectfont] at (132.0,135.0) {Free};
\draw[draw={rgb,255:red,99;green,115;blue,134},-{Stealth[length=4pt]}] (238,135) -- (258,135);
\draw[rounded corners=5pt,draw={rgb,255:red,25;green,118;blue,210},fill={rgb,255:red,25;green,118;blue,210},fill opacity=.075] (262,90) rectangle (466,180);
\node[anchor=center,text={rgb,255:red,35;green,52;blue,72},font=\fontsize{9.1}{11.0}\selectfont] at (364.0,135.0) {BeingWritten};
\draw[draw={rgb,255:red,99;green,115;blue,134},-{Stealth[length=4pt]}] (470,135) -- (490,135);
\draw[rounded corners=5pt,draw={rgb,255:red,88;green,167;blue,0},fill={rgb,255:red,88;green,167;blue,0},fill opacity=.075] (494,90) rectangle (698,180);
\node[anchor=center,text={rgb,255:red,35;green,52;blue,72},font=\fontsize{9.1}{11.0}\selectfont] at (596.0,135.0) {Complete};
\draw[draw={rgb,255:red,99;green,115;blue,134},-{Stealth[length=4pt]}] (702,135) -- (722,135);
\draw[rounded corners=5pt,draw={rgb,255:red,129;green,80;blue,181},fill={rgb,255:red,129;green,80;blue,181},fill opacity=.075] (726,90) rectangle (930,180);
\node[anchor=center,text={rgb,255:red,35;green,52;blue,72},font=\fontsize{9.1}{11.0}\selectfont] at (828.0,135.0) {BeingRead};
\draw[draw={rgb,255:red,99;green,115;blue,134}] (828,186) -- (828,230);
\draw[draw={rgb,255:red,99;green,115;blue,134}] (828,230) -- (132,230);
\draw[draw={rgb,255:red,99;green,115;blue,134},-{Stealth[length=4pt]}] (132,230) -- (132,186);
\node[anchor=center,text={rgb,255:red,35;green,52;blue,72},font=\fontsize{8.6}{10.4}\selectfont] at (480,282) {writer 写入与发布；service 拷贝后释放所有权};
\end{tikzpicture}
\captionof{figure}{chunk 的正常交接状态。典型路径是 Free → BeingWritten → Complete → BeingRead → Free。packet 可跨 chunk，长度补丁还可能延后完成；图未列出全部例外路径。}
\end{minipage}\end{center}

SMB 是生产者到服务之间的\textbf{暂存区}，不是整份 trace 的保留区。其页尺寸在 producer 连接期间固定，要求为 4 KiB 的倍数并有上限；页可按有限方案分成若干 chunk，ABI 最多支持每页 14 个 chunk。页头描述布局与 chunk 状态，chunk 头携带 writer、序号、分片等信息。不能把这里的页直接等同于操作系统缺页单位，也不能任意选择每页一百个小 chunk。\cite{pf-shmem}


一个 chunk 正常由单个 writer 独占写入，常对应一条线程本地事件流。writer 在已取得的 chunk 内前移写指针，像 arena 一样追加数据；用完后申请下一个 chunk，才需要进入生产者内部的协调路径。这个设计使常见写入避免每条事件争用全局锁，但并不意味着整个 SDK、IPC 或分配路径完全无锁。\cite{pf-shmem}\cite{pf-arbiter}\cite{pf-writer}


完成 chunk 的典型状态转移是 \texttt{Free → BeingWritten → Complete → BeingRead → Free}。writer 取得空闲所有权后写字节，发布完成状态；service 取得可读所有权并拷贝，随后释放。状态发布与获取承载跨线程、跨进程的可见性要求，不能只把它当作几个普通布尔字段。共享内存模拟不采用完全相同的 BeingRead 路径；某些补丁优化还会延长 BeingWritten 状态，因此该图表示正常交付的主路径，不是穷尽所有内部转移。\cite{pf-shmem}\cite{pf-arbiter}


TracePacket 与 chunk 没有一一对应关系。一个 chunk 可以包含多个小 packet，一个大 packet 也能跨多个 chunk，甚至在服务持续搬运且策略允许时大于整个 SMB。读取端必须收齐 packet 的有关片段才可交付；“已写出一个 chunk”不等于“已经获得一个独立可解析的事件”。这也解释了为什么丢失少量字节可能让比它更大的逻辑记录失效。\cite{pf-buffer-v1}\cite{pf-writer}


Protobuf 的长度字段造成一个微妙问题：开始写嵌套消息时还不知道最终长度，结束时才能回填。如果嵌套消息跨 chunk，前一个 chunk 可能已经交给 service，甚至其 SMB 位置已被复用。Perfetto 的 writer 为长度预留 varint 空间，记录补丁位置；安全时在本地直接补齐，否则通过 out-of-band patch 修改中央缓冲区中对应 chunk 的预留字节。这是序列化完整性协议，不是允许应用随意修改历史事件含义。\cite{pf-writer}\cite{pf-commit}


\texttt{CommitDataRequest} 的 \texttt{chunks\_\allowbreak{}to\_\allowbreak{}move} 用页号、chunk 号和目标缓冲区定位待搬运数据；\texttt{chunks\_\allowbreak{}to\_\allowbreak{}patch} 用 writer ID、chunk ID、偏移和新字节定位补丁。服务先处理 move，再处理 patch，因此同一次消息可包含同一个 chunk 的交付与回填。\texttt{has\_\allowbreak{}more\_\allowbreak{}patches} 表示是否仍有补丁待到达；未满足完整条件的 chunk 不能因“看上去已有字节”就被读成最终 packet。目标缓冲区、writer 与 chunk 的身份必须同时正确。\cite{pf-commit}\cite{pf-service-impl}


设 packet 的外层长度字段位于 chunk 41，后续内容延伸到 chunk 42。写完 41 时只能交付第一段并标记待补；写完 42 后知道总长度，把补丁发往 41，再允许组装读取。若中央 ring 已覆盖 41，补丁不能复活丢失内容；系统必须保留缺口或丢弃不完整 packet。此例说明 commit 是缓冲区交接，patch 是完成编码，二者都不是应用事务提交。\cite{pf-buffer-v1}\cite{pf-commit}


SMB 满时，writer 需要在数据损失和应用扰动之间选择。v58.2 包含 \texttt{kDrop}、\texttt{kStall} 和 \texttt{kStallThenDrop}：丢弃让应用继续但留下观测缺口；等待可降低瞬时损失却把诊断系统的拥塞传递给应用；先等后丢是折中。还必须区分等待策略的失败边界：这一固定实现的 stall 计数阈值为 300；持续耗尽达到阈值时，\texttt{kStall} 进入 \texttt{PERFETTO\_\allowbreak{}FATAL}，而 \texttt{kStallThenDrop} 返回无效 Chunk。因此 \texttt{kStall} 不承诺无限等待，应用扰动也不只限于延迟增加。该阈值是此版本的实现细节，不能当作稳定 API 或固定超时时间。源码特别处理等待时推进 pending commit 的问题，因为服务和 producer 的 IPC 线程若不能获得进展，简单循环等待空块并不安全。不能把“无丢失模式”当作没有代价的配置。\cite{pf-arbiter}


\section{中央缓冲区、配置与三层数据损失}

\Needspace{7.34cm}
\begin{center}\begin{minipage}{\linewidth}
\centering
\begin{tikzpicture}[x=0.016cm,y=-0.016cm]
\path[use as bounding box] (0,0) rectangle (960,340);
\node[anchor=east,text={rgb,255:red,35;green,52;blue,72},font=\fontsize{8.6}{10.4}\selectfont] at (160,100) {RING\_BUFFER};
\node[anchor=east,text={rgb,255:red,35;green,52;blue,72},font=\fontsize{8.6}{10.4}\selectfont] at (160,215) {DISCARD};
\draw[rounded corners=5pt,draw={rgb,255:red,88;green,167;blue,0},fill={rgb,255:red,88;green,167;blue,0},fill opacity=.075] (205,60) rectangle (360,140);
\node[anchor=center,text={rgb,255:red,35;green,52;blue,72},font=\fontsize{9.1}{11.0}\selectfont] at (282.5,100.0) {3};
\draw[rounded corners=5pt,draw={rgb,255:red,25;green,118;blue,210},fill={rgb,255:red,25;green,118;blue,210},fill opacity=.075] (205,175) rectangle (360,255);
\node[anchor=center,text={rgb,255:red,35;green,52;blue,72},font=\fontsize{9.1}{11.0}\selectfont] at (282.5,215.0) {1};
\draw[rounded corners=5pt,draw={rgb,255:red,88;green,167;blue,0},fill={rgb,255:red,88;green,167;blue,0},fill opacity=.075] (380,60) rectangle (535,140);
\node[anchor=center,text={rgb,255:red,35;green,52;blue,72},font=\fontsize{9.1}{11.0}\selectfont] at (457.5,100.0) {4};
\draw[rounded corners=5pt,draw={rgb,255:red,25;green,118;blue,210},fill={rgb,255:red,25;green,118;blue,210},fill opacity=.075] (380,175) rectangle (535,255);
\node[anchor=center,text={rgb,255:red,35;green,52;blue,72},font=\fontsize{9.1}{11.0}\selectfont] at (457.5,215.0) {2};
\draw[rounded corners=5pt,draw={rgb,255:red,88;green,167;blue,0},fill={rgb,255:red,88;green,167;blue,0},fill opacity=.075] (555,60) rectangle (710,140);
\node[anchor=center,text={rgb,255:red,35;green,52;blue,72},font=\fontsize{9.1}{11.0}\selectfont] at (632.5,100.0) {5};
\draw[rounded corners=5pt,draw={rgb,255:red,25;green,118;blue,210},fill={rgb,255:red,25;green,118;blue,210},fill opacity=.075] (555,175) rectangle (710,255);
\node[anchor=center,text={rgb,255:red,35;green,52;blue,72},font=\fontsize{9.1}{11.0}\selectfont] at (632.5,215.0) {3};
\draw[rounded corners=5pt,draw={rgb,255:red,88;green,167;blue,0},fill={rgb,255:red,88;green,167;blue,0},fill opacity=.075] (730,60) rectangle (885,140);
\node[anchor=center,text={rgb,255:red,35;green,52;blue,72},font=\fontsize{9.1}{11.0}\selectfont] at (807.5,100.0) {6};
\draw[rounded corners=5pt,draw={rgb,255:red,25;green,118;blue,210},fill={rgb,255:red,25;green,118;blue,210},fill opacity=.075] (730,175) rectangle (885,255);
\node[anchor=center,text={rgb,255:red,35;green,52;blue,72},font=\fontsize{9.1}{11.0}\selectfont] at (807.5,215.0) {4};
\node[anchor=center,text={rgb,255:red,35;green,52;blue,72},font=\fontsize{8.6}{10.4}\selectfont] at (480,302) {保留较新数据                         保留较早的未读数据};
\end{tikzpicture}
\captionof{figure}{相同容量，两种保留政策。容量假设为四个等宽块，输入为1至6。示意块不是实际packet；真实可读窗口还依赖分片完整性与增量状态。}
\end{minipage}\end{center}

v58.2 默认使用 TraceBufferV1；V2 是需显式选择的实验模式。本章依据 V1 分析中央缓冲区，不能把 V2 的新标志或恢复行为泛化到所有采集。V1 拷贝 chunk payload，重新打包带可信 producer 身份的头部，以 16 字节对齐存放，并用 \texttt{(ProducerID,\allowbreak{} WriterID,\allowbreak{} ChunkID)} 建立索引。这些额外结构会占空间，所以文件或缓冲区容量不能全部视为事件有效负载。\cite{pf-config}\cite{pf-buffer-v1}


\texttt{RING\_\allowbreak{}BUFFER} 保留较新的数据：写指针追上未读区时按完整 chunk 覆盖旧内容。\texttt{DISCARD} 保留较早的未读数据，空间不足后拒绝新写入。V1 还有容易漏掉的细节：一旦第一次 discard 打断序列，后续写入会持续成为 no-op，即使读者后来追上也不自动恢复为普通可写状态。这不是一个“读快一点就自然恢复”的通用有界队列。\cite{pf-buffer-v1}


中央存储按实际 commit 次序放置 chunk，不保证来自不同 writer 的全局时间顺序。读取时会按各 writer 的 chunk 序号重组 packet、检查碎片与缺口，保证单序列次序；跨 writer 的事件时间排序留给后续处理。某条序列有 1、2、5、6 时，缺口可阻止读者直接把后缀当作连续记录；旧前缀被读走或覆盖后，可读性也会变化。因此“缓冲区里还存在字节”和“现在能输出完整 packet”不是同一个问题。\cite{pf-buffer-v1}


配置应由问题反推。研究尾部卡顿倾向 ring；必须观察启动过程时可能选择 discard；高频调度与低频应用事件可分到不同中央缓冲区，避免前者冲掉后者。下面只展示缓冲区与增量状态政策，实际数据源仍须在相应平台注册并接收配置：



\begin{codeblock}
buffers { size_kb: 16384 fill_policy: RING_BUFFER }
buffers { size_kb: 4096  fill_policy: DISCARD }
duration_ms: 20000
incremental_state_config { clear_period_ms: 1000 }
\end{codeblock}

数据源通过 \texttt{target\_\allowbreak{}buffer} 或受支持的名称引用选择目的地。周期性增量清理只通知声明支持该能力的数据源，不能强迫任意 producer 重新发送所有历史状态。中央缓冲区大小也不是 producer SMB 大小，调整前者不能直接修复后者的瞬时拥塞。\cite{pf-config}


第一层损失发生在内核各 CPU 的 ftrace buffer：探针产生事件后，用户态 \texttt{traced\_\allowbreak{}probes} 来不及读取就可能覆盖。第二层发生在 producer SMB：服务调度延迟或突发写入超过暂存能力。第三层发生在中央缓冲区及其流式落盘过程：容量和读取频率不足，旧内容被覆盖或新内容被拒绝。三层处于不同位置，增大最后一层并不能追回第一层已丢失的 \texttt{sched\_\allowbreak{}switch}。\cite{pf-buffers}


容量估算要同时看平均速率、峰值和最长不可服务时间。中央 16 MiB、平均 2 MiB/s 时，粗略保留窗口为 \(16/2=8\) s，尚未扣除头部、填充与速率波动。若每 5 s 才排空一次，理论产生量为 10 MiB，实际预算还需余量。SMB 峰值 8 MiB/s、服务暂停 10 ms 时，需要约 \(8\times0.010=0.08\) MiB，即 81.92 KiB 的暂存余量；计算应使用峰值而不是全天平均。\cite{pf-buffers}


反例更能揭示风险：每 10 s 突发写入 2 MiB，平均仅 0.2 MiB/s，但 256 KiB SMB 只有单次突发的八分之一。若突发足够快、服务无法同步排空，仍会溢出。这里 256 KiB 是教学假设，不宣称是所有 SDK、后端或版本的默认值。不同官方文档对默认 SMB 的描述有差异，工程配置应读取实际连接参数，而不从一个文档段落猜测。\cite{pf-buffers}\cite{pf-service}


导入后应查询 \texttt{stats}。先检查 \texttt{severity='data\_\allowbreak{}loss' AND value!=0}，再单独关注 \texttt{ftrace\_\allowbreak{}cpu\_\allowbreak{}overrun\_\allowbreak{}end}、\texttt{traced\_\allowbreak{}buf\_\allowbreak{}trace\_\allowbreak{}writer\_\allowbreak{}packet\_\allowbreak{}loss}、\texttt{traced\_\allowbreak{}buf\_\allowbreak{}chunks\_\allowbreak{}overwritten}、\texttt{traced\_\allowbreak{}buf\_\allowbreak{}chunks\_\allowbreak{}discarded} 和增量状态失效统计。仅过滤 data\_\allowbreak{}loss 并不充分：官方文档示例中的 chunks\_\allowbreak{}discarded 为 info，而 ring 覆盖即使符合既定政策，也可能缩短所需分析窗口。统计字段还是计数、字节数或每 CPU 指标，也应分别解释，不能把不同单位相加成“总丢失率”。\cite{pf-buffers}\cite{pf-reader}


Flush 要求有关 producer 推进待提交数据并回报完成状态；它不是让所有应用在同一时刻冻结，也不是磁盘 fsync 的持久化承诺。超时、producer 退出后的 scraping 可以改善尾部可见性，却无法制造未记录的事件。最后应同时核对请求的配置、服务确认、Trace Processor 统计以及分析窗口是否闭合。\cite{pf-service-impl}


\section{TrackEvent、字典状态与时钟：时间线语义从哪里来}

\Needspace{7.34cm}
\begin{center}\begin{minipage}{\linewidth}
\centering
\begin{tikzpicture}[x=0.016cm,y=-0.016cm]
\path[use as bounding box] (0,0) rectangle (960,340);
\draw[rounded corners=5pt,draw={rgb,255:red,25;green,118;blue,210},fill={rgb,255:red,25;green,118;blue,210},fill opacity=.075] (40.0,75) rectangle (310.0,187);
\node[anchor=center,text={rgb,255:red,35;green,52;blue,72},font=\fontsize{9.1}{11.0}\selectfont] at (175.0,117.5) {序列字典};
\node[anchor=center,text={rgb,255:red,35;green,52;blue,72},font=\fontsize{8.2}{9.9}\selectfont] at (175.0,144.5) {iid 7 → 名字 Encode};
\draw[draw={rgb,255:red,99;green,115;blue,134},-{Stealth[length=4pt]}] (313.0,130) -- (341.0,130);
\draw[rounded corners=5pt,draw={rgb,255:red,129;green,80;blue,181},fill={rgb,255:red,129;green,80;blue,181},fill opacity=.075] (345.0,75) rectangle (615.0,187);
\node[anchor=center,text={rgb,255:red,35;green,52;blue,72},font=\fontsize{9.1}{11.0}\selectfont] at (480.0,117.5) {后续 TrackEvent};
\node[anchor=center,text={rgb,255:red,35;green,52;blue,72},font=\fontsize{8.2}{9.9}\selectfont] at (480.0,144.5) {引用 iid 7};
\draw[draw={rgb,255:red,99;green,115;blue,134},-{Stealth[length=4pt]}] (618.0,130) -- (646.0,130);
\draw[rounded corners=5pt,draw={rgb,255:red,88;green,167;blue,0},fill={rgb,255:red,88;green,167;blue,0},fill opacity=.075] (650.0,75) rectangle (920.0,187);
\node[anchor=center,text={rgb,255:red,35;green,52;blue,72},font=\fontsize{9.1}{11.0}\selectfont] at (785.0,117.5) {解码器状态};
\node[anchor=center,text={rgb,255:red,35;green,52;blue,72},font=\fontsize{8.2}{9.9}\selectfont] at (785.0,144.5) {序列有效才可解释};
\node[anchor=center,text={rgb,255:red,35;green,52;blue,72},font=\fontsize{8.6}{10.4}\selectfont] at (480,235) {同一个 iid 在不同序列中不保证相同含义};
\node[anchor=center,text={rgb,255:red,35;green,52;blue,72},font=\fontsize{8.6}{10.4}\selectfont] at (480,264) {文件仍可解析，不意味着每个事件都可完整恢复};
\end{tikzpicture}
\captionof{figure}{增量状态为压缩提供上下文。InternedData 的标识在 packet sequence 等上下文内解释。丢失依赖的状态包会使后续引用无法可靠解析，需等待有效重置。}
\end{minipage}\end{center}

\Needspace{7.34cm}
\begin{center}\begin{minipage}{\linewidth}
\centering
\begin{tikzpicture}[x=0.016cm,y=-0.016cm]
\path[use as bounding box] (0,0) rectangle (960,340);
\draw[rounded corners=5pt,draw={rgb,255:red,25;green,118;blue,210},fill={rgb,255:red,25;green,118;blue,210},fill opacity=.075] (40.0,75) rectangle (310.0,187);
\node[anchor=center,text={rgb,255:red,35;green,52;blue,72},font=\fontsize{9.1}{11.0}\selectfont] at (175.0,117.5) {时钟 A};
\node[anchor=center,text={rgb,255:red,35;green,52;blue,72},font=\fontsize{8.2}{9.9}\selectfont] at (175.0,144.5) {timestamp + clock id};
\draw[draw={rgb,255:red,99;green,115;blue,134},-{Stealth[length=4pt]}] (313.0,130) -- (341.0,130);
\draw[rounded corners=5pt,draw={rgb,255:red,129;green,80;blue,181},fill={rgb,255:red,129;green,80;blue,181},fill opacity=.075] (345.0,75) rectangle (615.0,187);
\node[anchor=center,text={rgb,255:red,35;green,52;blue,72},font=\fontsize{9.1}{11.0}\selectfont] at (480.0,117.5) {同步信息};
\node[anchor=center,text={rgb,255:red,35;green,52;blue,72},font=\fontsize{8.2}{9.9}\selectfont] at (480.0,144.5) {ClockSnapshot / 转换};
\draw[draw={rgb,255:red,99;green,115;blue,134},-{Stealth[length=4pt]}] (618.0,130) -- (646.0,130);
\draw[rounded corners=5pt,draw={rgb,255:red,88;green,167;blue,0},fill={rgb,255:red,88;green,167;blue,0},fill opacity=.075] (650.0,75) rectangle (920.0,187);
\node[anchor=center,text={rgb,255:red,35;green,52;blue,72},font=\fontsize{9.1}{11.0}\selectfont] at (785.0,117.5) {统一分析时间};
\node[anchor=center,text={rgb,255:red,35;green,52;blue,72},font=\fontsize{8.2}{9.9}\selectfont] at (785.0,144.5) {排序 / 区间 / 查询};
\node[anchor=center,text={rgb,255:red,35;green,52;blue,72},font=\fontsize{8.6}{10.4}\selectfont] at (480,235) {先说明时钟来源，再做跨来源差值};
\node[anchor=center,text={rgb,255:red,35;green,52;blue,72},font=\fontsize{8.6}{10.4}\selectfont] at (480,264) {观测间隔小于相对不确定度时，严格先后难以确定};
\end{tikzpicture}
\captionof{figure}{时间戳相同单位，不代表相同时间域。ClockSnapshot 提供不同时间域的关联依据。转换后仍需考虑采集精度、快照覆盖和误差；纳秒存储不等于纳秒准确度。}
\end{minipage}\end{center}

TrackEvent 的 track 表达一条逻辑事件流。线程 track 适合线程内同步嵌套；自定义 track 可以表达跨线程请求的整体生命期；counter 表达数值变化；flow 表达阶段间的关联。设计埋点时应先区分“请求存在了多久”和“线程何时执行哪个阶段”。跨线程请求可以有一条自定义请求 track，再在各线程 track 中记录执行阶段，以 flow 标记派发与交接。\cite{pf-trackevent}


把两个并行执行流放到同一条按栈嵌套的 track，可能制造并不存在的父子包含关系；把 begin 写到线程 A、end 写到线程 B，也不能默认得到正确同步嵌套。不同并行流通常应分轨，并明确关联标识的作用域。flow 箭头记录的是生产者声明的连接，单凭相同 ID 不能证明依赖完整、工作独占或该边一定在关键路径上。埋点语义是分析证据的一部分，不是画图以后自动产生的事实。\cite{pf-trackevent}


频繁重复的名字、类别、栈或默认字段不必每个 packet 都写全。producer 可以先 intern 字符串，后续用整数引用；也可使用增量时间和默认 track 等状态来减小编码成本。这里的 ID 作用于特定 packet sequence，不能把 writer A 的字符串 7 与 writer B 的字符串 7 当作同一个名字。同一 writer 清理增量状态后，也可能重新使用编号。\cite{pf-packet}\cite{pf-reader}


\texttt{trusted\_\allowbreak{}packet\_\allowbreak{}sequence\_\allowbreak{}id} 帮助识别序列，\texttt{SEQ\_\allowbreak{}INCREMENTAL\_\allowbreak{}STATE\_\allowbreak{}CLEARED} 表示状态重建边界，\texttt{SEQ\_\allowbreak{}NEEDS\_\allowbreak{}INCREMENTAL\_\allowbreak{}STATE} 表示该 packet 依赖之前状态。读者发现前序丢失后，会把相应状态标记为无效；需要该状态的后续 packet 可能被跳过，直到清理与重建恢复可解释性。\texttt{previous\_\allowbreak{}packet\_\allowbreak{}dropped} 在序列开始处也可能保守出现，不能把每次标记都机械计作一次已证实的物理丢包。\cite{pf-packet}\cite{pf-reader}


因此，压缩比和可恢复性存在直接权衡。字典引用越多，稳态编码越紧凑，但 ring 覆盖字典后会留下不可解释的后缀。若缓冲区保留原始字节 8 s，期间最近一次有效状态重建只覆盖后 6 s，可使用窗口就不应仍宣称为完整 8 s。周期清理能够缩短某些恢复等待，却增加重发成本，且依赖数据源实现；不能保证任何丢失都在一个周期内无条件恢复。\cite{pf-config}\cite{pf-reader}


时间戳也需要状态。不同源可能使用 BOOTTIME、MONOTONIC、REALTIME 或自定义时钟。MONOTONIC 与 BOOTTIME 对系统休眠的处理不同，REALTIME 可能被校时改变；两个主机的同名时钟也没有天然共同零点。Perfetto 用 ClockSnapshot 建立时钟之间的对应关系，内置 ID、自定义全局 ID 和 sequence-local ID 具有不同作用域。序列局部时钟号不能脱离序列身份解释。\cite{pf-clocks}\cite{pf-reader}


假设同一快照记录时钟 A 为 1000、B 为 2100，在适用区间内只按偏移换算，则 A 上 1104 对应 B 上 2204。这个四则运算示例不证明两个时钟永远同速，也不证明不同主机间已准确同步。漂移、快照稀疏、映射不连通与时钟跳变都会影响对齐。读者应检查可用同步关系与时间范围；源码对暂缺映射的 packet 可以延后处理，最终仍无法转换的情况会记录失败，不能悄悄把它当作纳秒零点相同。\cite{pf-clocks}\cite{pf-reader}


对时钟精度还应给出误差预算。若两个已换算时刻的误差上界分别为 \(\epsilon_1\)、\(\epsilon_2\)，其差值的保守误差上界为二者之和。一个 5 微秒的阶段，若同步误差也在数微秒量级，就不宜据此断言跨源事件的严格先后。数据库统一使用纳秒单位只规定数值尺度，不意味着采集、同步与重建都有一纳秒精度。


\section{Trace Processor：先恢复语义，再排序、建表与计算}

\Needspace{7.34cm}
\begin{center}\begin{minipage}{\linewidth}
\centering
\begin{tikzpicture}[x=0.016cm,y=-0.016cm]
\path[use as bounding box] (0,0) rectangle (960,340);
\draw[rounded corners=5pt,draw={rgb,255:red,25;green,118;blue,210},fill={rgb,255:red,25;green,118;blue,210},fill opacity=.075] (40.0,75) rectangle (310.0,187);
\node[anchor=center,text={rgb,255:red,35;green,52;blue,72},font=\fontsize{9.1}{11.0}\selectfont] at (175.0,117.5) {识别 / tokenize};
\node[anchor=center,text={rgb,255:red,35;green,52;blue,72},font=\fontsize{8.2}{9.9}\selectfont] at (175.0,144.5) {包边界 / 增量状态};
\draw[draw={rgb,255:red,99;green,115;blue,134},-{Stealth[length=4pt]}] (313.0,130) -- (341.0,130);
\draw[rounded corners=5pt,draw={rgb,255:red,129;green,80;blue,181},fill={rgb,255:red,129;green,80;blue,181},fill opacity=.075] (345.0,75) rectangle (615.0,187);
\node[anchor=center,text={rgb,255:red,35;green,52;blue,72},font=\fontsize{9.1}{11.0}\selectfont] at (480.0,117.5) {sort / parse};
\node[anchor=center,text={rgb,255:red,35;green,52;blue,72},font=\fontsize{8.2}{9.9}\selectfont] at (480.0,144.5) {时间顺序 / 语义重建};
\draw[draw={rgb,255:red,99;green,115;blue,134},-{Stealth[length=4pt]}] (618.0,130) -- (646.0,130);
\draw[rounded corners=5pt,draw={rgb,255:red,88;green,167;blue,0},fill={rgb,255:red,88;green,167;blue,0},fill opacity=.075] (650.0,75) rectangle (920.0,187);
\node[anchor=center,text={rgb,255:red,35;green,52;blue,72},font=\fontsize{9.1}{11.0}\selectfont] at (785.0,117.5) {storage / SQL};
\node[anchor=center,text={rgb,255:red,35;green,52;blue,72},font=\fontsize{8.2}{9.9}\selectfont] at (785.0,144.5) {slice / sched / counter};
\node[anchor=center,text={rgb,255:red,35;green,52;blue,72},font=\fontsize{8.6}{10.4}\selectfont] at (480,235) {通用 slice 与内核 sched\_slice 表达不同事实};
\node[anchor=center,text={rgb,255:red,35;green,52;blue,72},font=\fontsize{8.6}{10.4}\selectfont] at (480,264) {dur=-1 的未闭合区间要显式处理，不能当作负耗时};
\end{tikzpicture}
\captionof{figure}{Trace Processor 的处理流水线。不同输入由相应 tokenizer/importer 接受；时间排序和语义解析后落入列式存储，再通过 SQL 表查询。}
\end{minipage}\end{center}

Trace Processor 的输入处理不是“拿到 JSON 就插入 SQLite”。v58.2 的格式识别由 importer registry 管理，按检测优先级调用各导入器的 Sniff；具体可用集合还受构建配置和插件注册影响。压缩、容器和内部格式可能形成多层处理。官方架构文档能解释总体分工，但部分函数位置已变化，源码分析应沿本版本实际入口追踪，而不照抄旧图中的路径。\cite{pf-detection}\cite{pf-importers}\cite{pf-architecture}


对于 Perfetto protobuf，tokenizer 处理外层字节与解压等工作，\texttt{ProtoTraceReader} 负责关键的 packet 语义：识别 sequence、处理丢失与增量清理、更新 defaults 与 interned data、处理时钟快照，再把带相应状态代际的 token 交给排序器。不能先把所有裸 packet 按 timestamp 排序，再解释字典；文件顺序中出现的清理可能导致相同 ID 在后一个代际变成另一个字符串。\cite{pf-reader}


设 sequence S 第一代定义 \texttt{7=Parse}，事件 E 的 timestamp 为 20；随后清理并在第二代定义 \texttt{7=Encode}，事件 F 的 timestamp 为 10。不同时间来源或迟交付会让事件时间与输入顺序不完全一致。读者先按输入推进代际，把 E 绑定第一代、F 绑定第二代，再按时间输出 F、E。若只保存一个可变全局字典，E 可能被错误命名为 Encode。源码传递 \texttt{current\_\allowbreak{}generation()} 正是维持这种解释上下文。\cite{pf-reader}\cite{pf-sequence}


排序器按输入流维护队列，对需要的未排序后缀进行排序，再合并抽取合适的最早事件交给解析器。这样可以利用流内大体有序的性质，又允许跨流到达顺序不同；超出预期的乱序还需要被记录。最终 parser 或 tracker 根据事件重建 slice、进程线程身份、调度状态等语义记录。输入分块读取不意味着总内存恒定：排序可能需要暂存很多事件，导入后的列存储也会增长，pprof 导入器还会缓存整份 profile。\cite{pf-sorter}\cite{pf-pprof}


存储层将大量字段放入专用列式表，重复字符串通过 string pool 共享，再暴露给 SQL 分析层。这样适合按时间、身份、状态和数值进行扫描与连接，不应描述成每个原始字节都转成普通 SQLite 行，也不能据此承诺所有查询零拷贝。表名与视图随版本演进；本章实际 v58.2 查询使用 \texttt{sched\_\allowbreak{}slice}，另有 \texttt{slice}、\texttt{thread\_\allowbreak{}state}、\texttt{thread\_\allowbreak{}track}、\texttt{thread}、\texttt{process}、\texttt{counter}、\texttt{flow}、\texttt{args}、\texttt{stats} 等语义入口。\cite{pf-storage}\cite{pf-request}


进程线程关联应使用 \texttt{upid}、\texttt{utid} 等生命周期身份。原始 pid、tid 可能复用，跨越较长 trace 时按裸 tid 连接会把不同线程实例合并。自定义逻辑 track 也不必然对应一个 OS 线程；必须先通过 track 类型和相应映射确认。查询字段看似都是整数，只说明数据库类型相同，不说明它们属于相同身份空间。\cite{pf-storage}\cite{pf-thread-state}


调度状态来自事件重建：\texttt{sched\_\allowbreak{}switch} 结束旧线程的 Running，按 prev\_\allowbreak{}state 开始它的新状态，并为切入线程打开 Running；waking 将适用的阻塞状态转为 runnable。已经 running 或 runnable 的线程再次收到唤醒，不应简单重置成一次新的等待起点。缺失事件、首尾不闭合都可能造成未知状态，因此查询必须处理负持续时间或边界不完整，而不是无条件对全表求和。\cite{pf-thread-state}


下面是配套已执行查询的核心。先定位请求 slice，再将主线程状态与请求窗口相交；\texttt{dur>=0} 排除未完成区间，重叠谓词确保交集为正，纳秒除以 \(10^6\) 转为毫秒。在真实工作负载中，\texttt{tid=101} 这样的教学筛选应换成已验证的线程生命周期和阶段归属。\cite{pf-request}



\begin{codeblock}
WITH request_window AS (
  SELECT s.ts AS lo, s.ts+s.dur AS hi
  FROM slice s JOIN thread_track tt ON s.track_id=tt.id
  JOIN thread t ON tt.utid=t.utid
  WHERE s.name='request' AND t.tid=101 AND s.dur>=0
), clipped AS (
  SELECT st.state, MIN(st.ts+st.dur,w.hi)-MAX(st.ts,w.lo) AS ns
  FROM thread_state st JOIN thread t USING(utid)
  CROSS JOIN request_window w
  WHERE t.tid=101 AND st.dur>=0
    AND st.ts<w.hi AND st.ts+st.dur>w.lo
)
SELECT state,SUM(ns)/1e6 AS ms FROM clipped GROUP BY state;
\end{codeblock}

真实解析结果为 Running 8 ms、R 2 ms、S 10 ms，合计主线程窗口 20 ms。\texttt{R} 表示具备运行条件，不代表已经消耗 CPU；\texttt{S} 是本例的可中断睡眠状态，不是自动识别出的“等待哪个 Go channel”。应用的 \texttt{join\_\allowbreak{}wait} slice 为 10 ms，调度状态提供其睡眠事实，两者互相补充。Go goroutine 的 runnable 则属于运行时调度层，不能与 OS 线程 R 直接画等号。\cite{pf-thread-state}\cite{pf-request}


对 \texttt{sched\_\allowbreak{}slice} 使用同样相交公式，教学指定的 main 与 worker 分别为 8、15 ms，总和 23 ms。关键路径不能取它们的简单总和：worker 1—16 ms 与 main 前段存在并行，main 后段要等 worker 完成。仅靠线程状态没有完整业务依赖；本例由明确的派发、等待和完成语义补足。实际场景应为每个归属阶段定义窗口，与调度区间相交，并把无法覆盖的时间单独记为未知，而不是强行归到等待或 CPU。


区间连接还要防止重复计费。若同一请求有嵌套的外层阶段和内层阶段，将二者同时与调度表相交再求和，会把相同 CPU 片段计两次。计算请求总成本时应选择互斥阶段，或者先求已归属执行窗口的并集；计算阶段 inclusive 成本则可以保留嵌套，但必须声明不能再把各行直接相加。多条同名 request slice 也应按实例分组，而不是用名字把所有窗口混成一个集合。


缺失一次切换事件的影响也不局限于一行。状态机依赖相邻边界关闭区间，一个边界缺失可能让之后的区间含义不确定。把不完整行过滤掉能避免明显错误，却不会自动补齐覆盖。因此最好同时输出已知 Running、Runnable、Sleeping 的交集长度，以及分析窗口中未被有效状态覆盖的长度；只有覆盖关系成立时，才能要求它们相加恰等于墙钟。


\section{Go trace、Chrome JSON、pprof 与 traceconv 的兼容边界}

\Needspace{7.34cm}
\begin{center}\begin{minipage}{\linewidth}
\centering
\begin{tikzpicture}[x=0.016cm,y=-0.016cm]
\path[use as bounding box] (0,0) rectangle (960,340);
\draw[rounded corners=5pt,draw={rgb,255:red,25;green,118;blue,210},fill={rgb,255:red,25;green,118;blue,210},fill opacity=.075] (35,65) rectangle (300,170);
\node[anchor=center,text={rgb,255:red,35;green,52;blue,72},font=\fontsize{9.1}{11.0}\selectfont] at (167.5,104.0) {pprof 文件};
\node[anchor=center,text={rgb,255:red,35;green,52;blue,72},font=\fontsize{8.2}{9.9}\selectfont] at (167.5,131.0) {栈 + 权重};
\draw[rounded corners=5pt,draw={rgb,255:red,88;green,167;blue,0},fill={rgb,255:red,88;green,167;blue,0},fill opacity=.075] (355,65) rectangle (620,170);
\node[anchor=center,text={rgb,255:red,35;green,52;blue,72},font=\fontsize{9.1}{11.0}\selectfont] at (487.5,104.0) {aggregate\_sample};
\node[anchor=center,text={rgb,255:red,35;green,52;blue,72},font=\fontsize{8.2}{9.9}\selectfont] at (487.5,131.0) {4 / 15 / 4 ms};
\draw[rounded corners=5pt,draw={rgb,255:red,189;green,101;blue,0},fill={rgb,255:red,189;green,101;blue,0},fill opacity=.075] (675,65) rectangle (925,170);
\node[anchor=center,text={rgb,255:red,35;green,52;blue,72},font=\fontsize{9.1}{11.0}\selectfont] at (800.0,104.0) {时间线并未出现};
\node[anchor=center,text={rgb,255:red,35;green,52;blue,72},font=\fontsize{8.2}{9.9}\selectfont] at (800.0,131.0) {slice / sched = 0};
\draw[draw={rgb,255:red,99;green,115;blue,134},-{Stealth[length=4pt]}] (305,115) -- (350,115);
\draw[draw={rgb,255:red,99;green,115;blue,134},-{Stealth[length=4pt]}] (625,115) -- (670,115);
\node[anchor=center,text={rgb,255:red,35;green,52;blue,72},font=\fontsize{8.6}{10.4}\selectfont] at (480,245) {原生 Go trace → 本次 TP：Unknown trace type provided};
\end{tikzpicture}
\captionof{figure}{导入格式不等于恢复所有语义。Perfetto v58.2 可导入 pprof 聚合权重；本实验 slice 与 sched\_slice 仍为空。原生 Go trace 的直接导入失败。}
\end{minipage}\end{center}

Go runtime trace 具有自己的版本化二进制编码，Go 1.25 源码明确使用 \texttt{go 1.\allowbreak{}\%d trace\textbackslash{}x00\textbackslash{}x00\textbackslash{}x00} 文件头。它不是 Perfetto Trace protobuf，也不是 Chrome Trace Event JSON。v58.2 的内置 importer 工厂有 protobuf、JSON、systrace、perf.data、pprof 等入口，没有相应原生 Go runtime trace 工厂。浏览器最终能显示某些相似泳道，不代表底层文件可以直接互换。\cite{pf-go-header}\cite{pf-importers}


配套负例真正运行 Go 1.22.12，生成头部为 \texttt{go 1.\allowbreak{}22 trace\textbackslash{}x00\textbackslash{}x00\textbackslash{}x00}、大小 1727 字节的原生文件，再交给官方 v58.2 Trace Processor。程序退出码为 1，报 \texttt{Unknown trace type provided (ERR:fmt)}。这验证了本版本和本输入的直接导入不成立；不是以扩展名或静态猜测替代运行，也不声称证明所有第三方转换方案均不可用。\cite{pf-validation}


\texttt{go tool trace} 自己先解码 Go runtime 事件，再生成浏览器视图。Go 1.25 的 \texttt{/\allowbreak{}jsontrace} 处理器支持线程视图、goroutine/task 聚焦以及范围选择，输出采用 Chrome Trace Event 的 \texttt{traceEvents}、\texttt{ph}、\texttt{ts}、\texttt{dur}、\texttt{pid}、\texttt{tid} 等字段。若要走这一桥接，应记录生成所用 Go 版本、选择的视图与范围，获取实际 viewer JSON，再用目标 Perfetto 验证所保留事件，而不能把原始 \texttt{.\allowbreak{}trace} 改名 \texttt{.\allowbreak{}json}。\cite{pf-go-json}\cite{pf-go-format}


视图转换存在语义损失。Go viewer 可把 P、goroutine 或其他逻辑实体安排为泳道，JSON 的数字 \texttt{tid} 不必等于内核 OS TID；拿它与 \texttt{sched\_\allowbreak{}slice} 的原始 tid 直接连接，可能得到看似完整的错误结论。已转换的 slice/flow 也不自动保留原始 runtime trace 的全部调度约束、栈、任务关系和未展示事件。能够显示这些视图，不等于 Trace Processor 获得 Go 专用语义分析器。\cite{pf-go-json}\cite{pf-go-format}


Go 1.25 命令行的 \texttt{-\allowbreak{}pprof=net|sync|syscall|sched} 会从 trace 派生相应 pprof 类权重；它们不是通用 CPU profile，转换后也不再保留完整事件顺序。\texttt{-\allowbreak{}d} 的 wire、parsed、footprint 是调试输出模式，不是标准 Chrome JSON 导出开关。教程应区分“调试文本”“浏览器 JSON”与“派生聚合 profile”，避免只因三者都能重定向到文件就称为同一种 trace。\cite{pf-go-main}


另一方面，\textbf{Perfetto v58.2 的确能直接导入 pprof}。读取器解析映射、函数、位置和内联符号；为每个 sample type 建立 aggregate profile，按 pprof 的叶到根 location 顺序反转构建调用位置树，再将各 value 写入 aggregate sample。该路径保留类型与单位，不会读取 Period 后再次乘周期；它也不会由聚合值恢复逐事件时间或调度区间。\cite{pf-pprof}


实际教学 \texttt{request.\allowbreak{}pprof} 中的 cpu/nanoseconds 权重 4,000,000、15,000,000、4,000,000 被真实 v58.2 导入原值；查询 \texttt{\_\allowbreak{}\_\allowbreak{}intrinsic\_\allowbreak{}aggregate\_\allowbreak{}profile} 与 \texttt{\_\allowbreak{}\_\allowbreak{}intrinsic\_\allowbreak{}aggregate\_\allowbreak{}sample} 成功，而 \texttt{sched\_\allowbreak{}slice} 和 \texttt{slice} 行数均为零。内部表名适合本版本核验，不能作为永久稳定公共 API 的承诺。这个实验比“能不能打开”更强：它同时检查了保留的信息和没有产生的信息。\cite{pf-validation}


转换仍需审计：本版本 pprof 读取路径没有处理每个 sample 的 labels；将 int64 权重转换为 double，超过 \(2^{53}\) 后整数精确性可能下降。原文件中业务标签是否重要、整数是否可能很大，都影响这个分析入口能否满足任务。\texttt{traceconv json}、\texttt{traceconv profile} 是支持数据的输出转换模式，不是任意格式输入解码器，也不能把已经聚合的 pprof 逆变换成真实调用时间线。\cite{pf-pprof}\cite{pf-traceconv}


\section{可复现实验、误读检查与研究练习}

复现实验先固定源码与可执行文件身份。配套保存 v58.2 tag、完整提交、release 资料和下载校验；官方 Linux amd64 ZIP 的 SHA-256 为 \texttt{32c739f7\allowbreak{}1b2d3972\allowbreak{}1afd294c\allowbreak{}0b038f24\allowbreak{}99a44a71\allowbreak{}b5b6bcdd\allowbreak{}85b83fac\allowbreak{}a0b240b9}。运行 \texttt{-\allowbreak{}-\allowbreak{}version} 返回同一完整提交。一个值得警惕的实际发现是：此源码 tag 内的 Python \texttt{tools/\allowbreak{}trace\_\allowbreak{}processor} launcher 仍引用 v57.2 下载产物，因此“脚本来自 v58.2”不能替代对最终二进制的版本检查。\cite{pf-release}\cite{pf-launcher}\cite{pf-validation}


在材料包根目录执行 \texttt{python3 examples/\allowbreak{}request/\allowbreak{}verify\_\allowbreak{}perfetto.\allowbreak{}py}，脚本生成合成 ftrace，通过真实 Trace Processor 运行状态、CPU、slice 和 stats 查询，检查八项断言。\texttt{verification/\allowbreak{}request-\allowbreak{}perfetto.\allowbreak{}json} 保存版本、输入摘要、原始 SQL、退出状态、标准输出与解析结果；\texttt{verification/\allowbreak{}go-\allowbreak{}trace-\allowbreak{}import.\allowbreak{}json} 保存原生 Go 负例；\texttt{verification/\allowbreak{}pprof-\allowbreak{}perfetto-\allowbreak{}import.\allowbreak{}json} 保存聚合导入。这里验证的是解析与口径，不是生产环境开销，也没有实测 SMB 压力和内核过载。\cite{pf-request}\cite{pf-validation}


对真实服务重复实验时，应补三类证据：启用哪些数据源和类别、采集窗口与线程身份；观察到哪些丢失和状态恢复；应用阶段如何映射到请求。只有这些条件清楚，SQL 数字才可用于比较。比如优化后 CPU 降低而延迟上升，可能来自并发度、排队或工作量改变；应检查相同工作量下的分布和请求分组，不把一次时间线当作整个服务的统计结论。


常见误读可按证据边界逐项纠正：SMB 不是最终保留文件；减少 writer 拷贝不是端到端零拷贝；Flush 不是全局暂停或 fsync；已有完整字节不代表增量字典仍可解释；flow ID 不是独占归因证明；线程 R 不是 CPU 执行，也不是 Go G 的 runnable；总 CPU 大于墙钟不构成错误；Chrome JSON 的 tid 不保证是 OS TID；直接导入 pprof 不会生成调度时间线；过滤所有 error 后没有记录也不能证明完整无丢失。\cite{pf-service}\cite{pf-packet}\cite{pf-thread-state}\cite{pf-pprof}


采集开销本身也应实验检验。在相同负载、固定构建和相近环境下，分别运行关闭采集、仅启用必要类别、加入高频数据源三组，比较吞吐与延迟分布，同时记录输出速率、CPU 和丢失。等待型 SMB 策略可能减少缺口却增加业务停顿，丢弃型策略可能维持吞吐却降低解释覆盖；两者都需要纳入结论。只报告文件很小或界面很流畅，不能证明生产者热路径开销可以忽略。


练习一：把教学中央容量改为 24 MiB、速率改为 3 MiB/s，说明粗略窗口仍为 8 s；再加入每 2 s 一次 1 MiB 字典重发，计算长期平均速率变为 3.5 MiB/s，粗略窗口约 6.86 s，并指出实际瞬时拥塞仍需单独分析。练习二：给 SQL 增加一个 \texttt{dur=-\allowbreak{}1} 的尾部状态，解释为什么排除后应记录未覆盖时间，而不是把结果重新归一到 100\%。练习三：令同一 OS 工作线程在请求期间执行另一个请求 3 ms，指出原来的全窗口 CPU 查询会高估归因，应改成已识别阶段窗口求交集。


练习四：丢弃字典重建 packet，但保留随后三个引用事件，区分“物理丢失一条 packet”和“语义不可用多条事件”。练习五：给 pprof 加业务 label 和大于 \(2^{53}\) 的整数，比较本版本 aggregate 表与原文件，验证信息保留边界。练习六：为真实 Go 服务同时保存 runtime trace 和 OS 调度数据，设计显式标识与时钟对齐方案，列出仍无法证明的 goroutine 到请求映射，而不要凭名称相似直接关联。完成这些练习，才算掌握从记录协议到诊断结论之间的推理责任。


\chapter{同一请求，三种视角：可复现的跨工具实验}
\label{ch:21}

本章把原理缩小到可以手算的规模。输入是合成教学数据，数字由场景明确设定；解析和报告由真实工具执行。它验证数据格式、聚合和查询口径，不验证内核采样精度、eBPF 开销或生产环境尾延迟。配套目录同时保留生成脚本、输入文件、SQL 和原始输出。\cite{pp-teaching}\cite{pf-request}\cite{pf-validation}


\section{先写清请求、线程和时间}

\Needspace{7.98cm}
\begin{center}\begin{minipage}{\linewidth}
\centering
\begin{tikzpicture}[x=0.016cm,y=-0.016cm]
\path[use as bounding box] (0,0) rectangle (960,380);
\node[anchor=east,text={rgb,255:red,35;green,52;blue,72},font=\fontsize{9.1}{11.0}\selectfont] at (110,90) {main};
\node[anchor=east,text={rgb,255:red,35;green,52;blue,72},font=\fontsize{9.1}{11.0}\selectfont] at (110,195) {worker};
\draw[rounded corners=5pt,draw={rgb,255:red,88;green,167;blue,0},fill={rgb,255:red,88;green,167;blue,0},fill opacity=.075] (140,55) rectangle (292,130);
\node[anchor=center,text={rgb,255:red,35;green,52;blue,72},font=\fontsize{9.1}{11.0}\selectfont] at (216.0,92.5) {运行 4};
\draw[rounded corners=5pt,draw={rgb,255:red,189;green,101;blue,0},fill={rgb,255:red,189;green,101;blue,0},fill opacity=.075] (292,55) rectangle (368,130);
\node[anchor=center,text={rgb,255:red,35;green,52;blue,72},font=\fontsize{9.1}{11.0}\selectfont] at (330.0,92.5) {就绪 2};
\draw[rounded corners=5pt,draw={rgb,255:red,25;green,118;blue,210},fill={rgb,255:red,25;green,118;blue,210},fill opacity=.075] (368,55) rectangle (748,130);
\node[anchor=center,text={rgb,255:red,35;green,52;blue,72},font=\fontsize{9.1}{11.0}\selectfont] at (558.0,92.5) {睡眠 10};
\draw[rounded corners=5pt,draw={rgb,255:red,88;green,167;blue,0},fill={rgb,255:red,88;green,167;blue,0},fill opacity=.075] (748,55) rectangle (900,130);
\node[anchor=center,text={rgb,255:red,35;green,52;blue,72},font=\fontsize{9.1}{11.0}\selectfont] at (824.0,92.5) {运行 4};
\draw[rounded corners=5pt,draw={rgb,255:red,129;green,80;blue,181},fill={rgb,255:red,129;green,80;blue,181},fill opacity=.075] (178,160) rectangle (748,235);
\node[anchor=center,text={rgb,255:red,35;green,52;blue,72},font=\fontsize{9.1}{11.0}\selectfont] at (463.0,197.5) {运行 15 ms};
\draw[draw={rgb,255:red,99;green,115;blue,134}] (140,250) -- (140,260);
\node[anchor=center,text={rgb,255:red,35;green,52;blue,72},font=\fontsize{8.2}{9.9}\selectfont] at (140,285) {0};
\draw[draw={rgb,255:red,99;green,115;blue,134}] (178,250) -- (178,260);
\node[anchor=center,text={rgb,255:red,35;green,52;blue,72},font=\fontsize{8.2}{9.9}\selectfont] at (178,285) {1};
\draw[draw={rgb,255:red,99;green,115;blue,134}] (292,250) -- (292,260);
\node[anchor=center,text={rgb,255:red,35;green,52;blue,72},font=\fontsize{8.2}{9.9}\selectfont] at (292,285) {4};
\draw[draw={rgb,255:red,99;green,115;blue,134}] (368,250) -- (368,260);
\node[anchor=center,text={rgb,255:red,35;green,52;blue,72},font=\fontsize{8.2}{9.9}\selectfont] at (368,285) {6};
\draw[draw={rgb,255:red,99;green,115;blue,134}] (748,250) -- (748,260);
\node[anchor=center,text={rgb,255:red,35;green,52;blue,72},font=\fontsize{8.2}{9.9}\selectfont] at (748,285) {16};
\draw[draw={rgb,255:red,99;green,115;blue,134}] (900,250) -- (900,260);
\node[anchor=center,text={rgb,255:red,35;green,52;blue,72},font=\fontsize{8.2}{9.9}\selectfont] at (900,285) {20};
\draw[draw={rgb,255:red,99;green,115;blue,134}] (140,250) -- (900,250);
\node[anchor=center,text={rgb,255:red,35;green,52;blue,72},font=\fontsize{9.6}{11.6}\selectfont] at (480,332) {墙钟 20 ms；CPU 8 + 15 = 23 ms；平均容量 1.15 核};
\end{tikzpicture}
\captionof{figure}{同一请求的三种时间。合成模型，非内核实测。主线程运行8ms、可运行2ms、睡眠10ms；worker运行15ms。两线程已由模型明确归属于该请求。}
\end{minipage}\end{center}

场景有两个已明确归属于 request=demo 的 OS 线程。请求从相对 0 ms 开始，在 20 ms 完成。main 先执行 4 ms，然后处于 runnable 2 ms，随后 sleeping 10 ms，最后执行 4 ms；worker 在另一 CPU 上从 1 ms 执行到 16 ms。模型声明 main 等待 worker 完成后继续；这些业务归属与依赖是模型输入，不能仅凭 sched\_\allowbreak{}switch 自动发现。



{\small
\begin{longtable}{@{}p{3.665cm}@{\hspace{0.38cm}}p{3.665cm}@{\hspace{0.38cm}}p{3.665cm}@{\hspace{0.38cm}}p{3.665cm}@{}}
\toprule
对象 & 时间区间 & 状态或作用 & 时长 \\
\midrule\endhead
main & [0,4) ms & first\_\allowbreak{}work，运行 & 4 ms \\
main & [4,6) ms & 可运行，尚未调度 & 2 ms \\
main & [6,16) ms & join\_\allowbreak{}wait，睡眠 & 10 ms \\
main & [16,20) ms & finish\_\allowbreak{}work，运行 & 4 ms \\
worker & [1,16) ms & worker\_\allowbreak{}work，运行 & 15 ms \\
\bottomrule\end{longtable}}
手算得到 main CPU 为 8 ms，worker CPU 为 15 ms，请求总 CPU 为 23 ms；墙钟为 20 ms，平均 CPU 容量尺度为 1.15 核。main 的状态分解满足 8+2+10=20 ms。worker 的 15 ms 与 main 的前段运行重叠，因此不能将全部 CPU 相加作为请求延迟。这里所有时间都是教学数值，HTML 动画与 PDF 图共享相同场景。


\section{eBPF 应该怎样观察这个场景}

若通过 eBPF 观察切出事件，只记录当前 tid 和时间戳，可以统计切出频度，却不足以重建上述完整时间表。必须增加 next/prev 身份、切出状态、切入配对及相关唤醒事件，处理任务退出、TID 复用、丢失与边界。通过 perf 采样当前用户栈得到的是正在执行的栈；main 在睡眠时并没有一条相应的 CPU 样本。\cite{eb-sched-trace}\cite{eb-sched-core}\cite{eb-stack-helper}


若使用固定 CPU 时间周期，可在内核 map 中按栈累计采样数，再乘该周期生成适合 pprof 的 CPU 权重。频率反馈或可变周期采样应累计各事件实际代表的权重，不能只用栈计数乘一个全局平均周期；cycles 等硬件事件也不能未经模型换算就称为 CPU 纳秒。若选择向 ring buffer 输出带时钟与身份的逐事件记录，经过正确的协议转换可以供时序分析。两条路径保留的信息不同：前者节约数据量，后者能够支持事件配对，但要承担更高传输与恢复成本。本章没有运行特权 BPF 加载；配套 C/libbpf 示例展示切出观察和失败账本，其覆盖范围比完整调度追踪更窄。\cite{eb-ring-doc}\cite{eb-perf-abi}


\section{pprof 只接收它确实拥有的量}

\Needspace{7.34cm}
\begin{center}\begin{minipage}{\linewidth}
\centering
\begin{tikzpicture}[x=0.016cm,y=-0.016cm]
\path[use as bounding box] (0,0) rectangle (960,340);
\draw[rounded corners=5pt,draw={rgb,255:red,25;green,118;blue,210},fill={rgb,255:red,25;green,118;blue,210},fill opacity=.075] (40.0,75) rectangle (310.0,187);
\node[anchor=center,text={rgb,255:red,35;green,52;blue,72},font=\fontsize{9.1}{11.0}\selectfont] at (175.0,117.5) {场景定义};
\node[anchor=center,text={rgb,255:red,35;green,52;blue,72},font=\fontsize{8.2}{9.9}\selectfont] at (175.0,144.5) {确定归属 / 区间 / 栈};
\draw[draw={rgb,255:red,99;green,115;blue,134},-{Stealth[length=4pt]}] (313.0,130) -- (341.0,130);
\draw[rounded corners=5pt,draw={rgb,255:red,129;green,80;blue,181},fill={rgb,255:red,129;green,80;blue,181},fill opacity=.075] (345.0,75) rectangle (615.0,187);
\node[anchor=center,text={rgb,255:red,35;green,52;blue,72},font=\fontsize{9.1}{11.0}\selectfont] at (480.0,117.5) {两份合成输入};
\node[anchor=center,text={rgb,255:red,35;green,52;blue,72},font=\fontsize{8.2}{9.9}\selectfont] at (480.0,144.5) {request.pprof + ftrace};
\draw[draw={rgb,255:red,99;green,115;blue,134},-{Stealth[length=4pt]}] (618.0,130) -- (646.0,130);
\draw[rounded corners=5pt,draw={rgb,255:red,88;green,167;blue,0},fill={rgb,255:red,88;green,167;blue,0},fill opacity=.075] (650.0,75) rectangle (920.0,187);
\node[anchor=center,text={rgb,255:red,35;green,52;blue,72},font=\fontsize{9.1}{11.0}\selectfont] at (785.0,117.5) {真实工具验证};
\node[anchor=center,text={rgb,255:red,35;green,52;blue,72},font=\fontsize{8.2}{9.9}\selectfont] at (785.0,144.5) {pprof + Trace Processor};
\node[anchor=center,text={rgb,255:red,35;green,52;blue,72},font=\fontsize{8.6}{10.4}\selectfont] at (480,235) {pprof：main 8 / worker 15 ms；CPU 总量 23 ms};
\node[anchor=center,text={rgb,255:red,35;green,52;blue,72},font=\fontsize{8.6}{10.4}\selectfont] at (480,264) {Perfetto：主线程 Running 8 / R 2 / S 10 ms};
\end{tikzpicture}
\captionof{figure}{同一模型经过两条数据通路。两种文件共享教学场景，并非互相无损转换。pprof验证聚合算法；ftrace输入由Perfetto重建调度与区间。}
\end{minipage}\end{center}

生成器写出三条 Sample，唯一数值维度是 cpu/nanoseconds，值为 4,000,000、15,000,000、4,000,000。没有虚构样本次数，也没有写 period=10 ms；因为这三个精确数来自模型，并非 Go 默认 100 Hz 采样。文件级 duration 为 20 ms，三条均有 request=demo 标签。\cite{pp-proto}


主线程的 \texttt{main.\allowbreak{}first\_\allowbreak{}work} 和 \texttt{main.\allowbreak{}finish\_\allowbreak{}work} 都位于 \texttt{main.\allowbreak{}handle\_\allowbreak{}request} 栈下。工作线程的 \texttt{worker.\allowbreak{}work} 位于 \texttt{worker.\allowbreak{}loop} 栈下。工作线程不是主线程当前调用栈的一部分，所以没有人为增加共同的函数根。请求标签连接业务归属，栈仍表达各自执行上下文。



\begin{codeblock}
python3 examples/request/generate_pprof.py --verify
go tool pprof -symbolize=none -top examples/request/request.pprof
\end{codeblock}

本次实际客户端是 Go 1.22.12。查询显示 Duration 20 ms、Total samples 23 ms（115\%）；worker.work 的 flat 为 15 ms，main.first\_\allowbreak{}work 和 main.finish\_\allowbreak{}work 各 4 ms；main.handle\_\allowbreak{}request 的 cum 为 8 ms。按 request=demo 筛选仍保留全部 23 ms。三次查询成功、八项断言通过，证据位于 verification/request-pprof.json。这里关闭符号化，是因为文件只包含构造的逻辑函数与调用栈，未写入 Location.Address 或 Mapping，也没有对应的真实 ELF。


让 worker 在时间轴上平移而保持相同 CPU 权重，普通 profile 仍可能不变。因此 pprof 不能从这三条权重推回 2 ms runnable、10 ms sleeping 或 20 ms 内的执行顺序。文件级 duration 只提供窗口，不替代逐事件时间。


\section{Perfetto 从调度事件重建时间区间}

另一生成器输出符合 ftrace 文本格式的 sched\_\allowbreak{}switch、sched\_\allowbreak{}waking 和 tracing\_\allowbreak{}mark\_\allowbreak{}write。绝对起点选择 10 秒，正文统一显示相对毫秒。main 在 4 ms 切出时状态为 R，在 6 ms 短暂切入后立即以 S 切出；16 ms 被唤醒并重新运行。应用 B/E 标记定义 request、first\_\allowbreak{}work、join\_\allowbreak{}wait 和 finish\_\allowbreak{}work。输入属于合成调度记录，实际重建使用官方 Perfetto v58.2。\cite{pf-request}


下载脚本会校验官方 release ZIP 的 SHA-256 并核对最终二进制版本。材料包不内嵌大型二进制；具备网络时在材料目录执行：



\begin{codeblock}
python3 scripts/fetch_perfetto.py
python3 examples/request/verify_perfetto.py
\end{codeblock}

SQL 先确定请求窗口 \([a,b)\)，再把已完成的区间 \([t,t+d)\) 裁剪为 \(\max(0,\min(t+d,b)-\max(t,a))\)。它避免把请求前后运行算入窗口，但并不自动解决业务归属。配套查询中用 tid 101/102 是受控模型的简化；真实长期追踪应使用生命周期身份与任务阶段。


八项断言全部通过：request=20 ms、join\_\allowbreak{}wait=10 ms；main 的 Running=8 ms、R=2 ms、S=10 ms；sched\_\allowbreak{}slice 中 main CPU=8 ms、worker CPU=15 ms、两者共23 ms。产生的零时长状态不增加总量；dur=-1 的未闭合区间被排除。排除未知区间只避免误算，不能在真实输入中自动保证剩余覆盖完整。\cite{pf-request}


\section{两个兼容性对照实验}

第一个对照直接把 request.pprof 输入 Perfetto v58.2。实际导入成功，聚合表保留 cpu/nanoseconds 和三项权重；slice 与 sched\_\allowbreak{}slice 行数都为零。这说明它支持聚合格式，并没有从 profile 中恢复事件顺序。源码还表明这条导入路径没有处理每条 Sample 的标签，int64 被转为 double；使用前要评估这些信息是否重要。\cite{pf-pprof}\cite{pf-validation}


第二个对照用实际 Go 1.22.12 程序生成原生 execution trace，再直接输入同一 Trace Processor。退出码为 1，报 Unknown trace type provided (ERR:fmt)。Go 1.25 源码的 importer/输出路径阅读与这个运行结果共同支持本章的格式边界；本实验未运行 Go viewer JSON 的完整转换，因此没有把该转换列为已验证能力。\cite{pf-go-header}\cite{pf-go-json}\cite{pf-validation}


第19章还有一组独立的48条基线 Sample 实验，覆盖内联、递归、条件分母、差分、ASLR位置身份和归一化，共15次实际查询、14项语义断言。它与本章三条权重文件是两套独立教材输入，不能把其480 ms总量混入此处23 ms场景。\cite{pp-teaching}


\section{实验完成后应能解释什么}

同一个请求可以具有20 ms墙钟、23 ms CPU、主线程12 ms非运行时间。这三个数来自不同的求和集合，没有矛盾。eBPF可以提供部分系统事件或栈样本；pprof解释权重分布；Perfetto解释有证据支撑的状态与时间关系。任何一个工具都不能补回从未采集的请求身份、丢失的栈或被聚合删除的顺序。


复现时应首先检查 verification 中的版本与输入摘要，再看查询和输出。教学文件可以稳定复算；原生 Go trace 的大小和内容会随运行变化。若换用不同版本，先确认接口与表名，再将结果记录为新的验证，不应把本次结果改写成跨版本保证。


\chapter{组合诊断：让每一种证据回答合适的问题}
\label{ch:22}

\section{从问题选择采集器，而不是从工具寻找问题}

\Needspace{7.34cm}
\begin{center}\begin{minipage}{\linewidth}
\centering
\begin{tikzpicture}[x=0.016cm,y=-0.016cm]
\path[use as bounding box] (0,0) rectangle (960,340);
\draw[rounded corners=5pt,draw={rgb,255:red,25;green,118;blue,210},fill={rgb,255:red,25;green,118;blue,210},fill opacity=.075] (40.0,75) rectangle (310.0,187);
\node[anchor=center,text={rgb,255:red,35;green,52;blue,72},font=\fontsize{9.1}{11.0}\selectfont] at (175.0,117.5) {成本分布不明};
\node[anchor=center,text={rgb,255:red,35;green,52;blue,72},font=\fontsize{8.2}{9.9}\selectfont] at (175.0,144.5) {CPU / heap profile};
\draw[draw={rgb,255:red,99;green,115;blue,134},-{Stealth[length=4pt]}] (313.0,130) -- (341.0,130);
\draw[rounded corners=5pt,draw={rgb,255:red,129;green,80;blue,181},fill={rgb,255:red,129;green,80;blue,181},fill opacity=.075] (345.0,75) rectangle (615.0,187);
\node[anchor=center,text={rgb,255:red,35;green,52;blue,72},font=\fontsize{9.1}{11.0}\selectfont] at (480.0,117.5) {时间原因不明};
\node[anchor=center,text={rgb,255:red,35;green,52;blue,72},font=\fontsize{8.2}{9.9}\selectfont] at (480.0,144.5) {调度 / runtime trace};
\draw[draw={rgb,255:red,99;green,115;blue,134},-{Stealth[length=4pt]}] (618.0,130) -- (646.0,130);
\draw[rounded corners=5pt,draw={rgb,255:red,88;green,167;blue,0},fill={rgb,255:red,88;green,167;blue,0},fill opacity=.075] (650.0,75) rectangle (920.0,187);
\node[anchor=center,text={rgb,255:red,35;green,52;blue,72},font=\fontsize{9.1}{11.0}\selectfont] at (785.0,117.5) {内核边界不明};
\node[anchor=center,text={rgb,255:red,35;green,52;blue,72},font=\fontsize{8.2}{9.9}\selectfont] at (785.0,144.5) {有针对性的 eBPF};
\node[anchor=center,text={rgb,255:red,35;green,52;blue,72},font=\fontsize{8.6}{10.4}\selectfont] at (480,235) {保存身份、构建、单位、窗口、过滤条件和缺失计数};
\node[anchor=center,text={rgb,255:red,35;green,52;blue,72},font=\fontsize{8.6}{10.4}\selectfont] at (480,264) {采样降开销；聚合降传输；每种压缩都有信息代价};
\end{tikzpicture}
\captionof{figure}{按缺失的证据选择下一步。先写出尚不确定的问题，再选择能够增加所缺维度的采集。多个工具给出一致图形也不能替代研究设计。}
\end{minipage}\end{center}

若症状是总 CPU 很高，先固定请求量和输入分布，采集 CPU profile，检查总权重、未知归因、flat/cum 和热点路径。若症状是 CPU 不高而尾延迟变差，应补充请求区间、排队、线程调度或 Go execution trace，区分未获 CPU、同步阻塞、外部依赖与流控。若怀疑内核路径、系统调用、线程唤醒或跨语言栈，才把 eBPF 的观察点设计加入方案。\cite{pp-cpu-api}\cite{pp-graph}



{\small
\begin{longtable}{@{}p{5.013cm}@{\hspace{0.38cm}}p{5.013cm}@{\hspace{0.38cm}}p{5.013cm}@{}}
\toprule
要回答的问题 & 优先证据 & 仍需补足的语义 \\
\midrule\endhead
CPU 花在哪些执行路径 & 带匹配构建信息的 CPU profile & 稳定负载、采样覆盖、单位请求成本 \\
哪些分配路径推动堆增长 & GC 口径一致的 heap profile & 持有路径、RSS 与堆的区别 \\
为什么请求长时间未完成 & 请求区间、调度和运行时事件 & 请求身份、队列与依赖关系 \\
哪种系统行为频繁发生 & 明确 hook 与过滤的 eBPF 观测 & 缺失计数、开销、对象生命周期 \\
前后变化是否真由优化导致 & 对照实验与多种证据交叉检查 & 工作量、错误率、环境和统计波动 \\
\bottomrule\end{longtable}}
工具组合不必意味着所有采集器长期开到最高精度。可先保留低频汇总以发现异常，再触发一小段详细追踪。触发窗口是否覆盖原因、是否只保留故障结果、缓冲保留时长是否足够，都属于实验设计。需要复盘时保存构建身份、配置、时区/时钟说明、采集窗口、筛选条件、原始文件与工具版本，避免只剩无法复算的截图。


\section{端到端数据通路如何搭建}

一种 CPU 归因通路是：perf 事件触发 eBPF 程序 →\allowbreak{} 获取允许的栈信息 →\allowbreak{} 以进程/构建/栈/标签为键累加 →\allowbreak{} 用户态周期性取出计数 →\allowbreak{} 符号化 →\allowbreak{} 编码为 profile.proto →\allowbreak{} 用 pprof 查看。设计者必须选择实际事件来源和周期，解释权重单位，报告栈截断与读取失败。把任意计数写进 cpu/nanoseconds 并不会让它变成 CPU 时间。\cite{pp-proto}


另一种时序通路是：内核调度数据和应用埋点 →\allowbreak{} 事件传输与丢失统计 →\allowbreak{} 带身份及时间域的 trace 包 →\allowbreak{} Perfetto Trace Processor →\allowbreak{} SQL/时间线。若采集器使用 eBPF ring buffer，仍需自己完成到 Perfetto 协议或其支持输入格式的转换。CO-RE 不提供转换格式；ring buffer 不提供跨主机时钟同步；可导入文件也不保证其业务归因正确。


结合两条通路时，优先关联“同一构建、同一进程生命周期、同一时间窗口、同一工作类别”。火焰图中的 Encode 热点可以指导在时间线检查编码区间；时间线发现队列等待可以指导增加对应边界观测。两者是相互约束的解释。若只凭“某段区间里某个函数也很热”就认定该函数造成整段等待，便跨越了证据允许的范围。


\section{六种常见误判及其修正}

第一，把 flame graph 横轴看成时间。宽度表示所选权重，排列通常是布局；应到原始 trace 检查先后。第二，把 cum 当作 self CPU。cum 是含下游的累计归因；应该结合 flat 和调用上下文。第三，把 heap 的 inuse 当作操作系统驻留集。Go 堆剖析有采样和 GC 发布时间，并不覆盖全部 RSS。\cite{pp-flame}\cite{pp-graph}\cite{pp-mprof}


第四，把 off-CPU 全叫 I/O。线程可能等待锁、定时器、网络、用户态调度或其他条件；需要状态、唤醒源和程序边界共同解释。第五，把 verifier 接受当作程序测量正确。验证器关心其模型下的安全性，不能证明筛选条件、单位或统计归因满足研究问题。第六，把一份能解析的 trace 当作完整事实。丢包、增量状态失效、未闭合区间和观测边界都会改变可回答的问题。


每一种修正都要求检查被工具省略的维度：对象、时间、单位、路径上下文或覆盖集合。最有价值的阅读习惯是先找“不再保留什么”，再讨论图形怎样解释。这能把不同工具的界面还原为同一个测量设计问题。


\section{研读源码的顺序}

eBPF 路线从程序加载和 attach 开始，再看 verifier 的寄存器/指针状态与分支收敛，接着看 map 和 ringbuf 的同步协议，最后阅读 BTF/CO-RE 重定位与所用架构的栈路径。每一段都应回答输入、状态、失败分支和输出。仅阅读 README 很难发现 busy 记录阻塞消费、栈展开失败或 helper 可用性约束。


pprof 路线从 profile.proto 的字段和单位开始，回到 Go producer 确认权重如何产生，再看 merge/filter/report。可以使用本专题的合成 profile，逐次改变标签、地址、内联和递归，观察哪些报告变化、哪些不变。小而可手算的反例比大图更能说明算法。\cite{pp-proto}\cite{pp-merge}\cite{pp-filter}\cite{pp-total}


Perfetto 路线先看 producer—service—consumer 的通路，再看 TracePacket 的序列状态与时钟，最后跟进 tokenizer—sorter—parser—storage 的数据流。对每张 SQL 表问清行代表事件、区间还是采样；对每个 dur 检查负值和窗口裁剪。使用可控的小输入先验证自己的查询，再处理大规模生产 trace。


\section{练习与解题线索}

练习一：保持第21章两条线程的 CPU 分别为 8 和 15 ms，把 worker 推迟到主线程计算结束后才开始。CPU profile 是否必须变化？请求墙钟是否可能变化？解题线索：栈权重保留，先后与依赖可能变化；若请求等待 worker 完成，关键路径会改变。


练习二：某次采集 1\% 的事件无法展开栈，能否把所有已知函数权重统一乘以 1/0.99？解题线索：只有缺失对目标量近似独立且其他假设满足时才可能用这种估计。若缺失集中在某个原生库或深栈路径，简单扩张会把偏差传播到所有已知函数；应先保留 unknown 桶并研究缺失机制。


练习三：两个同时采集 10 秒的实例合并后，pprof 文件 duration 为 20 秒，能否说系统经历了 20 秒？解题线索：头部合并使用求和；真实时间并集、实例时间总和、CPU 总量是三种量。\cite{pp-merge}


练习四：调度 trace 中某线程在请求区间内执行 5 ms，能否把 5 ms 全归到该请求？解题线索：线程可能执行其他任务，G 与 M 的映射也可能变化；时间重叠只是候选关联，需要明确的任务边界或标签。


练习五：ring buffer 没有 reserve 失败，能否宣布零丢事件？解题线索：还需检查 hook 覆盖、过滤、采样、栈失败、用户态处理和文件写出等其他环节；“本环节没有报告失败”和“端到端完整”是不同命题。


练习六：同一请求 A 与 B 的调用栈汇总完全相同，能否从 profile 恢复两者的 p99？解题线索：聚合已丢掉请求级持续时间分布，无法从总权重反演分位数。必须保留请求级时间记录，或有另外经过验证的统计模型。


\clearpage
\endgroup
% END APPENDED MECHANISM CHAPTERS
\backmatter
\begin{thebibliography}{999}
\addcontentsline{toc}{chapter}{参考资料}
\bibitem{go-diagnostics} Diagnostics. 2026（读取年份）.\par
\url{https://go.dev/doc/diagnostics}\par
{\small official-doc；访问日期：2026-10-02。诊断工具的分类与组合；持续更新文档，读取年份不代表首次发表年份。}

\bibitem{go-pgo} Profile-guided optimization. 2026（读取年份）.\par
\url{https://go.dev/doc/pgo}\par
{\small official-doc；访问日期：2026-10-02。PGO 默认文件、剖析数据合并和工作负载代表性；持续更新文档。}

\bibitem{go-pgo-blog} Profile-guided optimization in Go 1.21. 2023.\par
\url{https://go.dev/blog/pgo}\par
{\small official-blog；访问日期：2026-10-02。Go 1.\allowbreak{}21 正式支持 PGO；文中的性能数字仅适用于原文评估条件。}

\bibitem{go-memory} A Guide to the Go Garbage Collector. 2026（读取年份）.\par
\url{https://go.dev/doc/gc-guide}\par
{\small official-doc；访问日期：2026-10-02。GC 成本模型、GOGC 和软内存限制；持续更新文档。}

\bibitem{go-testing} testing package (Go 1.25). 2025（Go 1.25.0）.\par
\url{https://pkg.go.dev/testing@go1.25.0}\par
{\small source；访问日期：2026-10-02。testing.\allowbreak{}B.\allowbreak{}Loop、基准计时和分配统计。}

\bibitem{go-proto} profile.proto. 2026（读取年份）.\par
\url{https://github.com/google/pprof/blob/main/proto/profile.proto}\par
{\small source；访问日期：2026-10-02。profile 值向量、sample\_\allowbreak{}type、location 和 string\_\allowbreak{}table；版本固定的实现分析见第十五章。}

\bibitem{go-scheduler} Go runtime scheduler. 2025（Go 1.25.0）.\par
\url{https://github.com/golang/go/blob/go1.25.0/src/runtime/proc.go}\par
{\small source；访问日期：2026-10-02。G/\allowbreak{}M/\allowbreak{}P 模型与工作窃取；属于所读版本的实现细节。}

\bibitem{go-race} Data Race Detector. 2026（读取年份）.\par
\url{https://go.dev/doc/articles/race_detector}\par
{\small official-doc；访问日期：2026-10-02。动态数据竞争检测只覆盖实际执行路径；带检测插桩的运行结果不作为普通性能基准。}

\bibitem{go-model} The Go Memory Model. 2025（Go 1.25.0 源码快照）.\par
\url{https://go.dev/ref/mem}\par
{\small official-doc；访问日期：2026-10-02。happens-\allowbreak{}before 关系与无数据竞争程序的保证；快照年份不代表内存模型首次发布年份。}

\bibitem{go-benchstat} benchstat. 2026（读取年份）.\par
\url{https://pkg.go.dev/golang.org/x/perf/cmd/benchstat}\par
{\small source；访问日期：2026-10-02。多次基准运行的比较与统计摘要；独立于标准库的 golang.\allowbreak{}org/\allowbreak{}x/\allowbreak{}perf 工具。}

\bibitem{go-http} net/http/pprof. 2025（Go 1.25.0）.\par
\url{https://pkg.go.dev/net/http/pprof@go1.25.0}\par
{\small source；访问日期：2026-10-02。HTTP GET 端点、seconds、gc=1 和处理器定义。}

\bibitem{go-pprof} runtime/pprof. 2025（Go 1.25.0）.\par
\url{https://pkg.go.dev/runtime/pprof@go1.25.0}\par
{\small source；访问日期：2026-10-02。CPU、heap、block、mutex 剖析的语义与采样机制。}

\bibitem{go-sync-pool} Go 1.25 sync.Pool：对象可随时移除的接口契约. 2025.\par
\url{https://github.com/golang/go/blob/go1.25.0/src/sync/pool.go#L14-L43}\par
{\small 官方版本源码；访问日期：2026-10-02。Pool用于临时对象复用；对象可随时被自动移除，不保证持久性；首次使用后不可复制。}

\bibitem{proj-pprof} google/pprof：固定提交的驱动与报告实现. 2026.\par
\url{https://github.com/google/pprof/tree/60bf690a9302b860b2110fbb4dccdfbbeb26ab89}\par
{\small primary-source-code；访问日期：2026-10-02。固定提交60bf690a93\allowbreak{}02b860b211\allowbreak{}0fbb4dccdf\allowbreak{}bbeb26ab89；提交时间经API核实为2026-\allowbreak{}10-\allowbreak{}01T06:\allowbreak{}43:\allowbreak{}31Z。v2审读profile/\allowbreak{}merge.\allowbreak{}go:\allowbreak{}135–574、internal/\allowbreak{}driver/\allowbreak{}fetch.\allowbreak{}go:\allowbreak{}37–89；合并错误、差分与输入缺失边界。不是最新版承诺。}

\bibitem{proj-pprof-model} pprof Profile 数据模型与归并实现. 2026.\par
\url{https://github.com/google/pprof/blob/60bf690a9302b860b2110fbb4dccdfbbeb26ab89/profile/profile.go}\par
{\small primary-source-code；访问日期：2026-10-02。审读Profile/\allowbreak{}Sample/\allowbreak{}Location/\allowbreak{}Mapping/\allowbreak{}Function模型、ParseData和CheckValid；对象图和本地ID是合并算法的前提。配合固定提交profile/\allowbreak{}merge.\allowbreak{}go交叉验证。}

\bibitem{proj-delve} Delve v1.24.2：调试器核心与Linux native后端. 2025.\par
\url{https://github.com/go-delve/delve/tree/v1.24.2}\par
{\small primary-source-code；访问日期：2026-10-02。固定v1.\allowbreak{}24.\allowbreak{}2；pkg/\allowbreak{}proc/\allowbreak{}target.\allowbreak{}go:\allowbreak{}65–68缓存失效；pkg/\allowbreak{}proc/\allowbreak{}native/\allowbreak{}proc.\allowbreak{}go:\allowbreak{}395–417固定OS线程ptrace；历史实现不作为Go1.\allowbreak{}25兼容性保证。}

\bibitem{proj-delve-stack} Delve v1.24.2：goroutine栈展开. 2025.\par
\url{https://github.com/go-delve/delve/blob/v1.24.2/pkg/proc/stack.go}\par
{\small primary-source-code；访问日期：2026-10-02。v1.\allowbreak{}24.\allowbreak{}2；实读130–179、380–553：线程/\allowbreak{}保存G寄存器、部分栈错误、内联栈和callpc修正、FDE/\allowbreak{}CFA寄存器恢复。}

\bibitem{proj-pyroscope} Grafana Pyroscope v1.14.0：Distributor写入路径. 2025.\par
\url{https://github.com/grafana/pyroscope/blob/v1.14.0/pkg/distributor/distributor.go}\par
{\small primary-source-code；访问日期：2026-10-02。v1.\allowbreak{}14.\allowbreak{}0；distributor.\allowbreak{}go约390–530：先限流、规范化、可选aggregate、ingester/\allowbreak{}segmentwri\allowbreak{}ter路由。best-\allowbreak{}effort/\allowbreak{}at-\allowbreak{}most-\allowbreak{}once只限定所读可选聚合路径。}

\bibitem{proj-pyroscope-store} Grafana Pyroscope v1.14.0：phlaredb、SymDB与查询. 2025.\par
\url{https://github.com/grafana/pyroscope/blob/v1.14.0/pkg/phlaredb/head.go}\par
{\small primary-source-code；访问日期：2026-10-02。v1.\allowbreak{}14.\allowbreak{}0；head.\allowbreak{}go:\allowbreak{}191–251；另深读symdb/\allowbreak{}partition\_\allowbreak{}memory.\allowbreak{}go:\allowbreak{}1–207、dedup\_\allowbreak{}slice.\allowbreak{}go:\allowbreak{}32–220、rewriter.\allowbreak{}go:\allowbreak{}36–217、stacktrace\_\allowbreak{}tree.\allowbreak{}go:\allowbreak{}16–164与querier/\allowbreak{}querier.\allowbreak{}go:\allowbreak{}578–710。区分符号DAG重写、结构去重、累计差分和展示裁剪。}

\bibitem{proj-parca} Parca v0.23.1：查询与符号化. 2025.\par
\url{https://github.com/parca-dev/parca/tree/v0.23.1}\par
{\small primary-source-code；访问日期：2026-10-02。v0.\allowbreak{}23.\allowbreak{}1；pkg/\allowbreak{}profile/\allowbreak{}schema.\allowbreak{}go:\allowbreak{}38–158物理布局；pkg/\allowbreak{}symbolizer/\allowbreak{}symbolizer.\allowbreak{}go:\allowbreak{}490–525符号后端选择；同时审读pkg/\allowbreak{}parcacol/\allowbreak{}querier.\allowbreak{}go。未部署或测定系统性能。}

\bibitem{proj-parca-store} Parca v0.23.1：ProfileColumnStore与列式schema. 2025.\par
\url{https://github.com/parca-dev/parca/blob/v0.23.1/pkg/profilestore/profilecolumnstore.go}\par
{\small primary-source-code；访问日期：2026-10-02。v0.\allowbreak{}23.\allowbreak{}1；WriteRaw\ensuremath{\rightarrow}\allowbreak{}\allowbreak{}normalizer.\allowbreak{}WriteRawRe\allowbreak{}questToArr\allowbreak{}owRecord\ensuremath{\rightarrow}\allowbreak{}\allowbreak{}ingester.\allowbreak{}Ingest；深读pkg/\allowbreak{}normalizer/\allowbreak{}normalizer.\allowbreak{}go:\allowbreak{}51–230及pkg/\allowbreak{}ingester/\allowbreak{}ingester.\allowbreak{}go:\allowbreak{}24–69，确认类型单位、动态标签与Arrow字典初始化。}

\bibitem{proj-parca-agent} Parca Agent v0.38.1：采集器装配与Arrow报告. 2025.\par
\url{https://github.com/parca-dev/parca-agent/tree/v0.38.1}\par
{\small primary-source-code；访问日期：2026-10-02。v0.\allowbreak{}38.\allowbreak{}1；reporter/\allowbreak{}parca\_\allowbreak{}reporter.\allowbreak{}go:\allowbreak{}205–264、532–548、969–1175、1280–1286；TraceHash/\allowbreak{}LRU、CPU count与offCPU时间、Arrow样本/\allowbreak{}按需补栈协议、is\_\allowbreak{}complete标记。}

\bibitem{proj-otel-ebpf} Parca使用的OpenTelemetry eBPF profiler固定fork. 2025.\par
\url{https://github.com/parca-dev/opentelemetry-ebpf-profiler/tree/edfa6253a8b2}\par
{\small primary-source-code；访问日期：2026-10-02。Agent go.\allowbreak{}mod替换到Parca fork edfa6253a8b2；审读interpreter/\allowbreak{}golang/\allowbreak{}golang.\allowbreak{}go:\allowbreak{}60–85及runtime\_\allowbreak{}data.\allowbreak{}go:\allowbreak{}1–35，标签latestVers\allowbreak{}ion=go1.\allowbreak{}24且未知版本回退告警。Go标签与native展开需独立验证。}

\bibitem{front-go121} Go 1.21 Release Notes：PGO正式可用. 2023.\par
\url{https://go.dev/doc/go1.21}\par
{\small official-release-notes；访问日期：2026-10-02。Go1.\allowbreak{}21发布说明：PGO正式支持、-\allowbreak{}pgo=auto/\allowbreak{}default.\allowbreak{}pgo、带保护去虚化等。工程发布材料，不与同行评审论文混为一类；本书Go迁移实验为设计。}

\bibitem{front-trace2024} More powerful Go execution traces. 2024.\par
\url{https://go.dev/blog/execution-traces-2024}\par
{\small official-engineering-article；访问日期：2026-10-02。官方文章2024-\allowbreak{}03-\allowbreak{}14：Go1.\allowbreak{}22追踪重写和generation/\allowbreak{}分段方向；文章当时go tool trace仍整体加载，不能把格式潜力当作当时全部工具能力。}

\bibitem{front-flight2025} Flight Recorder in Go 1.25. 2025.\par
\url{https://go.dev/blog/flight-recorder}\par
{\small official-engineering-article；访问日期：2026-10-02。官方文章2025-\allowbreak{}09-\allowbreak{}26：异常触发前保留trace窗口；与Go1.\allowbreak{}25.\allowbreak{}0源码交叉核实。选中工程资料最晚明确发布日期，不声称2026无其他研究。}

\bibitem{front-flight-source} Go 1.25.0 runtime/trace FlightRecorder实现与API契约. 2025.\par
\url{https://github.com/golang/go/blob/go1.25.0/src/runtime/trace/flightrecorder.go}\par
{\small primary-source-code；访问日期：2026-10-02。固定Go1.\allowbreak{}25.\allowbreak{}0；flightreco\allowbreak{}rder.\allowbreak{}go完整接口及157–175配置语义：MaxBytes优先MinAge但为hint，非WriteTo文件或RSS硬上限；单记录器、可与trace.\allowbreak{}Start并存、并发WriteTo受限。}

\bibitem{front-otel2024} OpenTelemetry announces support for profiling. 2024.\par
\url{https://opentelemetry.io/blog/2024/profiling/}\par
{\small official-project-announcement；访问日期：2026-10-02。官方2024-\allowbreak{}03-\allowbreak{}19介绍Profiles方向。区分计划/\allowbreak{}动机与后续规范完成程度，不据公告声称所有后端兼容。}

\bibitem{front-otel-spec} OpenTelemetry Profiles Specification. 2026.\par
\url{https://opentelemetry.io/docs/specs/otel/profiles/}\par
{\small living-specification-snapshot；访问日期：2026-10-02。2026-\allowbreak{}10-\allowbreak{}02实读规范快照（页面版本1.\allowbreak{}61.\allowbreak{}0）仍Alpha；Resource/\allowbreak{}Scope、字典、trace/\allowbreak{}span links、original\_\allowbreak{}payload。访问时间不等于论文发表时间。}

\bibitem{paper-scalene} Emery D. Berger, Sam Stern, Juan Altmayer Pizzorno. Triangulating Python Performance Issues with Scalene. OSDI 2023.. 2023.\par
\url{https://www.usenix.org/conference/osdi23/presentation/berger}\par
{\small peer-reviewed-conference-paper；访问日期：2026-10-02。OSDI2023原始PDF全文方法/\allowbreak{}实验/\allowbreak{}讨论已读，布局文本1–1030；重点§2–3/\allowbreak{}图3–4、§6/\allowbreak{}表1–3/\allowbreak{}图5–8、§7–8。表2样本中位18倍、pprint7976/\allowbreak{}23、sympy6757/\allowbreak{}10；表3完整模式1.\allowbreak{}32x中位、pprint4.\allowbreak{}03x。印刷页56泄漏公式与变量定义存在可演算一致性疑问，未宣称源码bug。所有性能数是作者原实验，不是本次复现。}

\bibitem{paper-hardware} Aarati Kakaraparthy, Jignesh M. Patel. Fine-Grained Hardware Profiling – Are You Using the Right Tools? SIGMOD Record 53(2), 2024.. 2024.\par
\url{https://doi.org/10.1145/3685980.3685986}\par
{\small research-journal-article；访问日期：2026-10-02。SIGMOD Record53(2)，2024年6月，pp38–43；DOI在线元数据2024-\allowbreak{}07-\allowbreak{}30。原PDF布局文本1–390实读，Listings2–4、图1–5。图1精度23.\allowbreak{}9/\allowbreak{}99.\allowbreak{}3/\allowbreak{}99.\allowbreak{}7\%；图2对无工具+40/\allowbreak{}+3/\allowbreak{}+2\%；图4对自身粗粒度+268/\allowbreak{}+8\%；图5对最佳策略+20/\allowbreak{}+25/\allowbreak{}+10\%，分母必须分别注明。}

\bibitem{paper-epping} Simon Sundberg, Anna Brunstrom, Simone Ferlin-Reiter, Toke Høiland-Jørgensen, Jesper Dangaard Brouer. Efficient Continuous Latency Monitoring with eBPF. PAM 2023.. 2023.\par
\url{https://doi.org/10.1007/978-3-031-28486-1_9}\par
{\small peer-reviewed-conference-paper；访问日期：2026-10-02。PAM2023 LNCS13882 pp191–208，在线2023-\allowbreak{}03-\allowbreak{}10。原PDF布局文本1–640实读含附录A、图8–9、Algorithm1。图4共同ACK子集；图5CPU47/\allowbreak{}57/\allowbreak{}77\%且PPing仅观测约6\%；图6单核1.\allowbreak{}53\ensuremath{\rightarrow}\allowbreak{}\allowbreak{}1.\allowbreak{}13Mpps；图7限流/\allowbreak{}直方图。Zenodo7555\allowbreak{}410在原文可见，但归档访问失败，未核验代码内容。}

\bibitem{proj-pprof-merge} google/pprof：语义合并与对象重编号源码. 2026.\par
\url{https://github.com/google/pprof/blob/60bf690a9302b860b2110fbb4dccdfbbeb26ab89/profile/merge.go}\par
{\small primary-source-code；访问日期：2026-10-02。实读135–574；mapSample155–193、sampleKey196–245、mapLocatio\allowbreak{}n284–337、mapMapping\allowbreak{}346–409、函数身份426–458、头部468–544、ID缓存549–574。教材代码片段逐字对照。}

\bibitem{proj-pyro-tree} Pyroscope v1.14.0：stacktrace tree. 2025.\par
\url{https://github.com/grafana/pyroscope/blob/v1.14.0/pkg/phlaredb/symdb/stacktrace_tree.go}\par
{\small primary-source-code；访问日期：2026-10-02。实读16–164：parent/\allowbreak{}reference/\allowbreak{}firstchild/\allowbreak{}nextsibling，反向插入，89–101父指针恢复，持久化只存p/\allowbreak{}r。对照partition\_\allowbreak{}memory.\allowbreak{}go:\allowbreak{}63–149锁、缓存与切片头。复杂度是本书根据控制流推导。}

\bibitem{proj-pyro-delta} Pyroscope v1.14.0：累计profile差分. 2025.\par
\url{https://github.com/grafana/pyroscope/blob/v1.14.0/pkg/phlaredb/delta.go}\par
{\small primary-source-code；访问日期：2026-10-02。实读21–99；首次基线为空、计数下降重置；仅memory alloc\_\allowbreak{}objects和alloc\_\allowbreak{}space的该路径。100\ensuremath{\rightarrow}\allowbreak{}\allowbreak{}160\ensuremath{\rightarrow}\allowbreak{}\allowbreak{}20\ensuremath{\rightarrow}\allowbreak{}\allowbreak{}45为教学输入。 多栈输入的原地相减后重置可污染重建基线：原始delta.\allowbreak{}go逐字复制的隔离Go测试已复现该序列；Samples方法保持原文，外围标签/\allowbreak{}类型为桩，未执行完整上游测试或生产部署。}

\bibitem{proj-parca-query} Parca v0.23.1：列式查询与count乘period. 2025.\par
\url{https://github.com/parca-dev/parca/blob/v0.23.1/pkg/parcacol/querier.go}\par
{\small primary-source-code；访问日期：2026-10-02。实读QueryRange\allowbreak{}341–381、queryRange\allowbreak{}Delta405–511及selectMerge/\allowbreak{}QueryMerge\allowbreak{}1367–1498。count逐行乘period再SUM；过滤时间边界因具体接口不同。}

\bibitem{proj-ebpf-unwind} OTel eBPF Parca fork：有界native栈展开. 2025.\par
\url{https://github.com/parca-dev/opentelemetry-ebpf-profiler/blob/edfa6253a8b2/support/ebpf/native_stack_trace.ebpf.c}\par
{\small primary-source-code；访问日期：2026-10-02。实读137–217、575–646；64KiB页、有界二分前驱查找、先压当前帧后展开、错误和tail call。非空栈不等于完整。}

\bibitem{proj-go-unwind} OTel eBPF Parca fork：Go pclntab转换为展开规则. 2025.\par
\url{https://github.com/parca-dev/opentelemetry-ebpf-profiler/blob/edfa6253a8b2/nativeunwind/elfunwindinfo/elfgopclntab.go}\par
{\small primary-source-code；访问日期：2026-10-02。实读480–645；Go函数表、PCSP、x86 FP/\allowbreak{}SP策略与ARM64历史实现注释。教材CFA例子为架构限定简化推导，不是Go ABI固定布局。}

\bibitem{tools-pprof} Go 1.25 runtime/pprof：CPU、heap、block、mutex 语义. 2025.\par
\url{https://github.com/golang/go/blob/go1.25.0/src/runtime/pprof/pprof.go}\par
{\small 官方版本源码；访问日期：2026-10-02。Go 1.\allowbreak{}25.\allowbreak{}0；CPU 名义采样频率为 100 Hz；heap 的四种样本类型；block 记录等待栈，mutex 记录释放栈；单个进程内同时只能运行一个 runtime/\allowbreak{}pprof CPU 采集器。}

\bibitem{tools-mprof} Go 1.25 runtime 内存与同步剖析实现. 2025.\par
\url{https://github.com/golang/go/blob/go1.25.0/src/runtime/mprof.go}\par
{\small 官方版本源码；访问日期：2026-10-02。MemProfile\allowbreak{}Rate 默认 512 KiB；采样率应尽早只改一次；MemProfile 最多滞后两个 GC 周期；block 按时长采样、mutex 按事件采样及权重校正。}

\bibitem{tools-http-pprof} Go 1.25 net/http/pprof 参数与处理器. 2025.\par
\url{https://github.com/golang/go/blob/go1.25.0/src/net/http/pprof/pprof.go}\par
{\small 官方版本源码；访问日期：2026-10-02。debug/\allowbreak{}gc/\allowbreak{}seconds 参数、CPU 与 trace 窗口、其他 profile 的差分；手动注册自定义 mux；默认 CPU 30 秒。}

\bibitem{tools-labels} Go 1.25 pprof Labels 与 WithLabels. 2025.\par
\url{https://github.com/golang/go/blob/go1.25.0/src/runtime/pprof/label.go}\par
{\small 官方版本源码；访问日期：2026-10-02。明确只有 CPU 与 goroutine profiles 使用 labels；WithLabels 只返回带标签的 context。}

\bibitem{tools-label-runtime} Go 1.25 pprof.Do 与 goroutine labels 继承. 2025.\par
\url{https://github.com/golang/go/blob/go1.25.0/src/runtime/pprof/runtime.go}\par
{\small 官方版本源码；访问日期：2026-10-02。Do 安装增强标签并在回调后恢复；新建 goroutine 继承创建时标签。}

\bibitem{tools-gc-guide} A Guide to the Go Garbage Collector. 2026.\par
\url{https://go.dev/doc/gc-guide}\par
{\small 官方持续更新指南；访问日期：2026-10-02。通过 golang/\allowbreak{}website 的 \_\allowbreak{}content/\allowbreak{}doc/\allowbreak{}gc-\allowbreak{}guide.\allowbreak{}html 核查；年份为读取版本年份，不代表首次发表。GC 目标公式包含根；GOMEMLIMIT 为软限制。}

\bibitem{tools-metrics} Go 1.25 runtime/metrics 指标定义. 2025.\par
\url{https://github.com/golang/go/blob/go1.25.0/src/runtime/metrics/description.go}\par
{\small 官方版本源码；访问日期：2026-10-02。名称、类型、Cumulative；调度延迟、锁等待与内存口径；CPU classes total 为 GOMAXPROCS 积分的估计可用容量。}

\bibitem{tools-metrics-read} Go 1.25 runtime/metrics.Read 与存储复用. 2025.\par
\url{https://github.com/golang/go/blob/go1.25.0/src/runtime/metrics/sample.go}\par
{\small 官方版本源码；访问日期：2026-10-02。未知指标返回 KindBad；重复 Read 可复用直方图存储；跨 goroutine 使用必须深拷贝相关数据。}

\bibitem{tools-go122} Go 1.22 Release Notes：mutex 与 execution tracer. 2024.\par
\url{https://go.dev/doc/go1.22}\par
{\small 官方发行说明；访问日期：2026-10-02。通过 golang/\allowbreak{}website 对应官方文档源码核查。mutex 按等待者数量缩放；新增内部锁竞争；trace 重写和代际分段。}

\bibitem{tools-go125} Go 1.25 Release Notes：FlightRecorder 与容器感知. 2025.\par
\url{https://go.dev/doc/go1.25}\par
{\small 官方发行说明；访问日期：2026-10-02。通过 golang/\allowbreak{}website 对应官方文档源码核查。标准库 FlightReco\allowbreak{}rder 与容器相关 GOMAXPROCS 改进；不将 1.\allowbreak{}25 声称为 2026 年最新版本。}

\bibitem{tools-trace2024} More powerful Go execution traces. 2024.\par
\url{https://go.dev/blog/execution-traces-2024}\par
{\small 官方技术文章；访问日期：2026-10-02。2024 年 trace 重写、开销改进、分段、实验 flight recorder 和 reader API。文中的性能数据仅代表其测试背景。}

\bibitem{tools-flight-blog} Flight Recorder in Go 1.25. 2025.\par
\url{https://go.dev/blog/flight-recorder}\par
{\small 官方技术文章；访问日期：2026-10-02。2025-\allowbreak{}09-\allowbreak{}26；Go 1.\allowbreak{}25 标准库飞行记录和异常触发案例。容量与并发细节以固定源码文档为准。}

\bibitem{tools-flight-source} Go 1.25 FlightRecorder 实现与并发契约. 2025.\par
\url{https://github.com/golang/go/blob/go1.25.0/src/runtime/trace/flightrecorder.go}\par
{\small 官方版本源码；访问日期：2026-10-02。一个 recorder 可与 trace.\allowbreak{}Start 并存；WriteTo 单写入者；MaxBytes 是提示，不保证内存或输出硬上限；容量优先于 MinAge。}

\bibitem{tools-trace-doc} Go 1.25 runtime/trace 文档与入口. 2025.\par
\url{https://github.com/golang/go/blob/go1.25.0/src/runtime/trace/trace.go}\par
{\small 官方版本源码；访问日期：2026-10-02。调度、syscall、GC 与用户注解事件；Start/\allowbreak{}Stop 及写出完成语义。}

\bibitem{tools-trace-annotation} Go 1.25 task、region、log 注解实现. 2025.\par
\url{https://github.com/golang/go/blob/go1.25.0/src/runtime/trace/annotation.go}\par
{\small 官方版本源码；访问日期：2026-10-02。Task 跨 goroutine；region 绑定调用 goroutine；类别应有限；Logf 内部格式化受 IsEnabled 保护。}

\bibitem{tools-trace-cmd} Go 1.25 go tool trace 命令实现. 2025.\par
\url{https://github.com/golang/go/blob/go1.25.0/src/cmd/trace/main.go}\par
{\small 官方版本源码；访问日期：2026-10-02。-\allowbreak{}pprof 支持 net、sync、syscall、sched；现代 trace 不必额外提供二进制；查看器浏览器支持说明。}

\bibitem{tools-delve-debug} Delve v1.25.2：dlv debug. 2025.\par
\url{https://github.com/go-delve/delve/blob/v1.25.2/Documentation/usage/dlv_debug.md}\par
{\small 项目官方版本文档；访问日期：2026-10-02。debug 构建关闭优化的程序；Go 版本检查；headless 与监听参数。}

\bibitem{tools-delve-exec} Delve v1.25.2：dlv exec. 2025.\par
\url{https://github.com/go-delve/delve/blob/v1.25.2/Documentation/usage/dlv_exec.md}\par
{\small 项目官方版本文档；访问日期：2026-10-02。已有二进制调试；建议本地调试构建 -\allowbreak{}gcflags=all=-\allowbreak{}N -\allowbreak{}l；默认 loopback 与 only-\allowbreak{}same-\allowbreak{}user。}

\bibitem{tools-delve-core} Delve v1.25.2：dlv core 支持矩阵. 2025.\par
\url{https://github.com/go-delve/delve/blob/v1.25.2/Documentation/usage/dlv_core.md}\par
{\small 项目官方版本文档；访问日期：2026-10-02。Linux amd64/\allowbreak{}arm64 core、Windows amd64 minidump 及 Delve dump；必须匹配可执行文件。}

\bibitem{tools-delve-cli} Delve v1.25.2 命令手册. 2025.\par
\url{https://github.com/go-delve/delve/blob/v1.25.2/Documentation/cli/README.md}\par
{\small 项目官方版本文档；访问日期：2026-10-02。condition、goroutine、frame、tracepoint、dump 等精确语法；goroutine 条件表达式示例；dump 的平台格式。}

\bibitem{tools-runtime-debug} Go 1.25 runtime 环境诊断参数. 2025.\par
\url{https://github.com/golang/go/blob/go1.25.0/src/runtime/extern.go}\par
{\small 官方版本源码；访问日期：2026-10-02。GODEBUG gctrace/\allowbreak{}schedtrace/\allowbreak{}scavtrace 与 GOTRACEBACK；crash 在 Unix 使用 SIGABRT；输出格式可能变化。}

\bibitem{tools-perf} Linux v6.14 perf record 官方手册. 2025.\par
\url{https://github.com/torvalds/linux/blob/v6.14/tools/perf/Documentation/perf-record.txt}\par
{\small 内核项目官方版本文档；访问日期：2026-10-02。进程/\allowbreak{}线程目标、采样频率、FP/\allowbreak{}DWARF/\allowbreak{}LBR 调用图模式及相应限制。}

\bibitem{tools-strace} strace v6.14 官方手册. 2025.\par
\url{https://github.com/strace/strace/blob/v6.14/doc/strace.1.in}\par
{\small 项目官方版本文档；访问日期：2026-10-02。-\allowbreak{}T syscall 持续时间；-\allowbreak{}c 系统时间估计与 -\allowbreak{}w 墙钟汇总差异；事件过滤。}

\bibitem{tools-ebpf} Linux v6.14 BPF Design Q\&A. 2025.\par
\url{https://github.com/torvalds/linux/blob/v6.14/Documentation/bpf/bpf_design_QA.rst}\par
{\small 内核项目官方版本文档；访问日期：2026-10-02。BPF 应用领域、内核执行与验证器边界；本章只使用通用设计事实，不引用文档中的历史限制为最新能力。}

\bibitem{tools-pprof-cli} google/pprof 官方使用文档. 2026.\par
\url{https://github.com/google/pprof/blob/main/doc/README.md}\par
{\small 项目官方持续更新文档；访问日期：2026-10-02。2026-\allowbreak{}10-\allowbreak{}02 读取；年份为读取版本年份。核查 flat/\allowbreak{}cum、递归去重、sample\_\allowbreak{}index、tagfocus、base/\allowbreak{}diff\_\allowbreak{}base 的分母与负值。}

\bibitem{tools-runtime-debug-api} Go 1.25 runtime/debug GC、内存限制与堆转储 API. 2025.\par
\url{https://github.com/golang/go/blob/go1.25.0/src/runtime/debug/garbage.go}\par
{\small 官方版本源码；访问日期：2026-10-02。SetGCPercent、SetMemoryL\allowbreak{}imit、FreeOSMemory、ReadGCStats、SetTraceback 与 WriteHeapD\allowbreak{}ump；负参数语义不同；完整堆转储暂停全部 goroutine。}

\bibitem{tools-fgprof} fgprof v0.9.5 官方 README. 2024.\par
\url{https://github.com/felixge/fgprof/blob/v0.9.5/README.md}\par
{\small 官方固定版本文档与源码；访问日期：2026-10-02。2024-\allowbreak{}08-\allowbreak{}30 发布；混合 CPU/\allowbreak{}I/\allowbreak{}O 的 wall-\allowbreak{}clock profiling；goroutine 数与开销、旧 Go 暂停和调度偏差限制。}

\bibitem{tools-fgprof-source} fgprof v0.9.5 采样实现. 2024.\par
\url{https://github.com/felixge/fgprof/blob/v0.9.5/fgprof.go}\par
{\small 官方固定版本文档与源码；访问日期：2026-10-02。固定版本；99 Hz ticker、runtime.\allowbreak{}GoroutineP\allowbreak{}rofile、实际采样率修正及 Go 1.\allowbreak{}23.\allowbreak{}0 兼容处理。}

\bibitem{tools-gops} google/gops v0.3.28 官方 README. 2023.\par
\url{https://github.com/google/gops/blob/v0.3.28/README.md}\par
{\small 官方固定版本文档与源码；访问日期：2026-10-02。2023-\allowbreak{}08-\allowbreak{}10 发布；进程发现与 agent 诊断区分，本地同用户模式、远程地址模式、gc/\allowbreak{}setgc/\allowbreak{}profile/\allowbreak{}trace 能力。}

\bibitem{tools-gops-agent} google/gops v0.3.28 agent 源码. 2023.\par
\url{https://github.com/google/gops/blob/v0.3.28/agent/agent.go}\par
{\small 官方固定版本文档与源码；访问日期：2026-10-02。默认 127.\allowbreak{}0.\allowbreak{}0.\allowbreak{}1:\allowbreak{}0；端口发现配置；TCP 端点可能被任何同机程序访问；不能将发现约定视为认证。}

\bibitem{tools-expvar} Go 1.25 expvar 官方实现. 2025.\par
\url{https://github.com/golang/go/blob/go1.25.0/src/expvar/expvar.go}\par
{\small 官方固定版本文档与源码；访问日期：2026-10-02。JSON /\allowbreak{}debug/\allowbreak{}vars；Go 1.\allowbreak{}22+ GET；cmdline/\allowbreak{}memstats 默认导出；原子变量和回调的边界。}

\bibitem{tools-nm} Go 1.25 go tool nm 官方手册. 2025.\par
\url{https://github.com/golang/go/blob/go1.25.0/src/cmd/nm/doc.go}\par
{\small 官方固定版本文档与源码；访问日期：2026-10-02。符号地址、类型、名称；-\allowbreak{}size 与 -\allowbreak{}sort size 从大到小排序。}

\bibitem{tools-objdump} Go 1.25 go tool objdump 官方实现与用法. 2025.\par
\url{https://github.com/golang/go/blob/go1.25.0/src/cmd/objdump/main.go}\par
{\small 官方固定版本文档与源码；访问日期：2026-10-02。反汇编可执行文件；-\allowbreak{}s 正则过滤符号、-\allowbreak{}S 显示 Go 源码。}

\bibitem{tools-crash-output} Go 1.23 SetCrashOutput 首发版本源码. 2024.\par
\url{https://github.com/golang/go/blob/go1.23.0/src/runtime/debug/stack.go}\par
{\small 官方固定版本文档与源码；访问日期：2026-10-02。单一额外 fatal 文本输出；保留 stderr；复制 FD；替换并发崩溃可能仍写旧目标；不是 core 或 recover。}

\bibitem{tools-delve-compat} Delve v1.25.2：Go 版本与 DWARF 兼容检查. 2025.\par
\url{https://github.com/go-delve/delve/blob/v1.25.2/pkg/goversion/compat.go}\par
{\small 项目官方版本源码；访问日期：2026-10-02。所读固定版本的常规 Go 版本检查范围为 1.\allowbreak{}22—1.\allowbreak{}25；riscv64 下限为 Go 1.\allowbreak{}24；读取 DWARF v5 要求构建 Delve 的 Go 工具链达到 1.\allowbreak{}25。版本检查范围不代表所有平台组合均已验证。}

\bibitem{eb-libbpf-overview} Linux: libbpf overview and CO-RE lifecycle.\par
\url{https://github.com/torvalds/linux/blob/adc218676eef25575469234709c2d87185ca223a/Documentation/bpf/libbpf/libbpf_overview.rst}\par
{\small\raggedright 固定版本：Linux v6.12; commit \texttt{adc21867\allowbreak{}6eef2557\allowbreak{}54692347\allowbreak{}09c2d871\allowbreak{}85ca223a}。\par
阅读范围：\nolinkurl{Documentation/bpf/libbpf/libbpf_overview.rst}：1–236 行。\par
证据边界：不把布局适配表述为程序语义兼容证明；不承诺所有目标内核具备相同能力。
}

\bibitem{eb-libbpf-load} libbpf: ELF collection, load ordering, section mapping and attach.\par
\url{https://github.com/libbpf/libbpf/blob/09b9e83102eb8ab9e540d36b4559c55f3bcdb95d/src/libbpf.c}\par
{\small\raggedright 固定版本：libbpf v1.5.0; commit \texttt{09b9e831\allowbreak{}02eb8ab9\allowbreak{}e540d36b\allowbreak{}4559c55f\allowbreak{}3bcdb95d}。\par
阅读范围：\nolinkurl{src/libbpf.c}：8035–8072 行；\nolinkurl{src/libbpf.c}：8550–8608 行；\nolinkurl{src/libbpf.c}：9456–9476 行；\nolinkurl{src/libbpf.c}：10820–10906 行；\nolinkurl{src/libbpf.c}：12127–12222 行；\nolinkurl{src/libbpf.c}：12460–12510 行；\nolinkurl{src/libbpf.c}：12665–12710 行。\par
证据边界：仅阅读列出的实际范围，不宣称逐行审完完整 libbpf.c；实际选择路径取决于目标内核支持。
}

\bibitem{eb-btf} Linux: BPF Type Format and BTF.ext.\par
\url{https://github.com/torvalds/linux/blob/adc218676eef25575469234709c2d87185ca223a/Documentation/bpf/btf.rst}\par
{\small\raggedright 固定版本：Linux v6.12; commit \texttt{adc21867\allowbreak{}6eef2557\allowbreak{}54692347\allowbreak{}09c2d871\allowbreak{}85ca223a}。\par
阅读范围：\nolinkurl{Documentation/bpf/btf.rst}：1–70 行；\nolinkurl{Documentation/bpf/btf.rst}：766–850 行；\nolinkurl{Documentation/bpf/btf.rst}：994–1085 行。\par
证据边界：文档的历史 pahole 能力描述不能泛化为现有所有版本；BTF 不表达全部业务语义。
}

\bibitem{eb-llvm-reloc} Linux: LLVM ELF relocations and CO-RE relocation records.\par
\url{https://github.com/torvalds/linux/blob/adc218676eef25575469234709c2d87185ca223a/Documentation/bpf/llvm_reloc.rst}\par
{\small\raggedright 固定版本：Linux v6.12; commit \texttt{adc21867\allowbreak{}6eef2557\allowbreak{}54692347\allowbreak{}09c2d871\allowbreak{}85ca223a}。\par
阅读范围：\nolinkurl{Documentation/bpf/llvm_reloc.rst}：1–215 行；\nolinkurl{Documentation/bpf/llvm_reloc.rst}：288–390 行。\par
证据边界：只展示固定版本定义，不推断未来重定位种类；示例结构偏移不是真实内核布局。
}

\bibitem{eb-core-relo} libbpf: calculating and patching CO-RE relocations.\par
\url{https://github.com/libbpf/libbpf/blob/09b9e83102eb8ab9e540d36b4559c55f3bcdb95d/src/relo_core.c}\par
{\small\raggedright 固定版本：libbpf v1.5.0; commit \texttt{09b9e831\allowbreak{}02eb8ab9\allowbreak{}e540d36b\allowbreak{}4559c55f\allowbreak{}3bcdb95d}。\par
阅读范围：\nolinkurl{src/relo_core.c}：879–960 行；\nolinkurl{src/relo_core.c}：1024–1108 行；\nolinkurl{src/relo_core.c}：1394–1422 行。\par
证据边界：结构重定位不保证单位、锁、生命周期或事件语义兼容；缺失字段的替代分支必须由程序设计提供。
}

\bibitem{eb-helpers} libbpf: BPF C section and map-definition helper macros.\par
\url{https://github.com/libbpf/libbpf/blob/09b9e83102eb8ab9e540d36b4559c55f3bcdb95d/src/bpf_helpers.h}\par
{\small\raggedright 固定版本：libbpf v1.5.0; commit \texttt{09b9e831\allowbreak{}02eb8ab9\allowbreak{}e540d36b\allowbreak{}4559c55f\allowbreak{}3bcdb95d}。\par
阅读范围：\nolinkurl{src/bpf_helpers.h}：1–95 行。\par
证据边界：宏本身不授予内核能力；当前环境没有安装可构建的 Clang/libbpf 依赖。
}

\bibitem{eb-verifier-doc} Linux: verifier register state, range tracking and pruning.\par
\url{https://github.com/torvalds/linux/blob/adc218676eef25575469234709c2d87185ca223a/Documentation/bpf/verifier.rst}\par
{\small\raggedright 固定版本：Linux v6.12; commit \texttt{adc21867\allowbreak{}6eef2557\allowbreak{}54692347\allowbreak{}09c2d871\allowbreak{}85ca223a}。\par
阅读范围：\nolinkurl{Documentation/bpf/verifier.rst}：1–180 行；\nolinkurl{Documentation/bpf/verifier.rst}：288–430 行。\par
证据边界：开头禁止循环的旧描述与同版本实际源码不一致，本文以源码为准；验证不是业务正确性或统计有效性证明。
}

\bibitem{eb-verifier-code} Linux: verifier state containment, liveness and loop exploration.\par
\url{https://github.com/torvalds/linux/blob/adc218676eef25575469234709c2d87185ca223a/kernel/bpf/verifier.c}\par
{\small\raggedright 固定版本：Linux v6.12; commit \texttt{adc21867\allowbreak{}6eef2557\allowbreak{}54692347\allowbreak{}09c2d871\allowbreak{}85ca223a}。\par
阅读范围：\nolinkurl{kernel/bpf/verifier.c}：6000–6120 行；\nolinkurl{kernel/bpf/verifier.c}：16300–16555 行；\nolinkurl{kernel/bpf/verifier.c}：17217–17335 行；\nolinkurl{kernel/bpf/verifier.c}：17577–17648 行；\nolinkurl{kernel/bpf/verifier.c}：17870–18015 行；\nolinkurl{kernel/bpf/verifier.c}：18312–18346 行。\par
证据边界：只对已阅读路径作机制解释，不宣称完整验证器形式化证明；具体接受集合受内核版本和程序类型影响。
}

\bibitem{eb-load-syscall} Linux: BPF program load permissions, verifier and runtime selection.\par
\url{https://github.com/torvalds/linux/blob/adc218676eef25575469234709c2d87185ca223a/kernel/bpf/syscall.c}\par
{\small\raggedright 固定版本：Linux v6.12; commit \texttt{adc21867\allowbreak{}6eef2557\allowbreak{}54692347\allowbreak{}09c2d871\allowbreak{}85ca223a}。\par
阅读范围：\nolinkurl{kernel/bpf/syscall.c}：2640–2710 行；\nolinkurl{kernel/bpf/syscall.c}：2835–2880 行。\par
证据边界：不把 root 或固定 capability 集合当作所有环境的充分条件；未执行特权加载。
}

\bibitem{eb-runtime-select} Linux: interpreter dispatch and JIT runtime selection.\par
\url{https://github.com/torvalds/linux/blob/adc218676eef25575469234709c2d87185ca223a/kernel/bpf/core.c}\par
{\small\raggedright 固定版本：Linux v6.12; commit \texttt{adc21867\allowbreak{}6eef2557\allowbreak{}54692347\allowbreak{}09c2d871\allowbreak{}85ca223a}。\par
阅读范围：\nolinkurl{kernel/bpf/core.c}：1708–1745 行；\nolinkurl{kernel/bpf/core.c}：2385–2449 行。\par
证据边界：不能承诺任意 JIT 失败都会回退解释器；未量化某平台执行开销。
}

\bibitem{eb-x86-jit} Linux x86: iterative BPF JIT instruction layout.\par
\url{https://github.com/torvalds/linux/blob/adc218676eef25575469234709c2d87185ca223a/arch/x86/net/bpf_jit_comp.c}\par
{\small\raggedright 固定版本：Linux v6.12; commit \texttt{adc21867\allowbreak{}6eef2557\allowbreak{}54692347\allowbreak{}09c2d871\allowbreak{}85ca223a}。\par
阅读范围：\nolinkurl{arch/x86/net/bpf_jit_comp.c}：3380–3455 行。\par
证据边界：架构特定实现，不能推广为所有 JIT 的相同算法；没有测量 JIT 与解释器性能差值。
}

\bibitem{eb-trampoline} Linux: BPF fentry/fexit trampoline construction and registration.\par
\url{https://github.com/torvalds/linux/blob/adc218676eef25575469234709c2d87185ca223a/kernel/bpf/trampoline.c}\par
{\small\raggedright 固定版本：Linux v6.12; commit \texttt{adc21867\allowbreak{}6eef2557\allowbreak{}54692347\allowbreak{}09c2d871\allowbreak{}85ca223a}。\par
阅读范围：\nolinkurl{kernel/bpf/trampoline.c}：1–135 行；\nolinkurl{kernel/bpf/trampoline.c}：200–231 行；\nolinkurl{kernel/bpf/trampoline.c}：393–477 行。\par
证据边界：具体支持依赖架构、内核和函数可挂载性；未给出跨机制统一开销排名。
}

\bibitem{eb-uprobe} Linux: uprobe inode/offset identity and software breakpoints.\par
\url{https://github.com/torvalds/linux/blob/adc218676eef25575469234709c2d87185ca223a/kernel/events/uprobes.c}\par
{\small\raggedright 固定版本：Linux v6.12; commit \texttt{adc21867\allowbreak{}6eef2557\allowbreak{}54692347\allowbreak{}09c2d871\allowbreak{}85ca223a}。\par
阅读范围：\nolinkurl{kernel/events/uprobes.c}：50–78 行；\nolinkurl{kernel/events/uprobes.c}：125–152 行；\nolinkurl{kernel/events/uprobes.c}：560–609 行。\par
证据边界：实际陷入与异位执行细节具有架构差异；不声称对 Go ABI 或任意返回路径均可直接套用 C 探针。
}

\bibitem{eb-map-array-doc} Linux: ARRAY/PERCPU\_ARRAY layout and concurrency.\par
\url{https://github.com/torvalds/linux/blob/adc218676eef25575469234709c2d87185ca223a/Documentation/bpf/map_array.rst}\par
{\small\raggedright 固定版本：Linux v6.12; commit \texttt{adc21867\allowbreak{}6eef2557\allowbreak{}54692347\allowbreak{}09c2d871\allowbreak{}85ca223a}。\par
阅读范围：\nolinkurl{Documentation/bpf/map_array.rst}：1–130 行。\par
证据边界：per-CPU 不自动提供多字段原子快照或全部上下文无竞态。
}

\bibitem{eb-map-array-code} Linux: per-CPU array lookup, copy and memory accounting.\par
\url{https://github.com/torvalds/linux/blob/adc218676eef25575469234709c2d87185ca223a/kernel/bpf/arraymap.c}\par
{\small\raggedright 固定版本：Linux v6.12; commit \texttt{adc21867\allowbreak{}6eef2557\allowbreak{}54692347\allowbreak{}09c2d871\allowbreak{}85ca223a}。\par
阅读范围：\nolinkurl{kernel/bpf/arraymap.c}：95–122 行；\nolinkurl{kernel/bpf/arraymap.c}：235–268 行；\nolinkurl{kernel/bpf/arraymap.c}：298–325 行；\nolinkurl{kernel/bpf/arraymap.c}：351–422 行；\nolinkurl{kernel/bpf/arraymap.c}：760–786 行。\par
证据边界：示例容量只计算 value 字节，不冒充完整分配开销；持续读出不是同时全局快照。
}

\bibitem{eb-ring-doc} Linux: BPF ring buffer reservation, ordering and wakeups.\par
\url{https://github.com/torvalds/linux/blob/adc218676eef25575469234709c2d87185ca223a/Documentation/bpf/ringbuf.rst}\par
{\small\raggedright 固定版本：Linux v6.12; commit \texttt{adc21867\allowbreak{}6eef2557\allowbreak{}54692347\allowbreak{}09c2d871\allowbreak{}85ca223a}。\par
阅读范围：\nolinkurl{Documentation/bpf/ringbuf.rst}：1–206 行。\par
证据边界：普通 reserve 的大小约束不代表 dynptr API 全部约束；保留顺序不是业务因果顺序；NMI 可能因锁竞争失败。
}

\bibitem{eb-ring-code} Linux: ring buffer capacity, pending position and publish protocol.\par
\url{https://github.com/torvalds/linux/blob/adc218676eef25575469234709c2d87185ca223a/kernel/bpf/ringbuf.c}\par
{\small\raggedright 固定版本：Linux v6.12; commit \texttt{adc21867\allowbreak{}6eef2557\allowbreak{}54692347\allowbreak{}09c2d871\allowbreak{}85ca223a}。\par
阅读范围：\nolinkurl{kernel/bpf/ringbuf.c}：118–154 行；\nolinkurl{kernel/bpf/ringbuf.c}：190–231 行；\nolinkurl{kernel/bpf/ringbuf.c}：408–538 行。\par
证据边界：容量例子未包含单独元数据页；不能把 reserve 失败全部解释为消费者已满；并发成本数字是教学假设。
}

\bibitem{eb-ring-consume} libbpf: acquire/consume/release ring consumer.\par
\url{https://github.com/libbpf/libbpf/blob/09b9e83102eb8ab9e540d36b4559c55f3bcdb95d/src/ringbuf.c}\par
{\small\raggedright 固定版本：libbpf v1.5.0; commit \texttt{09b9e831\allowbreak{}02eb8ab9\allowbreak{}e540d36b\allowbreak{}4559c55f\allowbreak{}3bcdb95d}。\par
阅读范围：\nolinkurl{src/ringbuf.c}：234–279 行。\par
证据边界：只分析 libbpf v1.5.0 消费协议；用户上传与持久化正确性仍需应用层处理。
}

\bibitem{eb-perf-abi} Linux perf ABI: sampling configuration and loss records.\par
\url{https://github.com/torvalds/linux/blob/adc218676eef25575469234709c2d87185ca223a/include/uapi/linux/perf_event.h}\par
{\small\raggedright 固定版本：Linux v6.12; commit \texttt{adc21867\allowbreak{}6eef2557\allowbreak{}54692347\allowbreak{}09c2d871\allowbreak{}85ca223a}。\par
阅读范围：\nolinkurl{include/uapi/linux/perf_event.h}：340–366 行；\nolinkurl{include/uapi/linux/perf_event.h}：392–443 行；\nolinkurl{include/uapi/linux/perf_event.h}：875–891 行；\nolinkurl{include/uapi/linux/perf_event.h}：1090–1105 行。\par
证据边界：perf 丢失记录不自动覆盖自建 BPF ring 的失败；计数缩放不能恢复未采调用栈。
}

\bibitem{eb-perf-core} Linux perf: interrupt throttling and sample frequency feedback.\par
\url{https://github.com/torvalds/linux/blob/adc218676eef25575469234709c2d87185ca223a/kernel/events/core.c}\par
{\small\raggedright 固定版本：Linux v6.12; commit \texttt{adc21867\allowbreak{}6eef2557\allowbreak{}54692347\allowbreak{}09c2d871\allowbreak{}85ca223a}。\par
阅读范围：\nolinkurl{kernel/events/core.c}：9652–9687 行。\par
证据边界：名义频率不是严格数量保证；没有测量当前主机节流状态。
}

\bibitem{eb-perf-record} Linux perf record: call graph modes and BPF off-CPU scope.\par
\url{https://github.com/torvalds/linux/blob/adc218676eef25575469234709c2d87185ca223a/tools/perf/Documentation/perf-record.txt}\par
{\small\raggedright 固定版本：Linux v6.12; commit \texttt{adc21867\allowbreak{}6eef2557\allowbreak{}54692347\allowbreak{}09c2d871\allowbreak{}85ca223a}。\par
阅读范围：\nolinkurl{tools/perf/Documentation/perf-record.txt}：278–323 行；\nolinkurl{tools/perf/Documentation/perf-record.txt}：809–837 行。\par
证据边界：不能推广为所有 eBPF profiler 均不支持 DWARF；用户与内核栈展开机制应分开。
}

\bibitem{eb-stack-helper} Linux BPF helper contract: context IDs, clocks and stack collection.\par
\url{https://github.com/torvalds/linux/blob/adc218676eef25575469234709c2d87185ca223a/include/uapi/linux/bpf.h}\par
{\small\raggedright 固定版本：Linux v6.12; commit \texttt{adc21867\allowbreak{}6eef2557\allowbreak{}54692347\allowbreak{}09c2d871\allowbreak{}85ca223a}。\par
阅读范围：\nolinkurl{include/uapi/linux/bpf.h}：55–90 行；\nolinkurl{include/uapi/linux/bpf.h}：1884–1902 行；\nolinkurl{include/uapi/linux/bpf.h}：1963–1987 行；\nolinkurl{include/uapi/linux/bpf.h}：2038–2088 行；\nolinkurl{include/uapi/linux/bpf.h}：2105–2125 行；\nolinkurl{include/uapi/linux/bpf.h}：2405–2465 行；\nolinkurl{include/uapi/linux/bpf.h}：3300–3355 行。\par
证据边界：文档建议提高系统限制不是本实验已执行动作；记录存在与栈完整是两个质量维度。
}

\bibitem{eb-stackmap} Linux: stack trace hashing, collision and helper call chain.\par
\url{https://github.com/torvalds/linux/blob/adc218676eef25575469234709c2d87185ca223a/kernel/bpf/stackmap.c}\par
{\small\raggedright 固定版本：Linux v6.12; commit \texttt{adc21867\allowbreak{}6eef2557\allowbreak{}54692347\allowbreak{}09c2d871\allowbreak{}85ca223a}。\par
阅读范围：\nolinkurl{kernel/bpf/stackmap.c}：227–326 行。\par
证据边界：ID 不能视为无限唯一永久身份；计数 map 与栈表轮转需协同。
}

\bibitem{eb-callchain} Linux perf: callchain selection and uretprobe return repair.\par
\url{https://github.com/torvalds/linux/blob/adc218676eef25575469234709c2d87185ca223a/kernel/events/callchain.c}\par
{\small\raggedright 固定版本：Linux v6.12; commit \texttt{adc21867\allowbreak{}6eef2557\allowbreak{}54692347\allowbreak{}09c2d871\allowbreak{}85ca223a}。\par
阅读范围：\nolinkurl{kernel/events/callchain.c}：190–268 行。\par
证据边界：跨任务和特殊上下文存在限制；不是 Go goroutine 全量快照。
}

\bibitem{eb-x86-callchain} Linux x86: user frame-pointer walk.\par
\url{https://github.com/torvalds/linux/blob/adc218676eef25575469234709c2d87185ca223a/arch/x86/events/core.c}\par
{\small\raggedright 固定版本：Linux v6.12; commit \texttt{adc21867\allowbreak{}6eef2557\allowbreak{}54692347\allowbreak{}09c2d871\allowbreak{}85ca223a}。\par
阅读范围：\nolinkurl{arch/x86/events/core.c}：2910–2975 行。\par
证据边界：这是固定 x86 内置路径，不代表所有架构或所有 profiler；不是通用 DWARF unwinder。
}

\bibitem{eb-go-runtime} Go runtime HACKING: G/M/P, user/system/signal stacks.\par
\url{https://github.com/golang/go/blob/6e676ab2b809d46623acb5988248d95d1eb7939c/src/runtime/HACKING.md}\par
{\small\raggedright 固定版本：Go go1.25.0; commit \texttt{6e676ab2\allowbreak{}b809d466\allowbreak{}23acb598\allowbreak{}8248d95d\allowbreak{}1eb7939c}。\par
阅读范围：\nolinkurl{src/runtime/HACKING.md}：1–110 行。\par
证据边界：runtime 文档自述可能落后，本文固定到明确提交；线程身份不能替代 goroutine 或请求身份。
}

\bibitem{eb-go-stack} Go runtime: copying stacks and pointer adjustment.\par
\url{https://github.com/golang/go/blob/6e676ab2b809d46623acb5988248d95d1eb7939c/src/runtime/stack.go}\par
{\small\raggedright 固定版本：Go go1.25.0; commit \texttt{6e676ab2\allowbreak{}b809d466\allowbreak{}23acb598\allowbreak{}8248d95d\allowbreak{}1eb7939c}。\par
阅读范围：\nolinkurl{src/runtime/stack.go}：899–989 行。\par
证据边界：不把 runtime 私有结构当跨版本稳定 ABI；未运行特定 Go unwinder 进行覆盖率验证。
}

\bibitem{eb-sched-trace} Linux: sched\_switch event schema.\par
\url{https://github.com/torvalds/linux/blob/adc218676eef25575469234709c2d87185ca223a/include/trace/events/sched.h}\par
{\small\raggedright 固定版本：Linux v6.12; commit \texttt{adc21867\allowbreak{}6eef2557\allowbreak{}54692347\allowbreak{}09c2d871\allowbreak{}85ca223a}。\par
阅读范围：\nolinkurl{include/trace/events/sched.h}：218–278 行。\par
证据边界：最小实验有意不读取上下文字段；线程事件不自动提供 goroutine 调度状态。
}

\bibitem{eb-sched-core} Linux scheduler: trace\_sched\_switch before context\_switch.\par
\url{https://github.com/torvalds/linux/blob/adc218676eef25575469234709c2d87185ca223a/kernel/sched/core.c}\par
{\small\raggedright 固定版本：Linux v6.12; commit \texttt{adc21867\allowbreak{}6eef2557\allowbreak{}54692347\allowbreak{}09c2d871\allowbreak{}85ca223a}。\par
阅读范围：\nolinkurl{kernel/sched/core.c}：6678–6697 行。\par
证据边界：固定 v6.12 调用点证据；不从此推断所有事件均能配成 off-CPU 时长。
}

\bibitem{eb-attach-api} libbpf public API: object, tracepoint, ring and possible CPUs.\par
\url{https://github.com/libbpf/libbpf/blob/09b9e83102eb8ab9e540d36b4559c55f3bcdb95d/src/libbpf.h}\par
{\small\raggedright 固定版本：libbpf v1.5.0; commit \texttt{09b9e831\allowbreak{}02eb8ab9\allowbreak{}e540d36b\allowbreak{}4559c55f\allowbreak{}3bcdb95d}。\par
阅读范围：\nolinkurl{src/libbpf.h}：210–239 行；\nolinkurl{src/libbpf.h}：754–774 行；\nolinkurl{src/libbpf.h}：1310–1337 行；\nolinkurl{src/libbpf.h}：1680–1696 行。\par
证据边界：只做源码 API 审查，未在当前环境编译或加载示例；实际错误与策略须在目标环境验证。
}

\bibitem{eb-isa} Linux BPF ISA: basic/wide encoding and program-local calls.\par
\url{https://github.com/torvalds/linux/blob/adc218676eef25575469234709c2d87185ca223a/Documentation/bpf/standardization/instruction-set.rst}\par
{\small\raggedright 固定版本：Linux v6.12; commit \texttt{adc21867\allowbreak{}6eef2557\allowbreak{}54692347\allowbreak{}09c2d871\allowbreak{}85ca223a}。\par
阅读范围：\nolinkurl{Documentation/bpf/standardization/instruction-set.rst}：1–140 行；\nolinkurl{Documentation/bpf/standardization/instruction-set.rst}：145–265 行；\nolinkurl{Documentation/bpf/standardization/instruction-set.rst}：490–580 行；\nolinkurl{Documentation/bpf/standardization/instruction-set.rst}：716–744 行。\par
证据边界：ISA 规范不等于每个内核与架构支持所有指令组；Linux helper 与 kfunc 的稳定性需看平台文档。
}

\bibitem{eb-design-qa} Linux BPF design Q\&A: calling convention and API stability.\par
\url{https://github.com/torvalds/linux/blob/adc218676eef25575469234709c2d87185ca223a/Documentation/bpf/bpf_design_QA.rst}\par
{\small\raggedright 固定版本：Linux v6.12; commit \texttt{adc21867\allowbreak{}6eef2557\allowbreak{}54692347\allowbreak{}09c2d871\allowbreak{}85ca223a}。\par
阅读范围：\nolinkurl{Documentation/bpf/bpf_design_QA.rst}：30–75 行；\nolinkurl{Documentation/bpf/bpf_design_QA.rst}：195–242 行。\par
证据边界：文档 zero overhead 措辞不作为测得的零开销承诺；512 B 栈上限需要结合实际调用链验证源码解释。
}

\bibitem{eb-kfuncs} Linux: kfunc registration, argument contracts and lifecycle.\par
\url{https://github.com/torvalds/linux/blob/adc218676eef25575469234709c2d87185ca223a/Documentation/bpf/kfuncs.rst}\par
{\small\raggedright 固定版本：Linux v6.12; commit \texttt{adc21867\allowbreak{}6eef2557\allowbreak{}54692347\allowbreak{}09c2d871\allowbreak{}85ca223a}。\par
阅读范围：\nolinkurl{Documentation/bpf/kfuncs.rst}：1–220 行；\nolinkurl{Documentation/bpf/kfuncs.rst}：290–445 行。\par
证据边界：存在名称或 BTF 不保证对当前程序类型可调用；需按目标版本和配置检测支持。
}

\bibitem{eb-verifier-state} Linux: verifier call-frame and stack-slot representation.\par
\url{https://github.com/torvalds/linux/blob/adc218676eef25575469234709c2d87185ca223a/include/linux/bpf_verifier.h}\par
{\small\raggedright 固定版本：Linux v6.12; commit \texttt{adc21867\allowbreak{}6eef2557\allowbreak{}54692347\allowbreak{}09c2d871\allowbreak{}85ca223a}。\par
阅读范围：\nolinkurl{include/linux/bpf_verifier.h}：335–365 行。\par
证据边界：调用帧常量不允许任意递归或每帧独占相同上限；tail call 混合调用还有更严格条件。
}

\bibitem{pp-scope} pprof 项目定位.\par
\url{https://github.com/google/pprof/blob/ebaad5f31b4d25e6f51abf4d91b9611cff92db7b/README.md}\par
{\small\raggedright 固定版本：google/pprof commit \texttt{ebaad5f3\allowbreak{}1b4d25e6\allowbreak{}f51abf4d\allowbreak{}91b9611c\allowbreak{}ff92db7b}。\par
阅读范围：\nolinkurl{README.md}：5–24 行。\par
证据边界：此项目自述明确不是Google官方产品；不把仓库所有者当作产品支持承诺。
}

\bibitem{pp-proto} profile.proto 数据契约.\par
\url{https://github.com/google/pprof/blob/ebaad5f31b4d25e6f51abf4d91b9611cff92db7b/proto/profile.proto}\par
{\small\raggedright 固定版本：google/pprof commit \texttt{ebaad5f3\allowbreak{}1b4d25e6\allowbreak{}f51abf4d\allowbreak{}91b9611c\allowbreak{}ff92db7b}。\par
阅读范围：\nolinkurl{proto/profile.proto}：20–37 行；\nolinkurl{proto/profile.proto}：46–233 行。\par
证据边界：格式只能描述字段，具体采样与估计语义仍由producer决定。
}

\bibitem{pp-cpu-api} Go CPU采集生命周期.\par
\url{https://github.com/golang/go/blob/6e676ab2b809d46623acb5988248d95d1eb7939c/src/runtime/pprof/pprof.go}\par
{\small\raggedright 固定版本：Go 1.25.0; commit \texttt{6e676ab2\allowbreak{}b809d466\allowbreak{}23acb598\allowbreak{}8248d95d\allowbreak{}1eb7939c}。\par
阅读范围：\nolinkurl{src/runtime/pprof/pprof.go}：806–894 行。\par
证据边界：针对Go1.25标准库API，不推广到其他语言或第三方producer。
}

\bibitem{pp-cpu-runtime} Go CPU信号入口与额外样本.\par
\url{https://github.com/golang/go/blob/6e676ab2b809d46623acb5988248d95d1eb7939c/src/runtime/cpuprof.go}\par
{\small\raggedright 固定版本：Go 1.25.0; commit \texttt{6e676ab2\allowbreak{}b809d466\allowbreak{}23acb598\allowbreak{}8248d95d\allowbreak{}1eb7939c}。\par
阅读范围：\nolinkurl{src/runtime/cpuprof.go}：5–195 行；\nolinkurl{src/runtime/cpuprof.go}：226–258 行。\par
证据边界：丢样计数不等于完整观测误差；无法恢复丢失的原栈。
}

\bibitem{pp-cpu-linux} Linux线程CPU时钟与去重.\par
\url{https://github.com/golang/go/blob/6e676ab2b809d46623acb5988248d95d1eb7939c/src/runtime/os_linux.go}\par
{\small\raggedright 固定版本：Go 1.25.0; commit \texttt{6e676ab2\allowbreak{}b809d466\allowbreak{}23acb598\allowbreak{}8248d95d\allowbreak{}1eb7939c}。\par
阅读范围：\nolinkurl{src/runtime/os_linux.go}：574–696 行。\par
证据边界：Linux实现细节；不宣称Windows/macOS采用相同信号机制。
}

\bibitem{pp-sigprof} 异步栈展开与标签读取.\par
\url{https://github.com/golang/go/blob/6e676ab2b809d46623acb5988248d95d1eb7939c/src/runtime/proc.go}\par
{\small\raggedright 固定版本：Go 1.25.0; commit \texttt{6e676ab2\allowbreak{}b809d466\allowbreak{}23acb598\allowbreak{}8248d95d\allowbreak{}1eb7939c}。\par
阅读范围：\nolinkurl{src/runtime/proc.go}：5574–5693 行。\par
证据边界：未知或抽象帧不自动定位到某个业务函数。
}

\bibitem{pp-buffer} CPU日志双环缓冲与溢出.\par
\url{https://github.com/golang/go/blob/6e676ab2b809d46623acb5988248d95d1eb7939c/src/runtime/profbuf.go}\par
{\small\raggedright 固定版本：Go 1.25.0; commit \texttt{6e676ab2\allowbreak{}b809d466\allowbreak{}23acb598\allowbreak{}8248d95d\allowbreak{}1eb7939c}。\par
阅读范围：\nolinkurl{src/runtime/profbuf.go}：12–86 行；\nolinkurl{src/runtime/profbuf.go}：303–328 行；\nolinkurl{src/runtime/profbuf.go}：429–475 行。\par
证据边界：writer在CPU入口还需序列化；不是所有runtime路径完全无锁。
}

\bibitem{pp-cpu-builder} 原始CPU记录到聚合值向量.\par
\url{https://github.com/golang/go/blob/6e676ab2b809d46623acb5988248d95d1eb7939c/src/runtime/pprof/proto.go}\par
{\small\raggedright 固定版本：Go 1.25.0; commit \texttt{6e676ab2\allowbreak{}b809d466\allowbreak{}23acb598\allowbreak{}8248d95d\allowbreak{}1eb7939c}。\par
阅读范围：\nolinkurl{src/runtime/pprof/proto.go}：274–393 行；\nolinkurl{src/runtime/pprof/proto.go}：397–483 行；\nolinkurl{src/runtime/pprof/proto.go}：585–645 行。\par
证据边界：不能据最终profile恢复事件顺序。
}

\bibitem{pp-heap-poisson} 堆采样的指数步长.\par
\url{https://github.com/golang/go/blob/6e676ab2b809d46623acb5988248d95d1eb7939c/src/runtime/malloc.go}\par
{\small\raggedright 固定版本：Go 1.25.0; commit \texttt{6e676ab2\allowbreak{}b809d466\allowbreak{}23acb598\allowbreak{}8248d95d\allowbreak{}1eb7939c}。\par
阅读范围：\nolinkurl{src/runtime/malloc.go}：1824–1884 行。\par
证据边界：数学推导采用理想Poisson模型；未测生产采样偏差。
}

\bibitem{pp-heap-coverage} 堆分配块、tiny allocator与内存槽覆盖.\par
\url{https://github.com/golang/go/blob/6e676ab2b809d46623acb5988248d95d1eb7939c/src/runtime/malloc.go}\par
{\small\raggedright 固定版本：Go 1.25.0; commit \texttt{6e676ab2\allowbreak{}b809d466\allowbreak{}23acb598\allowbreak{}8248d95d\allowbreak{}1eb7939c}。\par
阅读范围：\nolinkurl{src/runtime/malloc.go}：1124–1144 行；\nolinkurl{src/runtime/malloc.go}：1171–1179 行；\nolinkurl{src/runtime/malloc.go}：1229–1243 行。\par
证据边界：heap采集不等同所有语言级对象；不含栈上对象清单。
}

\bibitem{pp-profile-kinds} Go profile种类与归因契约.\par
\url{https://github.com/golang/go/blob/6e676ab2b809d46623acb5988248d95d1eb7939c/src/runtime/pprof/pprof.go}\par
{\small\raggedright 固定版本：Go 1.25.0; commit \texttt{6e676ab2\allowbreak{}b809d466\allowbreak{}23acb598\allowbreak{}8248d95d\allowbreak{}1eb7939c}。\par
阅读范围：\nolinkurl{src/runtime/pprof/pprof.go}：122–171 行；\nolinkurl{src/runtime/pprof/pprof.go}：594–626 行。\par
证据边界：文档层面的统计契约；具体边界需结合runtime实现。
}

\bibitem{pp-heap-scale} 堆样本逆概率外推.\par
\url{https://github.com/golang/go/blob/6e676ab2b809d46623acb5988248d95d1eb7939c/src/runtime/pprof/protomem.go}\par
{\small\raggedright 固定版本：Go 1.25.0; commit \texttt{6e676ab2\allowbreak{}b809d466\allowbreak{}23acb598\allowbreak{}8248d95d\allowbreak{}1eb7939c}。\par
阅读范围：\nolinkurl{src/runtime/pprof/protomem.go}：15–66 行；\nolinkurl{src/runtime/pprof/protomem.go}：69–93 行。\par
证据边界：不是精确逐对象清单；聚合和校正顺序不能任意交换。
}

\bibitem{pp-mprof} 堆与同步统计桶和发布周期.\par
\url{https://github.com/golang/go/blob/6e676ab2b809d46623acb5988248d95d1eb7939c/src/runtime/mprof.go}\par
{\small\raggedright 固定版本：Go 1.25.0; commit \texttt{6e676ab2\allowbreak{}b809d466\allowbreak{}23acb598\allowbreak{}8248d95d\allowbreak{}1eb7939c}。\par
阅读范围：\nolinkurl{src/runtime/mprof.go}：275–327 行；\nolinkurl{src/runtime/mprof.go}：364–479 行；\nolinkurl{src/runtime/mprof.go}：483–524 行；\nolinkurl{src/runtime/mprof.go}：795–880 行；\nolinkurl{src/runtime/mprof.go}：918–931 行。\par
证据边界：累计快照差分不自动按观测窗口裁剪事件。
}

\bibitem{pp-sema} 同步等待记录点与mutex近似收费.\par
\url{https://github.com/golang/go/blob/6e676ab2b809d46623acb5988248d95d1eb7939c/src/runtime/sema.go}\par
{\small\raggedright 固定版本：Go 1.25.0; commit \texttt{6e676ab2\allowbreak{}b809d466\allowbreak{}23acb598\allowbreak{}8248d95d\allowbreak{}1eb7939c}。\par
阅读范围：\nolinkurl{src/runtime/sema.go}：127–199 行；\nolinkurl{src/runtime/sema.go}：228–275 行。\par
证据边界：同步等待估计不是每锁每请求精确积分；需考虑再次睡眠和窗口边界。
}

\bibitem{pp-goroutine} 协作式goroutine快照.\par
\url{https://github.com/golang/go/blob/6e676ab2b809d46623acb5988248d95d1eb7939c/src/runtime/mprof.go}\par
{\small\raggedright 固定版本：Go 1.25.0; commit \texttt{6e676ab2\allowbreak{}b809d466\allowbreak{}23acb598\allowbreak{}8248d95d\allowbreak{}1eb7939c}。\par
阅读范围：\nolinkurl{src/runtime/mprof.go}：1254–1409 行；\nolinkurl{src/runtime/mprof.go}：1412–1508 行；\nolinkurl{src/runtime/mprof.go}：1630–1661 行；\nolinkurl{src/runtime/pprof/pprof.go}：743–769 行。\par
证据边界：不是所有系统G清单，也不是运行时间或请求身份表；debug>=2走runtime.Stack(buf,true)，在格式化全量traceback期间STW；不适用协作式短暂停顿解释。
}

\bibitem{pp-labels} Go标签安装和继承.\par
\url{https://github.com/golang/go/blob/6e676ab2b809d46623acb5988248d95d1eb7939c/src/runtime/pprof/runtime.go}\par
{\small\raggedright 固定版本：Go 1.25.0; commit \texttt{6e676ab2\allowbreak{}b809d466\allowbreak{}23acb598\allowbreak{}8248d95d\allowbreak{}1eb7939c}。\par
阅读范围：\nolinkurl{src/runtime/pprof/runtime.go}：32–51 行。\par
证据边界：现存worker需要在执行点安装标签。
}

\bibitem{pp-label-support} Go标签支持范围.\par
\url{https://github.com/golang/go/blob/6e676ab2b809d46623acb5988248d95d1eb7939c/src/runtime/pprof/label.go}\par
{\small\raggedright 固定版本：Go 1.25.0; commit \texttt{6e676ab2\allowbreak{}b809d466\allowbreak{}23acb598\allowbreak{}8248d95d\allowbreak{}1eb7939c}。\par
阅读范围：\nolinkurl{src/runtime/pprof/label.go}：44–59 行；\nolinkurl{src/runtime/pprof/label.go}：96–127 行。\par
证据边界：其他producer支持数字或其他标签不等于Go所有profile都支持。
}

\bibitem{pp-pc} return PC与内联帧恢复.\par
\url{https://github.com/golang/go/blob/6e676ab2b809d46623acb5988248d95d1eb7939c/src/runtime/symtab.go}\par
{\small\raggedright 固定版本：Go 1.25.0; commit \texttt{6e676ab2\allowbreak{}b809d466\allowbreak{}23acb598\allowbreak{}8248d95d\allowbreak{}1eb7939c}。\par
阅读范围：\nolinkurl{src/runtime/symtab.go}：95–167 行。\par
证据边界：不应对已归一化的Location.Address再次无条件减一。
}

\bibitem{pp-symbolizer} pprof本地符号化与Build ID核验.\par
\url{https://github.com/google/pprof/blob/ebaad5f31b4d25e6f51abf4d91b9611cff92db7b/internal/symbolizer/symbolizer.go}\par
{\small\raggedright 固定版本：google/pprof commit \texttt{ebaad5f3\allowbreak{}1b4d25e6\allowbreak{}f51abf4d\allowbreak{}91b9611c\allowbreak{}ff92db7b}。\par
阅读范围：\nolinkurl{internal/symbolizer/symbolizer.go}：129–244 行。\par
证据边界：缺失二进制或符号时不能保证完整源码行或原生栈。
}

\bibitem{pp-binutils} 进程地址到对象文件地址.\par
\url{https://github.com/google/pprof/blob/ebaad5f31b4d25e6f51abf4d91b9611cff92db7b/internal/binutils/binutils.go}\par
{\small\raggedright 固定版本：google/pprof commit \texttt{ebaad5f3\allowbreak{}1b4d25e6\allowbreak{}f51abf4d\allowbreak{}91b9611c\allowbreak{}ff92db7b}。\par
阅读范围：\nolinkurl{internal/binutils/binutils.go}：609–677 行。\par
证据边界：不能把任意ELF映射简单等同pc减mapping.start。
}

\bibitem{pp-merge} pprof合并、局部ID与归一化.\par
\url{https://github.com/google/pprof/blob/ebaad5f31b4d25e6f51abf4d91b9611cff92db7b/profile/merge.go}\par
{\small\raggedright 固定版本：google/pprof commit \texttt{ebaad5f3\allowbreak{}1b4d25e6\allowbreak{}f51abf4d\allowbreak{}91b9611c\allowbreak{}ff92db7b}。\par
阅读范围：\nolinkurl{profile/merge.go}：44–122 行；\nolinkurl{profile/merge.go}：155–245 行；\nolinkurl{profile/merge.go}：284–414 行；\nolinkurl{profile/merge.go}：426–465 行；\nolinkurl{profile/merge.go}：468–544 行。\par
证据边界：类型兼容检查不证明负载、版本和统计总体可比。
}

\bibitem{pp-filter} 样本筛选与帧隐藏.\par
\url{https://github.com/google/pprof/blob/ebaad5f31b4d25e6f51abf4d91b9611cff92db7b/profile/filter.go}\par
{\small\raggedright 固定版本：google/pprof commit \texttt{ebaad5f3\allowbreak{}1b4d25e6\allowbreak{}f51abf4d\allowbreak{}91b9611c\allowbreak{}ff92db7b}。\par
阅读范围：\nolinkurl{profile/filter.go}：20–78 行；\nolinkurl{profile/filter.go}：229–274 行。\par
证据边界：不能从隐藏后的图推断原始直接调用关系。
}

\bibitem{pp-scale} 粒度聚合与整数量化.\par
\url{https://github.com/google/pprof/blob/ebaad5f31b4d25e6f51abf4d91b9611cff92db7b/profile/profile.go}\par
{\small\raggedright 固定版本：google/pprof commit \texttt{ebaad5f3\allowbreak{}1b4d25e6\allowbreak{}f51abf4d\allowbreak{}91b9611c\allowbreak{}ff92db7b}。\par
阅读范围：\nolinkurl{profile/profile.go}：364–430 行；\nolinkurl{profile/profile.go}：445–489 行；\nolinkurl{profile/profile.go}：793–824 行。\par
证据边界：按更粗粒度聚合后不能逆向恢复原始记录；ScaleN的keepSample未计入ratio=1列，混合比例且其他列为零时可能丢弃仍有非零列的样本；独立反例已复现。
}

\bibitem{pp-diff} 差分的操作顺序.\par
\url{https://github.com/google/pprof/blob/ebaad5f31b4d25e6f51abf4d91b9611cff92db7b/internal/driver/fetch.go}\par
{\small\raggedright 固定版本：google/pprof commit \texttt{ebaad5f3\allowbreak{}1b4d25e6\allowbreak{}f51abf4d\allowbreak{}91b9611c\allowbreak{}ff92db7b}。\par
阅读范围：\nolinkurl{internal/driver/fetch.go}：58–89 行。\par
证据边界：普通差分保留正负变化；不是自动按请求量归一化。
}

\bibitem{pp-denominator} 报告分母与过滤顺序.\par
\url{https://github.com/google/pprof/blob/ebaad5f31b4d25e6f51abf4d91b9611cff92db7b/internal/driver/driver.go}\par
{\small\raggedright 固定版本：google/pprof commit \texttt{ebaad5f3\allowbreak{}1b4d25e6\allowbreak{}f51abf4d\allowbreak{}91b9611c\allowbreak{}ff92db7b}。\par
阅读范围：\nolinkurl{internal/driver/driver.go}：74–107 行。\par
证据边界：相同函数权重在两种分母下百分比不同，不意味着成本变化。
}

\bibitem{pp-total} 差分总量的绝对值与基线规则.\par
\url{https://github.com/google/pprof/blob/ebaad5f31b4d25e6f51abf4d91b9611cff92db7b/internal/report/report.go}\par
{\small\raggedright 固定版本：google/pprof commit \texttt{ebaad5f3\allowbreak{}1b4d25e6\allowbreak{}f51abf4d\allowbreak{}91b9611c\allowbreak{}ff92db7b}。\par
阅读范围：\nolinkurl{internal/report/report.go}：1279–1288 行；\nolinkurl{internal/report/report.go}：1307–1335 行。\par
证据边界：不能把net zero报告解读为无变化。
}

\bibitem{pp-percent} 百分比使用幅度.\par
\url{https://github.com/google/pprof/blob/ebaad5f31b4d25e6f51abf4d91b9611cff92db7b/internal/measurement/measurement.go}\par
{\small\raggedright 固定版本：google/pprof commit \texttt{ebaad5f3\allowbreak{}1b4d25e6\allowbreak{}f51abf4d\allowbreak{}91b9611c\allowbreak{}ff92db7b}。\par
阅读范围：\nolinkurl{internal/measurement/measurement.go}：165–179 行。\par
证据边界：差分的方向看带符号值，不能仅看flat\%。
}

\bibitem{pp-graph} flat/cum和调用图累积规则.\par
\url{https://github.com/google/pprof/blob/ebaad5f31b4d25e6f51abf4d91b9611cff92db7b/internal/graph/graph.go}\par
{\small\raggedright 固定版本：google/pprof commit \texttt{ebaad5f3\allowbreak{}1b4d25e6\allowbreak{}f51abf4d\allowbreak{}91b9611c\allowbreak{}ff92db7b}。\par
阅读范围：\nolinkurl{internal/graph/graph.go}：330–391 行；\nolinkurl{internal/graph/graph.go}：416–459 行；\nolinkurl{internal/graph/graph.go}：657–690 行。\par
证据边界：边权不是调用次数；普通调用图不完整保存路径或递归深度。
}

\bibitem{pp-flame} 火焰图栈数据与递归索引.\par
\url{https://github.com/google/pprof/blob/ebaad5f31b4d25e6f51abf4d91b9611cff92db7b/internal/report/stacks.go}\par
{\small\raggedright 固定版本：google/pprof commit \texttt{ebaad5f3\allowbreak{}1b4d25e6\allowbreak{}f51abf4d\allowbreak{}91b9611c\allowbreak{}ff92db7b}。\par
阅读范围：\nolinkurl{internal/report/stacks.go}：27–74 行；\nolinkurl{internal/report/stacks.go}：82–103 行；\nolinkurl{internal/report/stacks.go}：159–195 行。\par
证据边界：无逐事件时间信息，横向布局不是时间轴。
}

\bibitem{pp-teaching} 48条教学CPU样本的实际查询验证.\par
\nolinkurl{research/pprof/teaching/verification.json}\par
{\small\raggedright 固定版本：go version go1.22.12 linux/amd64; synthetic fixture v1。\par
阅读范围：\nolinkurl{research/pprof/teaching/verification.json}：all cases and checks。\par
证据边界：样本是教学构造；没有运行真实producer或模拟随机采样；未编译执行所读最新pprof提交；实际客户端是Go1.22.12内置版本；关闭符号化并使用虚构地址，不能据此宣称实际ELF符号化已验证。
}

\bibitem{pp-scale-repro} ScaleN混合比例边界的固定源码实际复现.\par
\nolinkurl{verification/pprof-review-repro/result.json}\par
{\small\raggedright 固定版本：google/pprof \texttt{ebaad5f3\allowbreak{}1b4d25e6\allowbreak{}f51abf4d\allowbreak{}91b9611c\allowbreak{}ff92db7b} profile subpackage; Go 1.22.12。\par
阅读范围：\nolinkurl{verification/pprof-review-repro/result.json}：full result；\nolinkurl{verification/pprof-review-repro/source-manifest.json}：unchanged copied source hashes。\par
证据边界：仅子包边界验证，不代表固定提交完整CLI已运行。
}

\bibitem{pf-release} Perfetto v58.2 release 与固定提交.\par
\url{https://github.com/google/perfetto/releases/tag/v58.2}\par
{\small\raggedright 固定版本：Perfetto v58.2; commit \texttt{add693d8\allowbreak{}b338ba95\allowbreak{}99dbcbc3\allowbreak{}e300b1ab\allowbreak{}8c000897}。\par
阅读范围：\nolinkurl{research/perfetto/release-latest.json}：tag\_name, published\_at, linux-amd64.zip asset digest；\nolinkurl{research/perfetto/pin.json}：git ls-remote tag and peeled commit。\par
证据边界：release元数据不替代运行时--version和下载字节摘要验证。
}

\bibitem{pf-service} Producer、consumer、service 与传输边界.\par
\url{https://github.com/google/perfetto/blob/add693d8b338ba9599dbcbc3e300b1ab8c000897/docs/concepts/service-model.md}\par
{\small\raggedright 固定版本：Perfetto v58.2; commit \texttt{add693d8\allowbreak{}b338ba95\allowbreak{}99dbcbc3\allowbreak{}e300b1ab\allowbreak{}8c000897}。\par
阅读范围：\nolinkurl{docs/concepts/service-model.md}：1–118 行。\par
证据边界：本文不采用文档中的默认SMB容量，另一官方页面默认值不同；通常进程与producer对应但并非强制一一对应。
}

\bibitem{pf-service-impl} Service 的信任边界、flush 与 scraping.\par
\url{https://github.com/google/perfetto/blob/add693d8b338ba9599dbcbc3e300b1ab8c000897/src/tracing/service/tracing_service_impl.cc}\par
{\small\raggedright 固定版本：Perfetto v58.2; commit \texttt{add693d8\allowbreak{}b338ba95\allowbreak{}99dbcbc3\allowbreak{}e300b1ab\allowbreak{}8c000897}。\par
阅读范围：\nolinkurl{src/tracing/service/tracing_service_impl.cc}：2117–2396 行；\nolinkurl{src/tracing/service/tracing_service_impl.cc}：2790–2810 行；\nolinkurl{src/tracing/service/tracing_service_impl.cc}：3412–3488 行；\nolinkurl{src/tracing/service/tracing_service_impl.cc}：3542–3664 行。\par
证据边界：不把flush解释为全局同时暂停或磁盘fsync；没有执行恶意producer或断连压力测试。
}

\bibitem{pf-shmem} 共享内存ABI、页面布局与chunk状态.\par
\url{https://github.com/google/perfetto/blob/add693d8b338ba9599dbcbc3e300b1ab8c000897/include/perfetto/ext/tracing/core/shared_memory_abi.h}\par
{\small\raggedright 固定版本：Perfetto v58.2; commit \texttt{add693d8\allowbreak{}b338ba95\allowbreak{}99dbcbc3\allowbreak{}e300b1ab\allowbreak{}8c000897}。\par
阅读范围：\nolinkurl{include/perfetto/ext/tracing/core/shared_memory_abi.h}：35–220 行。\par
证据边界：共享内存模拟和直接补丁优化存在不同状态路径；不是整个SDK全程无锁的证明。
}

\bibitem{pf-arbiter} SMB仲裁、回压策略与commit推进.\par
\url{https://github.com/google/perfetto/blob/add693d8b338ba9599dbcbc3e300b1ab8c000897/src/tracing/core/shared_memory_arbiter_impl.cc}\par
{\small\raggedright 固定版本：Perfetto v58.2; commit \texttt{add693d8\allowbreak{}b338ba95\allowbreak{}99dbcbc3\allowbreak{}e300b1ab\allowbreak{}8c000897}。\par
阅读范围：\nolinkurl{src/tracing/core/shared_memory_arbiter_impl.cc}：119–274 行；\nolinkurl{src/tracing/core/shared_memory_arbiter_impl.cc}：294–450 行。\par
证据边界：源码策略不等于特定工作负载的实际开销；未运行SMB溢出压力实验。
}

\bibitem{pf-writer} TraceWriter 分片序列化与延迟长度回填.\par
\url{https://github.com/google/perfetto/blob/add693d8b338ba9599dbcbc3e300b1ab8c000897/src/tracing/core/trace_writer_impl.cc}\par
{\small\raggedright 固定版本：Perfetto v58.2; commit \texttt{add693d8\allowbreak{}b338ba95\allowbreak{}99dbcbc3\allowbreak{}e300b1ab\allowbreak{}8c000897}。\par
阅读范围：\nolinkurl{src/tracing/core/trace_writer_impl.cc}：68–175 行；\nolinkurl{src/tracing/core/trace_writer_impl.cc}：230–390 行。\par
证据边界：packet大于SMB仍要求服务取得进展和策略允许；分片已交付不代表完整packet已可读取。
}

\bibitem{pf-commit} CommitDataRequest 协议.\par
\url{https://github.com/google/perfetto/blob/add693d8b338ba9599dbcbc3e300b1ab8c000897/protos/perfetto/common/commit_data_request.proto}\par
{\small\raggedright 固定版本：Perfetto v58.2; commit \texttt{add693d8\allowbreak{}b338ba95\allowbreak{}99dbcbc3\allowbreak{}e300b1ab\allowbreak{}8c000897}。\par
阅读范围：\nolinkurl{protos/perfetto/common/commit_data_request.proto}：20–94 行。\par
证据边界：commit不是应用事务提交；patch不能恢复已覆盖的旧内容。
}

\bibitem{pf-buffer-v1} 默认中央TraceBufferV1的数据结构与读取契约.\par
\url{https://github.com/google/perfetto/blob/add693d8b338ba9599dbcbc3e300b1ab8c000897/src/tracing/service/trace_buffer_v1.h}\par
{\small\raggedright 固定版本：Perfetto v58.2; commit \texttt{add693d8\allowbreak{}b338ba95\allowbreak{}99dbcbc3\allowbreak{}e300b1ab\allowbreak{}8c000897}。\par
阅读范围：\nolinkurl{src/tracing/service/trace_buffer_v1.h}：45–145 行。\par
证据边界：不能将此行为无条件推广到实验性V2；缓冲区含字节不保证全部可读。
}

\bibitem{pf-buffers} 三层缓冲、容量计算与丢失统计.\par
\url{https://github.com/google/perfetto/blob/add693d8b338ba9599dbcbc3e300b1ab8c000897/docs/concepts/buffers.md}\par
{\small\raggedright 固定版本：Perfetto v58.2; commit \texttt{add693d8\allowbreak{}b338ba95\allowbreak{}99dbcbc3\allowbreak{}e300b1ab\allowbreak{}8c000897}。\par
阅读范围：\nolinkurl{docs/concepts/buffers.md}：1–180 行；\nolinkurl{docs/concepts/buffers.md}：180–294 行；\nolinkurl{docs/concepts/buffers.md}：295–366 行。\par
证据边界：文档中的默认容量和严重级别不是跨版本承诺；本文容量数值属于明确给定假设的教学运算。
}

\bibitem{pf-config} TraceConfig缓冲政策、默认V1和增量清理.\par
\url{https://github.com/google/perfetto/blob/add693d8b338ba9599dbcbc3e300b1ab8c000897/protos/perfetto/config/trace_config.proto}\par
{\small\raggedright 固定版本：Perfetto v58.2; commit \texttt{add693d8\allowbreak{}b338ba95\allowbreak{}99dbcbc3\allowbreak{}e300b1ab\allowbreak{}8c000897}。\par
阅读范围：\nolinkurl{protos/perfetto/config/trace_config.proto}：35–110 行；\nolinkurl{protos/perfetto/config/trace_config.proto}：285–315 行；\nolinkurl{protos/perfetto/config/trace_config.proto}：467–489 行。\par
证据边界：配置能否生效还取决于已注册数据源和平台；中央buffer大小不等于producer SMB大小。
}

\bibitem{pf-trackevent} TrackEvent、自定义轨道与flow.\par
\url{https://github.com/google/perfetto/blob/add693d8b338ba9599dbcbc3e300b1ab8c000897/docs/instrumentation/track-events.md}\par
{\small\raggedright 固定版本：Perfetto v58.2; commit \texttt{add693d8\allowbreak{}b338ba95\allowbreak{}99dbcbc3\allowbreak{}e300b1ab\allowbreak{}8c000897}。\par
阅读范围：\nolinkurl{docs/instrumentation/track-events.md}：1–38 行；\nolinkurl{docs/instrumentation/track-events.md}：96–120 行；\nolinkurl{docs/instrumentation/track-events.md}：244–265 行；\nolinkurl{docs/instrumentation/track-events.md}：495–583 行。\par
证据边界：flow关联不自动证明独占归因或完整关键路径；文章提供设计原则，没有编译TrackEvent SDK埋点程序。
}

\bibitem{pf-packet} TracePacket序列身份与增量状态标志.\par
\url{https://github.com/google/perfetto/blob/add693d8b338ba9599dbcbc3e300b1ab8c000897/protos/perfetto/trace/trace_packet.proto}\par
{\small\raggedright 固定版本：Perfetto v58.2; commit \texttt{add693d8\allowbreak{}b338ba95\allowbreak{}99dbcbc3\allowbreak{}e300b1ab\allowbreak{}8c000897}。\par
阅读范围：\nolinkurl{protos/perfetto/trace/trace_packet.proto}：104–144 行；\nolinkurl{protos/perfetto/trace/trace_packet.proto}：350–475 行。\par
证据边界：不同sequence或代际中的相同整数ID不表示同一字符串；V2的损失原因位不能泛化到默认V1。
}

\bibitem{pf-reader} ProtoTraceReader的状态恢复、时钟处理与排序输入.\par
\url{https://github.com/google/perfetto/blob/add693d8b338ba9599dbcbc3e300b1ab8c000897/src/trace_processor/importers/proto/proto_trace_reader.cc}\par
{\small\raggedright 固定版本：Perfetto v58.2; commit \texttt{add693d8\allowbreak{}b338ba95\allowbreak{}99dbcbc3\allowbreak{}e300b1ab\allowbreak{}8c000897}。\par
阅读范围：\nolinkurl{src/trace_processor/importers/proto/proto_trace_reader.cc}：325–555 行；\nolinkurl{src/trace_processor/importers/proto/proto_trace_reader.cc}：1140–1200 行。\par
证据边界：本文没有演示所有多机clock同步模式；无解析错误不证明生产端没有漏报。
}

\bibitem{pf-clocks} 时钟域与ClockSnapshot换算.\par
\url{https://github.com/google/perfetto/blob/add693d8b338ba9599dbcbc3e300b1ab8c000897/docs/concepts/clock-sync.md}\par
{\small\raggedright 固定版本：Perfetto v58.2; commit \texttt{add693d8\allowbreak{}b338ba95\allowbreak{}99dbcbc3\allowbreak{}e300b1ab\allowbreak{}8c000897}。\par
阅读范围：\nolinkurl{docs/concepts/clock-sync.md}：1–254 行。\par
证据边界：文档的单个MONOTONIC数值描述有方向混淆，正文用独立偏移示例；时间单位不代表实际同步精度，未执行跨机时钟校准。
}

\bibitem{pf-architecture} Trace Processor总体架构.\par
\url{https://github.com/google/perfetto/blob/add693d8b338ba9599dbcbc3e300b1ab8c000897/docs/design-docs/trace-processor-architecture.md}\par
{\small\raggedright 固定版本：Perfetto v58.2; commit \texttt{add693d8\allowbreak{}b338ba95\allowbreak{}99dbcbc3\allowbreak{}e300b1ab\allowbreak{}8c000897}。\par
阅读范围：\nolinkurl{docs/design-docs/trace-processor-architecture.md}：1–120 行。\par
证据边界：文档部分入口路径已过时，v58.2实际入口以所引源码为准；流式输入不保证总内存为常量。
}

\bibitem{pf-detection} v58.2 importer registry 的格式检测.\par
\url{https://github.com/google/perfetto/blob/add693d8b338ba9599dbcbc3e300b1ab8c000897/src/trace_processor/util/trace_type.cc}\par
{\small\raggedright 固定版本：Perfetto v58.2; commit \texttt{add693d8\allowbreak{}b338ba95\allowbreak{}99dbcbc3\allowbreak{}e300b1ab\allowbreak{}8c000897}。\par
阅读范围：\nolinkurl{src/trace_processor/util/trace_type.cc}：70–96 行。\par
证据边界：此处不是固定万能魔数列表；可用导入器受构建/插件影响。
}

\bibitem{pf-importers} 内置导入器工厂清单.\par
\url{https://github.com/google/perfetto/blob/add693d8b338ba9599dbcbc3e300b1ab8c000897/src/trace_processor/importers/common/builtin_trace_importers.h}\par
{\small\raggedright 固定版本：Perfetto v58.2; commit \texttt{add693d8\allowbreak{}b338ba95\allowbreak{}99dbcbc3\allowbreak{}e300b1ab\allowbreak{}8c000897}。\par
阅读范围：\nolinkurl{src/trace_processor/importers/common/builtin_trace_importers.h}：20–54 行。\par
证据边界：清单缺失与真实负例结合才形成本文兼容性结论；不否认第三方转换工具或未来版本实现。
}

\bibitem{pf-sequence} PacketSequenceStateBuilder代际状态.\par
\url{https://github.com/google/perfetto/blob/add693d8b338ba9599dbcbc3e300b1ab8c000897/src/trace_processor/importers/proto/packet_sequence_state_builder.h}\par
{\small\raggedright 固定版本：Perfetto v58.2; commit \texttt{add693d8\allowbreak{}b338ba95\allowbreak{}99dbcbc3\allowbreak{}e300b1ab\allowbreak{}8c000897}。\par
阅读范围：\nolinkurl{src/trace_processor/importers/proto/packet_sequence_state_builder.h}：20–82 行。\par
证据边界：示例的编号复用用于解释状态绑定，不是新的采集性能测量。
}

\bibitem{pf-sorter} TraceSorter流内后缀排序与跨流合并.\par
\url{https://github.com/google/perfetto/blob/add693d8b338ba9599dbcbc3e300b1ab8c000897/src/trace_processor/sorter/trace_sorter.cc}\par
{\small\raggedright 固定版本：Perfetto v58.2; commit \texttt{add693d8\allowbreak{}b338ba95\allowbreak{}99dbcbc3\allowbreak{}e300b1ab\allowbreak{}8c000897}。\par
阅读范围：\nolinkurl{src/trace_processor/sorter/trace_sorter.cc}：72–110 行；\nolinkurl{src/trace_processor/sorter/trace_sorter.cc}：132–210 行。\par
证据边界：不能由分块读取推断常量内存；不同输入格式的排序策略可能不同。
}

\bibitem{pf-storage} TraceStorage列式表与字符串驻留.\par
\url{https://github.com/google/perfetto/blob/add693d8b338ba9599dbcbc3e300b1ab8c000897/src/trace_processor/storage/trace_storage.h}\par
{\small\raggedright 固定版本：Perfetto v58.2; commit \texttt{add693d8\allowbreak{}b338ba95\allowbreak{}99dbcbc3\allowbreak{}e300b1ab\allowbreak{}8c000897}。\par
阅读范围：\nolinkurl{src/trace_processor/storage/trace_storage.h}：90–210 行。\par
证据边界：不能将此简化成每个原始事件插入普通SQLite行；本文未基准测试SQL执行器或零拷贝性能。
}

\bibitem{pf-thread-state} 线程状态重建与重复唤醒处理.\par
\url{https://github.com/google/perfetto/blob/add693d8b338ba9599dbcbc3e300b1ab8c000897/src/trace_processor/importers/common/thread_state_tracker.cc}\par
{\small\raggedright 固定版本：Perfetto v58.2; commit \texttt{add693d8\allowbreak{}b338ba95\allowbreak{}99dbcbc3\allowbreak{}e300b1ab\allowbreak{}8c000897}。\par
阅读范围：\nolinkurl{src/trace_processor/importers/common/thread_state_tracker.cc}：38–104 行；\nolinkurl{src/trace_processor/importers/common/thread_state_tracker.cc}：135–210 行。\par
证据边界：OS线程状态不是Go goroutine状态；状态相交本身不证明请求归因。
}

\bibitem{pf-go-header} Go execution trace版本化文件头.\par
\url{https://github.com/golang/go/blob/6e676ab2b809d46623acb5988248d95d1eb7939c/src/internal/trace/version/version.go}\par
{\small\raggedright 固定版本：Go 1.25.0; commit \texttt{6e676ab2\allowbreak{}b809d466\allowbreak{}23acb598\allowbreak{}8248d95d\allowbreak{}1eb7939c}。\par
阅读范围：\nolinkurl{src/internal/trace/version/version.go}：61–80 行。\par
证据边界：源码为Go1.25，配套实际负例为Go1.22.12。
}

\bibitem{pf-go-main} go tool trace命令模式.\par
\url{https://github.com/golang/go/blob/6e676ab2b809d46623acb5988248d95d1eb7939c/src/cmd/trace/main.go}\par
{\small\raggedright 固定版本：Go 1.25.0; commit \texttt{6e676ab2\allowbreak{}b809d466\allowbreak{}23acb598\allowbreak{}8248d95d\allowbreak{}1eb7939c}。\par
阅读范围：\nolinkurl{src/cmd/trace/main.go}：30–65 行。\par
证据边界：固定Go1.25命令行，不跨版本推断参数。
}

\bibitem{pf-go-json} Go viewer JSON生成与视图选择.\par
\url{https://github.com/golang/go/blob/6e676ab2b809d46623acb5988248d95d1eb7939c/src/cmd/trace/jsontrace.go}\par
{\small\raggedright 固定版本：Go 1.25.0; commit \texttt{6e676ab2\allowbreak{}b809d466\allowbreak{}23acb598\allowbreak{}8248d95d\allowbreak{}1eb7939c}。\par
阅读范围：\nolinkurl{src/cmd/trace/jsontrace.go}：20–139 行。\par
证据边界：未将start/end参数擅自解释为时间单位；没有在配套中宣称完成Go viewer JSON到Perfetto的全语义验证。
}

\bibitem{pf-go-format} Go viewer使用的Chrome Trace Event数据结构.\par
\url{https://github.com/golang/go/blob/6e676ab2b809d46623acb5988248d95d1eb7939c/src/internal/trace/traceviewer/format/format.go}\par
{\small\raggedright 固定版本：Go 1.25.0; commit \texttt{6e676ab2\allowbreak{}b809d466\allowbreak{}23acb598\allowbreak{}8248d95d\allowbreak{}1eb7939c}。\par
阅读范围：\nolinkurl{src/internal/trace/traceviewer/format/format.go}：5–79 行。\par
证据边界：共享JSON字段不证明保留全部Go runtime语义。
}

\bibitem{pf-pprof} v58.2 pprof导入器的信息保留与损失.\par
\url{https://github.com/google/perfetto/blob/add693d8b338ba9599dbcbc3e300b1ab8c000897/src/trace_processor/importers/pprof/pprof_trace_reader.cc}\par
{\small\raggedright 固定版本：Perfetto v58.2; commit \texttt{add693d8\allowbreak{}b338ba95\allowbreak{}99dbcbc3\allowbreak{}e300b1ab\allowbreak{}8c000897}。\par
阅读范围：\nolinkurl{src/trace_processor/importers/pprof/pprof_trace_reader.cc}：80–110 行；\nolinkurl{src/trace_processor/importers/pprof/pprof_trace_reader.cc}：110–280 行；\nolinkurl{src/trace_processor/importers/pprof/pprof_trace_reader.cc}：285–380 行；\nolinkurl{src/trace_processor/importers/pprof/pprof_trace_reader.cc}：385–451 行。\par
证据边界：超过2\textasciicircum{}53整数可能失去精度；聚合导入不恢复时间、调度和请求因果；不承诺内部aggregate表名永久稳定。
}

\bibitem{pf-traceconv} traceconv输出模式与分发.\par
\url{https://github.com/google/perfetto/blob/add693d8b338ba9599dbcbc3e300b1ab8c000897/src/traceconv/traceconv.cc}\par
{\small\raggedright 固定版本：Perfetto v58.2; commit \texttt{add693d8\allowbreak{}b338ba95\allowbreak{}99dbcbc3\allowbreak{}e300b1ab\allowbreak{}8c000897}。\par
阅读范围：\nolinkurl{src/traceconv/traceconv.cc}：66–107 行；\nolinkurl{src/traceconv/traceconv.cc}：312–385 行。\par
证据边界：未执行全部转换模式；不能由聚合权重逆推真实调用时间线。
}

\bibitem{pf-launcher} v58.2源码中的launcher仍引用v57.2二进制.\par
\url{https://github.com/google/perfetto/blob/add693d8b338ba9599dbcbc3e300b1ab8c000897/tools/trace_processor}\par
{\small\raggedright 固定版本：Perfetto v58.2; commit \texttt{add693d8\allowbreak{}b338ba95\allowbreak{}99dbcbc3\allowbreak{}e300b1ab\allowbreak{}8c000897}。\par
阅读范围：\nolinkurl{research/perfetto/trace_processor-launcher.py}：30–76 行。\par
证据边界：本文选择v58.2官方release ZIP并校验，不运行此launcher下载。
}

\bibitem{pf-request} 合成ftrace经官方v58.2解析的8项断言.\par
\nolinkurl{verification/request-perfetto.json}\par
{\small\raggedright 固定版本：Perfetto v58.2; commit \texttt{add693d8\allowbreak{}b338ba95\allowbreak{}99dbcbc3\allowbreak{}e300b1ab\allowbreak{}8c000897}; synthetic request scenario v1。\par
阅读范围：\nolinkurl{examples/request/generate_trace.py}：full synthetic event scenario；\nolinkurl{examples/request/verify_perfetto.py}：full query and assertion protocol；\nolinkurl{verification/request-perfetto.json}：version, checks and raw SQL outputs。\par
证据边界：输入是合成ftrace，未从真实Linux内核采集；没有由时间重叠推断业务归因，线程归属由场景显式给定；不能据此评价生产环境采集开销或SMB丢失。
}

\bibitem{pf-validation} 实际二进制、原生Go负例与pprof导入证据.\par
\nolinkurl{verification/go-trace-import.json}\par
{\small\raggedright 固定版本：Perfetto v58.2; commit \texttt{add693d8\allowbreak{}b338ba95\allowbreak{}99dbcbc3\allowbreak{}e300b1ab\allowbreak{}8c000897}; Go runtime go1.22.12 linux/amd64。\par
阅读范围：\nolinkurl{verification/request-perfetto.json}：official binary full version；\nolinkurl{verification/go-trace-import.json}：full generated trace negative test；\nolinkurl{verification/pprof-perfetto-import.json}：full aggregate import and empty timing tables。\par
证据边界：负例固定输入和版本，不证明所有未来版本或第三方桥接均失败；pprof为教学构造文件，不是一次真实随机CPU采样；没有宣称完成Go viewer JSON桥接执行验证。
}

\end{thebibliography}
\end{document}
